\documentclass[11pt,a4paper,oneside]{book}

% ---------- packages ----------
\usepackage[utf8]{inputenc}
\usepackage[T1]{fontenc}
\usepackage{mathpazo}
\usepackage{amsmath,amssymb,amsthm,mathtools}
\usepackage[margin=2.6cm]{geometry}
\usepackage{microtype}
\usepackage{enumitem}
\usepackage{booktabs}
\usepackage{xcolor}
\usepackage{tikz}
\usetikzlibrary{arrows.meta,calc,patterns,positioning}
\usepackage[most]{tcolorbox}
\usepackage{fancyhdr}
\usepackage[hidelinks]{hyperref}

\setlength{\parskip}{0.4em}
\setlength{\parindent}{0pt}
\setlist{itemsep=0.2em,topsep=0.3em}

\pagestyle{fancy}
\fancyhf{}
\fancyhead[L]{\small\nouppercase{\leftmark}}
\fancyfoot[C]{\thepage}
\renewcommand{\headrulewidth}{0.3pt}

% ---------- colours and tikz styles ----------
\definecolor{ink}{HTML}{1F2933}
\definecolor{accent}{HTML}{B45309}
\definecolor{accent2}{HTML}{1F5FA8}
\definecolor{accent3}{HTML}{15803D}
\definecolor{mid}{HTML}{6B7280}
\tikzset{
  >=Stealth,
  axis/.style={-{Stealth[length=2mm]},thick,mid},
  vec/.style={-{Stealth[length=2.4mm]},very thick},
  guide/.style={dashed,mid,thin},
  pt/.style={circle,fill=ink,inner sep=1.3pt},
  lab/.style={font=\small,inner sep=1.5pt}
}

% ---------- macros ----------
\newcommand{\R}{\mathbb{R}}
\newcommand{\C}{\mathbb{C}}
\newcommand{\N}{\mathbb{N}}
\newcommand{\Z}{\mathbb{Z}}
\newcommand{\Q}{\mathbb{Q}}
\newcommand{\K}{\mathbb{K}}
\newcommand{\Prob}{\mathbb{P}}
\newcommand{\Exp}{\mathbb{E}}
\DeclareMathOperator{\Var}{Var}
\DeclareMathOperator{\Cov}{Cov}
\DeclareMathOperator{\Corr}{Corr}
\DeclareMathOperator{\se}{se}
\DeclareMathOperator{\bias}{bias}
\DeclareMathOperator{\MSE}{MSE}
\newcommand{\Normal}{\mathcal{N}}
\newcommand{\Fish}{\mathcal{I}}
\newcommand{\given}{\mid}
\DeclareMathOperator{\tr}{tr}
\DeclareMathOperator{\rank}{rank}
\DeclareMathOperator{\Null}{Null}
\DeclareMathOperator{\Col}{Col}
\DeclareMathOperator{\Row}{Row}
\DeclareMathOperator{\Span}{span}
\DeclareMathOperator{\diag}{diag}
\DeclareMathOperator{\proj}{proj}
\DeclareMathOperator{\sgn}{sgn}
\DeclareMathOperator{\adj}{adj}
\DeclareMathOperator{\im}{im}
\newcommand{\id}{\mathrm{id}}
\newcommand{\T}{^{\top}}
\newcommand{\herm}{^{*}}
\newcommand{\ip}[2]{\langle #1,#2\rangle}
\DeclarePairedDelimiter{\abs}{\lvert}{\rvert}
\DeclarePairedDelimiter{\norm}{\lVert}{\rVert}

% ---------- theorem environments ----------
\theoremstyle{plain}
\newtheorem{theorem}{Theorem}[chapter]
\newtheorem{lemma}[theorem]{Lemma}
\newtheorem{proposition}[theorem]{Proposition}
\newtheorem{corollary}[theorem]{Corollary}
\theoremstyle{definition}
\newtheorem{definition}[theorem]{Definition}
\newtheorem{example}[theorem]{Example}
\theoremstyle{remark}
\newtheorem*{remark}{Remark}
\newcounter{exercise}[chapter]
\newenvironment{exercise}{\refstepcounter{exercise}\par\medskip\noindent\textbf{Exercise \thechapter.\theexercise.}\ }{\par}

% ---------- call-out boxes ----------
\newtcolorbox{intuition}{breakable,colback=green!5,colframe=green!45!black,
  fonttitle=\bfseries,title={Intuition}}
\newtcolorbox{nuance}[1]{breakable,colback=orange!6,colframe=orange!75!black,
  fonttitle=\bfseries,title={Nuance: #1}}
\newtcolorbox{pitfall}[1]{breakable,colback=red!5,colframe=red!70!black,
  fonttitle=\bfseries,title={Common mistake: #1}}
\newtcolorbox{interview}[1]{breakable,colback=blue!5,colframe=blue!60!black,
  fonttitle=\bfseries,title={Interview lens: #1}}
\newtcolorbox{summarybox}{breakable,colback=black!4,colframe=black!60,
  fonttitle=\bfseries,title={Chapter summary}}

\title{\Huge\bfseries Inferential Statistics\\[0.3em]\Large From Estimation to Experiments and Regression}
\author{}
\date{}

\begin{document}
\frontmatter
\maketitle
\tableofcontents
\mainmatter

\part{Point Estimation and Statistical Intervals}
\chapter{Point Estimation of Parameters}

\section{Introduction}

A \emph{population} is described by a probability distribution that depends on one or more unknown constants called \emph{parameters}: the mean $\mu$ and variance $\sigma^2$ of a normal distribution, the rate $\lambda$ of an exponential distribution, the success probability $p$ of a Bernoulli trial. We cannot see the parameter. We see a \emph{random sample} $X_1,\dots,X_n$, independent random variables each with the population distribution, and we must use it to say something about the parameter.

\begin{definition}[Point estimator and estimate]
A \emph{point estimator} of a parameter $\theta$ is a statistic $\hat\Theta=h(X_1,\dots,X_n)$, a function of the sample whose value is intended to be close to $\theta$. After the data are observed, the number $\hat\theta=h(x_1,\dots,x_n)$ is the \emph{point estimate}.
\end{definition}

The distinction matters: the estimator is a random variable (it varies from sample to sample) and the estimate is one realised number. Everything we ask of an estimator, such as ``it is unbiased'' or ``it has small variance'', is a statement about the distribution of the random variable $\hat\Theta$.

\begin{example}
To estimate the mean $\mu$ of a population we may use the sample mean $\bar{X}$, the sample median $\tilde{X}$, or the midrange $(X_{\min}+X_{\max})/2$. All three are point estimators of $\mu$. Which is best depends on the population and on what we mean by ``best'', which is the subject of the next section.
\end{example}

\section{General Concepts of Point Estimation}

\subsection{Unbiased Estimators}

\begin{definition}[Unbiased estimator]
$\hat\Theta$ is an \emph{unbiased estimator} of $\theta$ if $\Exp[\hat\Theta]=\theta$ for every possible value of $\theta$. The \emph{bias} is $\bias(\hat\Theta)=\Exp[\hat\Theta]-\theta$.
\end{definition}

Unbiasedness says that the estimator is correct on average over repeated samples: it does not systematically overshoot or undershoot.

\begin{theorem}[Standard unbiased estimators]\label{thm:unbiased}
Let $X_1,\dots,X_n$ be a random sample from a distribution with mean $\mu$ and variance $\sigma^2$. Then
\begin{enumerate}[label=(\roman*)]
  \item $\bar{X}$ is an unbiased estimator of $\mu$;
  \item $S^2=\dfrac{1}{n-1}\sum_{i=1}^n(X_i-\bar{X})^2$ is an unbiased estimator of $\sigma^2$;
  \item if the data are Bernoulli trials with success probability $p$ and $X$ is the number of successes, $\hat{P}=X/n$ is an unbiased estimator of $p$.
\end{enumerate}
\end{theorem}

\begin{proof}
(i) $\Exp[\bar{X}]=\frac1n\sum\Exp[X_i]=\mu$. (iii) $\Exp[X]=np$, so $\Exp[\hat{P}]=p$. (ii) Expand the sum of squares about the true mean:
\[ \sum(X_i-\bar{X})^2=\sum(X_i-\mu)^2-n(\bar{X}-\mu)^2. \]
Taking expectations, $\Exp\Big[\sum(X_i-\bar{X})^2\Big]=n\sigma^2-n\cdot\frac{\sigma^2}n=(n-1)\sigma^2$, using $\Var[\bar{X}]=\sigma^2/n$. Dividing by $n-1$ gives $\Exp[S^2]=\sigma^2$.
\end{proof}

\begin{intuition}
The sum of squares around $\bar{X}$ is systematically smaller than the sum of squares around the true $\mu$, because $\bar{X}$ is the point that minimises the sum of squares for this very sample. Dividing by $n$ would therefore underestimate $\sigma^2$ on average, by the factor $(n-1)/n$. Dividing by $n-1$ exactly corrects this. The estimator $\frac1n\sum(X_i-\bar{X})^2$ is biased, with $\bias=-\sigma^2/n$.
\end{intuition}

\begin{example}
A sample of four measurements is $12.1,\ 11.8,\ 12.4,\ 12.3$. Then $\bar{x}=12.15$. The deviations $-0.05,-0.35,0.25,0.15$, squares $0.0025,0.1225,0.0625,0.0225$ with sum $\sum(x_i-\bar{x})^2=0.21$. So $s^2=0.21/3=0.07$ and $s=0.265$. These are point estimates of $\mu$ and $\sigma$.
\end{example}

\begin{nuance}{unbiased is not always sensible}
An unbiased estimator need not exist, and when it exists it need not be good. If $X$ is a single observation from a Poisson distribution, the only unbiased estimator of $e^{-2\lambda}$ is $(-1)^X$, which takes the values $\pm1$ for a quantity lying in $(0,1)$. Unbiasedness is a desirable property, not a requirement; Section~\ref{sec:mse} balances bias against variance.
\end{nuance}

\subsection{Proof that $S$ is a Biased Estimator of $\sigma$ (CD Only)}

Although $S^2$ is unbiased for $\sigma^2$, the square root $S$ is \emph{not} unbiased for $\sigma$.

\begin{theorem}\label{thm:sbiased}
If $\sigma^2>0$ and the sample is not degenerate, $\Exp[S]<\sigma$.
\end{theorem}

\begin{proof}
The square root is a strictly concave function, so Jensen's inequality gives $\Exp[\sqrt{S^2}]<\sqrt{\Exp[S^2]}=\sqrt{\sigma^2}=\sigma$ (strict inequality because $S^2$ is not constant).
\end{proof}

For a normal sample the bias is explicit: $\Exp[S]=c_4(n)\,\sigma$ with $c_4(n)=\sqrt{\dfrac2{n-1}}\,\dfrac{\Gamma(n/2)}{\Gamma((n-1)/2)}$; for example $c_4(5)\approx0.940$, $c_4(10)\approx0.973$, $c_4(25)\approx0.990$. The bias is small unless $n$ is small, and it is routinely ignored; control charts (Chapter~10) do use $c_4$ to correct it.

\subsection{Variance of a Point Estimator}

Among several unbiased estimators of the same parameter, we prefer the one with the smallest variance, because it varies least from sample to sample.

\begin{definition}[MVUE]
Among all unbiased estimators of $\theta$, the one with the smallest variance is the \emph{minimum variance unbiased estimator} (MVUE).
\end{definition}

\begin{example}[Mean versus median]
For a sample of size $n$ from $\Normal(\mu,\sigma^2)$ both $\bar{X}$ and the sample median $\tilde{X}$ are unbiased for $\mu$ (the normal is symmetric). But $\Var[\bar{X}]=\sigma^2/n$ while, for large $n$, $\Var[\tilde{X}]\approx\frac\pi2\cdot\frac{\sigma^2}n\approx1.57\,\sigma^2/n$. The median needs about $57\%$ more data to reach the same accuracy.
\end{example}

\begin{theorem}[$\bar{X}$ is the MVUE for the normal mean]\label{thm:mvue}
If $X_1,\dots,X_n$ is a random sample from $\Normal(\mu,\sigma^2)$, then $\bar{X}$ is the MVUE of $\mu$.
\end{theorem}

This is a consequence of the Cram\'er--Rao inequality, proved in Chapter~12. For non-normal populations the answer changes: for a uniform distribution on $(\mu-a,\mu+a)$ the midrange has variance of order $1/n^2$, far smaller than $\Var[\bar{X}]=a^2/(3n)$.

\subsection{Standard Error: Reporting a Point Estimator}

A point estimate alone is incomplete without some measure of its precision.

\begin{definition}[Standard error]
The \emph{standard error} of an estimator $\hat\Theta$ is its standard deviation, $\sigma_{\hat\Theta}=\sqrt{\Var[\hat\Theta]}$. If this depends on unknown parameters and these are replaced by estimates, the result is the \emph{estimated standard error}, written $\hat\sigma_{\hat\Theta}$ or $\se(\hat\Theta)$.
\end{definition}

For the sample mean, $\sigma_{\bar{X}}=\sigma/\sqrt n$, and since $\sigma$ is usually unknown the estimated standard error is $\se(\bar{X})=s/\sqrt n$. The standard practice is to report the estimate together with its standard error, as $\hat\theta\ \pm\ \se(\hat\theta)$.

\begin{example}
Thermal conductivity measurements on $n=10$ specimens give $\bar{x}=41.924$ and $s=0.284$. The estimated standard error of $\bar{x}$ is $s/\sqrt{10}=0.0898$, so we report $41.924\pm0.0898$. For a normal population, about $95\%$ of the time the estimate $\bar{X}$ is within $2\sigma_{\bar{X}}$ of $\mu$ (Section~\ref{sec:sampmeans}).
\end{example}

\begin{pitfall}{standard error is not standard deviation}
The standard deviation $s$ measures the spread of individual observations; the standard error $s/\sqrt n$ measures the spread of the \emph{mean} over repeated samples and shrinks as $n$ grows. Reporting $\pm s$ when one means $\pm\se$ overstates the uncertainty of the mean by a factor $\sqrt n$.
\end{pitfall}

\begin{example}[Textbook example: thermal conductivity of Armco iron]
Ten measurements of thermal conductivity at $100^\circ$F and $550$ W give $\bar{x}=41.924$ and $s=0.284$. The estimated standard error of $\bar{X}$ is $s/\sqrt{10}=0.0898$, about $0.2\%$ of the estimate, so the estimate is precise. If conductivity is normal, $2\,\widehat{\se}=0.1796$ and we are highly confident that the true mean lies in $41.924\pm0.1796$, i.e.\ $(41.744,\,42.104)$. Even if the estimator is not normal, as long as it is unbiased the estimate deviates from the true value by as much as four standard errors at most about $6\%$ of the time (Chebyshev), so a very conservative statement is that the parameter is within four standard errors of the estimate.
\end{example}

\subsection{Bootstrap Estimate of the Standard Error (CD Only)}

For many estimators there is no formula for the standard error. The \emph{bootstrap} estimates it by resampling the data.

\begin{definition}[Bootstrap standard error]
Given data $x_1,\dots,x_n$ and an estimator $\hat\theta=h(x_1,\dots,x_n)$:
\begin{enumerate}[label=(\arabic*)]
  \item draw a \emph{bootstrap sample} $x_1^*,\dots,x_n^*$ by sampling $n$ values \emph{with replacement} from the observed data;
  \item compute $\hat\theta^*=h(x_1^*,\dots,x_n^*)$;
  \item repeat steps 1--2 $B$ times, obtaining $\hat\theta^*_1,\dots,\hat\theta^*_B$;
  \item the bootstrap standard error is the sample standard deviation of these values,
  \[ \se_{\mathrm{boot}}(\hat\theta)=\sqrt{\frac1{B-1}\sum_{b=1}^B\big(\hat\theta^*_b-\bar{\theta}^*\big)^2},\qquad\bar{\theta}^*=\frac1B\sum_b\hat\theta^*_b. \]
\end{enumerate}
Typically $B$ between $200$ and $1000$ suffices for a standard error.
\end{definition}

The idea is that the observed data are our best picture of the population, so sampling from the data mimics sampling from the population. The method works for the median, the trimmed mean, a correlation coefficient, a ratio, or any statistic whatsoever. Chapter~13 develops the bootstrap in more detail.

\begin{example}
For the data $2,\ 5,\ 9$ the bootstrap sample $(5,5,9)$ gives the mean $6.33$ and $(2,2,2)$ gives $2$. Repeating many times, the standard deviation of the bootstrap means approaches $\sqrt{\sum(x_i-\bar{x})^2/n^2}$, which for $n=3$ is about $1.7$.
\end{example}

\subsection{Mean Square Error of an Estimator}\label{sec:mse}

Unbiased estimators are not always best: allowing a small bias can sharply reduce variance. We need a single criterion that accounts for both.

\begin{definition}[Mean square error]
$\MSE(\hat\Theta)=\Exp\big[(\hat\Theta-\theta)^2\big]$, the expected squared distance from the target.
\end{definition}

\begin{theorem}[Bias--variance decomposition]\label{thm:biasvar}
\[ \MSE(\hat\Theta)=\Var(\hat\Theta)+\big[\bias(\hat\Theta)\big]^2. \]
\end{theorem}

\begin{proof}
Write $\hat\Theta-\theta=(\hat\Theta-\Exp[\hat\Theta])+(\Exp[\hat\Theta]-\theta)$ and square. The cross term has expectation $2(\Exp[\hat\Theta]-\theta)\Exp(\hat\Theta-\Exp[\hat\Theta])=0$, leaving the variance plus the squared bias.
\end{proof}

For an unbiased estimator the MSE equals the variance. For biased ones the decomposition shows the trade-off.

\begin{example}[A biased estimator can beat an unbiased one]\label{ex:mseexample}
Let $X_1,\dots,X_n$ be normal with variance $\sigma^2$. Consider $\hat\sigma^2_c=c\sum(X_i-\bar{X})^2$. From $\sum(X_i-\bar{X})^2/\sigma^2\sim\chi^2_{n-1}$ (with mean $n-1$ and variance $2(n-1)$), the MSE is $\sigma^4\big[2c^2(n-1)+((n-1)c-1)^2\big]$, which is minimised at $c=1/(n+1)$. The unbiased choice $c=1/(n-1)$ has $\MSE=2\sigma^4/(n-1)$, while $c=1/(n+1)$ gives the smaller $\MSE=2\sigma^4/(n+1)$.
\end{example}

\begin{definition}[Relative efficiency, consistency]
The \emph{relative efficiency} of $\hat\Theta_2$ to $\hat\Theta_1$ is $\MSE(\hat\Theta_1)/\MSE(\hat\Theta_2)$; if it is less than $1$, $\hat\Theta_1$ is the more efficient estimator of $\theta$, in the sense that it has the smaller mean square error. An estimator is \emph{consistent} if $\hat\Theta_n\to\theta$ in probability as $n\to\infty$. A sufficient condition: $\bias\to0$ and $\Var\to0$ (then $\MSE\to0$).
\end{definition}

\section{Methods of Point Estimation}

So far we judged estimators that were proposed ad hoc. Two systematic recipes produce estimators from the distribution.

\subsection{Method of Moments}

\begin{definition}[Population and sample moments]
The $k$-th \emph{population moment} is $\Exp(X^k)$; the $k$-th \emph{sample moment} is $\frac1n\sum_{i=1}^nX_i^k$.
\end{definition}

\begin{definition}[Method of moments]
Suppose the distribution has $m$ unknown parameters $\theta_1,\dots,\theta_m$. Equate the first $m$ population moments (which are functions of the parameters) to the first $m$ sample moments, and solve the resulting $m$ equations for the parameters. The solutions are the \emph{moment estimators}.
\end{definition}

\begin{example}[Exponential]\label{ex:momexp}
For the exponential with rate $\lambda$, $\Exp[X]=1/\lambda$. Equating to $\bar{X}$ gives $\hat\lambda=1/\bar{X}$.
\end{example}

\begin{example}[Normal]
$\Exp[X]=\mu$ and $\Exp[X^2]=\sigma^2+\mu^2$. Equating to $\bar{X}$ and $\frac1n\sum X_i^2$: $\hat\mu=\bar{X}$ and $\hat\sigma^2=\frac1n\sum X_i^2-\bar{X}^2=\frac1n\sum(X_i-\bar{X})^2$, which is the \emph{biased} version of the sample variance.
\end{example}

\begin{example}[Gamma]
For a gamma distribution with shape $r$ and rate $\lambda$, $\Exp[X]=r/\lambda$ and $\Var[X]=r/\lambda^2$. Using the first two moments (equivalently the mean and the variance $\frac1n\sum(X_i-\bar{X})^2=:\hat v$): $\hat r=\bar{X}^2/\hat v$ and $\hat\lambda=\bar{X}/\hat v$.
\end{example}

\begin{example}[Uniform]
For a uniform distribution on $(0,\theta)$, $\Exp[X]=\theta/2$, so $\hat\theta=2\bar{X}$. This estimator is unbiased but can be silly: if the sample is $(1,2,9)$ then $\hat\theta=8$, which is smaller than the observed value $9$ and hence impossible. The maximum likelihood estimator below does not have this defect.
\end{example}

Moment estimators are simple and usually consistent, but they need not be efficient, and may fall outside the parameter space.

\begin{example}[Textbook example: failure-time data]
Eight electronic modules, tested at elevated temperature, have exponentially distributed failure times (hours) with $\bar{x}=21.65$ and $\sum x_i^2=6639.40$. The moment estimate of the exponential rate is $\hat\lambda=1/\bar{x}=0.0462$. Fitting a gamma distribution by moments gives $\hat{r}=1.30$ and $\hat\lambda=0.0599$; since $\hat{r}$ is slightly above $1$ (where the gamma reduces to the exponential) either distribution could be plausible.
\end{example}

\subsection{Method of Maximum Likelihood}

\begin{definition}[Likelihood function]
Let $X_1,\dots,X_n$ be a random sample from a distribution with probability mass or density function $f(x;\theta)$. The \emph{likelihood function} of the observed data is
\[ L(\theta)=f(x_1;\theta)\,f(x_2;\theta)\cdots f(x_n;\theta)=\prod_{i=1}^nf(x_i;\theta), \]
regarded as a function of $\theta$ with the data fixed. The \emph{maximum likelihood estimator} (MLE) is the value $\hat\theta$ that maximises $L(\theta)$.
\end{definition}

\begin{intuition}
For discrete data $L(\theta)$ is the probability, under parameter $\theta$, of obtaining exactly the data observed. The MLE chooses the parameter value under which the observed data were most probable. It is the value that makes the data least surprising.
\end{intuition}

Since the logarithm is increasing, maximising $L$ is the same as maximising the \emph{log-likelihood} $\ell(\theta)=\ln L(\theta)=\sum\ln f(x_i;\theta)$, which turns products into sums and is easier to differentiate. The usual recipe: write $\ell(\theta)$, solve $\ell'(\theta)=0$ (the \emph{likelihood equation}), and check that the solution is a maximum.

\begin{example}[Bernoulli]\label{ex:mlebern}
For $n$ trials with $x$ successes, $L(p)=p^x(1-p)^{n-x}$ and $\ell(p)=x\ln p+(n-x)\ln(1-p)$. Setting $\ell'(p)=x/p-(n-x)/(1-p)=0$ gives $\hat{p}=x/n$, the sample proportion.
\end{example}

\begin{example}[Exponential]\label{ex:mleexp}
For $f(x;\lambda)=\lambda e^{-\lambda x}$, $\ell(\lambda)=n\ln\lambda-\lambda\sum x_i$, and $\ell'(\lambda)=n/\lambda-\sum x_i=0$ gives $\hat\lambda=n/\sum x_i=1/\bar{x}$, the same as the moment estimator.
\end{example}

\begin{example}[Normal]\label{ex:mlenormal}
For $\Normal(\mu,\sigma^2)$, $\ell(\mu,\sigma^2)=-\frac n2\ln(2\pi\sigma^2)-\frac1{2\sigma^2}\sum(x_i-\mu)^2$. Setting both partial derivatives to zero: $\partial\ell/\partial\mu=\frac1{\sigma^2}\sum(x_i-\mu)=0$ gives $\hat\mu=\bar{x}$, and $\partial\ell/\partial\sigma^2=-\frac n{2\sigma^2}+\frac1{2\sigma^4}\sum(x_i-\mu)^2=0$ gives $\hat\sigma^2=\frac1n\sum(x_i-\bar{x})^2$. Thus the MLE of the variance is the \emph{biased} estimator with divisor $n$.
\end{example}

\begin{example}[Poisson]
For $f(x;\lambda)=e^{-\lambda}\lambda^x/x!$, $\ell(\lambda)=-n\lambda+(\sum x_i)\ln\lambda-\sum\ln x_i!$, so $\hat\lambda=\bar{x}$.
\end{example}

\begin{example}[Uniform: a likelihood without a stationary point]\label{ex:mleunif}
For $f(x;\theta)=1/\theta$ on $0\le x\le\theta$: $L(\theta)=\theta^{-n}$ provided $\theta\ge\max x_i$ and $0$ otherwise. The function $\theta^{-n}$ is decreasing, so the maximum is at the smallest admissible value, $\hat\theta=\max(x_1,\dots,x_n)$. No calculus is possible here; one must look at the shape of $L$ directly.
\end{example}

\begin{theorem}[Invariance property]\label{thm:invariance}
If $\hat\theta$ is the MLE of $\theta$, then for any function $g$ the MLE of $g(\theta)$ is $g(\hat\theta)$.
\end{theorem}

\begin{example}
Since $\hat\sigma^2=\frac1n\sum(x_i-\bar{x})^2$ is the MLE of $\sigma^2$ for a normal sample, the MLE of $\sigma$ is $\hat\sigma=\sqrt{\hat\sigma^2}$. For the exponential, the MLE of the mean $1/\lambda$ is $\bar{x}$ and of the probability $\Prob(X>t)=e^{-\lambda t}$ is $e^{-t/\bar{x}}$.
\end{example}

\begin{theorem}[Large-sample properties of the MLE]\label{thm:mleasymp}
Under mild regularity conditions, as $n\to\infty$ the MLE $\hat\Theta$ is
\begin{enumerate}[label=(\roman*)]
  \item asymptotically unbiased: $\Exp[\hat\Theta]\to\theta$;
  \item consistent;
  \item asymptotically normal, $\hat\Theta\approx\Normal\big(\theta,\ 1/(n\Fish(\theta))\big)$, with variance equal to the Cram\'er--Rao lower bound, where $\Fish(\theta)=\Exp\big[-\partial^2\ln f(X;\theta)/\partial\theta^2\big]$ is the Fisher information of one observation;
  \item consequently asymptotically the MVUE (no consistent estimator has smaller large-sample variance).
\end{enumerate}
\end{theorem}

The Fisher information and the Cram\'er--Rao bound are developed in Chapter~12. Practical consequence: the standard error of an MLE can be estimated by $1/\sqrt{n\Fish(\hat\theta)}$, the reciprocal square root of the curvature of the log-likelihood at its peak.

\begin{example}
For Bernoulli data, $\Fish(p)=1/[p(1-p)]$, so $\se(\hat{p})\approx\sqrt{\hat{p}(1-\hat{p})/n}$: the familiar standard error of a proportion. For the exponential, $\Fish(\lambda)=1/\lambda^2$ and $\se(\hat\lambda)\approx\hat\lambda/\sqrt n$.
\end{example}

\begin{interview}{when the MLE is awkward}
Likelihood equations often have no closed form (gamma shape, logistic regression) and are solved numerically by Newton's method or EM. The MLE can be biased in small samples (normal variance), can fail to be a stationary point (uniform), and can be non-unique or even infinite (logistic regression with perfectly separated classes). Know the Bernoulli, exponential, Poisson and normal cases by heart.
\end{interview}

\subsection{Bayesian Estimation of Parameters (CD Only)}

The methods so far treat $\theta$ as an unknown constant. The Bayesian approach treats $\theta$ as a random variable with a \emph{prior} distribution $f(\theta)$ expressing what is known before the data. Bayes' theorem updates it to the \emph{posterior} distribution after observing $x_1,\dots,x_n$:
\[ f(\theta\mid x_1,\dots,x_n)=\frac{f(x_1,\dots,x_n\mid\theta)\,f(\theta)}{\int f(x_1,\dots,x_n\mid\theta')f(\theta')\,d\theta'}\ \propto\ L(\theta)\,f(\theta). \]
The \emph{Bayes estimator} of $\theta$ is the mean of the posterior distribution, which minimises the expected squared error.

\begin{example}[Bernoulli with a beta prior]\label{ex:betabin}
Let $X\sim\mathrm{Binomial}(n,p)$ and let the prior on $p$ be $\mathrm{Beta}(a,b)$ with density $\propto p^{a-1}(1-p)^{b-1}$. Multiplying by the likelihood $p^x(1-p)^{n-x}$ shows that the posterior is $\mathrm{Beta}(a+x,\ b+n-x)$. The Bayes estimator is the posterior mean
\[ \hat{p}_B=\frac{a+x}{a+b+n}=\frac{n}{a+b+n}\cdot\frac xn+\frac{a+b}{a+b+n}\cdot\frac a{a+b}, \]
a weighted average of the sample proportion and the prior mean. As $n$ grows the data dominate and $\hat{p}_B\to x/n$.
\end{example}

\begin{example}[Normal mean with a normal prior]
Let $X_1,\dots,X_n$ be $\Normal(\mu,\sigma^2)$ with $\sigma^2$ known and prior $\mu\sim\Normal(\mu_0,\sigma_0^2)$. The posterior is normal with mean $\dfrac{(\sigma^2/n)\mu_0+\sigma_0^2\bar{x}}{\sigma^2/n+\sigma_0^2}$ and variance $\dfrac{(\sigma^2/n)\sigma_0^2}{\sigma^2/n+\sigma_0^2}$. The estimate pulls $\bar{x}$ toward the prior mean, more strongly when the prior is precise (small $\sigma_0^2$) or the data are few.
\end{example}

\section{Sampling Distributions}

\begin{definition}[Sampling distribution]
A statistic is a function of the sample and is therefore a random variable; its probability distribution is called its \emph{sampling distribution}.
\end{definition}

Every estimator and test statistic in this book has a sampling distribution, and the whole of inference rests on knowing it.

\begin{example}
A population consists of the equally likely values $1,3,5$. For random samples of size $2$ with replacement there are nine equally likely pairs, and the sample mean takes the following values: $(1,1)\to1$, $(1,3),(3,1)\to2$, $(1,5),(5,1),(3,3)\to3$, $(3,5),(5,3)\to4$, $(5,5)\to5$. So $\bar{X}$ has probabilities $1/9,\ 2/9,\ 3/9,\ 2/9,\ 1/9$ on $1,2,3,4,5$, with mean $3$ (equal to the population mean) and variance $4/3=\sigma^2/2$, where $\sigma^2=8/3$.
\end{example}

\section{Sampling Distribution of Means}\label{sec:sampmeans}

\begin{theorem}[Distribution of $\bar{X}$]\label{thm:dist-xbar}
Let $X_1,\dots,X_n$ be a random sample with mean $\mu$ and variance $\sigma^2$. Then $\Exp[\bar{X}]=\mu$ and $\Var[\bar{X}]=\sigma^2/n$. If the population is normal, $\bar{X}$ is exactly $\Normal(\mu,\sigma^2/n)$.
\end{theorem}

\begin{theorem}[Central limit theorem]\label{thm:clt-text}
If $X_1,\dots,X_n$ is a random sample from a population with mean $\mu$ and finite variance $\sigma^2$, then for large $n$
\[ Z=\frac{\bar{X}-\mu}{\sigma/\sqrt n}\ \text{ is approximately }\Normal(0,1),\qquad\text{i.e. }\bar{X}\approx\Normal\Big(\mu,\frac{\sigma^2}n\Big). \]
The approximation improves as $n$ grows. For most populations it is good when $n\ge30$; for a nearly symmetric population smaller $n$ suffices, and for a strongly skewed population larger $n$ may be required.
\end{theorem}

\begin{example}
An electronics company makes resistors with mean $\mu=100\ \Omega$ and standard deviation $\sigma=10\ \Omega$. For a sample of $n=25$, what is the probability that the average resistance is less than $95\ \Omega$? Since $\sigma_{\bar{X}}=10/\sqrt{25}=2$, $\Prob(\bar{X}<95)\approx\Prob\big(Z<(95-100)/2\big)=\Prob(Z<-2.5)=0.0062$.
\end{example}

\begin{theorem}[Difference of two sample means]\label{thm:diffmeans}
Let independent random samples of sizes $n_1,n_2$ come from populations with means $\mu_1,\mu_2$ and variances $\sigma_1^2,\sigma_2^2$. Then $\bar{X}_1-\bar{X}_2$ has mean $\mu_1-\mu_2$ and variance $\sigma_1^2/n_1+\sigma_2^2/n_2$, and
\[ Z=\frac{(\bar{X}_1-\bar{X}_2)-(\mu_1-\mu_2)}{\sqrt{\sigma_1^2/n_1+\sigma_2^2/n_2}} \]
is exactly $\Normal(0,1)$ if both populations are normal, and approximately so for large $n_1,n_2$.
\end{theorem}

\begin{example}
Two machines fill cartons. Machine 1: $\mu_1=500$, $\sigma_1=6$; machine 2: $\mu_2=497$, $\sigma_2=5$. For $n_1=n_2=36$, $\bar{X}_1-\bar{X}_2$ has mean $3$ and standard deviation $\sqrt{36/36+25/36}=1.30$. The probability that machine 1's sample mean exceeds machine 2's is $\Prob(Z>(0-3)/1.30)=\Prob(Z>-2.30)=0.989$.
\end{example}

\begin{summarybox}
\begin{itemize}
  \item An estimator is a statistic; its value is an estimate. Judge estimators by bias, variance and $\MSE=\Var+\bias^2$.
  \item $\bar{X}$, $S^2$ and $\hat{P}$ are unbiased; $S$ is slightly biased low; the mean beats the median for normal data (the median has about $57\%$ larger variance).
  \item Report estimates with their standard errors; the bootstrap gives standard errors without a formula.
  \item Method of moments equates sample and population moments; maximum likelihood maximises $L(\theta)=\prod f(x_i;\theta)$, is invariant under transformations and is asymptotically unbiased, normal and efficient.
  \item Bayesian estimation combines a prior with the likelihood; the posterior mean is a weighted average of the data estimate and the prior mean.
  \item $\bar{X}\approx\Normal(\mu,\sigma^2/n)$ (exactly for normal populations, by the CLT otherwise); $\bar{X}_1-\bar{X}_2$ has variance $\sigma_1^2/n_1+\sigma_2^2/n_2$.
\end{itemize}
\end{summarybox}

\begin{example}[Textbook examples: sampling distributions]
\emph{Uniform.} If $X\sim U(4,6)$ then $\mu=5$, $\sigma^2=(6-4)^2/12=1/3$, and for $n=40$ the central limit theorem gives $\bar{X}\approx\Normal(5,\ \sigma^2/n=1/120)$.
\emph{Two samples.} A jet-engine component has life with mean $5000$ h and $\sigma=40$ (old process), and after improvement mean $5050$ h and $\sigma=30$ (new process); samples $n_1=16$ old and $n_2=25$ new are taken. $\bar{X}_2-\bar{X}_1\sim\Normal\big(50,\ 30^2/25+40^2/16=36+100=136\big)$, so $\Prob(\bar{X}_2-\bar{X}_1\ge25)=\Prob\big(Z\ge(25-50)/\sqrt{136}=-2.14\big)=0.9838$.
\emph{Likelihood and sample size.} For the exponential failure data the log-likelihood is maximised at $\hat\lambda=0.0462$ but is quite flat for $n=8$, which means $\lambda$ is not estimated very precisely; keeping $\bar{x}=21.65$ fixed, the log-likelihood becomes steeper (more sharply peaked) for $n=20$ and $n=40$. For the gamma model the maximum likelihood estimates, found numerically, are about $\hat{r}=1.75$ and $\hat\lambda=0.08$.
\end{example}

\section{Exercises}

\begin{exercise}
Let $X_1,X_2,\dots,X_7$ be a random sample from a population with mean $\mu$ and variance $\sigma^2$. Consider $\hat\Theta_1=(X_1+X_2+\dots+X_7)/7$ and $\hat\Theta_2=(2X_1-X_6+X_4)/2$. Is either biased? Find the variance of each and say which is preferable.
\end{exercise}

\begin{exercise}
Show that $\sum_{i}(X_i-\bar{X})^2=\sum_i X_i^2-n\bar{X}^2$ and that $\frac1n\sum(X_i-\bar{X})^2$ has bias $-\sigma^2/n$. Find the MSE of this biased estimator for a normal sample and compare it with the MSE of $S^2$.
\end{exercise}

\begin{exercise}
For a random sample of size $n$ from a Poisson distribution with parameter $\lambda$, find the method of moments estimator and the maximum likelihood estimator of $\lambda$, their bias and variance, and the estimated standard error.
\end{exercise}

\begin{exercise}
Find the maximum likelihood estimator of $\theta$ for the density $f(x;\theta)=\theta x^{\theta-1}$, $0<x<1$, $\theta>0$, and the moment estimator. Which would you prefer when $n=3$ and the data are $0.2,0.5,0.9$?
\end{exercise}

\begin{exercise}
A sample of size $n=20$ from an exponential distribution has $\sum x_i=64$. Find the MLE of $\lambda$, of the mean $1/\lambda$ and of $\Prob(X>5)$ (use invariance), and the estimated standard error of $\hat\lambda$.
\end{exercise}

\begin{exercise}
Show that for $\hat\sigma^2_c=c\sum(X_i-\bar{X})^2$ the MSE is $\sigma^4[2c^2(n-1)+((n-1)c-1)^2]$ and verify the minimiser $c=1/(n+1)$.
\end{exercise}

\begin{exercise}
In a quality study $x=12$ defectives are found in $n=100$ items. Using a $\mathrm{Beta}(2,18)$ prior, compute the posterior distribution of $p$ and the Bayes estimate. Compare with the MLE.
\end{exercise}

\begin{exercise}
The scores of a test have $\mu=75$, $\sigma=12$. For a random sample of $36$ students find $\Prob(\bar{X}>78)$, $\Prob(72<\bar{X}<77)$, and the value $c$ such that $\Prob(\bar{X}>c)=0.05$.
\end{exercise}

\begin{exercise}
Write a short program that draws $B=1000$ bootstrap samples from the data $14,\ 18,\ 11,\ 16,\ 25,\ 9,\ 13$ and computes the bootstrap standard error of (a) the mean and (b) the median. Compare (a) with $s/\sqrt n$.
\end{exercise}

\chapter{Statistical Intervals for a Single Sample}

\section{Introduction}

An estimate $\hat\theta$ will not equal $\theta$ exactly. A \emph{statistical interval} says how far off it may be. Three kinds are used.

\begin{itemize}
  \item A \emph{confidence interval} contains a population parameter (such as $\mu$) with a stated probability, over repeated sampling.
  \item A \emph{prediction interval} contains a single future observation with a stated probability.
  \item A \emph{tolerance interval} contains at least a stated proportion of the population, with a stated confidence.
\end{itemize}

For the same data the three have different widths:
\begin{itemize}
  \item the confidence interval is narrowest. It concerns only the mean, whose uncertainty shrinks like $1/\sqrt n$;
  \item the prediction interval is wider. It must also cover the variability of one new value;
  \item the tolerance interval is widest.
\end{itemize}

\section{Confidence Interval on the Mean of a Normal Distribution, Variance Known}

\subsection{Development of the Confidence Interval and Its Basic Properties}

Let $X_1,\dots,X_n$ be a random sample from $\Normal(\mu,\sigma^2)$ with $\sigma^2$ known and $\mu$ unknown. From Chapter~1,
\[ Z=\frac{\bar{X}-\mu}{\sigma/\sqrt n}\sim\Normal(0,1). \]
Let $z_{\alpha/2}$ cut off a right-tail area $\alpha/2$ under the standard normal curve. ($z_{0.025}=1.96$, $z_{0.05}=1.645$, $z_{0.005}=2.576$.) Then $\Prob(-z_{\alpha/2}\le Z\le z_{\alpha/2})=1-\alpha$. Solve the inequalities for $\mu$:
\[ \Prob\Big(\bar{X}-z_{\alpha/2}\frac{\sigma}{\sqrt n}\ \le\ \mu\ \le\ \bar{X}+z_{\alpha/2}\frac\sigma{\sqrt n}\Big)=1-\alpha. \]

\begin{definition}[Confidence interval on the mean, $\sigma$ known]
If $\bar{x}$ is the mean of a random sample of size $n$ from a normal population with known variance $\sigma^2$, a $100(1-\alpha)\%$ confidence interval for $\mu$ is
\[ \bar{x}-z_{\alpha/2}\frac\sigma{\sqrt n}\ \le\ \mu\ \le\ \bar{x}+z_{\alpha/2}\frac\sigma{\sqrt n}. \]
$1-\alpha$ is the \emph{confidence level}. The two endpoints are the \emph{lower and upper confidence limits}.
\end{definition}

\begin{example}\label{ex:ci-known}
Ten measurements of the impact energy of a steel have $\bar{x}=64.46$ J. $\sigma=1$ J is known.

A $95\%$ interval: $64.46\pm1.96\cdot1/\sqrt{10}=64.46\pm0.62$, that is $63.84\le\mu\le65.08$.
\end{example}

\begin{pitfall}{what ``95\% confidence'' means}
The $95\%$ describes the \emph{procedure}, not the particular interval.
\begin{itemize}
  \item Before the data are collected, the random interval $\bar{X}\pm z_{\alpha/2}\sigma/\sqrt n$ contains the fixed $\mu$ with probability $0.95$.
  \item After the data are in, the computed interval either contains $\mu$ or it does not. You cannot say ``the probability that $\mu$ lies in $(63.84,65.08)$ is $0.95$''. In the frequentist view $\mu$ is not random.
  \item Correct reading: if you repeated the procedure on many independent samples, about $95\%$ of the intervals would cover $\mu$.
\end{itemize}
\end{pitfall}

\begin{figure}[h]
\centering
\begin{tikzpicture}[scale=1.0]
  \draw[thick,mid] (0,-0.3)--(0,4.3); \node[lab] at (0,4.55){$\mu$};
  \foreach \i/\c/\w/\col in {0/0.3/0.9/accent2,1/-0.4/0.9/accent2,2/0.1/0.9/accent2,3/1.3/0.9/accent,4/-0.2/0.9/accent2,5/0.5/0.9/accent2,6/-1.2/0.9/accent,7/0.0/0.9/accent2,8/-0.6/0.9/accent2,9/0.7/0.9/accent2}{
    \draw[thick,\col] ({\c-\w},{4.1-0.42*\i})--({\c+\w},{4.1-0.42*\i});
    \node[pt,\col] at (\c,{4.1-0.42*\i}){};}
\end{tikzpicture}
\caption{Ten confidence intervals from ten samples. The vertical line is the true $\mu$; the two intervals that miss it are drawn in orange. Over many samples, $95\%$ of $95\%$ intervals cover $\mu$.}
\end{figure}

\subsection*{Properties}

\begin{itemize}
  \item \textbf{Width.} The length is $2z_{\alpha/2}\sigma/\sqrt n$, called the \emph{precision}. It shrinks with larger $n$ and larger $\alpha$ (lower confidence). It grows with larger $\sigma$ and higher confidence. $E=z_{\alpha/2}\sigma/\sqrt n$ is the \emph{error} (margin of error) in estimating $\mu$ by $\bar{x}$.
  \item \textbf{Confidence against precision.} More confidence costs width. Going from $95\%$ to $99\%$ replaces $1.96$ by $2.576$, a $31\%$ longer interval.
\end{itemize}

\subsection{Choice of Sample Size}

With confidence $100(1-\alpha)\%$ we are sure the error $\abs{\bar{x}-\mu}$ does not exceed $E=z_{\alpha/2}\sigma/\sqrt n$. Solve for $n$:

\begin{theorem}[Sample size for estimating a mean]\label{thm:samplesize-mean}
To be $100(1-\alpha)\%$ confident that the error in estimating $\mu$ by $\bar{x}$ is at most $E$, use
\[ n=\left(\frac{z_{\alpha/2}\,\sigma}{E}\right)^2,\qquad\text{rounded up to the next integer.} \]
\end{theorem}

\begin{example}
In Example~\ref{ex:ci-known}, how large must $n$ be for a $95\%$ error of at most $0.5$ J? $n=(1.96\cdot1/0.5)^2=15.37$, so $n=16$.
\end{example}

$n$ varies as $1/E^2$. Halving the error multiplies the sample size by four.

\subsection{One-sided Confidence Bounds}

Sometimes only an upper or lower limit matters (a maximum emission level, a minimum strength). The two kinds of statement are built differently. Table~\ref{tab:onetwo} lists both for every parameter in this chapter.

\begin{definition}[One-sided confidence bounds]
\begin{itemize}
  \item $100(1-\alpha)\%$ \emph{upper confidence bound}: $\mu\le\bar{x}+z_\alpha\sigma/\sqrt n$.
  \item $100(1-\alpha)\%$ \emph{lower confidence bound}: $\bar{x}-z_\alpha\sigma/\sqrt n\le\mu$.
\end{itemize}
Note $z_\alpha$, not $z_{\alpha/2}$. All of the error probability $\alpha$ sits in \emph{one} tail, because only one end is claimed. Each bound alone has coverage $1-\alpha$.
\end{definition}

\begin{example}
Same data as above.
\begin{itemize}
  \item $95\%$ lower bound: $64.46-1.645/\sqrt{10}=63.94$. Tighter than the two-sided lower limit $63.84=64.46-1.96/\sqrt{10}$.
  \item $95\%$ upper bound: $64.46+0.52=64.98$. The two-sided upper limit is $65.08$.
\end{itemize}
\end{example}

\begin{nuance}{one-sided versus two-sided: the rules}
\begin{enumerate}
  \item \textbf{Two-sided} $100(1-\alpha)\%$ interval: $\alpha/2$ in each tail. The critical value has subscript $\alpha/2$ ($z_{\alpha/2}$, $t_{\alpha/2,\nu}$, $\chi^2_{\alpha/2,\nu}$ and $\chi^2_{1-\alpha/2,\nu}$, $f_{\alpha/2}$ and $f_{1-\alpha/2}$).
  \item \textbf{One-sided} $100(1-\alpha)\%$ bound: all of $\alpha$ in one tail. The critical value has subscript $\alpha$ ($z_\alpha$, $t_{\alpha,\nu}$, $\chi^2_{\alpha,\nu}$ or $\chi^2_{1-\alpha,\nu}$, $f_\alpha$ or $f_{1-\alpha}$). Write \emph{only one} limit. The other side stays open ($-\infty$ or $+\infty$, or $0$/$1$ for a variance or proportion that cannot go further).
  \item \textbf{They are not the same confidence level.} A one-sided $95\%$ upper bound is \emph{not} the upper end of a two-sided $95\%$ interval. The upper end of a two-sided $95\%$ interval is a one-sided $97.5\%$ bound. The upper end of a two-sided $90\%$ interval is a one-sided $95\%$ bound.
  \item \textbf{Tests.} A one-sided test at level $\alpha$ matches a one-sided $100(1-\alpha)\%$ bound. A two-sided test at level $\alpha$ matches a two-sided $100(1-\alpha)\%$ interval. Reject $H_0:\theta=\theta_0$ against $H_1:\theta>\theta_0$ exactly when the lower bound exceeds $\theta_0$. Reject against $H_1:\theta<\theta_0$ exactly when the upper bound is below $\theta_0$.
\end{enumerate}
\end{nuance}

\begin{table}
\caption{One-sided and two-sided $100(1-\alpha)\%$ limits. In each row the two-sided interval uses both limits with $\alpha/2$, the upper bound uses only the upper limit with $\alpha$, the lower bound only the lower limit with $\alpha$.}\label{tab:onetwo}
\centering\small
\begin{tabular}{p{2.6cm}p{3.6cm}p{3.9cm}p{3.9cm}}
\toprule
Parameter & Two-sided ($\alpha/2$) & Upper bound ($\le$, uses $\alpha$) & Lower bound ($\ge$, uses $\alpha$) \\
\midrule
$\mu$, $\sigma$ known & $\bar{x}\pm z_{\alpha/2}\dfrac{\sigma}{\sqrt n}$ & $\mu\le\bar{x}+z_\alpha\dfrac{\sigma}{\sqrt n}$ & $\mu\ge\bar{x}-z_\alpha\dfrac{\sigma}{\sqrt n}$ \\
$\mu$, $\sigma$ unknown & $\bar{x}\pm t_{\alpha/2,n-1}\dfrac{s}{\sqrt n}$ & $\mu\le\bar{x}+t_{\alpha,n-1}\dfrac{s}{\sqrt n}$ & $\mu\ge\bar{x}-t_{\alpha,n-1}\dfrac{s}{\sqrt n}$ \\
$\sigma^2$ & $\dfrac{(n-1)s^2}{\chi^2_{\alpha/2,n-1}}\le\sigma^2\le\dfrac{(n-1)s^2}{\chi^2_{1-\alpha/2,n-1}}$ & $\sigma^2\le\dfrac{(n-1)s^2}{\chi^2_{1-\alpha,n-1}}$ & $\sigma^2\ge\dfrac{(n-1)s^2}{\chi^2_{\alpha,n-1}}$ \\
$p$ (large $n$) & $\hat{p}\pm z_{\alpha/2}\sqrt{\dfrac{\hat{p}(1-\hat{p})}{n}}$ & $p\le\hat{p}+z_\alpha\sqrt{\dfrac{\hat{p}(1-\hat{p})}{n}}$ & $p\ge\hat{p}-z_\alpha\sqrt{\dfrac{\hat{p}(1-\hat{p})}{n}}$ \\
\bottomrule
\end{tabular}
\end{table}

Multipliers for a $95\%$ statement:
\begin{itemize}
  \item $z$: $z_{0.025}=1.960$ (two-sided) against $z_{0.05}=1.645$ (one-sided).
  \item $t$ with $9$ degrees of freedom: $t_{0.025,9}=2.262$ against $t_{0.05,9}=1.833$.
\end{itemize}

\emph{Sample size} follows the same rule. A two-sided interval with error at most $E$ needs $n=(z_{\alpha/2}\sigma/E)^2$. A one-sided bound at distance at most $E$ from $\bar{x}$ needs the smaller $n=(z_\alpha\sigma/E)^2$.

\subsection{General Method to Derive a Confidence Interval}

All the intervals in this chapter come from one idea.

\begin{definition}[Pivotal quantity]
A \emph{pivotal quantity} (pivot) is a function $g(X_1,\dots,X_n;\theta)$ of the sample and the unknown parameter whose distribution does not depend on $\theta$ (or on any other unknown parameter).
\end{definition}

Recipe:
\begin{enumerate}
  \item Find a pivot $g$ for $\theta$ with a known distribution.
  \item Find constants $a,b$ with $\Prob(a\le g\le b)=1-\alpha$.
  \item Rearrange $a\le g(X;\theta)\le b$ into $L(X)\le\theta\le U(X)$.
\end{enumerate}

\begin{example}
Pivots used in this chapter:
\begin{itemize}
  \item $\mu$, $\sigma$ known: $g=(\bar{X}-\mu)/(\sigma/\sqrt n)\sim\Normal(0,1)$.
  \item $\mu$, $\sigma$ unknown: $T=(\bar{X}-\mu)/(S/\sqrt n)\sim t_{n-1}$.
  \item $\sigma^2$: $(n-1)S^2/\sigma^2\sim\chi^2_{n-1}$.
\end{itemize}
\end{example}

\subsection{A Large-Sample Confidence Interval for $\mu$}

If the population is not normal (or not assumed so) but $n$ is large, the CLT gives $\bar{X}\approx\Normal(\mu,\sigma^2/n)$. If $\sigma$ is unknown, replacing it by $S$ costs little, because $S$ is close to $\sigma$ for large $n$.

\begin{theorem}[Large-sample interval]\label{thm:largesample}
If $n$ is large (say $n\ge40$), an approximate $100(1-\alpha)\%$ interval for the mean $\mu$ of any population with finite variance is
\[ \bar{x}-z_{\alpha/2}\frac{s}{\sqrt n}\ \le\ \mu\ \le\ \bar{x}+z_{\alpha/2}\frac s{\sqrt n}. \]
\end{theorem}

\begin{example}
$n=50$, $\bar{x}=23.4$, $s=3.1$. A $90\%$ interval: $23.4\pm1.645\cdot3.1/\sqrt{50}=23.4\pm0.72$.
\end{example}

\subsection{Bootstrap Confidence Intervals (CD Only)}

For statistics like the median or a correlation, no pivot with a known distribution exists. The bootstrap of Chapter~1 gives the \emph{percentile interval}.

\begin{definition}[Bootstrap percentile interval]
Generate $B$ bootstrap estimates $\hat\theta^*_1,\dots,\hat\theta^*_B$ by resampling the data with replacement, and sort them. A $100(1-\alpha)\%$ percentile interval is
\[ \big[\ \hat\theta^*_{(\alpha/2)},\ \hat\theta^*_{(1-\alpha/2)}\ \big], \]
the $\alpha/2$ and $1-\alpha/2$ sample quantiles of the bootstrap distribution. With $B=1000$ and $\alpha=0.05$, these are the $25$th and $975$th ordered values.
\end{definition}

A second version, the \emph{bootstrap standard-error interval}, is $\hat\theta\pm z_{\alpha/2}\se_{\mathrm{boot}}$. It is reasonable when the bootstrap distribution is roughly normal. The bootstrap needs no distributional assumption, but it needs a reasonably large and representative sample. See Chapter~13.

\section{Confidence Interval on the Mean of a Normal Distribution, Variance Unknown}

\subsection{The $t$ Distribution}

If $\sigma$ is unknown, replacing it by $S$ changes the sampling distribution of the pivot.

\begin{theorem}[The $t$ statistic]\label{thm:t-pivot}
If $X_1,\dots,X_n$ is a random sample from a \emph{normal} distribution with mean $\mu$, then
\[ T=\frac{\bar{X}-\mu}{S/\sqrt n}\sim t_{n-1}. \]
\end{theorem}

About the $t$ density:
\begin{itemize}
  \item Bell-shaped and symmetric about $0$, like the standard normal, but with heavier tails. The shape depends on the degrees of freedom $k=n-1$.
  \item $\Var[T]=k/(k-2)$ for $k>2$.
  \item As $k\to\infty$, $t_k\to\Normal(0,1)$.
  \item $t_{\alpha,k}$ is the point with right-tail area $\alpha$. For example $t_{0.025,9}=2.262$ and $t_{0.025,\infty}=1.960$.
\end{itemize}

\subsection{Development of the $t$ Distribution (CD Only)}

Write
\[ T=\frac{(\bar{X}-\mu)/(\sigma/\sqrt n)}{\sqrt{S^2/\sigma^2}}=\frac{Z}{\sqrt{V/(n-1)}},\qquad Z\sim\Normal(0,1),\quad V=\frac{(n-1)S^2}{\sigma^2}\sim\chi^2_{n-1}, \]
with $Z$ and $V$ independent. That is the definition of the $t$ distribution with $n-1$ degrees of freedom. Chapter~11 (Theorem~\ref{thm:cochran}) proves the independence of $\bar{X}$ and $S^2$ and the $\chi^2$ law of $V$, and writes down the density of $t_k$.

\subsection{The $t$ Confidence Interval on $\mu$}

Use $\Prob(-t_{\alpha/2,n-1}\le T\le t_{\alpha/2,n-1})=1-\alpha$ and rearrange:

\begin{theorem}[$t$ interval]\label{thm:t-ci}
If $\bar{x}$ and $s$ are the mean and standard deviation of a random sample from a normal distribution with unknown variance, a $100(1-\alpha)\%$ interval on $\mu$ is
\[ \bar{x}-t_{\alpha/2,n-1}\frac s{\sqrt n}\ \le\ \mu\ \le\ \bar{x}+t_{\alpha/2,n-1}\frac s{\sqrt n}. \]
\emph{One-sided $100(1-\alpha)\%$ bounds:} upper $\mu\le\bar{x}+t_{\alpha,n-1}\,s/\sqrt n$; lower $\bar{x}-t_{\alpha,n-1}\,s/\sqrt n\le\mu$. They use $t_{\alpha,n-1}$, not $t_{\alpha/2,n-1}$, and give one limit only.
\end{theorem}

\begin{example}\label{ex:t-ci}
Ten measurements of the load at failure of a component: $\bar{x}=22.4$ kN, $s=3.0$ kN.
\begin{itemize}
  \item $95\%$ interval: $t_{0.025,9}=2.262$, so $22.4\pm2.262\cdot3.0/\sqrt{10}=22.4\pm2.15$, that is $20.25\le\mu\le24.55$.
  \item One-sided $95\%$ bounds use $t_{0.05,9}=1.833$: $\mu\le22.4+1.833(0.9487)=24.14$ and $\mu\ge22.4-1.739=20.66$.
\end{itemize}
Each bound is tighter than the matching end of the two-sided interval.
\end{example}

\begin{example}[Tensile adhesion of U-700 alloy]
Load at failure for $n=22$ specimens: $\bar{x}=13.71$, $s=3.55$. (The box plot and normal plot support normality.)
\begin{itemize}
  \item $21$ degrees of freedom, $t_{0.025,21}=2.080$.
  \item $95\%$ interval: $13.71\pm2.080(3.55)/\sqrt{22}=13.71\pm1.57$, that is $12.14\le\mu\le15.28$.
  \item The interval is fairly wide because the variability is large.
\end{itemize}
You cannot solve directly for the sample size that gives a target length. The length involves $s$ and $t_{\alpha/2,n-1}$, which depend on the data. Choose $n$ by trial and error from a prior guess of $\sigma$.
\end{example}

\begin{pitfall}{using $z$ with a small sample}
With $n=10$, the normal multiplier $1.96$ gives $22.4\pm1.86$. That is too short and covers $\mu$ less than $95\%$ of the time. The $t$ multiplier corrects for the extra uncertainty from estimating $\sigma$. The correction is large for $n<30$ and negligible for $n>100$.
\end{pitfall}

The $t$ interval assumes a normal population. It is fairly robust to moderate non-normality for $n\ge15$ if the data are roughly symmetric. It is unreliable for small samples from skewed or heavy-tailed populations. Check with a normal probability plot.

\section{Confidence Interval on the Variance and Standard Deviation of a Normal Distribution}

Let $X_1,\dots,X_n$ be a random sample from $\Normal(\mu,\sigma^2)$ with sample variance $S^2$. The pivot is
\[ \chi^2=\frac{(n-1)S^2}{\sigma^2}\sim\chi^2_{n-1}. \]
Let $\chi^2_{\alpha,k}$ be the point with right-tail area $\alpha$. Then $\Prob\big(\chi^2_{1-\alpha/2,n-1}\le(n-1)S^2/\sigma^2\le\chi^2_{\alpha/2,n-1}\big)=1-\alpha$.

\begin{theorem}[Interval for $\sigma^2$ and $\sigma$]\label{thm:var-ci}
A $100(1-\alpha)\%$ interval for $\sigma^2$ is
\[ \frac{(n-1)s^2}{\chi^2_{\alpha/2,n-1}}\ \le\ \sigma^2\ \le\ \frac{(n-1)s^2}{\chi^2_{1-\alpha/2,n-1}}. \]
Take square roots for an interval on $\sigma$.

One-sided $100(1-\alpha)\%$ bounds use $\alpha$ (not $\alpha/2$) and one limit: upper $\sigma^2\le(n-1)s^2/\chi^2_{1-\alpha,n-1}$; lower $\sigma^2\ge(n-1)s^2/\chi^2_{\alpha,n-1}$. The upper bound divides by the \emph{small} chi-square point. The lower bound divides by the \emph{large} one.
\end{theorem}

\begin{example}
$n=20$ fill volumes, $s^2=0.0153$.
\begin{itemize}
  \item $95\%$ two-sided: $\chi^2_{0.025,19}=32.85$ and $\chi^2_{0.975,19}=8.91$. So $19(0.0153)/32.85=0.00885\le\sigma^2\le19(0.0153)/8.91=0.0326$, and $0.094\le\sigma\le0.181$.
  \item One-sided $95\%$: $\chi^2_{0.95,19}=10.12$ and $\chi^2_{0.05,19}=30.14$. So $\sigma^2\le19(0.0153)/10.12=0.0287$ and $\sigma^2\ge19(0.0153)/30.14=0.00965$.
\end{itemize}
\end{example}

\begin{nuance}{asymmetry and non-normality}
The $\chi^2$ distribution is skewed, so the interval is not symmetric about $s^2$. The point estimate lies closer to the lower limit.

Unlike the $t$ interval for the mean, this interval is \emph{very sensitive} to non-normality, because the distribution of $S^2$ depends on the fourth moment. Use it only when normality is credible.
\end{nuance}

\section{A Large-Sample Confidence Interval for a Population Proportion}

Let $X$ be the number of successes in $n$ independent trials with success probability $p$, and $\hat{P}=X/n$. For large $n$, the normal approximation to the binomial gives $\hat{P}\approx\Normal\big(p,\ p(1-p)/n\big)$. Replace $p$ in the variance by $\hat{p}$:

\begin{theorem}[Interval for a proportion]\label{thm:prop-ci}
If $\hat{p}$ is the proportion of successes in a random sample of size $n$, an approximate $100(1-\alpha)\%$ interval for $p$ is
\[ \hat{p}-z_{\alpha/2}\sqrt{\frac{\hat{p}(1-\hat{p})}n}\ \le\ p\ \le\ \hat{p}+z_{\alpha/2}\sqrt{\frac{\hat{p}(1-\hat{p})}n}. \]
The approximation is acceptable when $n\hat{p}$ and $n(1-\hat{p})$ are both at least about $5$ to $10$.

One-sided $100(1-\alpha)\%$ bounds use $z_\alpha$ and one limit: $p\le\hat{p}+z_\alpha\sqrt{\hat{p}(1-\hat{p})/n}$ (upper), $\hat{p}-z_\alpha\sqrt{\hat{p}(1-\hat{p})/n}\le p$ (lower).
\end{theorem}

\begin{example}
Among $n=85$ automobile crankshaft bearings, $10$ have rough surfaces. $\hat{p}=10/85=0.1176$.
\begin{itemize}
  \item $95\%$ interval: $0.1176\pm1.96\sqrt{0.1176\cdot0.8824/85}=0.1176\pm0.0685$, that is $0.049\le p\le0.186$.
  \item One-sided $95\%$ upper bound: $0.1176+1.645(0.0349)=0.175$. Lower bound: $0.060$. These use $1.645$, not $1.96$.
\end{itemize}
\end{example}

\begin{theorem}[Sample size for a proportion]\label{thm:samplesize-prop}
To be $100(1-\alpha)\%$ confident that the error $\abs{\hat{p}-p}$ is at most $E$, use
\[ n=\left(\frac{z_{\alpha/2}}E\right)^2p(1-p). \]
Since $p$ is unknown, use a preliminary estimate, or the worst case $p=0.5$ (which maximises $p(1-p)=0.25$). That gives $n=0.25\,(z_{\alpha/2}/E)^2$.
\end{theorem}

\begin{example}
A poll with a $\pm3\%$ margin at $95\%$ confidence and no prior information needs $n=0.25(1.96/0.03)^2=1067.1$, so $n=1068$.
\end{example}

\begin{interview}{better intervals for proportions}
The interval of Theorem~\ref{thm:prop-ci} (the Wald interval) can have poor coverage when $p$ is near $0$ or $1$, or $n$ is small. Two better options:
\begin{itemize}
  \item The \emph{Wilson score interval}: solve $\abs{\hat{p}-p}\le z_{\alpha/2}\sqrt{p(1-p)/n}$ for $p$ exactly, instead of plugging in $\hat{p}$.
  \item The ``add two successes and two failures'' adjustment: use $\tilde{p}=(x+2)/(n+4)$ in the formula above.
\end{itemize}
Both have much better coverage. For $x=0$ successes, the $95\%$ upper bound is about $3/n$ (the ``rule of three'').
\end{interview}

\section{A Prediction Interval for a Future Observation}

A confidence interval is about a fixed parameter. Often we want a range for the next individual value $X_{n+1}$, drawn from the same population.

\begin{theorem}[Prediction interval]\label{thm:pi}
Let $X_1,\dots,X_n$ be a random sample from $\Normal(\mu,\sigma^2)$ and $X_{n+1}$ an independent future observation. A $100(1-\alpha)\%$ prediction interval for $X_{n+1}$ is
\[ \bar{x}-t_{\alpha/2,n-1}\,s\sqrt{1+\frac1n}\ \le\ X_{n+1}\ \le\ \bar{x}+t_{\alpha/2,n-1}\,s\sqrt{1+\frac1n}. \]
\end{theorem}

\begin{proof}
$X_{n+1}-\bar{X}$ is normal with mean $0$ and variance $\sigma^2+\sigma^2/n=\sigma^2(1+1/n)$, because the new observation is independent of $\bar{X}$. So
\[ \frac{X_{n+1}-\bar{X}}{S\sqrt{1+1/n}}=\frac{(X_{n+1}-\bar{X})/(\sigma\sqrt{1+1/n})}{\sqrt{S^2/\sigma^2}}\sim t_{n-1}, \]
and the interval follows as in Theorem~\ref{thm:t-ci}.
\end{proof}

\begin{example}
Continue Example~\ref{ex:t-ci} ($n=10$, $\bar{x}=22.4$, $s=3.0$). The $95\%$ prediction interval for the next component is $22.4\pm2.262\cdot3.0\sqrt{1.1}=22.4\pm7.12$, that is $15.3$ to $29.5$. The confidence interval for $\mu$, $20.25$ to $24.55$, is much narrower.
\end{example}

As $n\to\infty$, the prediction interval tends to $\mu\pm z_{\alpha/2}\sigma$. It does not shrink to zero: no amount of data removes the variability of a single future value. The confidence interval shrinks to a point.

\begin{example}[Prediction interval, same alloy data]
For a twenty-third specimen, the $95\%$ prediction interval is $13.71\pm2.080(3.55)\sqrt{1+1/22}=13.71\pm7.55$, that is $6.16$ to $21.26$. This is much longer than the interval $(12.14,\,15.28)$ for the mean. There is always uncertainty about the next observation, even when $\mu$ is known.
\end{example}

\section{Tolerance Intervals for a Normal Distribution}

A \emph{tolerance interval} aims to capture a stated proportion of the whole population, with a stated confidence that it really does.

\begin{definition}[Tolerance interval]
A \emph{tolerance interval for capturing at least $100\gamma\%$ of the values of a normal distribution with confidence level $100(1-\alpha)\%$} is
\[ \bar{x}-k\,s\ \le\ x\ \le\ \bar{x}+k\,s. \]
The constant $k=k(n,\gamma,\alpha)$ comes from tables or software. It is chosen so that, over repeated samples, a proportion $1-\alpha$ of such intervals contain at least a proportion $\gamma$ of the population.
\end{definition}

A good approximation for the two-sided factor (Howe's formula) is
\[ k\approx z_{(1+\gamma)/2}\sqrt{\frac{(n-1)\,(1+1/n)}{\chi^2_{1-\alpha,n-1}}}, \]
where $\chi^2_{1-\alpha,n-1}$ is the point with \emph{left} tail area $\alpha$ (right-tail area $1-\alpha$).

\begin{example}
$n=10$, $\gamma=0.95$, $1-\alpha=0.95$.
\begin{itemize}
  \item $z_{0.975}=1.96$ and $\chi^2_{0.95,9}=3.325$.
  \item $k\approx1.96\sqrt{9(1.1)/3.325}=1.96\sqrt{2.977}=3.38$. The exact tabled value is $3.38$.
  \item With $\bar{x}=22.4$, $s=3.0$: $22.4\pm3.38\cdot3.0$, that is $12.3$ to $32.5$. Wider than the prediction interval $15.3$ to $29.5$.
\end{itemize}
\end{example}

\begin{example}[Tolerance interval, same alloy data]
To include $90\%$ of the population with $95\%$ confidence for $n=22$, Table XI gives $k=2.264$. The interval is $\bar{x}\pm2.264s=\bar{x}\pm8.04$. With $\bar{x}=13.71$ this is $(5.67,\,21.75)$. The width, $16.08$, is what matters.
\begin{itemize}
  \item As $n\to\infty$, $k$ approaches $z_{0.05}=1.645$ for a $90\%$ two-sided tolerance interval. A $100(1-\alpha)\%$ prediction interval tends to a tolerance interval containing $100(1-\alpha)\%$ of the distribution.
  \item If $\mu=600$ and $\sigma=30$ are known, $600\pm1.96(30)=(541.2,\,658.8)$ captures $95\%$ of the population.
  \item With $\bar{x}$ and $s$ in place of $\mu,\sigma$, the plain interval $\bar{x}\pm1.96s$ is likely to contain less than $95\%$. That is why $1.96$ is replaced by the larger factor $k$.
\end{itemize}
\end{example}

\begin{interview}{three intervals, three questions}
\begin{itemize}
  \item \emph{Confidence interval:} where is the population mean? Narrowest. Shrinks with $n$.
  \item \emph{Prediction interval:} where will the next single observation fall? Limit $\mu\pm z\sigma$.
  \item \emph{Tolerance interval:} what range covers $95\%$ of all units, with $95\%$ confidence that the claim is right? Widest. Limit $\mu\pm z_{(1+\gamma)/2}\sigma$ from above.
\end{itemize}
Quality engineers use tolerance intervals for specification limits and confidence intervals for process averages.
\end{interview}

\begin{summarybox}
\begin{itemize}
  \item $\sigma$ known: $\bar{x}\pm z_{\alpha/2}\sigma/\sqrt n$. Sample size $n=(z_{\alpha/2}\sigma/E)^2$. One-sided $100(1-\alpha)\%$ bounds: $\bar{x}+z_\alpha\sigma/\sqrt n$ (upper) or $\bar{x}-z_\alpha\sigma/\sqrt n$ (lower). One limit only, and $\alpha$ not $\alpha/2$. The end of a two-sided $95\%$ interval is a one-sided $97.5\%$ bound.
  \item $\sigma$ unknown (normal data): $\bar{x}\pm t_{\alpha/2,n-1}s/\sqrt n$. Large samples (any population): $\bar{x}\pm z_{\alpha/2}s/\sqrt n$.
  \item Variance: $\dfrac{(n-1)s^2}{\chi^2_{\alpha/2,n-1}}\le\sigma^2\le\dfrac{(n-1)s^2}{\chi^2_{1-\alpha/2,n-1}}$. Sensitive to non-normality.
  \item Proportion: $\hat{p}\pm z_{\alpha/2}\sqrt{\hat{p}(1-\hat{p})/n}$. Sample size $n=(z_{\alpha/2}/E)^2p(1-p)\le0.25(z_{\alpha/2}/E)^2$.
  \item Prediction interval: $\bar{x}\pm t_{\alpha/2,n-1}s\sqrt{1+1/n}$. Tolerance interval: $\bar{x}\pm ks$.
  \item The confidence level is the long-run coverage of the procedure. It is not the probability that one computed interval contains $\mu$.
\end{itemize}
\end{summarybox}

\section{Exercises}

\begin{exercise}
The tensile strength of a fiber is normal with $\sigma=2$ MPa. A sample of $9$ has mean $98.5$. Find $90\%$, $95\%$ and $99\%$ confidence intervals for $\mu$ and the sample size needed to estimate $\mu$ within $0.5$ MPa with $95\%$ confidence.
\end{exercise}

\begin{exercise}
A $95\%$ confidence interval for a mean based on $n=25$ and known $\sigma$ has length $4$. What would its length be for $n=100$? For $99\%$ confidence with $n=25$?
\end{exercise}

\begin{exercise}
Data on the time (minutes) for a chemical reaction: $\bar{x}=45.2$, $s=4.1$, $n=12$ from a normal population. Find the $95\%$ confidence interval for $\mu$, the $95\%$ upper confidence bound for $\mu$, and the $95\%$ confidence interval for $\sigma$.
\end{exercise}

\begin{exercise}
In a survey, $x=216$ of $n=400$ respondents prefer product A. Find the $95\%$ confidence interval for the proportion in the population. How large a sample would give a margin of error of $2\%$ with $99\%$ confidence?
\end{exercise}

\begin{exercise}
Explain, with a short numerical example, why a $95\%$ prediction interval is wider than a $95\%$ confidence interval from the same data, and compute both for $n=15$, $\bar{x}=50$, $s=6$.
\end{exercise}

\begin{exercise}
Use Howe's formula to compute the tolerance factor $k$ for $n=20$, $\gamma=0.99$, confidence $0.95$ (use $z_{0.995}=2.576$ and $\chi^2_{0.95,19}=10.12$).
\end{exercise}

\begin{exercise}
Show that the width of the large-sample confidence interval for a proportion is largest at $p=0.5$, and derive the sample size formula of Theorem~\ref{thm:samplesize-prop}.
\end{exercise}

\begin{exercise}
Draw $B=1000$ bootstrap samples from the data $4.1,\ 5.3,\ 3.9,\ 6.2,\ 4.8,\ 5.5,\ 4.4$ and form a bootstrap percentile $95\%$ interval for the median. Compare it with the $t$ interval for the mean.
\end{exercise}


\part{Tests of Hypotheses}
\chapter{Tests of Hypotheses for a Single Sample}

\section{Hypothesis Testing}

\subsection{Statistical Hypotheses}

\begin{definition}[Statistical hypothesis]
A \emph{statistical hypothesis} is a statement about the parameters of one or more populations (or about the form of the distribution). A \emph{hypothesis test} is a rule, based on a random sample, for deciding between two hypotheses:
\begin{itemize}
  \item the \emph{null hypothesis} $H_0$, the claim assumed true unless the data contradict it, written as an equality such as $H_0:\mu=50$;
  \item the \emph{alternative hypothesis} $H_1$, the claim accepted if $H_0$ is rejected, such as $H_1:\mu\ne50$.
\end{itemize}
\end{definition}

The two hypotheses are not treated symmetrically. We begin by assuming $H_0$ is true, and reject it only if the sample would be very unlikely under $H_0$. This is a proof by contradiction made probabilistic. For that reason the claim we hope to \emph{establish} is placed in $H_1$.

\begin{example}
A propellant must have a mean burning rate of $50$ cm/s. To test whether a new batch meets this we take $H_0:\mu=50$ against $H_1:\mu\ne50$. A manufacturer who wants evidence that a new process has \emph{reduced} the defect rate below $5\%$ would test $H_0:p=0.05$ against $H_1:p<0.05$.
\end{example}

\subsection{Tests of Statistical Hypotheses}

A test uses a \emph{test statistic} (a function of the sample) and a \emph{rejection region} (critical region): the set of values of the statistic for which $H_0$ is rejected. The boundary values are the \emph{critical values}.

\begin{definition}[Type I and Type II errors]
\begin{center}
\begin{tabular}{lcc}
\toprule
 & $H_0$ true & $H_0$ false \\
\midrule
Fail to reject $H_0$ & correct & \textbf{Type II error} \\
Reject $H_0$ & \textbf{Type I error} & correct \\
\bottomrule
\end{tabular}
\end{center}
$\alpha=\Prob(\text{Type I error})=\Prob(\text{reject }H_0\mid H_0\text{ true})$ is the \emph{significance level} (or size) of the test. $\beta=\Prob(\text{Type II error})=\Prob(\text{fail to reject }H_0\mid H_0\text{ false})$. The \emph{power} of the test is $1-\beta=\Prob(\text{reject }H_0\mid H_0\text{ false})$. Note that $\beta$ depends on which specific alternative value of the parameter is true.
\end{definition}

\begin{intuition}
$\alpha$ is the false-alarm rate and $1-\beta$ is the detection rate. For a fixed sample size the two errors trade off: shrinking the rejection region lowers $\alpha$ but raises $\beta$. Only a larger sample can reduce both. Because $\alpha$ is chosen directly while $\beta$ is not, rejecting $H_0$ is a strong conclusion (error rate controlled at $\alpha$), while failing to reject is a weak one (it may be due to low power). That is why one says ``fail to reject $H_0$'' rather than ``accept $H_0$''.
\end{intuition}

\begin{figure}[h]
\centering
\begin{tikzpicture}[scale=1.0]
  \draw[axis] (-3.4,0)--(7.6,0); \draw[axis] (0,-0.2)--(0,2.2);
  \draw[thick,accent2,domain=-3:3,samples=60] plot (\x,{2*exp(-\x*\x/2)});
  \draw[thick,accent,domain=-0.4:6.8,samples=70] plot (\x,{2*exp(-(\x-3.2)*(\x-3.2)/2)});
  \fill[accent2!25,domain=1.64:3,samples=40] (1.64,0)--plot (\x,{2*exp(-\x*\x/2)})--(3,0)--cycle;
  \fill[accent!22,domain=-0.4:1.64,samples=40] (-0.4,0)--plot (\x,{2*exp(-(\x-3.2)*(\x-3.2)/2)})--(1.64,0)--cycle;
  \draw[dashed,mid] (1.64,0)--(1.64,2.1); \node[lab] at (1.64,2.35){critical value};
  \node[lab,accent2] at (-1.6,1.6){$H_0$}; \node[lab,accent] at (5.3,1.6){$H_1$};
  \node[lab] at (2.5,0.25){$\alpha$};
\end{tikzpicture}
\caption{Distribution of the test statistic under $H_0$ and under one alternative. The tail of the $H_0$ curve beyond the critical value is $\alpha$; the part of the $H_1$ curve to the left of it is $\beta$.}
\end{figure}

\subsection{One-Sided and Two-Sided Hypotheses}

\begin{definition}[Forms of the alternative]
\begin{itemize}
  \item \emph{Two-sided}: $H_0:\mu=\mu_0$, $H_1:\mu\ne\mu_0$; the rejection region is in both tails.
  \item \emph{One-sided} (upper): $H_0:\mu=\mu_0$, $H_1:\mu>\mu_0$; rejection region in the upper tail.
  \item \emph{One-sided} (lower): $H_0:\mu=\mu_0$, $H_1:\mu<\mu_0$; rejection region in the lower tail.
\end{itemize}
\end{definition}

\begin{example}[Which hypothesis goes in $H_1$?]
A bottler requires a mean bursting strength of at least $200$ psi. Formulation (a): $H_0:\mu=200$ against $H_1:\mu>200$. Rejecting $H_0$ is a strong conclusion, so the manufacturer must \emph{demonstrate} that the bottles exceed the specification; bottles are judged unsatisfactory unless the data prove otherwise, and there is some probability that good bottles (true mean slightly above $200$) are not accepted. Formulation (b): $H_0:\mu=200$ against $H_1:\mu<200$. Now the bottles are accepted unless there is strong evidence that they fall short, and some probability remains that bottles slightly below $200$ are judged satisfactory. Which is correct depends on the situation: (a) suits a supplier with a history of problems or safety-critical use; (b) suits a supplier with a good record where small deviations are harmless. In either case \emph{put the statement about which a strong conclusion is wanted in the alternative hypothesis}.
\end{example}

The alternative must be chosen from the scientific question \emph{before} seeing the data. Choosing a one-sided alternative after looking at which way the data lean doubles the true $\alpha$.

\subsection{General Procedure for Hypothesis Testing}

\begin{definition}[The eight-step procedure]
\begin{enumerate}[label=\arabic*.]
  \item From the problem context, identify the parameter of interest.
  \item State the null hypothesis $H_0$.
  \item Specify an appropriate alternative hypothesis $H_1$.
  \item Choose the significance level $\alpha$.
  \item Determine an appropriate test statistic.
  \item State the rejection region for the statistic.
  \item Compute the sample quantities, substitute into the test statistic and find its value (and the $P$-value).
  \item Decide whether $H_0$ should be rejected and report that in the problem context.
\end{enumerate}
\end{definition}

\section{Tests on the Mean of a Normal Distribution, Variance Known}

\subsection{Hypothesis Tests on the Mean}

Let $X_1,\dots,X_n$ be a random sample from $\Normal(\mu,\sigma^2)$ with $\sigma^2$ known. Test $H_0:\mu=\mu_0$. If $H_0$ is true, then $\bar{X}\sim\Normal(\mu_0,\sigma^2/n)$, so
\[ Z_0=\frac{\bar{X}-\mu_0}{\sigma/\sqrt n}\sim\Normal(0,1). \]

\begin{theorem}[$z$-test]\label{thm:ztest}
For a given $\alpha$, reject $H_0:\mu=\mu_0$ in favour of
\begin{itemize}
  \item $H_1:\mu\ne\mu_0$ if $\abs{z_0}>z_{\alpha/2}$;
  \item $H_1:\mu>\mu_0$ if $z_0>z_\alpha$;
  \item $H_1:\mu<\mu_0$ if $z_0<-z_\alpha$.
\end{itemize}
Each test has Type I error exactly $\alpha$.
\end{theorem}

\begin{example}\label{ex:ztest}
The burning rate of a propellant has $\sigma=2$ cm/s. For a sample of $n=25$, $\bar{x}=51.3$. Test $H_0:\mu=50$ against $H_1:\mu\ne50$ at $\alpha=0.05$. The statistic is $z_0=(51.3-50)/(2/\sqrt{25})=3.25$. Since $\abs{3.25}>z_{0.025}=1.96$, reject $H_0$: there is evidence that the mean burning rate differs from $50$ cm/s.
\end{example}

\subsection{$P$-Values in Hypothesis Tests}

Reporting only ``reject'' or ``fail to reject'' at a fixed $\alpha$ hides how strong the evidence is.

\begin{definition}[$P$-value]
The \emph{$P$-value} is the smallest level of significance at which the observed value of the test statistic would lead to rejection of $H_0$; equivalently, the probability, computed assuming $H_0$ is true, of obtaining a test statistic value at least as extreme (in the direction of $H_1$) as the one observed.
\end{definition}

For the $z$-test with observed $z_0$: $P=2\big[1-\Phi(\abs{z_0})\big]$ for the two-sided alternative; $P=1-\Phi(z_0)$ for $H_1:\mu>\mu_0$; $P=\Phi(z_0)$ for $H_1:\mu<\mu_0$. The rule: \emph{reject $H_0$ if $P<\alpha$.}

\begin{example}
In Example~\ref{ex:ztest}, $P=2[1-\Phi(3.25)]=2(0.00058)=0.0012$. The data would be very unusual if $\mu$ were $50$; $H_0$ would be rejected at every conventional level.
\end{example}

\begin{pitfall}{what a $P$-value is not}
The $P$-value is \emph{not} the probability that $H_0$ is true, and $1-P$ is not the probability that $H_1$ is true. It is the probability of data this extreme \emph{if $H_0$ were true}. A small $P$ also does not measure the size of the effect: with a huge sample a trivial difference gives a tiny $P$ (see Section~\ref{sec:practical}).
\end{pitfall}

\subsection{Connection Between Hypothesis Tests and Confidence Intervals}

\begin{theorem}[Duality]\label{thm:duality}
The two-sided level-$\alpha$ test of $H_0:\mu=\mu_0$ rejects if and only if $\mu_0$ lies outside the $100(1-\alpha)\%$ confidence interval for $\mu$. The one-sided tests correspond to one-sided bounds, at the \emph{same} $\alpha$:
\begin{itemize}
  \item $H_1:\mu>\mu_0$ is rejected at level $\alpha$ iff $\mu_0<\bar{x}-z_\alpha\sigma/\sqrt n$, i.e.\ iff $\mu_0$ lies below the $100(1-\alpha)\%$ \emph{lower} confidence bound.
  \item $H_1:\mu<\mu_0$ is rejected at level $\alpha$ iff $\mu_0>\bar{x}+z_\alpha\sigma/\sqrt n$, i.e.\ iff $\mu_0$ lies above the $100(1-\alpha)\%$ \emph{upper} confidence bound.
\end{itemize}
\end{theorem}

\begin{proof}
$\abs{\bar{x}-\mu_0}>z_{\alpha/2}\sigma/\sqrt n$ is exactly the statement that $\mu_0\notin\big[\bar{x}-z_{\alpha/2}\sigma/\sqrt n,\ \bar{x}+z_{\alpha/2}\sigma/\sqrt n\big]$.
\end{proof}

So a confidence interval answers every possible two-sided hypothesis test at once: it is the set of all values $\mu_0$ that would not be rejected.

\begin{pitfall}{do not mix one-sided tests with two-sided intervals}
A one-sided test at $\alpha=0.05$ uses $z_{0.05}=1.645$ and matches a one-sided $95\%$ bound. A two-sided $95\%$ interval uses $z_{0.025}=1.96$ and matches a two-sided test at $\alpha=0.05$. If a one-sided test at $\alpha=0.05$ rejects but $\mu_0$ is still inside the two-sided $95\%$ interval, there is no contradiction: the interval corresponds to the two-sided test, which has a stricter cutoff. To compare a one-sided test with an interval, use the two-sided $90\%$ interval (whose end is a one-sided $95\%$ bound) or the one-sided $95\%$ bound. The same holds for $t$, $\chi^2$ and $F$ procedures: one-sided uses $\alpha$ in a single tail, two-sided uses $\alpha/2$ in each.
\end{pitfall}

\subsection{Type II Error and Choice of Sample Size}

Suppose the true mean is $\mu=\mu_0+\delta$ with $\delta\ne0$. Then $\bar{X}\sim\Normal(\mu_0+\delta,\sigma^2/n)$ and $Z_0=\dfrac{\bar{X}-\mu_0}{\sigma/\sqrt n}\sim\Normal\Big(\dfrac{\delta\sqrt n}\sigma,\ 1\Big)$.

\begin{theorem}[Type II error of the $z$-test]\label{thm:beta-z}
For the two-sided test,
\[ \beta=\Phi\Big(z_{\alpha/2}-\frac{\delta\sqrt n}\sigma\Big)-\Phi\Big(-z_{\alpha/2}-\frac{\delta\sqrt n}\sigma\Big). \]
For the one-sided alternative $H_1:\mu>\mu_0$ (with $\delta>0$), $\beta=\Phi\big(z_\alpha-\delta\sqrt n/\sigma\big)$; for $H_1:\mu<\mu_0$ ($\delta<0$), $\beta=1-\Phi\big(-z_\alpha-\delta\sqrt n/\sigma\big)$.
\end{theorem}

\begin{proof}
$\beta=\Prob(-z_{\alpha/2}\le Z_0\le z_{\alpha/2})$ with $Z_0\sim\Normal(\delta\sqrt n/\sigma,1)$; subtract the mean and use $\Phi$.
\end{proof}

\begin{example}\label{ex:beta}
In Example~\ref{ex:ztest} suppose the true mean is $51$, so $\delta=1$ and $\delta\sqrt n/\sigma=1\cdot5/2=2.5$. With $\alpha=0.05$, $\beta=\Phi(1.96-2.5)-\Phi(-1.96-2.5)=\Phi(-0.54)-\Phi(-4.46)=0.2946-0.0000=0.295$. The test misses a shift of $1$ cm/s about $29.5\%$ of the time; its power there is about $70.5\%$.
\end{example}

The properties of $\beta$: it decreases as $\abs{\delta}$ increases (large departures are easier to detect), as $n$ increases and as $\alpha$ increases, and it increases with $\sigma$. A plot of $\beta$ (or the power) against $\delta$ is the \emph{operating characteristic (OC) curve} (or power curve); textbooks tabulate OC curves in terms of the standardised difference $d=\abs{\delta}/\sigma$.

\begin{theorem}[Sample size for the $z$-test]\label{thm:samplesize-test}
To achieve Type II error $\beta$ at a specified true difference $\delta$ with a level-$\alpha$ test, use approximately
\[ n\approx\frac{(z_{\alpha/2}+z_\beta)^2\sigma^2}{\delta^2}\quad\text{(two-sided)},\qquad n=\frac{(z_\alpha+z_\beta)^2\sigma^2}{\delta^2}\quad\text{(one-sided)}. \]
(The two-sided formula neglects the term $\Phi(-z_{\alpha/2}-\delta\sqrt n/\sigma)$, which is negligible when $\beta$ is small.)
\end{theorem}

\begin{example}
For Example~\ref{ex:beta}, how large must $n$ be to detect $\delta=1$ with probability $0.90$ ($\beta=0.10$) at $\alpha=0.05$? $n=(z_{0.025}+z_{0.10})^2\sigma^2/\delta^2=(1.96+1.28)^2(2)^2/1^2=42$. (The textbook's operating characteristic charts with $d=\delta/\sigma=0.5$ give a very similar value, about $40$.)
\end{example}

\begin{example}[Textbook examples: propellant burning rate, continued]
\emph{Strong versus weak conclusions.} If we wish to show with a \emph{strong} conclusion that the burning rate is below $50$ cm/s, we take $H_0:\mu=50$, $H_1:\mu<50$, since rejecting $H_0$ is always the strong conclusion. Failing to reject does not mean $\mu=50$ exactly, only that there is no strong evidence for $H_1$ ($H_0$ is understood to include every value not covered by $H_1$).
\emph{Test.} $\sigma=2$, $n=25$, $\bar{x}=51.3$: $z_0=(51.3-50)/(2/\sqrt{25})=3.25>1.96$, so $H_0:\mu=50$ is rejected at $\alpha=0.05$ (two-sided $P=2[1-\Phi(3.25)]\approx0.0012$), with strong evidence that the mean exceeds $50$.
\emph{Power.} If the true mean is $51$ ($\delta=1$, $d=\delta/\sigma=0.5$), with $n=25$: $\beta=\Phi(1.96-2.5)-\Phi(-1.96-2.5)\approx0.30$, i.e.\ about a $30\%$ chance of missing a shift of $1$ cm/s. To detect a shift of $1$ cm/s with probability $0.90$ ($\beta=0.10$) we need $n=(z_{\alpha/2}+z_\beta)^2\sigma^2/\delta^2=(1.96+1.28)^2(4)/1=42$; the OC chart gives about $40$, in close agreement. OC curves exist for the two-sided case (charts VIa, VIb) and the one-sided cases (VIc, VId); for given $n$ and $\alpha$, larger departures are detected more easily, and for a given departure a larger $n$ gives a smaller $\beta$.
\end{example}

\subsection{Large Sample Test}

If $n$ is large (say $n\ge40$), the central limit theorem justifies the $z$-test even without normality, and replacing $\sigma$ by $s$ is harmless:
\[ Z_0=\frac{\bar{X}-\mu_0}{S/\sqrt n}\ \text{ is approximately }\Normal(0,1)\text{ when }H_0\text{ is true}. \]
All the decision rules of Theorem~\ref{thm:ztest} then apply.

\subsection{Some Practical Comments on Hypothesis Tests}\label{sec:practical}

\begin{itemize}
  \item \textbf{Statistical versus practical significance.} A test can detect a difference too small to matter if $n$ is large, and can miss an important difference if $n$ is small. Always report the estimated size of the effect with a confidence interval, not just $P$.
  \item \textbf{$\alpha$ is chosen by the analyst}, conventionally $0.05$ or $0.01$; the $P$-value lets the reader judge for themselves.
  \item \textbf{Assumptions matter.} Independence and (for small $n$) normality should be checked.
  \item \textbf{Do not choose the hypotheses after looking at the data.} Testing many hypotheses and reporting the significant ones inflates the false-positive rate (Chapter~13).
\end{itemize}

\section{Tests on the Mean of a Normal Distribution, Variance Unknown}

\subsection{Hypothesis Tests on the Mean}

When $\sigma$ is unknown and the population is normal, replace $\sigma$ by $S$; the statistic then follows a $t$ distribution.

\begin{theorem}[One-sample $t$-test]\label{thm:ttest}
To test $H_0:\mu=\mu_0$ use $T_0=\dfrac{\bar{X}-\mu_0}{S/\sqrt n}$, which has a $t_{n-1}$ distribution when $H_0$ is true. Reject $H_0$ in favour of
\begin{itemize}
  \item $H_1:\mu\ne\mu_0$ if $\abs{t_0}>t_{\alpha/2,n-1}$;
  \item $H_1:\mu>\mu_0$ if $t_0>t_{\alpha,n-1}$;
  \item $H_1:\mu<\mu_0$ if $t_0<-t_{\alpha,n-1}$.
\end{itemize}
\end{theorem}

\begin{example}\label{ex:ttest}
A drilled shaft is specified to have a mean diameter $2.5$ mm. A sample of $n=16$ has $\bar{x}=2.58$ and $s=0.12$. Test $H_0:\mu=2.5$ against $H_1:\mu\ne2.5$ at $\alpha=0.05$. $t_0=(2.58-2.5)/(0.12/4)=2.667$. With $15$ degrees of freedom $t_{0.025,15}=2.131$. Since $2.667>2.131$, reject $H_0$.
\end{example}

\subsection{$P$-Value for a $t$-Test}

The $P$-value is a tail area under $t_{n-1}$. Tables give only a few critical points, so the $P$-value is usually bracketed. In Example~\ref{ex:ttest}, $t_{0.01,15}=2.602$ and $t_{0.005,15}=2.947$, so the one-sided tail area is between $0.005$ and $0.01$ and the two-sided $0.01<P<0.02$ (software gives $P=0.0176$).

\subsection{Choice of Sample Size}

If the true mean is $\mu_0+\delta$, the statistic $T_0$ has a \emph{noncentral} $t$ distribution with noncentrality $\delta\sqrt n/\sigma$, and $\beta$ must be computed from it (software, or the OC curves in terms of $d=\abs{\delta}/\sigma$). A good approximation for planning is the normal-theory formula of Theorem~\ref{thm:samplesize-test} plus a small correction:
\[ n\approx\frac{(z_{\alpha/2}+z_\beta)^2\sigma^2}{\delta^2}+\frac{z_{\alpha/2}^2}2. \]
Since $\sigma$ is unknown, a rough estimate from a pilot study or past data is needed; the answer is only as good as that estimate.

\begin{example}[Textbook example: golf club coefficient of restitution]
$n=15$ drivers, $\bar{x}=0.83725$, $s=0.02456$. To show the mean exceeds $0.82$ at $\alpha=0.05$: $H_0:\mu=0.82$, $H_1:\mu>0.82$. $t_0=(0.83725-0.82)/(0.02456/\sqrt{15})=2.72>t_{0.05,14}=1.761$: reject $H_0$ ($P\approx0.008$, just under $0.01$). The normal probability plot supports normality, and since the aim is to show the mean \emph{exceeds} $0.82$, a one-sided alternative is appropriate. \emph{Sample size:} if the mean exceeds $0.82$ by $0.02$, then $d=0.02/0.02456=0.81$; the OC chart with $n=15$ gives $\beta\approx0.10$, so $n=15$ is adequate for power at least $0.8$.
\end{example}

\subsection{Likelihood Ratio Approach to Development of Test Procedures (CD Only)}

The tests above are all instances of one general construction.

\begin{definition}[Likelihood ratio test]
Let $L(\theta)$ be the likelihood, $\Theta_0$ the set of parameter values allowed by $H_0$ and $\Theta$ the full parameter space. The \emph{likelihood ratio} is
\[ \Lambda=\frac{\max_{\theta\in\Theta_0}L(\theta)}{\max_{\theta\in\Theta}L(\theta)}\in[0,1], \]
and the test rejects $H_0$ for small values of $\Lambda$ (data much more probable under the best unrestricted parameter than under the best parameter allowed by $H_0$).
\end{definition}

\begin{example}[The $t$-test is a likelihood ratio test]
For a normal sample, the unrestricted maximum of the likelihood uses $\hat\mu=\bar{x}$, $\hat\sigma^2=\frac1n\sum(x_i-\bar{x})^2$; under $H_0:\mu=\mu_0$ it uses $\hat\sigma_0^2=\frac1n\sum(x_i-\mu_0)^2=\hat\sigma^2+(\bar{x}-\mu_0)^2$. Then $\Lambda=(\hat\sigma^2/\hat\sigma_0^2)^{n/2}=\big[1+\tfrac{t_0^2}{n-1}\big]^{-n/2}$, a decreasing function of $\abs{t_0}$. Rejecting for small $\Lambda$ is the same as rejecting for large $\abs{t_0}$.
\end{example}

\begin{theorem}[Wilks' theorem]\label{thm:wilks}
Under $H_0$ and regularity conditions, $-2\ln\Lambda$ has approximately a chi-square distribution as $n\to\infty$, with degrees of freedom equal to the number of parameters restricted by $H_0$. So reject if $-2\ln\Lambda>\chi^2_{\alpha,r}$.
\end{theorem}

\section{Tests on the Variance and Standard Deviation of a Normal Distribution}

\subsection{The Hypothesis Testing Procedures}

Test $H_0:\sigma^2=\sigma_0^2$ for a normal population. The statistic is
\[ \chi_0^2=\frac{(n-1)S^2}{\sigma_0^2},\qquad\chi_0^2\sim\chi^2_{n-1}\text{ when }H_0\text{ is true.} \]

\begin{theorem}[$\chi^2$-test for a variance]\label{thm:chisqtest}
Reject $H_0:\sigma^2=\sigma_0^2$ in favour of
\begin{itemize}
  \item $H_1:\sigma^2\ne\sigma_0^2$ if $\chi_0^2>\chi^2_{\alpha/2,n-1}$ or $\chi_0^2<\chi^2_{1-\alpha/2,n-1}$;
  \item $H_1:\sigma^2>\sigma_0^2$ if $\chi_0^2>\chi^2_{\alpha,n-1}$;
  \item $H_1:\sigma^2<\sigma_0^2$ if $\chi_0^2<\chi^2_{1-\alpha,n-1}$.
\end{itemize}
Tests on $\sigma$ are the same with $\sigma_0^2$ replaced by its square.
\end{theorem}

\begin{example}\label{ex:vartest}
A filling machine is acceptable if $\sigma^2\le0.01$. A sample of $n=20$ has $s^2=0.0153$. Test $H_0:\sigma^2=0.01$ against $H_1:\sigma^2>0.01$ at $\alpha=0.05$. $\chi_0^2=19(0.0153)/0.01=29.07$ and $\chi^2_{0.05,19}=30.14$. Since $29.07<30.14$, do not reject $H_0$ ($P\approx0.065$): the evidence of a variance exceeding the limit is suggestive but not conclusive at the $5\%$ level.
\end{example}

\begin{example}[Textbook example: bottle filling]
$n=20$ bottles, $s^2=0.0153$; is the variance above $0.01$ (fluid ounces)$^2$? $H_0:\sigma^2=0.01$, $H_1:\sigma^2>0.01$, $\alpha=0.05$. $\chi_0^2=(n-1)s^2/\sigma_0^2=19(0.0153)/0.01=29.07<\chi^2_{0.05,19}=30.14$: do not reject, no strong evidence of a problem. From the table, $27.20<29.07<30.14$ gives $0.05<P<0.10$ (exact $P=0.0649$). One-sided tests use $\chi^2_{\alpha,n-1}$ in the upper tail (for $H_1:\sigma^2>\sigma_0^2$) or $\chi^2_{1-\alpha,n-1}$ in the lower tail (for $H_1:\sigma^2<\sigma_0^2$); two-sided uses $\chi^2_{\alpha/2,n-1}$ and $\chi^2_{1-\alpha/2,n-1}$. \emph{Power:} here $\sigma_0=0.10$; if the true $\sigma$ is $25\%$ larger ($\lambda=1.25$) the OC chart with $n=20$ gives $\beta\approx0.6$, so only about a $40\%$ chance of rejecting; to reach $\beta=0.20$ at $\lambda=1.25$ the chart gives $n\approx75$ bottles.
\end{example}

\subsection{$\beta$-Error and Choice of Sample Size}

If the true variance is $\sigma^2=\lambda^2\sigma_0^2$, then $(n-1)S^2/\sigma_0^2=\lambda^2V$ with $V\sim\chi^2_{n-1}$. For the two-sided test with critical values $c_L=\chi^2_{1-\alpha/2,n-1}$, $c_U=\chi^2_{\alpha/2,n-1}$,
\[ \beta=\Prob\big(c_L\le\lambda^2V\le c_U\big)=\Prob\Big(\frac{c_L}{\lambda^2}\le V\le\frac{c_U}{\lambda^2}\Big), \]
and for the upper one-sided test $\beta=\Prob\big(V\le\chi^2_{\alpha,n-1}/\lambda^2\big)$. OC curves in terms of $\lambda=\sigma/\sigma_0$ tabulate these. A large $\lambda$ (a big increase in the standard deviation) is detected with moderate $n$, but small departures from $\sigma_0$ need large samples.

\section{Tests on a Population Proportion}

\subsection{Large-Sample Tests on a Proportion}

Let $X$ be the number of successes in $n$ trials, $\hat{P}=X/n$, and test $H_0:p=p_0$. Under $H_0$, $X$ is binomial$(n,p_0)$, approximately normal with mean $np_0$ and variance $np_0(1-p_0)$ when $np_0\ge5$ and $n(1-p_0)\ge5$:
\[ Z_0=\frac{X-np_0}{\sqrt{np_0(1-p_0)}}=\frac{\hat{P}-p_0}{\sqrt{p_0(1-p_0)/n}}\approx\Normal(0,1). \]
The decision rules are those of Theorem~\ref{thm:ztest}.

\begin{example}\label{ex:proptest}
A process makes $5\%$ defectives. After a change, $x=4$ defectives are found in $n=200$ items. Test $H_0:p=0.05$ against $H_1:p<0.05$ at $\alpha=0.05$. $z_0=(4-200\cdot0.05)/\sqrt{200\cdot0.05\cdot0.95}=-6/3.082=-1.947<-1.645$, so reject $H_0$ ($P=\Phi(-1.947)=0.0258$): the defect rate has been reduced.
\end{example}

\begin{example}[Textbook example: process fallout]
A customer requires fraction defective at most $0.05$. To demonstrate capability the manufacturer takes $n=200$ devices and finds $4$ defective. The claim to be \emph{demonstrated} goes in $H_1$: $H_0:p=0.05$, $H_1:p<0.05$, $\alpha=0.05$. $z_0=(4-200(0.05))/\sqrt{200(0.05)(0.95)}=-6/3.082=-1.95<-1.645$: reject $H_0$ ($P=0.0256$); the process is capable. Software also gives the \emph{one-sided} $95\%$ upper confidence bound $0.0363<0.05$, consistent with rejection: a one-sided test goes with a one-sided bound, here an \emph{upper} bound because $H_1:p<p_0$. (A \emph{two-sided} $95\%$ interval would use $z_{0.025}=1.96$ and have a larger upper limit.) \emph{Power:} if the true $p$ is $0.03$, $\beta\approx0.7$ (power only $0.3$) with $n=200$; to allow $\beta=0.10$ at $p=0.03$ with $\alpha=0.05$ the sample size needed is $n\approx832$, large because $0.05$ versus $0.03$ is a small difference.
\end{example}

\subsection{Small-Sample Tests on a Proportion (CD Only)}

When the normal approximation is poor, use the exact binomial distribution. For $H_1:p>p_0$ the $P$-value is $\Prob(X\ge x\mid p=p_0)=\sum_{k\ge x}\binom nkp_0^k(1-p_0)^{n-k}$; for $H_1:p<p_0$ it is $\Prob(X\le x\mid p=p_0)$; for a two-sided alternative one doubles the smaller tail (capped at $1$).

\begin{example}
A machine should succeed $30\%$ of the time. If it succeeds $x=11$ times in $n=20$ trials, is its success rate higher ($H_0:p=0.3$ against $H_1:p>0.3$)? $P=\Prob(X\ge11\mid n=20,p=0.3)=0.0171$, so $H_0$ is rejected at $\alpha=0.05$.
\end{example}

\subsection{Type II Error and Choice of Sample Size}

If the true proportion is $p\ne p_0$, then $\hat{P}\approx\Normal\big(p,p(1-p)/n\big)$ and, for the two-sided test,
\[ \beta\approx\Phi\Big(\frac{p_0-p+z_{\alpha/2}\sqrt{p_0(1-p_0)/n}}{\sqrt{p(1-p)/n}}\Big)-\Phi\Big(\frac{p_0-p-z_{\alpha/2}\sqrt{p_0(1-p_0)/n}}{\sqrt{p(1-p)/n}}\Big). \]
Solving for $n$ (two-sided, ignoring the second term; replace $z_{\alpha/2}$ by $z_\alpha$ for one-sided):
\[ n\approx\left[\frac{z_{\alpha/2}\sqrt{p_0(1-p_0)}+z_\beta\sqrt{p(1-p)}}{p-p_0}\right]^2. \]

\begin{example}
To detect an increase of the defect rate from $p_0=0.05$ to $p=0.07$ with power $0.9$ ($\beta=0.1$) at $\alpha=0.05$ (one-sided): $n=\big[(1.645\sqrt{0.0475}+1.2816\sqrt{0.0651})/0.02\big]^2=\big[(0.3585+0.3270)/0.02\big]^2=1175$ items.
\end{example}

\section{Summary of Inference Procedures for a Single Sample}

\begin{center}
\small
\begin{tabular}{p{2.4cm}p{3.2cm}p{3.6cm}p{3.4cm}}
\toprule
Parameter and assumptions & Test statistic & Reject $H_0$ (two-sided) & Confidence interval \\
\midrule
$\mu$, $\sigma$ known (normal or $n$ large) & $z_0=\dfrac{\bar{x}-\mu_0}{\sigma/\sqrt n}$ & $\abs{z_0}>z_{\alpha/2}$ & $\bar{x}\pm z_{\alpha/2}\sigma/\sqrt n$ \\
$\mu$, $\sigma$ unknown (normal) & $t_0=\dfrac{\bar{x}-\mu_0}{s/\sqrt n}$ & $\abs{t_0}>t_{\alpha/2,n-1}$ & $\bar{x}\pm t_{\alpha/2,n-1}s/\sqrt n$ \\
$\sigma^2$ (normal) & $\chi_0^2=\dfrac{(n-1)s^2}{\sigma_0^2}$ & $\chi_0^2>\chi^2_{\alpha/2,n-1}$ or $<\chi^2_{1-\alpha/2,n-1}$ & $\dfrac{(n-1)s^2}{\chi^2_{\alpha/2}}\le\sigma^2\le\dfrac{(n-1)s^2}{\chi^2_{1-\alpha/2}}$ \\
$p$ (large $n$) & $z_0=\dfrac{x-np_0}{\sqrt{np_0(1-p_0)}}$ & $\abs{z_0}>z_{\alpha/2}$ & $\hat{p}\pm z_{\alpha/2}\sqrt{\hat{p}(1-\hat{p})/n}$ \\
\bottomrule
\end{tabular}
\end{center}

The table above is for \emph{two-sided} alternatives. The one-sided versions, with the matching one-sided $100(1-\alpha)\%$ bound, are:

\begin{center}
\small
\begin{tabular}{p{2.2cm}p{3.5cm}p{3.5cm}p{3.5cm}}
\toprule
Parameter & $H_1$ and rejection region (level $\alpha$) & Matching bound ($\alpha$, one limit) & Critical value \\
\midrule
$\mu$, $\sigma$ known & $\mu>\mu_0$: $z_0>z_\alpha$ ;\quad $\mu<\mu_0$: $z_0<-z_\alpha$ & lower: $\bar{x}-z_\alpha\sigma/\sqrt n\le\mu$ ;\quad upper: $\mu\le\bar{x}+z_\alpha\sigma/\sqrt n$ & $z_\alpha$ \\
$\mu$, $\sigma$ unknown & $\mu>\mu_0$: $t_0>t_{\alpha,n-1}$ ;\quad $\mu<\mu_0$: $t_0<-t_{\alpha,n-1}$ & lower: $\bar{x}-t_{\alpha,n-1}s/\sqrt n\le\mu$ ;\quad upper: $\mu\le\bar{x}+t_{\alpha,n-1}s/\sqrt n$ & $t_{\alpha,n-1}$ \\
$\sigma^2$ & $\sigma^2>\sigma_0^2$: $\chi_0^2>\chi^2_{\alpha,n-1}$ ;\quad $\sigma^2<\sigma_0^2$: $\chi_0^2<\chi^2_{1-\alpha,n-1}$ & lower: $\sigma^2\ge(n-1)s^2/\chi^2_{\alpha,n-1}$ ;\quad upper: $\sigma^2\le(n-1)s^2/\chi^2_{1-\alpha,n-1}$ & $\chi^2_{\alpha}$ or $\chi^2_{1-\alpha}$ \\
$p$ (large $n$) & $p>p_0$: $z_0>z_\alpha$ ;\quad $p<p_0$: $z_0<-z_\alpha$ & lower: $\hat{p}-z_\alpha\widehat{\se}\le p$ ;\quad upper: $p\le\hat{p}+z_\alpha\widehat{\se}$ & $z_\alpha$ \\
\bottomrule
\end{tabular}
\end{center}

Reading rule: a ``greater than'' alternative is paired with a \emph{lower} bound (the data must push the parameter above $\mu_0$, so we ask whether the lowest plausible value is still above it); a ``less than'' alternative is paired with an \emph{upper} bound.

\section{Testing for Goodness of Fit}

A \emph{goodness-of-fit test} checks whether the population follows a specified distribution. Divide the range of the variable into $k$ classes (bins). Let $O_i$ be the observed frequency in class $i$ and $E_i=np_i$ the expected frequency, where $p_i$ is the probability of class $i$ under the hypothesised distribution.

\begin{theorem}[Chi-square goodness-of-fit test]\label{thm:gof}
Under $H_0$ (the data come from the specified distribution), the statistic
\[ X_0^2=\sum_{i=1}^k\frac{(O_i-E_i)^2}{E_i} \]
has approximately a chi-square distribution with $k-p-1$ degrees of freedom, where $p$ is the number of parameters of the distribution that were \emph{estimated} from the data. Reject $H_0$ if $x_0^2>\chi^2_{\alpha,k-p-1}$. The approximation requires every $E_i\ge5$ (combine adjacent classes if necessary).
\end{theorem}

\begin{example}[Fair die]
A die is rolled $60$ times, with counts $8,12,9,11,14,6$ for faces $1$ to $6$. Under $H_0$ (fair die) $E_i=10$. $x_0^2=(4+4+1+1+16+16)/10=4.2$ with $k-1=5$ degrees of freedom ($p=0$). Since $4.2<\chi^2_{0.05,5}=11.07$, we do not reject fairness.
\end{example}

\begin{example}[Poisson with estimated mean]
Suppose the number of defects per unit is to be tested against a Poisson distribution whose mean is estimated by $\bar{x}$ from the same data. With $k=5$ classes after combining, $p=1$ parameter is estimated, so the degrees of freedom are $k-p-1=3$. Estimating a parameter uses up one degree of freedom because it makes the fit look better than it would for a fixed specified distribution.
\end{example}

\begin{example}[Textbook examples: goodness of fit]
\emph{Poisson.} $n=60$ printed circuit boards with $32, 15, 9, 4$ boards having $0,1,2,3$ defects: the estimated mean is $\hat\lambda=(0\cdot32+1\cdot15+2\cdot9+3\cdot4)/60=0.75$. Expected counts are $60p_i$ with $p_i$ from Poisson($0.75$); the last expected count is below $3$, so the last two cells are combined, leaving $k=3$ cells, and the degrees of freedom are $k-p-1=3-1-1=1$ because one parameter was estimated. $\chi_0^2=2.94<\chi^2_{0.05,1}=3.84$: cannot reject the Poisson model ($P=0.0864$). Minimum expected frequencies of $3$, $4$ or $5$ are widely used; cells with smaller expected counts are combined with neighbours and $k$ is reduced.
\emph{Normal, equiprobable cells.} For $n=100$ power-supply voltages ($\bar{x}=5.04$, $s=0.08$) choose $k=8$ cells with equal probability $0.125$, so every expected count is $np_i=12.5$; for the standard normal the boundaries are $0,\ \pm0.32,\ \pm0.675,\ \pm1.15$ (times $s$, plus $\bar{x}$). Two parameters are estimated, so the degrees of freedom are $8-2-1=5$; the test statistic is $0.64$, against $\chi^2_{0.05,5}=11.07$ (the textbook compares with $11.07$): not rejected ($P=0.9861$).
\end{example}

\section{Contingency Table Tests}

Two classifications of the same units can be summarised in an $r\times c$ \emph{contingency table}, where $O_{ij}$ is the number of units in row category $i$ and column category $j$. We ask whether the two classifications are \emph{independent}.

\begin{theorem}[Chi-square test of independence]\label{thm:indep}
Let $R_i$ be the row totals, $C_j$ the column totals and $n$ the grand total. Under $H_0$: ``the row and column classifications are independent'', the expected counts are $E_{ij}=\dfrac{R_iC_j}n$, and
\[ X_0^2=\sum_{i=1}^r\sum_{j=1}^c\frac{(O_{ij}-E_{ij})^2}{E_{ij}}\ \approx\ \chi^2_{(r-1)(c-1)}. \]
Reject $H_0$ if $x_0^2>\chi^2_{\alpha,(r-1)(c-1)}$ (again requiring $E_{ij}\ge5$ roughly).
\end{theorem}

\begin{proof}[Why the degrees of freedom are $(r-1)(c-1)$]
There are $rc$ cells. The row and column probabilities $p_{i\cdot}$ and $p_{\cdot j}$ must be estimated from the data; there are $(r-1)+(c-1)$ free estimates (each set sums to $1$). The degrees of freedom are $rc-1-[(r-1)+(c-1)]=(r-1)(c-1)$.
\end{proof}

\begin{example}\label{ex:cont}
Two treatments are compared: treatment A has $40$ successes and $10$ failures, treatment B has $30$ successes and $20$ failures. The row totals are $50,50$, column totals $70,30$ and $n=100$. The expected counts are $35,15,35,15$. $x_0^2=\dfrac{25}{35}+\dfrac{25}{15}+\dfrac{25}{35}+\dfrac{25}{15}=4.76$ with $1$ degree of freedom. Since $4.76>\chi^2_{0.05,1}=3.84$, reject independence ($P=0.029$): the success rate depends on the treatment.
\end{example}

The same test applies when the data are a set of random samples from $r$ populations classified into $c$ categories (a \emph{test of homogeneity}): $H_0$ is that all populations have the same category probabilities, and the statistic and degrees of freedom are identical.

\begin{interview}{common slips with chi-square tests}
Counts, never percentages, go into $O_i$ and $E_i$. The statistic is always upper-tailed (large values indicate poor fit). Small expected counts invalidate the approximation. Degrees of freedom are $k-1$ minus the number of estimated parameters for goodness of fit and $(r-1)(c-1)$ for independence. A significant result shows association, not causation, and says nothing about its strength: use a measure such as Cram\'er's $V=\sqrt{x_0^2/(n\min(r-1,c-1))}$.
\end{interview}

\begin{summarybox}
\begin{itemize}
  \item $H_0$ is assumed true until the data contradict it; $\alpha=\Prob(\text{reject }H_0\mid H_0)$, $\beta=\Prob(\text{fail to reject}\mid H_0\text{ false})$, power $=1-\beta$. Rejecting is strong; failing to reject is weak.
  \item Follow the eight steps. The $P$-value is the probability of a result at least as extreme if $H_0$ were true; reject if $P<\alpha$. A two-sided test rejects iff the value $\mu_0$ falls outside the $100(1-\alpha)\%$ confidence interval.
  \item $z$-test (known $\sigma$): $\beta=\Phi(z_{\alpha/2}-\delta\sqrt n/\sigma)-\Phi(-z_{\alpha/2}-\delta\sqrt n/\sigma)$, $n\approx(z_{\alpha/2}+z_\beta)^2\sigma^2/\delta^2$. $t$-test: $t_0=(\bar{x}-\mu_0)/(s/\sqrt n)\sim t_{n-1}$.
  \item Variance: $\chi_0^2=(n-1)s^2/\sigma_0^2\sim\chi^2_{n-1}$. Proportion: $z_0=(x-np_0)/\sqrt{np_0(1-p_0)}$, exact binomial for small $n$.
  \item Likelihood ratio tests generate the standard tests; $-2\ln\Lambda\approx\chi^2$ (Wilks).
  \item Goodness of fit: $\sum(O-E)^2/E\sim\chi^2_{k-p-1}$; independence in an $r\times c$ table: $\chi^2_{(r-1)(c-1)}$ with $E_{ij}=R_iC_j/n$.
\end{itemize}
\end{summarybox}

\section{Exercises}

\begin{exercise}
The mean breaking strength of a cable is claimed to be at least $1800$ lb. A sample of $50$ cables has $\bar{x}=1792$ with $\sigma=15$ known. Test the claim at $\alpha=0.01$ with the eight-step procedure and find the $P$-value.
\end{exercise}

\begin{exercise}
For the test in Exercise~1, compute $\beta$ when the true mean is $1785$ and the sample size needed for $\beta=0.05$ at that alternative.
\end{exercise}

\begin{exercise}
Twelve measurements of the octane rating of a gasoline have $\bar{x}=88.3$ and $s=0.9$. Test $H_0:\mu=87.5$ against $H_1:\mu\ne87.5$ at $\alpha=0.05$ and compare your decision with the $95\%$ confidence interval.
\end{exercise}

\begin{exercise}
Explain why the sentence ``the probability that the null hypothesis is true is $0.03$'' is not a correct reading of $P=0.03$. Give a correct reading.
\end{exercise}

\begin{exercise}
A sample of $n=25$ has $s^2=18.5$. Test $H_0:\sigma^2=12$ against $H_1:\sigma^2>12$ at $\alpha=0.05$ and compute $\beta$ when the true variance is $24$.
\end{exercise}

\begin{exercise}
In a poll of $n=500$, $265$ favour a proposal. Test $H_0:p=0.5$ against $H_1:p>0.5$ at $\alpha=0.05$. How large a sample would detect $p=0.55$ with power $0.8$?
\end{exercise}

\begin{exercise}
The numbers of accidents per week over $50$ weeks are: $0$: 17 weeks, $1$: 19, $2$: 9, $3$: 4, $4$ or more: 1. Test whether a Poisson distribution with $\lambda$ estimated from the data fits ($\alpha=0.05$; combine classes as needed).
\end{exercise}

\begin{exercise}
A survey cross-classifies $200$ employees by shift (day/night) and quality rating (high/medium/low): day: $50,\ 30,\ 20$; night: $20,\ 35,\ 45$. Test independence at $\alpha=0.05$ and describe the association.
\end{exercise}

\chapter{Statistical Inference for Two Samples}

\section{Introduction}

Let $X_{11},\dots,X_{1n_1}$ be a random sample from population~1 (mean $\mu_1$, variance $\sigma_1^2$) and $X_{21},\dots,X_{2n_2}$ an \emph{independent} random sample from population~2 (mean $\mu_2$, variance $\sigma_2^2$). Throughout, the two samples are independent of one another, except in Section~\ref{sec:paired} where the observations come in matched pairs. Typical questions: Is catalyst A better than catalyst B? Do the two suppliers have equal variability? Is the defect rate lower on the new line?

\section{Inference For a Difference in Means of Two Normal Distributions, Variances Known}

Assume both populations are normal (or the samples large) and $\sigma_1^2,\sigma_2^2$ are known. The sampling distribution (Chapter~1, Theorem~\ref{thm:diffmeans}) is
\[ \bar{X}_1-\bar{X}_2\sim\Normal\Big(\mu_1-\mu_2,\ \frac{\sigma_1^2}{n_1}+\frac{\sigma_2^2}{n_2}\Big). \]

\subsection{Hypothesis Tests for a Difference in Means, Variances Known}

To test $H_0:\mu_1-\mu_2=\Delta_0$ (usually $\Delta_0=0$) use
\[ Z_0=\frac{\bar{X}_1-\bar{X}_2-\Delta_0}{\sqrt{\sigma_1^2/n_1+\sigma_2^2/n_2}}\sim\Normal(0,1)\ \text{ under }H_0. \]

\begin{theorem}[Two-sample $z$-test]\label{thm:twoz}
Reject $H_0:\mu_1-\mu_2=\Delta_0$ if $\abs{z_0}>z_{\alpha/2}$ (for $H_1:\mu_1-\mu_2\ne\Delta_0$), if $z_0>z_\alpha$ (for $H_1:\mu_1-\mu_2>\Delta_0$) or if $z_0<-z_\alpha$ (for $H_1:\mu_1-\mu_2<\Delta_0$). The $P$-values are $2[1-\Phi(\abs{z_0})]$, $1-\Phi(z_0)$ and $\Phi(z_0)$.
\end{theorem}

\begin{example}\label{ex:twoz}
Sample~1: $n_1=12$, $\bar{x}_1=75.2$, $\sigma_1=3$. Sample~2: $n_2=15$, $\bar{x}_2=72.8$, $\sigma_2=4$. Test $H_0:\mu_1=\mu_2$ against $H_1:\mu_1>\mu_2$ at $\alpha=0.05$. $\sigma_D=\sqrt{9/12+16/15}=\sqrt{1.8167}=1.348$ and $z_0=2.4/1.348=1.78>1.645$, so reject $H_0$ ($P=0.0375$). For the two-sided alternative $P=0.075$ and $H_0$ would not be rejected at $0.05$.
\end{example}

\subsection{Choice of Sample Size}

Let the true difference be $\mu_1-\mu_2=\Delta$ and write $\sigma_D=\sqrt{\sigma_1^2/n_1+\sigma_2^2/n_2}$. For the two-sided test
\[ \beta=\Phi\Big(z_{\alpha/2}-\frac{\Delta-\Delta_0}{\sigma_D}\Big)-\Phi\Big(-z_{\alpha/2}-\frac{\Delta-\Delta_0}{\sigma_D}\Big), \]
and one-sided versions follow as in Chapter~3. For equal sample sizes $n_1=n_2=n$:

\begin{theorem}[Sample size for the two-sample $z$-test]\label{thm:twoz-n}
To detect a true difference $\Delta-\Delta_0$ with Type II error $\beta$ at level $\alpha$,
\[ n\approx\frac{(z_{\alpha/2}+z_\beta)^2(\sigma_1^2+\sigma_2^2)}{(\Delta-\Delta_0)^2}\quad\text{(two-sided)},\qquad n=\frac{(z_\alpha+z_\beta)^2(\sigma_1^2+\sigma_2^2)}{(\Delta-\Delta_0)^2}\quad\text{(one-sided)}. \]
\end{theorem}

\begin{nuance}{operating-characteristic (OC) curves and the equivalent sample size}
The textbook reads $\beta$ from OC curves (Appendix charts VIIa--VIIf) using the parameter
\[ d=\frac{\abs{\Delta-\Delta_0}}{\sqrt{\sigma_1^2+\sigma_2^2}}. \]
When $n_1\ne n_2$ or $\sigma_1\ne\sigma_2$ the chart is entered with the \emph{equivalent} common sample size
\[ n=\frac{\sigma_1^2+\sigma_2^2}{\sigma_1^2/n_1+\sigma_2^2/n_2}. \]
Large samples: if $n_1,n_2>40$, the sample standard deviations may be used for $\sigma_1,\sigma_2$ and the $z$-test is applied even when the populations are not normal (central limit theorem).
\end{nuance}

\begin{example}
With $\sigma_1=3$, $\sigma_2=4$, to detect a difference of $2$ with power $0.9$ at $\alpha=0.05$ (two-sided): $n=(1.96+1.2816)^2(9+16)/4=65.7$, so $66$ per group.
\end{example}

\subsection{Identifying Cause and Effect}

A significant difference between two samples shows that the populations differ; it does not by itself show \emph{why}.

\begin{itemize}
  \item In a \emph{designed experiment}, the experimenter assigns units to treatments \emph{at random}. Randomisation spreads all other factors (known and unknown) evenly across groups on average, so a significant difference can be attributed to the treatment: cause and effect.
  \item In an \emph{observational study} the groups form themselves (smokers versus non-smokers, users versus non-users). A difference may be due to a \emph{confounding} variable related both to group membership and to the response. Statistical significance then supports association only.
\end{itemize}

\begin{example}[Simpson's paradox]
Treatment~A cures $93\%$ of small stones (81 of 87) and $73\%$ of large stones (192 of 263); treatment~B cures $87\%$ of small stones (234 of 270) and $69\%$ of large stones (55 of 80). A is better within each stone size, yet B looks better overall ($83\%$ versus $78\%$) because B was given mostly to the easier small stones. Without randomisation the size of the stone confounds the comparison.
\end{example}

\subsection{Confidence Interval on a Difference in Means, Variances Known}

\begin{theorem}[Interval for $\mu_1-\mu_2$, $\sigma_1,\sigma_2$ known]\label{thm:twoz-ci}
A $100(1-\alpha)\%$ confidence interval is
\[ (\bar{x}_1-\bar{x}_2)-z_{\alpha/2}\sqrt{\frac{\sigma_1^2}{n_1}+\frac{\sigma_2^2}{n_2}}\ \le\ \mu_1-\mu_2\ \le\ (\bar{x}_1-\bar{x}_2)+z_{\alpha/2}\sqrt{\frac{\sigma_1^2}{n_1}+\frac{\sigma_2^2}{n_2}}. \]
\end{theorem}

\begin{theorem}[One-sided bounds for $\mu_1-\mu_2$, $\sigma_1,\sigma_2$ known]\label{thm:twoz-onesided}
With $\sigma_D=\sqrt{\sigma_1^2/n_1+\sigma_2^2/n_2}$, a $100(1-\alpha)\%$ \emph{upper} confidence bound and a $100(1-\alpha)\%$ \emph{lower} confidence bound are
\[ \mu_1-\mu_2\ \le\ (\bar{x}_1-\bar{x}_2)+z_{\alpha}\,\sigma_D, \qquad (\bar{x}_1-\bar{x}_2)-z_{\alpha}\,\sigma_D\ \le\ \mu_1-\mu_2. \]
Each bound uses $z_\alpha$ (all of $\alpha$ in one tail), not $z_{\alpha/2}$.
\end{theorem}

\begin{nuance}{one-sided versus two-sided, stated explicitly}
\begin{itemize}
  \item A \emph{two-sided} $100(1-\alpha)\%$ interval puts $\alpha/2$ in each tail and uses the critical value with $\alpha/2$ ($z_{\alpha/2}$ or $t_{\alpha/2,\nu}$).
  \item A \emph{one-sided} $100(1-\alpha)\%$ bound puts all of $\alpha$ in one tail and uses the critical value with $\alpha$ ($z_{\alpha}$ or $t_{\alpha,\nu}$).
  \item The two are \emph{different} statements. The upper end of a two-sided $95\%$ interval is a one-sided $97.5\%$ upper bound; a one-sided $95\%$ upper bound is tighter (it uses $z_{0.05}=1.645$, not $z_{0.025}=1.96$).
  \item Pairing with tests: $H_1:\mu_1-\mu_2>\Delta_0$ is matched by the \emph{lower} bound (reject iff $\Delta_0$ lies below it); $H_1:\mu_1-\mu_2<\Delta_0$ is matched by the \emph{upper} bound; $H_1:\mu_1-\mu_2\ne\Delta_0$ is matched by the two-sided interval, at the same $\alpha$.
\end{itemize}
\end{nuance}

\begin{example}[sample size for a confidence interval]\label{ex:twoz-ci-n}
If $n_1=n_2=n$, the half-width of the two-sided $100(1-\alpha)\%$ interval is $E=z_{\alpha/2}\sqrt{(\sigma_1^2+\sigma_2^2)/n}$, so to have error at most $E$,
\[ n=\Big(\frac{z_{\alpha/2}}{E}\Big)^2(\sigma_1^2+\sigma_2^2). \]
For a one-sided bound replace $z_{\alpha/2}$ by $z_\alpha$. With $\sigma_1=3$, $\sigma_2=4$, $E=1$ and $95\%$ two-sided: $n=1.96^2(25)=96.04$, so $97$ per group.
\end{example}

\begin{example}
For Example~\ref{ex:twoz}: $2.4\pm1.96(1.348)=2.4\pm2.64$, i.e.\ $-0.24\le\mu_1-\mu_2\le5.04$. The interval contains $0$, matching the two-sided non-rejection at $\alpha=0.05$ (duality, Chapter~3).
\end{example}

\section{Inference For a Difference in Means of Two Normal Distributions, Variances Unknown}

\subsection{Hypothesis Tests for a Difference in Means, Variances Unknown}

\subsubsection*{Case 1: equal variances, $\sigma_1^2=\sigma_2^2=\sigma^2$}

Both sample variances estimate the same $\sigma^2$, so we pool them.

\begin{definition}[Pooled variance]
$S_p^2=\dfrac{(n_1-1)S_1^2+(n_2-1)S_2^2}{n_1+n_2-2}$.
\end{definition}

\begin{theorem}[Pooled two-sample $t$-test]\label{thm:pooledt}
If the populations are normal with equal variances, then under $H_0:\mu_1-\mu_2=\Delta_0$
\[ T_0=\frac{\bar{X}_1-\bar{X}_2-\Delta_0}{S_p\sqrt{1/n_1+1/n_2}}\sim t_{n_1+n_2-2}. \]
Reject $H_0$ if $\abs{t_0}>t_{\alpha/2,\,n_1+n_2-2}$ (two-sided), $t_0>t_{\alpha,n_1+n_2-2}$ ($H_1:\mu_1-\mu_2>\Delta_0$), or $t_0<-t_{\alpha,n_1+n_2-2}$.
\end{theorem}

\begin{proof}
$(\bar{X}_1-\bar{X}_2-\Delta_0)/\big(\sigma\sqrt{1/n_1+1/n_2}\big)\sim\Normal(0,1)$, and $(n_1+n_2-2)S_p^2/\sigma^2=(n_1-1)S_1^2/\sigma^2+(n_2-1)S_2^2/\sigma^2$ is a sum of two independent $\chi^2$ variables, hence $\chi^2_{n_1+n_2-2}$, independent of the difference of means. Their ratio is the $t$ statistic.
\end{proof}

\begin{example}\label{ex:pooled}
Method~A: $n_1=10$, $\bar{x}_1=48.2$, $s_1=3.1$. Method~B: $n_2=12$, $\bar{x}_2=45.1$, $s_2=3.7$. Test $H_0:\mu_1=\mu_2$ against $H_1:\mu_1\ne\mu_2$ at $\alpha=0.05$ assuming equal variances. $s_p^2=(9\cdot9.61+11\cdot13.69)/20=11.854$, $s_p=3.443$, and the standard error is $3.443\sqrt{1/10+1/12}=1.474$. $t_0=3.1/1.474=2.10$ with $20$ degrees of freedom; $t_{0.025,20}=2.086$. Since $2.10>2.086$, reject $H_0$ (just barely; $P\approx0.048$).
\end{example}

\subsubsection*{Case 2: unequal variances, $\sigma_1^2\ne\sigma_2^2$}

When the variances may differ, use the unpooled statistic. It is only approximately $t$-distributed (the Behrens--Fisher problem).

\begin{theorem}[Welch's $t$-test]\label{thm:welch}
Let $T_0^*=\dfrac{\bar{X}_1-\bar{X}_2-\Delta_0}{\sqrt{S_1^2/n_1+S_2^2/n_2}}$. Under $H_0$, $T_0^*$ is approximately $t_\nu$ with (Satterthwaite)
\[ \nu=\frac{\big(s_1^2/n_1+s_2^2/n_2\big)^2}{\dfrac{(s_1^2/n_1)^2}{n_1-1}+\dfrac{(s_2^2/n_2)^2}{n_2-1}}, \]
rounded down to an integer. The rejection regions are those of Theorem~\ref{thm:pooledt} with $\nu$ degrees of freedom.
\end{theorem}

\begin{example}
For the data of Example~\ref{ex:pooled}: $\sqrt{s_1^2/n_1+s_2^2/n_2}=\sqrt{0.961+1.141}=1.450$, $t_0^*=3.1/1.450=2.14$ and $\nu=(2.102)^2/[(0.961)^2/9+(1.141)^2/11]=4.418/0.2209\approx19.99$, so $\nu=19$ degrees of freedom. With $t_{0.025,19}=2.093$ the conclusion is the same.
\end{example}

\subsection{More About the Equal Variance Assumption (CD Only)}

The pooled test is slightly more powerful if $\sigma_1^2=\sigma_2^2$, but it can be misleading if the variances differ \emph{and} the sample sizes are unequal: when the smaller sample comes from the more variable population, the pooled test's true Type I error can far exceed $\alpha$. With equal sample sizes $n_1=n_2$ the two statistics are numerically identical ($t_0=t_0^*$), and the pooled test is reasonably robust. The pooled test is an exact likelihood-ratio test when normality, independence and equal variances hold, whereas Welch's test is only approximate; but the pooled denominator effectively divides each variance by the wrong sample size when $\sigma_1^2\ne\sigma_2^2$ and $n_1\ne n_2$, which is why choosing $n_1=n_2$ is good practice. A practical check is the normal probability plot of both samples: clearly different slopes of the two straight lines indicate unequal variances, and Welch's procedure is then preferred. Welch's test loses very little when the variances are in fact equal, which is why many analysts use it by default and avoid the preliminary $F$-test of Section~\ref{sec:ftest}, which is itself sensitive to non-normality and affects the overall error rate.

\subsection{Choice of Sample Size}

The power of the pooled test depends on the noncentrality $\abs{\Delta-\Delta_0}/\big(\sigma\sqrt{1/n_1+1/n_2}\big)$. For $n_1=n_2=n$, power is read from OC curves in terms of $d=\abs{\Delta-\Delta_0}/(2\sigma)$ and $n^*=2n-1$. A normal-theory approximation, adequate for planning, is
\[ n\approx\frac{2(z_{\alpha/2}+z_\beta)^2\sigma^2}{(\Delta-\Delta_0)^2}+\frac{z_{\alpha/2}^2}4 \]
per group, where $\sigma$ is a prior estimate of the common standard deviation.

\subsection{Confidence Interval on a Difference in Means, Variances Unknown}

\begin{theorem}[$t$ intervals for $\mu_1-\mu_2$]\label{thm:twot-ci}
Equal variances: $(\bar{x}_1-\bar{x}_2)\pm t_{\alpha/2,n_1+n_2-2}\,s_p\sqrt{1/n_1+1/n_2}$. Unequal variances: $(\bar{x}_1-\bar{x}_2)\pm t_{\alpha/2,\nu}\sqrt{s_1^2/n_1+s_2^2/n_2}$ with $\nu$ from Theorem~\ref{thm:welch}.
\end{theorem}

\begin{theorem}[One-sided $t$ bounds for $\mu_1-\mu_2$]
Replace $t_{\alpha/2,\cdot}$ by $t_{\alpha,\cdot}$ and keep only one end. Equal variances: $\mu_1-\mu_2\le(\bar{x}_1-\bar{x}_2)+t_{\alpha,n_1+n_2-2}\,s_p\sqrt{1/n_1+1/n_2}$ (upper) and $(\bar{x}_1-\bar{x}_2)-t_{\alpha,n_1+n_2-2}\,s_p\sqrt{1/n_1+1/n_2}\le\mu_1-\mu_2$ (lower). Unequal variances: the same with $t_{\alpha,\nu}\sqrt{s_1^2/n_1+s_2^2/n_2}$.
\end{theorem}

\begin{example}
For Example~\ref{ex:pooled}: $3.1\pm2.086(1.474)=3.1\pm3.07$, i.e.\ $0.03\le\mu_1-\mu_2\le6.17$. The lower limit is barely above $0$, consistent with the marginal rejection.
\end{example}

\section{Paired $t$-Test}\label{sec:paired}

Sometimes observations come in natural pairs: the same subject before and after a treatment, two methods applied to the same specimen. The two measurements in a pair are positively correlated, which an independent-samples analysis ignores.

\begin{definition}[Paired design]
Observe pairs $(X_{1j},X_{2j})$, $j=1,\dots,n$, where the differences $D_j=X_{1j}-X_{2j}$ are independent and normal with mean $\mu_D=\mu_1-\mu_2$ and variance $\sigma_D^2$.
\end{definition}

\begin{theorem}[Paired $t$-test and interval]\label{thm:pairedt}
Let $\bar{D}$ and $S_D$ be the sample mean and standard deviation of the $n$ differences. To test $H_0:\mu_D=\Delta_0$ use
\[ T_0=\frac{\bar{D}-\Delta_0}{S_D/\sqrt n}\sim t_{n-1}\quad\text{under }H_0, \]
rejecting for $\abs{t_0}>t_{\alpha/2,n-1}$ (two-sided). A two-sided $100(1-\alpha)\%$ confidence interval for $\mu_D$ is $\bar{d}\pm t_{\alpha/2,n-1}s_D/\sqrt n$. One-sided: reject for $t_0>t_{\alpha,n-1}$ ($H_1:\mu_D>\Delta_0$) or $t_0<-t_{\alpha,n-1}$ ($H_1:\mu_D<\Delta_0$); the $100(1-\alpha)\%$ upper bound is $\mu_D\le\bar{d}+t_{\alpha,n-1}s_D/\sqrt n$ and the lower bound is $\bar{d}-t_{\alpha,n-1}s_D/\sqrt n\le\mu_D$.
\end{theorem}

The paired test is just the one-sample $t$-test (Chapter~3) applied to the differences.

\begin{example}\label{ex:paired}
Eight subjects' times were measured with and without a device, giving differences $2.1,\ 1.4,\ -0.3,\ 3.0,\ 1.8,\ 0.9,\ 2.5,\ 1.2$. $\bar{d}=1.575$, $s_d=1.025$, $\se=1.025/\sqrt8=0.362$, $t_0=1.575/0.362=4.35$ with $7$ degrees of freedom. Since $t_{0.005,7}=3.499<4.35$, reject $H_0:\mu_D=0$ ($P\approx0.003$); the $95\%$ interval is $1.575\pm2.365(0.362)=1.575\pm0.857$.
\end{example}

\begin{interview}{paired versus independent samples}
Pairing removes the variability between pairs (differences among subjects) and focuses on the within-pair difference. If the correlation within pairs is $\rho$, $\Var(D)=2\sigma^2(1-\rho)$, so pairing helps when $\rho>0$ and the variance of the mean difference shrinks by the factor $1-\rho$. The cost is that the degrees of freedom drop from $2n-2$ to $n-1$, which matters only if $\rho$ is small. \emph{Never} use an independent-samples test on paired data: it discards the pairing and usually loses a great deal of power. Conversely, never pair arbitrarily; the pairing must be built into the design (the same subject, matched batches).
\end{interview}

The textbook's guidelines for choosing between the two designs are:
\begin{enumerate}
  \item If the experimental units are relatively homogeneous (small $\sigma$) and the correlation within pairs is small, the gain from pairing is offset by the loss of degrees of freedom, so use independent samples.
  \item If the units are heterogeneous (large $\sigma$) and the positive correlation within pairs is large, use the paired design. This typically happens when the same units are used for both treatments.
\end{enumerate}
If the degrees of freedom are large (say $40$ or $50$) the loss from pairing is not serious; if they are small (say $10$ or $20$), losing half of them is serious unless pairing buys enough precision. The paired interval remains valid when $\sigma_1^2\ne\sigma_2^2$ (since $s_D^2$ estimates $\Var(X_1-X_2)$), and for large samples (about $30$ pairs or more) normality of the differences is not needed by the central limit theorem.

\section{Inference on the Variances of Two Normal Distributions}\label{sec:ftest}

\subsection{The $F$ Distribution}

\begin{definition}[$F$ distribution]
If $U\sim\chi^2_u$ and $V\sim\chi^2_v$ are independent, then $F=\dfrac{U/u}{V/v}$ has an $F$ distribution with $u$ numerator and $v$ denominator degrees of freedom, $F\sim F_{u,v}$. The distribution is skewed right, takes positive values and depends on both degrees of freedom. Let $f_{\alpha,u,v}$ be the point with right-tail area $\alpha$.
\end{definition}

A useful identity: $f_{1-\alpha,u,v}=\dfrac1{f_{\alpha,v,u}}$, so only upper-tail points need be tabulated (Table V gives upper points only). For example $f_{0.05,5,10}=3.33$ and $f_{0.95,5,10}=1/f_{0.05,10,5}=1/4.74=0.211$. The mean is $v/(v-2)$ for $v>2$ and the variance is $\dfrac{2v^2(u+v-2)}{u(v-2)^2(v-4)}$ for $v>4$.

\subsection{Development of the $F$ Distribution (CD Only)}

For independent normal samples, $(n_1-1)S_1^2/\sigma_1^2\sim\chi^2_{n_1-1}$ and $(n_2-1)S_2^2/\sigma_2^2\sim\chi^2_{n_2-1}$ are independent (Chapter~11, Theorem~\ref{thm:cochran}). Dividing each by its degrees of freedom,
\[ F=\frac{S_1^2/\sigma_1^2}{S_2^2/\sigma_2^2}\sim F_{n_1-1,\,n_2-1}. \]

\subsection{Hypothesis Tests on the Ratio of Two Variances}

\begin{theorem}[$F$-test for equality of variances]\label{thm:ftest}
To test $H_0:\sigma_1^2=\sigma_2^2$, use $F_0=S_1^2/S_2^2\sim F_{n_1-1,n_2-1}$ under $H_0$. Reject $H_0$ in favour of $H_1:\sigma_1^2\ne\sigma_2^2$ if $f_0>f_{\alpha/2,n_1-1,n_2-1}$ or $f_0<f_{1-\alpha/2,n_1-1,n_2-1}$; in favour of $H_1:\sigma_1^2>\sigma_2^2$ if $f_0>f_{\alpha,n_1-1,n_2-1}$; and in favour of $H_1:\sigma_1^2<\sigma_2^2$ if $f_0<f_{1-\alpha,n_1-1,n_2-1}$.
\end{theorem}

\begin{example}
For Example~\ref{ex:pooled}: $s_1^2=9.61$ ($n_1=10$) and $s_2^2=13.69$ ($n_2=12$), so $f_0=0.702$. With $9$ and $11$ degrees of freedom, $f_{0.025,9,11}=3.59$ and $f_{0.975,9,11}=1/f_{0.025,11,9}=1/3.91=0.256$. Since $0.256<0.702<3.59$, there is no evidence that the variances differ.

\emph{$P$-value for a two-sided $F$-test.} Since it is arbitrary which population is labelled~1, compute the tail probability on the side where $f_0$ falls and add the matching tail for $1/f_0$: the lower-tail probability of $f_0$ plus the upper-tail probability of $1/f_0$. (Textbook example: $n_1=n_2=20$, $s_1=1.96$, $s_2=2.13$, $f_0=0.85$; $\Prob(F_{19,19}<0.85)=0.3634$ and $\Prob(F_{19,19}>1.18)=0.3610$, so $P=0.7244$ and $H_0$ is not rejected.)
\end{example}

\subsection{$\beta$-Error and Choice of Sample Size}

If the true ratio is $\lambda^2=\sigma_1^2/\sigma_2^2$ then $F_0/\lambda^2\sim F_{n_1-1,n_2-1}$, so for the two-sided test
\[ \beta=\Prob\Big(\frac{f_{1-\alpha/2,n_1-1,n_2-1}}{\lambda^2}\le F\le\frac{f_{\alpha/2,n_1-1,n_2-1}}{\lambda^2}\Big),\qquad F\sim F_{n_1-1,n_2-1}. \]
OC curves (Appendix charts VIo--VIr, for $\alpha=0.05$ and $0.01$ with $n_1=n_2=n$; two-sided and one-sided versions) tabulate $\beta$ against the parameter $\lambda=\sigma_1/\sigma_2$. For example, to detect with probability at least $0.80$ that one standard deviation is half the other ($\lambda=2$) with $n_1=n_2=20$ at $\alpha=0.05$, the chart gives $\beta\approx0.20$, so power $0.80$ and $n=20$ is adequate.

\subsection{Confidence Interval on the Ratio of Two Variances}

\begin{theorem}[Interval for $\sigma_1^2/\sigma_2^2$]\label{thm:f-ci}
A $100(1-\alpha)\%$ confidence interval is
\[ \frac{s_1^2}{s_2^2}\,f_{1-\alpha/2,\,n_2-1,\,n_1-1}\ \le\ \frac{\sigma_1^2}{\sigma_2^2}\ \le\ \frac{s_1^2}{s_2^2}\,f_{\alpha/2,\,n_2-1,\,n_1-1}. \]

One-sided $100(1-\alpha)\%$ bounds use $f_\alpha$ alone: upper bound $\sigma_1^2/\sigma_2^2\le\dfrac{s_1^2}{s_2^2}f_{\alpha,n_2-1,n_1-1}$ and lower bound $\dfrac{s_1^2}{s_2^2}f_{1-\alpha,n_2-1,n_1-1}\le\sigma_1^2/\sigma_2^2$. They are not the ends of the two-sided interval above, which uses $\alpha/2$.
\end{theorem}

\begin{pitfall}{the $F$-test is not robust}
Like the $\chi^2$ test for a single variance, the $F$-test assumes normal populations and is badly affected by heavy tails: with non-normal data its true Type I error can be several times the nominal $\alpha$. Use Levene's or the Brown--Forsythe test (an ANOVA on absolute deviations from the group means or medians) when normality is doubtful.
\end{pitfall}

\section{Inference on Two Population Proportions}

\subsection{Large-Sample Test for $H_0:p_1=p_2$}

Let $X_1$ of $n_1$ trials and $X_2$ of $n_2$ trials be successes with sample proportions $\hat{P}_1=X_1/n_1$, $\hat{P}_2=X_2/n_2$. Then $\hat{P}_1-\hat{P}_2\approx\Normal\big(p_1-p_2,\ p_1(1-p_1)/n_1+p_2(1-p_2)/n_2\big)$. Under $H_0:p_1=p_2=p$ the common $p$ is estimated by the \emph{pooled proportion} $\hat{P}=(X_1+X_2)/(n_1+n_2)$.

\begin{theorem}[Two-proportion $z$-test]\label{thm:twoprop}
The statistic
\[ Z_0=\frac{\hat{P}_1-\hat{P}_2}{\sqrt{\hat{P}(1-\hat{P})\Big(\dfrac1{n_1}+\dfrac1{n_2}\Big)}}\ \approx\ \Normal(0,1)\quad\text{under }H_0:p_1=p_2. \]
The rejection regions are those of Theorem~\ref{thm:twoz}. The approximation needs $n_i\hat{p}_i\ge5$ and $n_i(1-\hat{p}_i)\ge5$.
\end{theorem}

\begin{example}\label{ex:twoprop}
Line~1: $15$ defectives in $300$ units ($\hat{p}_1=0.05$); line~2: $20$ in $250$ ($\hat{p}_2=0.08$). The pooled proportion is $35/550=0.0636$, the standard error is $\sqrt{0.0636\cdot0.9364\cdot(1/300+1/250)}=0.0209$, and $z_0=-0.03/0.0209=-1.435$. For the two-sided test $P=0.151$: no evidence of a difference at $\alpha=0.05$. (Note this is exactly the chi-square test of independence for the $2\times2$ table, with $z_0^2=\chi_0^2$.)
\end{example}

\subsection{Small Sample Test for $H_0:p_1=p_2$ (CD Only)}

When counts are small, use \emph{Fisher's exact test}. Given the table margins (total successes $s=x_1+x_2$ out of $N=n_1+n_2$), under $H_0$ the count $X_1$ has a hypergeometric distribution: $\Prob(X_1=k)=\dbinom{n_1}k\dbinom{n_2}{s-k}\Big/\dbinom Ns$. The $P$-value sums the probabilities of tables at least as extreme.

\begin{example}
Group 1: $5$ successes of $6$; group 2: $2$ of $6$. For $H_1:p_1>p_2$ with $N=12$, $s=7$: $P=\Prob(X_1\ge5)=\dfrac{\binom65\binom62+\binom66\binom61}{\binom{12}7}=\dfrac{90+6}{792}=0.121$, so there is no significant evidence at $0.05$.
\end{example}

\subsection{$\beta$-Error and Choice of Sample Size}

For the two-sided test with true proportions $p_1\ne p_2$ and $\bar{p}=(p_1+p_2)/2$, $\bar{q}=1-\bar{p}$, $\beta\approx\Phi\Big(\dfrac{z_{\alpha/2}\sqrt{\bar{p}\bar{q}(1/n_1+1/n_2)}-(p_1-p_2)}{\sigma_{\hat{P}_1-\hat{P}_2}}\Big)-\Phi\Big(\dfrac{-z_{\alpha/2}\sqrt{\bar{p}\bar{q}(1/n_1+1/n_2)}-(p_1-p_2)}{\sigma_{\hat{P}_1-\hat{P}_2}}\Big)$. For equal sample sizes, solving for $n$ gives
\[ n\approx\frac{\Big(z_{\alpha/2}\sqrt{(p_1+p_2)(q_1+q_2)/2}+z_\beta\sqrt{p_1q_1+p_2q_2}\Big)^2}{(p_1-p_2)^2}, \]
with $q_i=1-p_i$ (replace $z_{\alpha/2}$ by $z_\alpha$ for a one-sided test).

\subsection{Confidence Interval for $p_1-p_2$}

\begin{theorem}[Interval for a difference of proportions]\label{thm:twoprop-ci}
An approximate $100(1-\alpha)\%$ interval is
\[ (\hat{p}_1-\hat{p}_2)\pm z_{\alpha/2}\sqrt{\frac{\hat{p}_1(1-\hat{p}_1)}{n_1}+\frac{\hat{p}_2(1-\hat{p}_2)}{n_2}}. \]
\end{theorem}

\begin{example}
For Example~\ref{ex:twoprop}: $-0.03\pm1.96\sqrt{0.00016+0.00029}=-0.03\pm0.042$, i.e.\ $-0.072\le p_1-p_2\le0.012$. The interval uses the \emph{unpooled} standard error, unlike the test.
\end{example}

One-sided $100(1-\alpha)\%$ bounds use $z_\alpha$: $p_1-p_2\le(\hat{p}_1-\hat{p}_2)+z_\alpha\,\widehat{\se}$ (upper) and $(\hat{p}_1-\hat{p}_2)-z_\alpha\,\widehat{\se}\le p_1-p_2$ (lower), with the same unpooled $\widehat{\se}$. For Example~\ref{ex:twoprop} at $95\%$: upper bound $-0.03+1.645(0.0212)=0.005$.

\section{Summary Table for Inference Procedures for Two Samples}

\begin{center}
\small
\begin{tabular}{p{3.1cm}p{3.6cm}p{3.8cm}p{3.0cm}}
\toprule
Problem & Test statistic & Reject ($H_1$ two-sided) & Distribution under $H_0$ \\
\midrule
$\mu_1-\mu_2$, $\sigma$'s known & $z_0=\dfrac{\bar{x}_1-\bar{x}_2-\Delta_0}{\sqrt{\sigma_1^2/n_1+\sigma_2^2/n_2}}$ & $\abs{z_0}>z_{\alpha/2}$ & $\Normal(0,1)$ \\
$\mu_1-\mu_2$, $\sigma_1^2=\sigma_2^2$ unknown & $t_0=\dfrac{\bar{x}_1-\bar{x}_2-\Delta_0}{s_p\sqrt{1/n_1+1/n_2}}$ & $\abs{t_0}>t_{\alpha/2,n_1+n_2-2}$ & $t_{n_1+n_2-2}$ \\
$\mu_1-\mu_2$, $\sigma_1^2\ne\sigma_2^2$ & $t_0^*=\dfrac{\bar{x}_1-\bar{x}_2-\Delta_0}{\sqrt{s_1^2/n_1+s_2^2/n_2}}$ & $\abs{t_0^*}>t_{\alpha/2,\nu}$ & $\approx t_\nu$ \\
Paired data & $t_0=\dfrac{\bar{d}-\Delta_0}{s_d/\sqrt n}$ & $\abs{t_0}>t_{\alpha/2,n-1}$ & $t_{n-1}$ \\
$\sigma_1^2/\sigma_2^2$ & $f_0=s_1^2/s_2^2$ & $f_0>f_{\alpha/2}$ or $<f_{1-\alpha/2}$ & $F_{n_1-1,n_2-1}$ \\
$p_1-p_2$ & $z_0=\dfrac{\hat{p}_1-\hat{p}_2}{\sqrt{\hat{p}(1-\hat{p})(1/n_1+1/n_2)}}$ & $\abs{z_0}>z_{\alpha/2}$ & $\approx\Normal(0,1)$ \\
\bottomrule
\end{tabular}
\end{center}

\begin{summarybox}
\begin{itemize}
  \item Independent samples: $\bar{X}_1-\bar{X}_2$ has variance $\sigma_1^2/n_1+\sigma_2^2/n_2$; with $\sigma$'s known use $z$; with equal unknown variances pool ($s_p^2$, $n_1+n_2-2$ df); otherwise Welch's test with Satterthwaite degrees of freedom.
  \item Matched pairs: take differences and do a one-sample $t$-test with $n-1$ df; pairing helps when the within-pair correlation is positive.
  \item Randomisation is what supports cause-and-effect conclusions; observational comparisons can be confounded.
  \item $F=S_1^2/S_2^2\sim F_{n_1-1,n_2-1}$ compares variances; $f_{1-\alpha,u,v}=1/f_{\alpha,v,u}$; the test is sensitive to non-normality.
  \item Two proportions: pool under $H_0$ for the test, do not pool for the interval; Fisher's exact test for small counts.
  \item Always accompany a test with the confidence interval for the difference.
\end{itemize}
\end{summarybox}

\section{Exercises}

\begin{exercise}
Two machines fill cans with known standard deviations $\sigma_1=0.8$, $\sigma_2=1.1$. Samples of $n_1=30$, $n_2=36$ give means $355.4$ and $354.6$. Test $H_0:\mu_1=\mu_2$ against the two-sided alternative at $\alpha=0.05$, give the $P$-value and a $95\%$ confidence interval, and find the sample size per group to detect a difference of $0.5$ with power $0.9$.
\end{exercise}

\begin{exercise}
Two catalysts are compared. Catalyst~1: $n=8$, $\bar{x}=91.5$, $s=3.9$. Catalyst~2: $n=8$, $\bar{x}=89.2$, $s=4.1$. Assuming equal variances, test whether they differ ($\alpha=0.05$), compute a $95\%$ confidence interval for the difference, and check the equal-variance assumption with an $F$-test.
\end{exercise}

\begin{exercise}
Repeat the previous test with Welch's method for samples $n_1=6$, $s_1=2$ and $n_2=14$, $s_2=6$ with $\bar{x}_1-\bar{x}_2=5$. Compare the degrees of freedom with the pooled test and explain which test you trust.
\end{exercise}

\begin{exercise}
Ten runners' times (seconds) before and after a training program are: before $12.9,13.4,12.5,14.1,13.0,12.8,13.7,13.3,12.6,13.9$; after $12.5,13.1,12.6,13.6,12.7,12.4,13.2,13.1,12.2,13.5$. Test for improvement using a paired $t$-test and compare with the (incorrect) independent-samples $t$-test.
\end{exercise}

\begin{exercise}
Show that for paired data with common variance $\sigma^2$ and correlation $\rho$ within pairs, $\Var(\bar{D})=2\sigma^2(1-\rho)/n$. When does pairing improve on independent samples?
\end{exercise}

\begin{exercise}
In sample~1, $x_1=63$ of $n_1=200$ respond; in sample~2, $x_2=48$ of $n_2=180$. Test $H_0:p_1=p_2$, form a $95\%$ confidence interval for $p_1-p_2$, and verify that the squared $z$ statistic equals the chi-square statistic for the $2\times2$ table.
\end{exercise}

\begin{exercise}
Prove the identity $f_{1-\alpha,u,v}=1/f_{\alpha,v,u}$ and use it to find the lower critical point $f_{0.975,8,15}$ from $f_{0.025,15,8}=4.10$.
\end{exercise}

\begin{exercise}
How large must $n_1=n_2=n$ be to detect $p_1=0.10$ versus $p_2=0.15$ with power $0.8$ at $\alpha=0.05$ (two-sided)? Use the formula in this chapter.
\end{exercise}


\part{Regression and Correlation}
\chapter{Simple Linear Regression and Correlation}

\begin{nuance}{how to use this chapter}
This chapter follows Chapter~11 of the course textbook section by section, with the same titles and order. Regression builds an empirical model for how a response $Y$ changes with a regressor $x$, then uses the machinery of Chapters~2 and~3 (confidence intervals and tests) on the fitted coefficients. The worked example below is carried through every section so that all the formulas can be checked against one set of numbers. Numbering follows the textbook with the chapter number lowered by $6$: Section $5.2$ of this book is Section $11$-$2$ of the textbook, and Section $5.2.1$ is its Section $11$-$2.1$.
\end{nuance}

\section{Empirical Models}

Many problems involve the relationship between two or more variables: the strength of a weld and its temperature, the fuel use of a car and its weight. An \emph{empirical model} is an equation fitted to observed data that describes this relationship without being derived from first principles. Typically the response $Y$ is random and the regressor $x$ is a quantity whose value we choose or measure with negligible error.

\begin{definition}[Simple linear regression model]
The \emph{simple linear regression model} is
\[ Y=\beta_0+\beta_1x+\varepsilon, \]
where $\beta_0$ (the intercept) and $\beta_1$ (the slope) are unknown \emph{regression coefficients} and $\varepsilon$ is a random error with $\Exp\varepsilon=0$ and $\Var\varepsilon=\sigma^2$. For data $(x_i,y_i)$, $i=1,\dots,n$: $Y_i=\beta_0+\beta_1x_i+\varepsilon_i$ with the $\varepsilon_i$ independent. For inference we also assume the errors are normal, $\varepsilon_i\sim\Normal(0,\sigma^2)$.
\end{definition}

Then for fixed $x$, $\Exp(Y\mid x)=\beta_0+\beta_1x$ (the mean response lies on a straight line), $\Var(Y\mid x)=\sigma^2$ (constant variance), and $\beta_1$ is the change in the mean response per unit change in $x$.

\begin{example}\label{ex:regdata}
Throughout this chapter we use the following $n=8$ observations.
\begin{center}
\begin{tabular}{c|cccccccc}
\toprule
$x$ & 1 & 2 & 3 & 4 & 5 & 6 & 7 & 8\\
$y$ & 2.1 & 3.9 & 6.2 & 7.8 & 10.1 & 12.2 & 13.8 & 16.1\\
\bottomrule
\end{tabular}
\end{center}
\end{example}

\section{Simple Linear Regression}

\begin{figure}[h]
\centering
\begin{tikzpicture}[scale=0.55]
  \draw[axis] (0,0)--(9.4,0) node[right,lab]{$x$}; \draw[axis] (0,0)--(0,9.4) node[above,lab]{$y$};
  \draw[thick,accent2] (0.3,0.35)--(9,9.0);
  \foreach \x/\y in {1/1.05,2/1.95,3/3.1,4/3.9,5/5.05,6/6.1,7/6.9,8/8.05}{
    \node[pt] at (\x,\y){}; \draw[guide] (\x,\y)--(\x,{0.0357*0.5+\x*0.99881}); }
\end{tikzpicture}
\caption{The least squares line minimises the sum of squared vertical distances (dashed) from the data points to the line.}
\end{figure}

\subsection*{The method of least squares}

We estimate $\beta_0,\beta_1$ by choosing the line that makes the data as close as possible, in the sense of minimising the sum of squared vertical deviations
\[ L(\beta_0,\beta_1)=\sum_{i=1}^n\big(y_i-\beta_0-\beta_1x_i\big)^2. \]
Setting $\partial L/\partial\beta_0=0$ and $\partial L/\partial\beta_1=0$ gives the \emph{least squares normal equations}
\[ n\hat\beta_0+\hat\beta_1\sum x_i=\sum y_i,\qquad\hat\beta_0\sum x_i+\hat\beta_1\sum x_i^2=\sum x_iy_i. \]

\begin{definition}[Corrected sums of squares]
$S_{xx}=\sum(x_i-\bar{x})^2=\sum x_i^2-\dfrac{(\sum x_i)^2}n$, $\ S_{xy}=\sum(x_i-\bar{x})(y_i-\bar{y})=\sum x_iy_i-\dfrac{(\sum x_i)(\sum y_i)}n$, and $S_{yy}=\sum(y_i-\bar{y})^2$.
\end{definition}

\begin{theorem}[Least squares estimates]\label{thm:ls}
\[ \hat\beta_1=\frac{S_{xy}}{S_{xx}},\qquad\hat\beta_0=\bar{y}-\hat\beta_1\bar{x}. \]
The fitted line is $\hat{y}=\hat\beta_0+\hat\beta_1x$ and the \emph{residuals} are $e_i=y_i-\hat{y}_i$. The line passes through $(\bar{x},\bar{y})$.
\end{theorem}

\begin{proof}
The first normal equation gives $\hat\beta_0=\bar{y}-\hat\beta_1\bar{x}$. Substituting into the second: $\hat\beta_1\big(\sum x_i^2-n\bar{x}^2\big)=\sum x_iy_i-n\bar{x}\bar{y}$, i.e.\ $\hat\beta_1S_{xx}=S_{xy}$. The second derivatives show that this is a minimum ($L$ is a convex quadratic).
\end{proof}

\begin{example}[Continuing Example~\ref{ex:regdata}]\label{ex:regfit}
$\bar{x}=4.5$, $\bar{y}=9.025$, $S_{xx}=42$, $\sum xy=408.8$ so $S_{xy}=408.8-8(4.5)(9.025)=83.9$ and $S_{yy}=819.4-8(9.025)^2=167.795$. Hence $\hat\beta_1=83.9/42=1.9976$ and $\hat\beta_0=9.025-1.9976(4.5)=0.0357$. The fitted line is $\hat{y}=0.0357+1.9976\,x$: each unit increase in $x$ raises the mean response by about $2.0$.
\end{example}

\subsection*{Estimating $\sigma^2$}

The error variance is estimated from the residuals.

\begin{definition}[Error sum of squares, $\hat\sigma^2$]
$SS_E=\sum_{i=1}^ne_i^2=\sum(y_i-\hat{y}_i)^2$, and
\[ \hat\sigma^2=MS_E=\frac{SS_E}{n-2}. \]
A convenient computing formula is $SS_E=S_{yy}-\hat\beta_1S_{xy}$.
\end{definition}

The divisor $n-2$ is the degrees of freedom: two parameters ($\beta_0,\beta_1$) were estimated from the data, so two degrees of freedom are lost, and with it $\Exp[\hat\sigma^2]=\sigma^2$ (shown below). In Example~\ref{ex:regfit}: $SS_E=167.795-1.9976(83.9)=0.195$ and $\hat\sigma^2=0.195/6=0.0325$, so $\hat\sigma=0.180$.

\section{Properties of the Least Squares Estimators}

Write $\hat\beta_1=\sum c_iY_i$ with $c_i=(x_i-\bar{x})/S_{xx}$; the $x_i$ are fixed, so $\hat\beta_1$ is a linear combination of the independent $Y_i$.

\begin{theorem}[Sampling properties]\label{thm:lsprops}
Under the model assumptions:
\begin{enumerate}[label=(\roman*)]
  \item $\Exp\hat\beta_1=\beta_1$ and $\Var\hat\beta_1=\sigma^2/S_{xx}$;
  \item $\Exp\hat\beta_0=\beta_0$ and $\Var\hat\beta_0=\sigma^2\Big(\dfrac1n+\dfrac{\bar{x}^2}{S_{xx}}\Big)$;
  \item $\Exp\hat\sigma^2=\sigma^2$;
  \item if the errors are normal, $\hat\beta_0,\hat\beta_1$ are normal, and $(n-2)\hat\sigma^2/\sigma^2\sim\chi^2_{n-2}$ independently of $(\hat\beta_0,\hat\beta_1)$.
\end{enumerate}
The estimated standard errors are $\se(\hat\beta_1)=\sqrt{\hat\sigma^2/S_{xx}}$ and $\se(\hat\beta_0)=\sqrt{\hat\sigma^2(1/n+\bar{x}^2/S_{xx})}$.
\end{theorem}

\begin{proof}
(i) Since $\sum c_i=0$ and $\sum c_ix_i=1$: $\Exp\hat\beta_1=\sum c_i(\beta_0+\beta_1x_i)=\beta_1$, and $\Var\hat\beta_1=\sigma^2\sum c_i^2=\sigma^2/S_{xx}$. (ii) $\hat\beta_0=\bar{Y}-\hat\beta_1\bar{x}$; the terms $\bar{Y}$ and $\hat\beta_1$ are uncorrelated (because $\sum c_i=0$), so $\Var\hat\beta_0=\sigma^2/n+\bar{x}^2\sigma^2/S_{xx}$. (iii) and (iv) are proved in the matrix setting in Chapter~6 (Theorem~\ref{thm:mlr-dist}): the residual vector lies in an $(n-2)$-dimensional space.
\end{proof}

\begin{theorem}[Gauss--Markov]\label{thm:gm1}
Among all \emph{linear unbiased} estimators of $\beta_1$ (estimators of the form $\sum a_iY_i$), $\hat\beta_1$ has the smallest variance; likewise for $\hat\beta_0$. The least squares estimators are the \emph{best linear unbiased estimators} (BLUE), requiring only uncorrelated errors with equal variance, not normality.
\end{theorem}

\begin{interview}{what makes the slope precise?}
$\Var\hat\beta_1=\sigma^2/S_{xx}$ shrinks when the noise $\sigma$ is small, when $n$ is large, and when the $x$-values are \emph{spread out} (large $S_{xx}$). A designed experiment should place regressor values at the extremes of the range: half at each end maximises $S_{xx}$ for given $n$ and range (though it cannot detect curvature, so a few middle points are added). Extrapolating beyond the range of the data is dangerous, since the fitted line says nothing about whether the relationship stays linear there.
\end{interview}

\section{Some Comments on Uses of Regression (CD Only)}

Regression is used for \emph{description} (summarising a trend), \emph{prediction} (estimating the response at a new $x$ within the range of the data) and \emph{control} (choosing $x$ to achieve a target $y$). Cautions: (1) a fitted line is valid only in the region of the observed $x$, (2) a strong fit shows association, not causation, (3) the regression of $y$ on $x$ is not the regression of $x$ on $y$ (the two lines differ unless the correlation is $\pm1$), (4) the model assumes $x$ is measured without error, and (5) a few influential points can control the whole fit.

\section{Hypothesis Tests in Simple Linear Regression}

\subsection{Use of $t$-Tests}

With normal errors, standardising each coefficient with its estimated standard error gives a $t$ statistic.

\begin{theorem}[$t$-tests on the coefficients]\label{thm:regt}
$T=\dfrac{\hat\beta_1-\beta_{1,0}}{\se(\hat\beta_1)}\sim t_{n-2}$ when $\beta_1=\beta_{1,0}$, and similarly $T=\dfrac{\hat\beta_0-\beta_{0,0}}{\se(\hat\beta_0)}\sim t_{n-2}$. To test $H_0:\beta_1=\beta_{1,0}$ against $H_1:\beta_1\ne\beta_{1,0}$, reject if $\abs{t_0}>t_{\alpha/2,n-2}$.
\end{theorem}

The most important special case is $H_0:\beta_1=0$ against $H_1:\beta_1\ne0$, the \emph{test for significance of regression}. Failing to reject means there is no evidence of a linear relationship between $x$ and $Y$ (either $x$ is of little value in explaining $Y$, or the true relationship is not linear); rejecting means $x$ is useful in explaining $Y$ (the line is a better model than a constant, though a curve might be better still).

\begin{example}
For Example~\ref{ex:regfit}: $\se(\hat\beta_1)=\sqrt{0.0325/42}=0.0278$ and $t_0=1.9976/0.0278=71.8$ with $6$ degrees of freedom, far beyond $t_{0.005,6}=3.707$: reject $H_0:\beta_1=0$ ($P<0.0001$). For the intercept, $\se(\hat\beta_0)=\sqrt{0.0325(1/8+20.25/42)}=0.1405$ and $t_0=0.0357/0.1405=0.25$: no evidence that the intercept differs from $0$.
\end{example}

\subsection{Analysis of Variance Approach to Test Significance of Regression}

The total variability of the responses splits into a part explained by the regression and a part left over.

\begin{theorem}[Analysis of variance identity]\label{thm:anova-reg}
\[ \underbrace{\sum(y_i-\bar{y})^2}_{SS_T=S_{yy}}=\underbrace{\sum(\hat{y}_i-\bar{y})^2}_{SS_R=\hat\beta_1S_{xy}}+\underbrace{\sum(y_i-\hat{y}_i)^2}_{SS_E}, \]
with degrees of freedom $n-1=1+(n-2)$.
\end{theorem}

\begin{proof}
Write $y_i-\bar{y}=(\hat{y}_i-\bar{y})+(y_i-\hat{y}_i)$ and square; the cross term is $2\sum(\hat{y}_i-\bar{y})e_i=2\hat\beta_1\sum(x_i-\bar{x})e_i=0$, because the normal equations make the residuals sum to zero and orthogonal to the $x_i$.
\end{proof}

The ratio of mean squares gives an $F$ test equivalent to the $t$-test.

\begin{theorem}[$F$-test for significance of regression]\label{thm:regF}
Let $MS_R=SS_R/1$ and $MS_E=SS_E/(n-2)$. Under $H_0:\beta_1=0$, $F_0=MS_R/MS_E\sim F_{1,n-2}$, and $H_0$ is rejected if $f_0>f_{\alpha,1,n-2}$. Moreover $F_0=T_0^2$, where $T_0=\hat\beta_1/\se(\hat\beta_1)$.
\end{theorem}

\begin{center}
\begin{tabular}{lcccc}
\toprule
Source & Sum of squares & df & Mean square & $F_0$ \\
\midrule
Regression & $SS_R$ & $1$ & $MS_R$ & $MS_R/MS_E$ \\
Error & $SS_E$ & $n-2$ & $MS_E=\hat\sigma^2$ & \\
Total & $SS_T$ & $n-1$ & & \\
\bottomrule
\end{tabular}
\end{center}

\begin{example}
$SS_T=167.795$, $SS_R=1.9976(83.9)=167.600$, $SS_E=0.195$; $f_0=167.600/0.0325=5157$, which equals $t_0^2=71.8^2$ and exceeds $f_{0.05,1,6}=5.99$.
\end{example}

\section{Confidence Intervals}

\subsection{Confidence Intervals on the Slope and Intercept}

\begin{theorem}[Intervals for the coefficients]\label{thm:regci}
A $100(1-\alpha)\%$ confidence interval on $\beta_1$ is $\hat\beta_1\pm t_{\alpha/2,n-2}\sqrt{\hat\sigma^2/S_{xx}}$, and on $\beta_0$ is $\hat\beta_0\pm t_{\alpha/2,n-2}\sqrt{\hat\sigma^2(1/n+\bar{x}^2/S_{xx})}$. A confidence interval on $\sigma^2$ is $\dfrac{(n-2)\hat\sigma^2}{\chi^2_{\alpha/2,n-2}}\le\sigma^2\le\dfrac{(n-2)\hat\sigma^2}{\chi^2_{1-\alpha/2,n-2}}$.
\end{theorem}

\begin{example}
For the slope: $1.9976\pm2.447(0.0278)=1.9976\pm0.068$, i.e.\ $1.930\le\beta_1\le2.066$.
\end{example}

\subsection{Confidence Interval on the Mean Response}

The \emph{mean response} at $x=x_0$ is $\mu_{Y|x_0}=\beta_0+\beta_1x_0$, estimated by $\hat\mu_{Y|x_0}=\hat\beta_0+\hat\beta_1x_0$, which is normal and unbiased with variance $\sigma^2\Big[\dfrac1n+\dfrac{(x_0-\bar{x})^2}{S_{xx}}\Big]$.

\begin{theorem}[Interval on the mean response]\label{thm:meanresp}
A $100(1-\alpha)\%$ confidence interval on $\mu_{Y|x_0}$ is
\[ \hat\mu_{Y|x_0}\pm t_{\alpha/2,n-2}\sqrt{\hat\sigma^2\Big[\frac1n+\frac{(x_0-\bar{x})^2}{S_{xx}}\Big]}. \]
\end{theorem}

The interval is narrowest at $x_0=\bar{x}$ and widens as $x_0$ moves away from the centre of the data: our knowledge of the line is best in the middle.

\begin{example}
At $x_0=5$: $\hat\mu=0.0357+1.9976(5)=10.024$; the standard error is $\sqrt{0.0325(1/8+0.25/42)}=0.0652$, so the $95\%$ interval is $10.024\pm2.447(0.0652)=10.024\pm0.160$.
\end{example}

\section{Prediction of New Observations}

A new individual observation $Y_0$ at $x_0$ equals the mean response plus a new error $\varepsilon_0$ independent of the data.

\begin{theorem}[Prediction interval]\label{thm:regpi}
A $100(1-\alpha)\%$ prediction interval for a future observation $Y_0$ at $x_0$ is
\[ \hat{y}_0\pm t_{\alpha/2,n-2}\sqrt{\hat\sigma^2\Big[1+\frac1n+\frac{(x_0-\bar{x})^2}{S_{xx}}\Big]}. \]
\end{theorem}

\begin{proof}
$Y_0-\hat{Y}_0$ has mean $0$ and variance $\sigma^2+\sigma^2[1/n+(x_0-\bar{x})^2/S_{xx}]$ (the two parts are independent), and standardising with $\hat\sigma$ gives a $t_{n-2}$ pivot.
\end{proof}

\begin{example}
At $x_0=5$: $\sqrt{0.0325(1+0.1310)}=0.1917$, so the prediction interval is $10.024\pm2.447(0.1917)=10.024\pm0.469$, much wider than the interval $\pm0.160$ for the mean response: even with the line known exactly, a single new observation still scatters by $\sigma$.
\end{example}

\section{Adequacy of the Regression Model}

The tests and intervals are valid only if the model assumptions hold: linearity, independent errors with constant variance and (for small $n$) normality.

\subsection{Residual Analysis}

\begin{definition}[Residuals]
$e_i=y_i-\hat{y}_i$. The \emph{standardised residuals} are $d_i=e_i/\hat\sigma$; about $95\%$ of them should lie within $\pm2$ if the model is correct.
\end{definition}

Graphical checks:
\begin{itemize}
  \item \textbf{Normal probability plot of residuals:} points near a straight line support normality; curvature indicates skew or heavy tails.
  \item \textbf{Residuals versus fitted values (or versus $x$):} a random horizontal band is good; a curved pattern means the straight line is the wrong shape (try a quadratic term or a transformation); a funnel (spread increasing with $\hat{y}$) indicates non-constant variance (try transforming $y$ or weighted least squares).
  \item \textbf{Residuals versus time order:} runs of same-sign residuals indicate correlated errors.
\end{itemize}

\subsection{Coefficient of Determination ($R^2$)}\label{sec:regr2}

\begin{definition}[$R^2$]
$R^2=\dfrac{SS_R}{SS_T}=1-\dfrac{SS_E}{SS_T}$, the proportion of the variability in $y$ explained by the regression, with $0\le R^2\le1$.
\end{definition}

For the data above $R^2=167.600/167.795=0.9988$. Cautions: a large $R^2$ does not prove that the linear model is correct (a curved relationship over a limited range can give $R^2$ near $1$), and a small $R^2$ is not necessarily a bad model if the noise is naturally large. $R^2$ never decreases when regressors are added, so it cannot be used to compare models with different numbers of regressors (see Chapter~6).

\subsection{Lack-of-Fit Test (CD Only)}

If there are \emph{repeated observations} at some values of $x$ (replicates), the error sum of squares can be split into pure error (variation among replicates) and lack of fit:
\[ SS_E=SS_{PE}+SS_{LOF},\qquad SS_{PE}=\sum_i\sum_j(y_{ij}-\bar{y}_i)^2\ (n-m\text{ df}),\quad SS_{LOF}=SS_E-SS_{PE}\ (m-2\text{ df}), \]
where $m$ is the number of distinct $x$-values. Under the hypothesis that the straight line is adequate, $F_0=\dfrac{SS_{LOF}/(m-2)}{SS_{PE}/(n-m)}\sim F_{m-2,n-m}$; a large $F_0$ is evidence of curvature.

\section{Transformations to a Straight Line}

If a scatter plot is clearly curved, a transformation of $x$, $y$ or both may make the relationship linear, after which the methods above apply.

\begin{center}
\begin{tabular}{llll}
\toprule
Nonlinear model & Transformation & Linear form & Back-transform \\
\midrule
$y=\beta_0e^{\beta_1x}$ & $y^*=\ln y$ & $y^*=\ln\beta_0+\beta_1x$ & $\beta_0=e^{\text{intercept}}$ \\
$y=\beta_0x^{\beta_1}$ & $y^*=\ln y,\ x^*=\ln x$ & $y^*=\ln\beta_0+\beta_1x^*$ & $\beta_0=e^{\text{intercept}}$ \\
$y=\beta_0+\beta_1\ln x$ & $x^*=\ln x$ & $y=\beta_0+\beta_1x^*$ & --- \\
$y=1/(\beta_0+\beta_1x)$ & $y^*=1/y$ & $y^*=\beta_0+\beta_1x$ & --- \\
$y=\dfrac x{\beta_0x-\beta_1}$ & $y^*=1/y,\ x^*=1/x$ & $y^*=\beta_0-\beta_1x^*$ & --- \\
\bottomrule
\end{tabular}
\end{center}

\section{More About Transformations (CD Only)}

Two cautions. First, a transformation changes the error structure along with the mean: if $y=\beta_0e^{\beta_1x}\varepsilon$ with a multiplicative error $\varepsilon$, then $\ln y$ has an additive error and the log-linear fit is correct; if the original error is additive the log-transformed errors are not constant-variance. Always re-check the residuals \emph{after} transforming. Second, predictions made on the transformed scale (a prediction of $\ln y$) must be back-transformed, and the back-transformed mean of $\ln y$ estimates the \emph{median} of $y$, not its mean. For stabilising variance a common choice is $\sqrt y$ for counts and $\ln y$ when the standard deviation is proportional to the mean; the Box--Cox family $y^{(\lambda)}=(y^\lambda-1)/\lambda$ chooses $\lambda$ by maximum likelihood.

\section{Correlation}

So far $x$ was a fixed regressor. Suppose instead that $(X,Y)$ are \emph{both} random, with a bivariate normal distribution, and we are interested in their strength of linear association.

\begin{definition}[Correlation coefficient]
The population correlation is $\rho=\dfrac{\Cov(X,Y)}{\sigma_X\sigma_Y}\in[-1,1]$. For a sample the \emph{sample correlation coefficient} is
\[ R=\frac{S_{xy}}{\sqrt{S_{xx}S_{yy}}}=\hat\beta_1\sqrt{\frac{S_{xx}}{S_{yy}}}. \]
\end{definition}

Relations: $R^2$ (the square of $R$) equals the coefficient of determination of Section~\ref{sec:regr2}; $R$ has the same sign as $\hat\beta_1$; $R=\pm1$ if and only if all points lie exactly on a line; and $\hat\beta_1=R\,s_y/s_x$.

\begin{theorem}[Test and interval for $\rho$]\label{thm:corr}
If $(X,Y)$ is bivariate normal:
\begin{enumerate}[label=(\roman*)]
  \item to test $H_0:\rho=0$ use $T_0=\dfrac{R\sqrt{n-2}}{\sqrt{1-R^2}}\sim t_{n-2}$ under $H_0$ (identical to the $t$-test of $H_0:\beta_1=0$);
  \item for other hypotheses and intervals use Fisher's transformation: $Z=\tfrac12\ln\dfrac{1+R}{1-R}=\operatorname{arctanh}R$ is approximately $\Normal\Big(\operatorname{arctanh}\rho,\ \dfrac1{n-3}\Big)$. A $100(1-\alpha)\%$ interval for $\rho$ is $\tanh\big(\operatorname{arctanh}r\pm z_{\alpha/2}/\sqrt{n-3}\big)$.
\end{enumerate}
\end{theorem}

\begin{example}
For the example data $r=\sqrt{0.9988}=0.9994$; $t_0=0.9994\sqrt6/\sqrt{1-0.9988}=70.7$ (equals the regression $t$ up to rounding). For an interval with $n=20$ and $r=0.6$: $\operatorname{arctanh}0.6=0.693$, so $0.693\pm1.96/\sqrt{17}=0.693\pm0.475$, and $\rho\in(\tanh0.218,\tanh1.168)=(0.21,0.82)$.
\end{example}

\begin{pitfall}{correlation is not causation, and zero correlation is not independence}
A high correlation can come from a common cause or pure coincidence. Correlation measures \emph{linear} association only: $Y=X^2$ with $X$ symmetric about zero has $\rho=0$ yet $Y$ is a deterministic function of $X$. And a single outlier can create or destroy a large correlation; always look at the scatter plot (Anscombe's quartet has four very different data sets with identical $r=0.816$ and identical regression lines).
\end{pitfall}


\begin{summarybox}
\begin{itemize}
  \item Model $Y=\beta_0+\beta_1x+\varepsilon$, $\varepsilon\sim\Normal(0,\sigma^2)$. Least squares: $\hat\beta_1=S_{xy}/S_{xx}$, $\hat\beta_0=\bar{y}-\hat\beta_1\bar{x}$, $\hat\sigma^2=SS_E/(n-2)$ with $SS_E=S_{yy}-\hat\beta_1S_{xy}$.
  \item $\Var\hat\beta_1=\sigma^2/S_{xx}$, $\Var\hat\beta_0=\sigma^2(1/n+\bar{x}^2/S_{xx})$; $\hat\beta$'s are BLUE (Gauss--Markov).
  \item Tests and intervals use $t_{n-2}$; $H_0:\beta_1=0$ is also tested by $F_0=MS_R/MS_E=t_0^2$ with $SS_T=SS_R+SS_E$.
  \item Mean response: $\hat{y}_0\pm t\sqrt{\hat\sigma^2[1/n+(x_0-\bar{x})^2/S_{xx}]}$; new observation: add $1$ inside the bracket; both are narrowest at $\bar{x}$.
  \item Check the residuals (normal plot, versus $\hat{y}$); $R^2=SS_R/SS_T$; lack of fit needs replicates; transformations linearise some nonlinear models.
  \item $r=S_{xy}/\sqrt{S_{xx}S_{yy}}$, $t_0=r\sqrt{n-2}/\sqrt{1-r^2}$; Fisher's $z$ gives intervals for $\rho$.
\end{itemize}
\end{summarybox}

\section{Exercises}

\begin{exercise}
For the data $x$: $2,4,5,7,9,10$ and $y$: $5.1,8.3,9.8,14.2,18.0,19.6$, find the least squares line, $\hat\sigma^2$, the $95\%$ confidence interval for the slope, and test $H_0:\beta_1=2$ against $\beta_1\ne2$.
\end{exercise}

\begin{exercise}
Construct the analysis of variance table for the data of Exercise~1 and test the significance of regression with $\alpha=0.05$. Verify $F_0=t_0^2$ and compute $R^2$.
\end{exercise}

\begin{exercise}
Using Exercise~1, find a $95\%$ confidence interval on the mean response at $x_0=6$ and a $95\%$ prediction interval for a new observation at $x_0=6$. Why is the second wider? What happens at $x_0=20$?
\end{exercise}

\begin{exercise}
Prove that the residuals of a least squares line with an intercept satisfy $\sum e_i=0$ and $\sum x_ie_i=0$, and deduce that the fitted line passes through $(\bar{x},\bar{y})$.
\end{exercise}

\begin{exercise}
Show that $\hat\beta_1=R\,s_y/s_x$ and $R^2=SS_R/SS_T$. If $R=0.8$ and $s_y=10$, $s_x=2$, find $\hat\beta_1$ and the proportion of variance explained.
\end{exercise}

\begin{exercise}
The data below follow an approximate relationship $y=\beta_0e^{\beta_1x}$: $x=1,2,3,4,5$; $y=2.7,7.1,20.5,54.9,148.0$. Linearise, fit by least squares and estimate $\beta_0$ and $\beta_1$.
\end{exercise}

\begin{exercise}
In a sample of $n=27$ pairs $r=-0.45$. Test $H_0:\rho=0$ and find a $95\%$ confidence interval for $\rho$ using Fisher's transformation.
\end{exercise}

\begin{exercise}
Explain, with the help of a residual plot sketch, how to recognise (a) a nonlinear relationship, (b) non-constant variance and (c) a single outlier, and what you would do about each.
\end{exercise}

\chapter{Multiple Linear Regression}

\section{Multiple Linear Regression Model}

\subsection{Introduction}

\begin{definition}[Multiple linear regression model]
With $k$ regressors $x_1,\dots,x_k$,
\[ Y=\beta_0+\beta_1x_1+\beta_2x_2+\dots+\beta_kx_k+\varepsilon,\qquad\varepsilon\sim\Normal(0,\sigma^2). \]
The model is \emph{linear} because it is linear in the parameters $\beta_j$; the regressors themselves may be nonlinear functions (squares, products, logarithms). The coefficient $\beta_j$ is the change in the mean response per unit change in $x_j$ \emph{with the other regressors held fixed}.
\end{definition}

Models that look more complicated are still multiple linear regression models. A cubic in one regressor, $Y=\beta_0+\beta_1x+\beta_2x^2+\beta_3x^3+\varepsilon$, becomes linear on setting $x_1=x$, $x_2=x^2$, $x_3=x^3$. An interaction $Y=\beta_0+\beta_1x_1+\beta_2x_2+\beta_{12}x_1x_2+\varepsilon$ is linear with $x_3=x_1x_2$, $\beta_3=\beta_{12}$. The second-order model with interaction, $Y=\beta_0+\beta_1x_1+\beta_2x_2+\beta_{11}x_1^2+\beta_{22}x_2^2+\beta_{12}x_1x_2+\varepsilon$, is likewise linear with five regressors; its surface can be a mound, a saddle or a ridge depending on the coefficients. Any model linear in the $\beta$'s is a linear regression model, \emph{regardless of the shape of the surface} it generates. \emph{Interaction} means that the effect of changing $x_1$ depends on the level of $x_2$: in the model with $\beta_{12}$, $\partial\,\Exp(Y)/\partial x_1=\beta_1+\beta_{12}x_2$.

\begin{example}[Textbook example: wire bond pull strength]
For $25$ wire bonds, regress pull strength $y$ on wire length $x_1$ and die height $x_2$. The normal equations (with the sums from the data table) are
\[ 25\hat\beta_0+206\hat\beta_1+8294\hat\beta_2=725.82,\quad 206\hat\beta_0+2396\hat\beta_1+77{,}177\hat\beta_2=8008.37,\quad 8294\hat\beta_0+77{,}177\hat\beta_1+3{,}531{,}848\hat\beta_2=274{,}811.31, \]
with solution $\hat\beta_0=2.26379$, $\hat\beta_1=2.74427$, $\hat\beta_2=0.01253$. The matrix route $\hat{\boldsymbol\beta}=(X\T X)^{-1}X\T\mathbf{y}$ gives the same values. The estimated error variance is $\hat\sigma^2=5.2364$ and the standard errors are $\se(\hat\beta_0)=1.060$, $\se(\hat\beta_1)=0.09352$, $\se(\hat\beta_2)=0.002798$ (precisely estimated slopes; the intercept is relatively less precise).
\end{example}

\begin{example}[The running example]\label{ex:mlrdata}
Eight observations on a response $y$ with two coded regressors $x_1,x_2\in\{-1,+1\}$:
\begin{center}
\begin{tabular}{c|cccccccc}
\toprule
$x_1$ & $-1$ & $-1$ & $-1$ & $-1$ & $+1$ & $+1$ & $+1$ & $+1$\\
$x_2$ & $-1$ & $-1$ & $+1$ & $+1$ & $-1$ & $-1$ & $+1$ & $+1$\\
$y$ & 12.1 & 13.0 & 15.2 & 16.1 & 17.9 & 19.2 & 21.0 & 22.3\\
\bottomrule
\end{tabular}
\end{center}
\end{example}

\subsection{Least Squares Estimation of the Parameters}

Given $n$ observations $(y_i,x_{i1},\dots,x_{ik})$, the least squares estimates minimise
\[ L=\sum_{i=1}^n\Big(y_i-\beta_0-\sum_{j=1}^k\beta_jx_{ij}\Big)^2. \]
Setting each partial derivative to zero gives $p=k+1$ \emph{normal equations} in the $p$ unknowns. In matrix form they are very compact.

\subsection{Matrix Approach to Multiple Linear Regression}

Write
\[ \mathbf{y}=\begin{pmatrix}y_1\\\vdots\\y_n\end{pmatrix},\quad X=\begin{pmatrix}1&x_{11}&\cdots&x_{1k}\\\vdots&\vdots&&\vdots\\1&x_{n1}&\cdots&x_{nk}\end{pmatrix},\quad\boldsymbol\beta=\begin{pmatrix}\beta_0\\\vdots\\\beta_k\end{pmatrix},\quad\boldsymbol\varepsilon=\begin{pmatrix}\varepsilon_1\\\vdots\\\varepsilon_n\end{pmatrix}, \]
so that $\mathbf{y}=X\boldsymbol\beta+\boldsymbol\varepsilon$, with $X$ the $n\times p$ \emph{design matrix}.

\begin{theorem}[Least squares estimator]\label{thm:mlr-ls}
If $X\T X$ is invertible (the columns of $X$ are linearly independent),
\[ \hat{\boldsymbol\beta}=(X\T X)^{-1}X\T\mathbf{y}, \]
the unique solution of the normal equations $X\T X\hat{\boldsymbol\beta}=X\T\mathbf{y}$. The fitted values are $\hat{\mathbf{y}}=X\hat{\boldsymbol\beta}=H\mathbf{y}$ with the \emph{hat matrix} $H=X(X\T X)^{-1}X\T$, and the residuals are $\mathbf{e}=\mathbf{y}-\hat{\mathbf{y}}=(I-H)\mathbf{y}$.
\end{theorem}

\begin{proof}
$L=(\mathbf{y}-X\boldsymbol\beta)\T(\mathbf{y}-X\boldsymbol\beta)$ has gradient $-2X\T(\mathbf{y}-X\boldsymbol\beta)$, which vanishes at the stated $\hat{\boldsymbol\beta}$, and the Hessian $2X\T X$ is positive definite, so this is the unique minimum.
\end{proof}

\begin{theorem}[Properties of $H$ (geometry)]\label{thm:hat}
$H$ is symmetric and idempotent ($H^2=H$) with $\tr H=p$; it is the orthogonal projection of $\R^n$ onto the column space of $X$. Consequently $X\T\mathbf{e}=\mathbf{0}$: the residual vector is orthogonal to every column of $X$, and with an intercept column, $\sum e_i=0$.
\end{theorem}

\begin{definition}[Error variance estimate]
$SS_E=\mathbf{e}\T\mathbf{e}=\mathbf{y}\T\mathbf{y}-\hat{\boldsymbol\beta}\T X\T\mathbf{y}$ and $\hat\sigma^2=MS_E=\dfrac{SS_E}{n-p}$, with $p=k+1$ parameters.
\end{definition}

\begin{example}[Running example]\label{ex:mlrfit}
Here $X\T X=\diag(8,8,8)$, since the columns $\mathbf{1}$, $x_1$, $x_2$ are mutually orthogonal, so $(X\T X)^{-1}=\tfrac18I$ and each coefficient is a simple average: $\hat\beta_0=\bar{y}=17.1$, $\hat\beta_1=\tfrac18\sum x_{i1}y_i=24.0/8=3.0$, $\hat\beta_2=\tfrac18\sum x_{i2}y_i=12.4/8=1.55$. The fitted equation is $\hat{y}=17.1+3.0x_1+1.55x_2$. The fitted values are $12.55,\,12.55,\,15.65,\,15.65,\,18.55,\,18.55,\,21.65,\,21.65$, the residuals are $-0.45,\,0.45,\,-0.45,\,0.45,\,-0.65,\,0.65,\,-0.65,\,0.65$, $SS_E=4(0.2025)+4(0.4225)=2.5$ and $\hat\sigma^2=2.5/(8-3)=0.5$.
\end{example}

\subsection{Properties of the Least Squares Estimators}

\begin{theorem}[Sampling properties]\label{thm:mlr-dist}
Under the model with $\Exp[\boldsymbol\varepsilon]=\mathbf{0}$ and $\Cov[\boldsymbol\varepsilon]=\sigma^2I$:
\begin{enumerate}[label=(\roman*)]
  \item $\Exp[\hat{\boldsymbol\beta}]=\boldsymbol\beta$ and $\Cov[\hat{\boldsymbol\beta}]=\sigma^2(X\T X)^{-1}$; write $C_{jj}$ for the $j$-th diagonal entry of $(X\T X)^{-1}$, so $\Var[\hat\beta_j]=\sigma^2C_{jj}$ and $\se(\hat\beta_j)=\sqrt{\hat\sigma^2C_{jj}}$;
  \item $\Exp[\hat\sigma^2]=\sigma^2$;
  \item (Gauss--Markov) among all linear unbiased estimators, $\hat\beta_j$ has minimum variance;
  \item if $\boldsymbol\varepsilon\sim\Normal(\mathbf{0},\sigma^2I)$, then $\hat{\boldsymbol\beta}\sim\Normal\big(\boldsymbol\beta,\sigma^2(X\T X)^{-1}\big)$, $\ (n-p)\hat\sigma^2/\sigma^2\sim\chi^2_{n-p}$, and $\hat{\boldsymbol\beta}$ is independent of $\hat\sigma^2$.
\end{enumerate}
\end{theorem}

\begin{proof}
(i) $\hat{\boldsymbol\beta}=(X\T X)^{-1}X\T(X\boldsymbol\beta+\boldsymbol\varepsilon)=\boldsymbol\beta+(X\T X)^{-1}X\T\boldsymbol\varepsilon$, so $\Exp[\hat{\boldsymbol\beta}]=\boldsymbol\beta$ and $\Cov[\hat{\boldsymbol\beta}]=(X\T X)^{-1}X\T\,\sigma^2I\,X(X\T X)^{-1}=\sigma^2(X\T X)^{-1}$. (ii) $\mathbf{e}=(I-H)\boldsymbol\varepsilon$ because $(I-H)X=0$; then $\Exp(\mathbf{e}\T\mathbf{e})=\Exp[\boldsymbol\varepsilon\T(I-H)\boldsymbol\varepsilon]=\sigma^2\tr(I-H)=\sigma^2(n-p)$. (iv) $(I-H)$ is the projection onto a subspace of dimension $n-p$, so with orthonormal coordinates adapted to that subspace $\boldsymbol\varepsilon\T(I-H)\boldsymbol\varepsilon/\sigma^2$ is a sum of $n-p$ squared independent standard normals; independence of $\hat{\boldsymbol\beta}$ and $\mathbf{e}$ follows because $\hat{\boldsymbol\beta}-\boldsymbol\beta=(X\T X)^{-1}X\T\boldsymbol\varepsilon$ and $\mathbf{e}=(I-H)\boldsymbol\varepsilon$ are jointly normal and uncorrelated ($X\T(I-H)=0$). (iii) Any other linear unbiased estimator $A\mathbf{y}$ has $AX=I$; write $A=(X\T X)^{-1}X\T+D$ with $DX=0$, so $\Cov(A\mathbf{y})=\sigma^2(X\T X)^{-1}+\sigma^2DD\T\succeq\sigma^2(X\T X)^{-1}$.
\end{proof}

\begin{example}
In the running example $C_{jj}=1/8$ for each coefficient, so $\se(\hat\beta_j)=\sqrt{0.5/8}=0.25$ for $j=0,1,2$, and $\hat\beta_1,\hat\beta_2$ are uncorrelated (the off-diagonal entries of $(X\T X)^{-1}$ are $0$). The hat matrix diagonal is $h_{ii}=\tfrac38$ for every observation, which sums to $p=3$.
\end{example}

\section{Hypothesis Tests in Multiple Linear Regression}

\subsection{Test for Significance of Regression}

\begin{definition}[Sums of squares]
$SS_T=\sum(y_i-\bar{y})^2=\mathbf{y}\T\mathbf{y}-n\bar{y}^2$, $\ SS_R=\hat{\boldsymbol\beta}\T X\T\mathbf{y}-n\bar{y}^2$ and $SS_E=SS_T-SS_R$, with degrees of freedom $n-1=k+(n-p)$.
\end{definition}

\begin{theorem}[$F$-test for significance of regression]\label{thm:mlr-F}
To test $H_0:\beta_1=\beta_2=\dots=\beta_k=0$ against $H_1:$ at least one $\beta_j\ne0$, use
\[ F_0=\frac{SS_R/k}{SS_E/(n-p)}=\frac{MS_R}{MS_E}\sim F_{k,\,n-p}\ \text{ under }H_0. \]
Reject $H_0$ if $f_0>f_{\alpha,k,n-p}$.
\end{theorem}

\begin{definition}[$R^2$ and adjusted $R^2$]
$R^2=SS_R/SS_T$ and
\[ R^2_{\mathrm{adj}}=1-\frac{SS_E/(n-p)}{SS_T/(n-1)}. \]
$R^2$ never decreases when a regressor is added; $R^2_{\mathrm{adj}}$ penalises extra regressors (it can decrease) and is preferred for comparing models of different size.
\end{definition}

\begin{example}
For the running example: $SS_R=8(3.0)^2+8(1.55)^2=72+19.22=91.22$, $SS_E=2.5$, $SS_T=93.72$. $f_0=(91.22/2)/0.5=91.2>f_{0.05,2,5}=5.79$: reject $H_0$; at least one regressor is useful. $R^2=0.973$ and $R^2_{\mathrm{adj}}=1-(0.5)/(93.72/7)=0.963$.
\end{example}

\subsection{Tests on Individual Regression Coefficients and Subsets of Coefficients}

\begin{theorem}[$t$-test on a single coefficient]\label{thm:mlr-t}
To test $H_0:\beta_j=0$ against $H_1:\beta_j\ne0$ use $T_0=\hat\beta_j/\se(\hat\beta_j)\sim t_{n-p}$ under $H_0$, and reject if $\abs{t_0}>t_{\alpha/2,n-p}$. For the one-sided alternatives, reject $H_1:\beta_j>0$ if $t_0>t_{\alpha,n-p}$ and $H_1:\beta_j<0$ if $t_0<-t_{\alpha,n-p}$. This is a \emph{marginal} test: it asks whether $x_j$ adds anything given that all the other regressors are already in the model.
\end{theorem}

\begin{example}
$t_0=3.0/0.25=12.0$ for $\beta_1$ and $1.55/0.25=6.2$ for $\beta_2$, both exceeding $t_{0.025,5}=2.571$: both regressors are significant.
\end{example}

A \emph{subset} of coefficients is tested by the \emph{extra sum of squares} method. Partition $\boldsymbol\beta=(\boldsymbol\beta_1,\boldsymbol\beta_2)$, with $\boldsymbol\beta_2$ containing $r$ coefficients, and test $H_0:\boldsymbol\beta_2=\mathbf{0}$ against $H_1:\boldsymbol\beta_2\ne\mathbf{0}$. Fit the \emph{full model} and the \emph{reduced model} (without $\boldsymbol\beta_2$); the improvement in the regression sum of squares is
\[ SS_R(\boldsymbol\beta_2\mid\boldsymbol\beta_1)=SS_R(\text{full})-SS_R(\text{reduced})=SS_E(\text{reduced})-SS_E(\text{full}). \]

\begin{theorem}[Partial $F$-test]\label{thm:partialF}
Under $H_0:\boldsymbol\beta_2=\mathbf{0}$,
\[ F_0=\frac{SS_R(\boldsymbol\beta_2\mid\boldsymbol\beta_1)/r}{MS_E(\text{full})}\sim F_{r,\,n-p}, \]
and $H_0$ is rejected if $f_0>f_{\alpha,r,n-p}$. When $r=1$, $F_0=T_0^2$ for the $t$-test of that coefficient.
\end{theorem}

\begin{example}
Test whether $x_2$ adds anything to a model already containing $x_1$: the reduced model $\hat{y}=17.1+3.0x_1$ has $SS_R=72$, so $SS_R(\beta_2\mid\beta_1)=91.22-72=19.22$, and $F_0=19.22/0.5=38.4=6.2^2$ (compare $f_{0.05,1,5}=6.61$): significant.
\end{example}

\begin{example}[Textbook example: wire bond data]
For the wire bond fit ($n=25$, $p=3$, so $22$ error degrees of freedom): $SS_T=6105.9447$, $SS_R=5990.7712$, $SS_E=115.1735$. Then $f_0\approx572\gg f_{0.05,2,22}=3.44$, so the regression is significant: pull strength is linearly related to wire length, die height, or both. This does not by itself show the model is adequate for prediction, which needs the checks of Section~\ref{sec:madeq}. $R^2=5990.7712/6105.9447=0.9811$ and $R^2_{\mathrm{adj}}=0.979$, against $R^2=0.9640$ and $R^2_{\mathrm{adj}}=0.962$ for the one-variable model with wire length alone: adding die height raises $R^2$ by only $0.0171$ but also raises $R^2_{\mathrm{adj}}$, so it is a genuine improvement. For $H_0:\beta_2=0$: $t_0=4.4767>t_{0.025,22}=2.074$ ($P=0.0002$), so die height contributes given wire length. The partial $F$-statistic is $t_0^2\approx20.0>f_{0.05,1,22}=4.30$, the same conclusion.
\end{example}

\subsection{More About the Extra Sum of Squares Method (CD Only)}

With an orthogonal design (as above), $SS_R(\beta_2\mid\beta_1)=SS_R(\beta_2)$: the contribution of each regressor does not depend on the order in which they are entered. In non-orthogonal data it does, and the contributions generally do not add to $SS_R$. The sequential (``Type I'') sums of squares depend on the order of entry, the partial (``Type III'') sums of squares measure each regressor's contribution given all the others and are equivalent to the $t$-tests. The extra sum of squares framework also tests any linear restriction $C\boldsymbol\beta=\mathbf{0}$ via $F_0=\dfrac{(SS_E(\text{restricted})-SS_E(\text{full}))/r}{MS_E(\text{full})}$.

\section{Confidence Intervals in Multiple Linear Regression}

\subsection{Confidence Intervals on Individual Regression Coefficients}

\begin{theorem}[Interval on $\beta_j$]\label{thm:mlr-ci}
$\hat\beta_j-t_{\alpha/2,n-p}\sqrt{\hat\sigma^2C_{jj}}\le\beta_j\le\hat\beta_j+t_{\alpha/2,n-p}\sqrt{\hat\sigma^2C_{jj}}$ (two-sided, $\alpha/2$ in each tail). One-sided $100(1-\alpha)\%$ bounds use $t_{\alpha,n-p}$ and one limit: $\beta_j\le\hat\beta_j+t_{\alpha,n-p}\sqrt{\hat\sigma^2C_{jj}}$ (upper) or $\hat\beta_j-t_{\alpha,n-p}\sqrt{\hat\sigma^2C_{jj}}\le\beta_j$ (lower). Likewise for the mean-response and prediction intervals below: replace $t_{\alpha/2,n-p}$ by $t_{\alpha,n-p}$ and keep one limit.
\end{theorem}

\begin{example}
$3.0\pm2.571(0.25)=3.0\pm0.643$, i.e.\ $2.36\le\beta_1\le3.64$.
\end{example}

\subsection{Confidence Interval on the Mean Response}

At a point $\mathbf{x}_0=(1,x_{01},\dots,x_{0k})\T$ the mean response is $\mu_{Y|\mathbf{x}_0}=\mathbf{x}_0\T\boldsymbol\beta$, estimated by $\hat\mu=\mathbf{x}_0\T\hat{\boldsymbol\beta}$ with $\Var[\hat\mu]=\sigma^2\mathbf{x}_0\T(X\T X)^{-1}\mathbf{x}_0$.

\begin{theorem}[Interval on the mean response]\label{thm:mlr-meanresp}
$\hat\mu_{Y|\mathbf{x}_0}\pm t_{\alpha/2,n-p}\sqrt{\hat\sigma^2\,\mathbf{x}_0\T(X\T X)^{-1}\mathbf{x}_0}$.
\end{theorem}

\begin{example}
At $\mathbf{x}_0=(1,1,1)\T$: $\hat\mu=17.1+3.0+1.55=21.65$ and $\mathbf{x}_0\T(X\T X)^{-1}\mathbf{x}_0=3/8$, so the standard error is $\sqrt{0.5\cdot3/8}=0.433$ and the $95\%$ interval is $21.65\pm2.571(0.433)=21.65\pm1.11$.
\end{example}

\section{Prediction of New Observations}

The prediction interval is always wider than the confidence interval on the mean response at the same point, because it adds the inherent variability of a single new $Y$ to the error of estimating the mean.

\begin{theorem}[Prediction interval]\label{thm:mlr-pi}
A $100(1-\alpha)\%$ prediction interval for a new observation at $\mathbf{x}_0$ is $\hat{y}_0\pm t_{\alpha/2,n-p}\sqrt{\hat\sigma^2\big(1+\mathbf{x}_0\T(X\T X)^{-1}\mathbf{x}_0\big)}$.
\end{theorem}

\begin{example}
At $\mathbf{x}_0=(1,1,1)\T$: $21.65\pm2.571\sqrt{0.5(1+0.375)}=21.65\pm2.13$.
\end{example}

\begin{pitfall}{hidden extrapolation}
In several dimensions a point can lie inside the range of each regressor separately yet far from the cloud of data jointly (a \emph{hidden extrapolation}). The quantity $h_{00}=\mathbf{x}_0\T(X\T X)^{-1}\mathbf{x}_0$ detects this: if it exceeds $\max_ih_{ii}$ the point lies outside the \emph{regressor variable hull} and predictions there are unreliable.
\end{pitfall}

\section{Model Adequacy Checking}\label{sec:madeq}

\subsection{Residual Analysis}

The residuals $\mathbf{e}=(I-H)\mathbf{y}$ have $\Cov(\mathbf{e})=\sigma^2(I-H)$, so $\Var[e_i]=\sigma^2(1-h_{ii})$: residuals at high-leverage points are artificially small. Several scalings are used.

\begin{definition}[Scaled residuals]
Standardised residual $d_i=e_i/\hat\sigma$; \emph{studentised} residual $r_i=\dfrac{e_i}{\sqrt{\hat\sigma^2(1-h_{ii})}}$, which has approximately unit variance for every observation. Values with $\abs{r_i}>3$ are potential outliers.
\end{definition}

The plots of Chapter~5 apply: normal probability plot of the residuals, residuals versus $\hat{y}_i$, residuals versus each regressor $x_j$ (curvature suggests a missing quadratic term) and residuals versus time order. The \emph{prediction residual} $e_{(i)}=e_i/(1-h_{ii})$ is the error when observation $i$ is left out of the fit, and $PRESS=\sum e_{(i)}^2$ measures predictive ability.

\subsection{Influential Observations}

\begin{definition}[Leverage and Cook's distance]
The \emph{leverage} of observation $i$ is $h_{ii}=\mathbf{x}_i\T(X\T X)^{-1}\mathbf{x}_i$, the weight that $y_i$ has on its own fitted value ($\hat{y}_i=\sum_jh_{ij}y_j$); $0\le h_{ii}\le1$ and $\sum h_{ii}=p$, so the average leverage is $p/n$. Observations with $h_{ii}>2p/n$ have high leverage. \emph{Cook's distance}
\[ D_i=\frac{r_i^2}p\cdot\frac{h_{ii}}{1-h_{ii}} \]
measures how much all fitted values change when observation $i$ is deleted; $D_i>1$ marks an influential point.
\end{definition}

An observation is influential when it combines high leverage (unusual $x$) with a large residual; a high-leverage point lying exactly on the trend merely sharpens the fit, while an outlier at a low-leverage point shifts the line little. Investigate (do not automatically delete) influential points.

\section{Aspects of Multiple Regression Modeling}

\subsection{Polynomial Regression Models}

A curved relationship can be fitted by including powers of a regressor: $Y=\beta_0+\beta_1x+\beta_2x^2+\varepsilon$ (quadratic) is a linear model in $(x,x^2)$, so all of the above applies with $x_1=x$, $x_2=x^2$. Guidelines: keep the order low (degree $\le2$ or $3$); retain lower-order terms when a higher-order term is included; \emph{centre} $x$ (use $x-\bar{x}$) to reduce the strong correlation between $x$ and $x^2$ that makes $X\T X$ ill-conditioned; and extrapolate with great caution since polynomials behave wildly outside the data range. Interaction terms $x_1x_2$ let the effect of $x_1$ depend on the level of $x_2$.

\subsection{Categorical Regressors and Indicator Variables}

A qualitative factor with $m$ levels is represented by $m-1$ \emph{indicator (dummy) variables}. For a two-level factor (e.g.\ machine A or B) let $x=0$ for A and $x=1$ for B. In $Y=\beta_0+\beta_1x_1+\beta_2x+\varepsilon$, the coefficient $\beta_2$ is the difference between B and A at a fixed $x_1$: machine A has intercept $\beta_0$ and machine B intercept $\beta_0+\beta_2$, with parallel lines. Adding the product $x_1x$ allows different slopes. With $m=3$ levels use two indicators; using $m$ indicators together with an intercept makes $X\T X$ singular (the dummy-variable trap).

\subsection{Selection of Variables and Model Building}

With many candidate regressors we want a model that is simple yet predicts well. Adding regressors always reduces $SS_E$ but inflates the variance of predictions and risks fitting noise (overfitting).

\begin{itemize}
  \item \textbf{All possible regressions} for $k$ candidates ($2^k$ models) compared by $R^2_{\mathrm{adj}}$, $MS_E$, the PRESS statistic, Mallows' $C_p=\dfrac{SS_E(p)}{\hat\sigma^2_{\mathrm{full}}}-n+2p$ (a good model has $C_p\approx p$) or AIC$=n\ln(SS_E/n)+2p$.
  \item \textbf{Stepwise methods} use partial $F$-statistics with an entry threshold $f_{\mathrm{in}}$ and a removal threshold $f_{\mathrm{out}}$, $f_{\mathrm{in}}\ge f_{\mathrm{out}}$ (software often uses $\alpha=0.15$ to enter and remove, $0.10$ for backward elimination). \emph{Stepwise regression} starts with the regressor most highly correlated with $y$ (largest $F$), then at each step adds the candidate with the largest partial $F$ if it exceeds $f_{\mathrm{in}}$, and afterwards deletes any included regressor whose partial $F$ falls below $f_{\mathrm{out}}$; it stops when nothing can be added or removed. \emph{Forward selection} only adds, never testing earlier entries for removal, which is its weakness. \emph{Backward elimination} starts with all $K$ candidates and repeatedly removes the one with the smallest partial $F$ if it is below $f_{\mathrm{out}}$. Because the variable chosen is the one that maximises $F$, the advertised significance level is not the true level, and the degrees of freedom of $MS_E$ change from step to step.
  \item \textbf{Textbook recommendation.} If the number of candidates is not too large, use all possible regressions with the minimum $MS_E$ and $C_p$ criteria: it finds the best equation for those criteria, which stepwise methods do not guarantee, and it is not distorted by dependencies among regressors. $C_p$ uses $MS_E$ of the full model as the estimate of $\sigma^2$; a model with negligible bias has $C_p\approx p$, and one chooses the minimum $C_p$ or a model with a slightly larger $C_p$ that is not much biased. A final model should always be checked with residuals, lack of fit and influence, and several methods should be compared since no single ``best'' equation need exist.
\end{itemize}

\begin{pitfall}{variable selection and inference}
Selecting a model with the same data used to test it makes the reported $t$ and $F$ statistics and $P$-values too optimistic, because the winner of many comparisons was chosen \emph{because} it looked good. Validate the final model on fresh data or by cross-validation (Chapter~14), and prefer subject-matter judgement to blind automation.
\end{pitfall}

\subsection{Multicollinearity}

\begin{definition}[Multicollinearity and VIF]
\emph{Multicollinearity} means that some regressors are nearly linear combinations of others, so $X\T X$ is nearly singular. The \emph{variance inflation factor} of $x_j$ is
\[ \mathrm{VIF}_j=\frac1{1-R_j^2}, \]
where $R_j^2$ is the $R^2$ from regressing $x_j$ on the remaining regressors.
\end{definition}

One can show $\Var[\hat\beta_j]=\dfrac{\sigma^2}{S_{jj}}\mathrm{VIF}_j$ (with $S_{jj}=\sum(x_{ij}-\bar{x}_j)^2$), so VIF is the factor by which collinearity inflates the variance. The textbook rule: $\mathrm{VIF}_j>10$ ($R_j^2>0.9$) signals a problem, though some authors consider $10$ too liberal and use $4$ or $5$ as the limit. Symptoms: large standard errors and unstable coefficients (signs that flip with small changes in the data), a significant overall $F$ with no significant individual $t$'s, and an $X\T X$ with a very small eigenvalue (condition number $\kappa(X\T X)=\lambda_{\max}/\lambda_{\min}>1000$). Remedies: remove or combine redundant regressors, collect data that spread the regressors, centre variables, or use ridge regression or principal components regression. Predictions \emph{within} the range of the data can still be fine; it is the interpretation of individual coefficients that suffers.

Remedies listed in the textbook: augment the data with new observations designed to break the near-dependencies (often impossible for cost or physical reasons); delete regressors (at the price of discarding their information); or use an estimator less sensitive to collinearity, such as ridge regression. Even with collinearity the fitted equation can still predict well by \emph{interpolation} inside the original region of $x$-space, since $\mathbf{x}\T\hat{\boldsymbol\beta}$ may be well estimated though individual $\hat\beta_j$ are not; extrapolation needs good individual estimates and generally fails.

\subsection{Ridge Regression (CD Only)}

\begin{definition}[Ridge estimator]
For $\lambda\ge0$, $\hat{\boldsymbol\beta}_R=(X\T X+\lambda I)^{-1}X\T\mathbf{y}$ (with the regressors standardised and the intercept left unpenalised). It minimises $SS_E+\lambda\sum_{j\ge1}\beta_j^2$.
\end{definition}

In practice $\lambda$ is chosen by solving for several values (on standardised data, usually $0<\lambda<1$) and inspecting a plot of the coefficients against $\lambda$, the \emph{ridge trace}: take $\lambda$ as small as possible that still stabilises the estimates, because variance decreases and squared bias increases with $\lambda$. The ridge estimator is biased but has smaller variance, and for a suitable $\lambda>0$ a smaller mean square error than least squares when collinearity is severe. It shrinks the coefficients toward zero, most strongly along directions in which the data carry little information (the small singular values of $X$; see Linear Algebra, Chapter~17). The value of $\lambda$ is chosen by cross-validation or from the ridge trace (plot of the coefficients against $\lambda$).

\subsection{Nonlinear Regression Models (CD Only)}

Some models are \emph{intrinsically nonlinear} in the parameters, for example $Y=\theta_1\big(1-e^{-\theta_2x}\big)+\varepsilon$. The least squares estimates minimise $\sum\big(y_i-f(x_i;\boldsymbol\theta)\big)^2$ but there is no closed-form solution; one iterates, for example by the \emph{Gauss--Newton} method: linearise $f$ about the current estimate, $f(x_i;\boldsymbol\theta)\approx f(x_i;\boldsymbol\theta^{(m)})+J_i(\boldsymbol\theta-\boldsymbol\theta^{(m)})$, solve the resulting linear least squares problem for the step, and repeat until convergence. Good starting values (often from a transformation or from physical insight) are important, and the usual $t$ and $F$ tests hold only approximately (for large $n$).

\begin{summarybox}
\begin{itemize}
  \item Model $\mathbf{y}=X\boldsymbol\beta+\boldsymbol\varepsilon$; $\hat{\boldsymbol\beta}=(X\T X)^{-1}X\T\mathbf{y}$; $\hat{\mathbf{y}}=H\mathbf{y}$ is the orthogonal projection onto $\Col X$; $\hat\sigma^2=SS_E/(n-p)$, $p=k+1$.
  \item $\Cov[\hat{\boldsymbol\beta}]=\sigma^2(X\T X)^{-1}$; $\se(\hat\beta_j)=\sqrt{\hat\sigma^2C_{jj}}$; Gauss--Markov; normal errors give $t$ and $\chi^2$ distributions.
  \item Overall test $F_0=MS_R/MS_E\sim F_{k,n-p}$; individual tests $t_0=\hat\beta_j/\se(\hat\beta_j)\sim t_{n-p}$; subsets by partial $F$ with the extra sum of squares; $R^2_{\mathrm{adj}}$ for comparing models.
  \item Mean response $\mathbf{x}_0\T\hat{\boldsymbol\beta}\pm t\sqrt{\hat\sigma^2\mathbf{x}_0\T(X\T X)^{-1}\mathbf{x}_0}$; prediction adds $1$ under the root.
  \item Diagnostics: studentised residuals, leverage $h_{ii}>2p/n$, Cook's $D_i>1$, PRESS; check for hidden extrapolation.
  \item Modeling: polynomial terms (centre them), indicator variables ($m-1$ for $m$ levels), selection by $C_p$/AIC/stepwise (beware optimism), multicollinearity via $\mathrm{VIF}>10$, ridge as a remedy, Gauss--Newton for nonlinear models.
\end{itemize}
\end{summarybox}

\section{Exercises}

\begin{exercise}
For the running example, construct the analysis of variance table, test $H_0:\beta_1=\beta_2=0$ at $\alpha=0.05$ and compute $R^2$ and $R^2_{\mathrm{adj}}$. Find a $95\%$ confidence interval for $\beta_2$.
\end{exercise}

\begin{exercise}
For the running example, find the $95\%$ confidence interval for the mean response and the $95\%$ prediction interval at $(x_1,x_2)=(-1,1)$ and at $(0,0)$. Why is the second interval narrower?
\end{exercise}

\begin{exercise}
Show that $H=X(X\T X)^{-1}X\T$ is symmetric and idempotent, that $\tr H=p$, and that $0\le h_{ii}\le1$. Show also that $(I-H)X=0$.
\end{exercise}

\begin{exercise}
A regression with $n=25$ and $k=4$ regressors gives $SS_R=420$ and $SS_E=180$. Test the significance of regression, compute $R^2$ and $R^2_{\mathrm{adj}}$. Adding two more regressors raises $SS_R$ to $438$; test whether the two added regressors are jointly significant.
\end{exercise}

\begin{exercise}
Regressing $x_3$ on $x_1,x_2$ gives $R_3^2=0.96$. Find $\mathrm{VIF}_3$ and explain what it implies for the standard error of $\hat\beta_3$ relative to uncorrelated regressors.
\end{exercise}

\begin{exercise}
Write the design matrix for a model with one quantitative regressor $x$ and a three-level factor (levels A, B, C) with indicator variables, for the six observations $(x,\text{level})$: $(1,A),(2,A),(3,B),(4,B),(5,C),(6,C)$. Interpret each coefficient.
\end{exercise}

\begin{exercise}
In a regression with $n=20$, $p=3$ one observation has $h_{ii}=0.62$ and studentised residual $r_i=2.1$. Compute Cook's distance and decide whether it is influential; compare $h_{ii}$ with $2p/n$.
\end{exercise}

\begin{exercise}
Show that for ridge regression with $X\T X=\diag(\lambda_1,\dots,\lambda_p)$ each coefficient is the least squares coefficient multiplied by $\lambda_j/(\lambda_j+\lambda)$, and describe what this does to coefficients whose regressors carry little information.
\end{exercise}


\part{Designed Experiments}
\chapter{Design and Analysis of Single-Factor Experiments: The Analysis of Variance}

\section{Designing Engineering Experiments}

An \emph{experiment} changes controllable inputs on purpose and watches the response.
\begin{itemize}
  \item A \emph{factor} is an input we control. It takes several \emph{levels}.
  \item A \emph{treatment} is one combination of factor levels.
  \item Statistics helps in three places: choosing the design, analysing the data, and stating conclusions with known error rates.
\end{itemize}

Three principles of design:
\begin{itemize}
  \item \textbf{Replication.} Repeat each treatment. This gives an estimate of experimental error and improves precision.
  \item \textbf{Randomisation.} Assign treatments to units, and run the trials, in random order. This protects against hidden systematic effects and justifies the independence assumption. As in Chapter~4, it also supports cause-and-effect conclusions.
  \item \textbf{Blocking.} Group similar units and compare treatments within each group. This removes variation that would otherwise sit in the error (Section~\ref{sec:rcbd}).
\end{itemize}
Before collecting data, state the problem, the response, the factors and ranges, and the planned analysis.

\section{The Completely Randomized Single-Factor Experiment}

\subsection{An Example}

Suppose $a$ treatments (levels of one factor) are compared, and units are assigned to them completely at random. This is a \emph{completely randomized design} (CRD). Let $y_{ij}$ be the $j$-th observation under treatment $i$, $j=1,\dots,n_i$.

\begin{example}\label{ex:anovadata}
Three machine settings (A, B, C) are tried on four randomly assigned runs each. The response is strength.
\begin{center}
\begin{tabular}{lcccc|cc}
\toprule
Setting & \multicolumn{4}{c|}{Observations} & Total $y_{i\cdot}$ & Mean $\bar{y}_{i\cdot}$\\
\midrule
A & 12 & 14 & 13 & 15 & 54 & 13.5\\
B & 16 & 18 & 17 & 19 & 70 & 17.5\\
C & 13 & 15 & 14 & 16 & 58 & 14.5\\
\bottomrule
\end{tabular}
\end{center}
The grand total is $y_{\cdot\cdot}=182$, the grand mean is $\bar{y}_{\cdot\cdot}=15.17$, and $N=12$.
\end{example}

\subsection{The Analysis of Variance}

\begin{definition}[Single-factor model]
\[ Y_{ij}=\mu+\tau_i+\varepsilon_{ij},\qquad i=1,\dots,a;\ j=1,\dots,n_i. \]
\begin{itemize}
  \item $\mu$ is the overall mean and $\tau_i$ is the effect of treatment $i$, with $\sum n_i\tau_i=0$.
  \item The errors $\varepsilon_{ij}$ are independent $\Normal(0,\sigma^2)$.
  \item The treatment means are $\mu_i=\mu+\tau_i$.
\end{itemize}
\end{definition}

The test is
\[ H_0:\mu_1=\mu_2=\dots=\mu_a\ (\text{all }\tau_i=0)\qquad\text{against}\qquad H_1:\mu_i\ne\mu_j\text{ for at least one pair.} \]

\begin{theorem}[Partition of the total variability]\label{thm:anova-ss}
\[ \underbrace{\sum_{i=1}^a\sum_{j=1}^{n_i}(y_{ij}-\bar{y}_{\cdot\cdot})^2}_{SS_T}=\underbrace{\sum_{i=1}^an_i(\bar{y}_{i\cdot}-\bar{y}_{\cdot\cdot})^2}_{SS_{\mathrm{Treatments}}}+\underbrace{\sum_{i=1}^a\sum_{j=1}^{n_i}(y_{ij}-\bar{y}_{i\cdot})^2}_{SS_E}, \]
with degrees of freedom $N-1=(a-1)+(N-a)$.
\end{theorem}

\begin{proof}
Write $y_{ij}-\bar{y}_{\cdot\cdot}=(\bar{y}_{i\cdot}-\bar{y}_{\cdot\cdot})+(y_{ij}-\bar{y}_{i\cdot})$ and square. The cross term is $2\sum_i(\bar{y}_{i\cdot}-\bar{y}_{\cdot\cdot})\sum_j(y_{ij}-\bar{y}_{i\cdot})=0$, because deviations from a group mean sum to zero within each group.
\end{proof}

The \emph{mean squares} are $MS_{\mathrm{Tr}}=SS_{\mathrm{Tr}}/(a-1)$ and $MS_E=SS_E/(N-a)$. Their expected values explain the test:
\[ \Exp[MS_E]=\sigma^2,\qquad\Exp[MS_{\mathrm{Tr}}]=\sigma^2+\frac{\sum n_i\tau_i^2}{a-1}. \]
\begin{itemize}
  \item If $H_0$ is true, both estimate $\sigma^2$, so their ratio is near $1$.
  \item If the means differ, $MS_{\mathrm{Tr}}$ is inflated and the ratio is large.
\end{itemize}

\begin{theorem}[The $F$-test of ANOVA]\label{thm:anova-F}
Under $H_0$ and the model assumptions, $SS_{\mathrm{Tr}}/\sigma^2\sim\chi^2_{a-1}$ and $SS_E/\sigma^2\sim\chi^2_{N-a}$, independently. So
\[ F_0=\frac{MS_{\mathrm{Tr}}}{MS_E}\sim F_{a-1,\,N-a}. \]
Reject $H_0$ if $f_0>f_{\alpha,\,a-1,\,N-a}$ (upper tail only).
\end{theorem}

\begin{center}
\begin{tabular}{lcccc}
\toprule
Source & Sum of squares & df & Mean square & $F_0$ \\
\midrule
Treatments & $SS_{\mathrm{Tr}}$ & $a-1$ & $MS_{\mathrm{Tr}}$ & $MS_{\mathrm{Tr}}/MS_E$ \\
Error & $SS_E$ & $N-a$ & $MS_E$ & \\
Total & $SS_T$ & $N-1$ & & \\
\bottomrule
\end{tabular}
\end{center}

\emph{Computing formulas} (with $y_{i\cdot}$ the treatment totals):
\[ SS_T=\sum\sum y_{ij}^2-\frac{y_{\cdot\cdot}^2}N,\qquad SS_{\mathrm{Tr}}=\sum\frac{y_{i\cdot}^2}{n_i}-\frac{y_{\cdot\cdot}^2}N,\qquad SS_E=SS_T-SS_{\mathrm{Tr}}. \]

\begin{example}[Continuing Example~\ref{ex:anovadata}]\label{ex:anovacalc}
\begin{itemize}
  \item $SS_{\mathrm{Tr}}=4\big[(13.5-15.17)^2+(17.5-15.17)^2+(14.5-15.17)^2\big]=34.67$.
  \item Within each setting the deviations from the mean are $(-1.5,0.5,-0.5,1.5)$. Their squares sum to $5$, so $SS_E=3(5)=15.0$.
  \item $SS_T=49.67$. Degrees of freedom: $2$ and $9$. $MS_{\mathrm{Tr}}=17.33$, $MS_E=1.667$.
  \item $f_0=10.4>f_{0.05,2,9}=4.26$, so reject $H_0$ ($P\approx0.005$). The settings do not all give the same mean strength.
\end{itemize}
\end{example}

The formulas hold for unequal $n_i$. Still, prefer a balanced design ($n_i=n$): the test is more robust to unequal variances, and power is largest for a given $N$.

\begin{example}[Paper tensile strength]
A paper maker tests four hardwood concentrations ($5,10,15,20\%$) with six specimens each, in random order. So $a=4$, $n=6$, $N=24$. The treatment totals are $60,94,102,127$ (means $10.0,\,15.67,\,17.0,\,21.17$) and the grand total is $383$.
\begin{itemize}
  \item $SS_T=512.96$, $SS_{\text{Treatments}}=382.79$, $SS_E=130.17$, with $23$, $3$ and $20$ degrees of freedom.
  \item $MS_{\text{Treatments}}=127.60$, $MS_E=6.51$, $f_0=19.60$.
  \item At $\alpha=0.01$, $f_{0.01,3,20}=4.94$. Reject $H_0$ ($P=3.59\times10^{-6}$).
\end{itemize}
The rejection region is the \emph{upper} tail only, because under $H_1$ the expected value of $MS_{\text{Treatments}}$ exceeds $\sigma^2$.

Intervals: for one treatment mean, $\bar{y}_{i\cdot}\pm t_{\alpha/2,\,a(n-1)}\sqrt{MS_E/n}$. For a difference, $\bar{y}_{i\cdot}-\bar{y}_{j\cdot}\pm t_{\alpha/2,\,a(n-1)}\sqrt{2MS_E/n}$. If the interval for a difference contains $0$, no difference is detected.

Fisher's LSD: $t_{0.025,20}=2.086$, so $\mathrm{LSD}=2.086\sqrt{2(6.51)/6}=3.07$. Every pair of means differs by more than $3.07$, except $10\%$ and $15\%$ (means $15.67$ and $17.0$). So $20\%$ hardwood gives the greatest strength.
\end{example}

\subsection{Multiple Comparisons Following the ANOVA}

A significant $F$ says that \emph{some} means differ. It does not say which.

\begin{itemize}
  \item Comparing all $\binom a2$ pairs with $t$-tests at level $\alpha$ each makes the chance of at least one false positive much larger than $\alpha$. This chance is the \emph{experimentwise error rate}.
  \item For $a=5$ there are $10$ pairs. With $\alpha=0.05$ each, the experimentwise rate is about $0.40$.
\end{itemize}
Two procedures control it.

\begin{definition}[Tukey's method]
In a balanced design, declare $\mu_i\ne\mu_j$ if $\abs{\bar{y}_{i\cdot}-\bar{y}_{j\cdot}}>T_\alpha$, where
\[ T_\alpha=q_\alpha(a,f)\sqrt{\frac{MS_E}n},\qquad f=N-a. \]
$q_\alpha(a,f)$ is the upper $\alpha$ point of the studentised range distribution (tabulated).
\begin{itemize}
  \item Tukey keeps the \emph{overall} error rate for all pairwise comparisons at exactly $\alpha$ (balanced design).
  \item It gives simultaneous intervals $\bar{y}_{i\cdot}-\bar{y}_{j\cdot}\pm T_\alpha$ for all $\mu_i-\mu_j$.
  \item By contrast, the ``family error rate'' that software quotes for unprotected LSD comparisons is the chance of at least one Type~I error among all the comparisons. It is about $0.19$ for $a=4$ (six comparisons).
\end{itemize}
\end{definition}

\begin{definition}[Fisher's least significant difference (LSD)]
Declare $\mu_i\ne\mu_j$ if
\[ \abs{\bar{y}_{i\cdot}-\bar{y}_{j\cdot}}>\mathrm{LSD}=t_{\alpha/2,\,N-a}\sqrt{MS_E\Big(\frac1{n_i}+\frac1{n_j}\Big)}. \]
\begin{itemize}
  \item Apply LSD only \emph{after} a significant $F$-test (``protected LSD''). This protects the error rate.
  \item LSD and Tukey are \emph{two-sided} comparisons, with $\alpha/2$ in each tail. For a one-sided comparison $H_1:\mu_i>\mu_j$, use $t_{\alpha,\,N-a}$ instead of $t_{\alpha/2,\,N-a}$.
  \item The ANOVA $F$-test itself rejects only in the upper tail, $f_0>f_{\alpha,\,a-1,\,N-a}$. The ratio $MS_{\text{Tr}}/MS_E$ becomes large for \emph{any} departure from equal means. There is no ``two-sided $F$-test'' here.
\end{itemize}
\end{definition}

\begin{example}[Continuing Example~\ref{ex:anovacalc}]
\begin{itemize}
  \item Tukey: $q_{0.05}(3,9)=3.95$, so $T_{0.05}=3.95\sqrt{1.667/4}=2.55$.
  \item Differences: $B-A=4.0$, $B-C=3.0$, $C-A=1.0$. Setting B differs from both others. A and C are not distinguished.
  \item LSD: $2.262\sqrt{1.667\cdot2/4}=2.262(0.913)=2.07$. Same conclusions here.
\end{itemize}
\end{example}

\subsection{More About Multiple Comparisons (CD Only)}

Fisher's LSD is easy and very common, but it is liberal. Each test has level $\alpha$, yet the experimentwise Type~I error is much larger than $\alpha$. Three other approaches follow.

\subsubsection*{Graphical comparison of means}

Each treatment average $\bar{y}_{i\cdot}$ has standard deviation $\sigma/\sqrt n$. If all treatment means were equal, the averages would look like a random sample from a normal curve with that spread, estimated by $\sqrt{MS_E/n}$.
\begin{itemize}
  \item Imagine this curve on a movable transparency. Slide it along the axis where the $a$ means are plotted.
  \item If some position makes all the means look like typical draws, there is no evidence of differences.
  \item Means that do not fit under the curve belong to treatments with different mean responses.
\end{itemize}
This is a rough but effective picture of the ANOVA result. (The $t$ distribution would be exact, but the normal curve is close.)

\subsubsection*{Orthogonal contrasts}

\begin{definition}[Contrast]
A \emph{contrast} is a combination of treatment totals $C=\sum_{i=1}^ac_iy_{i\cdot}$ with $\sum c_i=0$ (unbalanced: $\sum n_ic_i=0$). It expresses one specific comparison, chosen \emph{before} the experiment.
\end{definition}
Examples: $H_0:\mu_3=\mu_4$ has coefficients $(0,0,1,-1)$. $H_0:\mu_1+\mu_3=\mu_2+\mu_4$ has coefficients $(1,-1,1,-1)$.

The sum of squares of a contrast has one degree of freedom:
\[ SS_C=\frac{\big(\sum_ic_iy_{i\cdot}\big)^2}{n\sum_ic_i^2}\qquad\Big(\text{unbalanced: }\frac{\big(\sum c_iy_{i\cdot}\big)^2}{\sum n_ic_i^2}\Big). \]
It is tested by $F_0=SS_C/MS_E\sim F_{1,\,N-a}$.

\begin{itemize}
  \item Contrasts $\{c_i\}$ and $\{d_i\}$ are \emph{orthogonal} if $\sum c_id_i=0$ (unbalanced: $\sum n_ic_id_i=0$).
  \item A set of $a-1$ mutually orthogonal contrasts splits $SS_{\mathrm{Tr}}$ into $a-1$ independent one-degree-of-freedom pieces. The tests are independent.
  \item With $a=3$ and treatment 1 a control, $(-2,1,1)$ (control against the average of the others) and $(0,1,-1)$ (one level against the other) are orthogonal.
  \item Choose contrasts \emph{before} looking at the data. Picking the largest observed differences afterwards inflates the Type~I error.
\end{itemize}

\begin{example}
For Example~\ref{ex:anovacalc}: $y_{1\cdot}=54$, $y_{2\cdot}=70$, $y_{3\cdot}=58$, $n=4$.
\begin{itemize}
  \item $c=(-2,1,1)$: $C=-108+70+58=20$ and $SS_C=20^2/(4\cdot6)=16.67$.
  \item $d=(0,1,-1)$: $C=12$ and $SS_d=12^2/(4\cdot2)=18.00$.
  \item $c\cdot d=0$, so they are orthogonal, and $SS_c+SS_d=34.67=SS_{\mathrm{Tr}}$.
  \item Against $MS_E=1.667$: $F_c=10.0$ and $F_d=10.8$. Compare each with $f_{0.05,1,9}=5.12$. Both are significant.
\end{itemize}
\end{example}

\subsubsection*{Tukey's method}

Tukey's procedure compares all pairs using the studentised range
\[ Q=\frac{\bar{Y}_{\max}-\bar{Y}_{\min}}{\sqrt{MS_E/n}}, \]
where $\bar{Y}_{\max}$ and $\bar{Y}_{\min}$ are the largest and smallest of the $a$ sample means. For equal sample sizes, declare $\mu_i\ne\mu_j$ if $\abs{\bar{y}_{i\cdot}-\bar{y}_{j\cdot}}>T_\alpha=q_\alpha(a,f)\sqrt{MS_E/n}$, with $f=N-a$. The overall error rate for all pairwise comparisons is exactly $\alpha$.

\begin{interview}{beyond the textbook}
Other procedures you will meet: Bonferroni (test each of $m$ comparisons at $\alpha/m$), Scheff\'e (all possible contrasts; most conservative), Dunnett (each treatment against a control), and the Holm and Benjamini--Hochberg corrections (Chapter~13). The textbook covers Fisher's LSD, the graphical method, orthogonal contrasts and Tukey's method.
\end{interview}

\subsection{Residual Analysis and Model Checking}

The residuals are $e_{ij}=y_{ij}-\bar{y}_{i\cdot}$. Check the model as in regression.
\begin{itemize}
  \item \textbf{Normal probability plot of residuals.} Checks normality. The $F$-test is fairly robust to moderate departures.
  \item \textbf{Residuals against fitted values $\bar{y}_{i\cdot}$.} A funnel shape means the variance grows with the mean. Try $\ln y$ or $\sqrt y$.
  \item \textbf{Residuals against run order.} A trend suggests a time effect, or that randomisation or independence failed.
  \item \textbf{Equal variances (beyond the textbook).} \emph{Bartlett's test} of $H_0:\sigma_1^2=\dots=\sigma_a^2$ uses $\chi_0^2=2.3026\,q/c$, with
  \[ q=(N-a)\log_{10}S_p^2-\sum(n_i-1)\log_{10}S_i^2,\qquad c=1+\frac1{3(a-1)}\Big(\sum\frac1{n_i-1}-\frac1{N-a}\Big). \]
  Compare with $\chi^2_{\alpha,\,a-1}$. It is sensitive to non-normality, so Levene's test is often preferred.
\end{itemize}

\subsection{Determining Sample Size}

The power of the $F$-test depends on a noncentrality parameter. For the balanced case define
\[ \Phi^2=\frac{n\sum_{i=1}^a\tau_i^2}{a\sigma^2}. \]
Read the Type~II error $\beta$ from OC curves with $a-1$ and $a(n-1)$ degrees of freedom. (Or compute power from the noncentral $F$ with noncentrality $\lambda=a\Phi^2$.)

To plan the experiment:
\begin{enumerate}
  \item Choose $D$, the smallest range of treatment means (largest minus smallest) that matters in practice.
  \item The worst case for a given $D$ puts two means at the extremes and the rest at the centre. Then $\Phi^2=nD^2/(2a\sigma^2)$.
  \item Use a prior estimate of $\sigma$. Choose $n$ so the power is adequate (say $0.8$ or $0.9$).
\end{enumerate}
Larger $n$, larger $D$ and smaller $\sigma$ all increase power.

\begin{example}
Five means, $\alpha=0.01$. We want to reject $H_0$ with probability at least $0.90$ when $\sum\tau_i^2=5.0\sigma^2$.
\begin{itemize}
  \item $\Phi^2=n\sum\tau_i^2/(a\sigma^2)=n$, with $\nu_1=4$ and $\nu_2=5(n-1)$.
  \item Try $n=4$: $\Phi^2=4$, $\Phi=2$, $\nu_2=15$. The OC curve gives $\beta\approx0.38$, so power $\approx0.62<0.90$. Not enough.
  \item Continue the same way. At least $n=6$ replicates are needed.
\end{itemize}
\end{example}

\subsection{Technical Details about the Analysis of Variance (CD Only)}

Write the observations as a vector. Each sum of squares is a quadratic form $\mathbf{y}\T A_k\mathbf{y}$, where the $A_k$ are symmetric idempotent matrices with $\sum A_k=I$ and $A_jA_k=0$.

By Cochran's theorem, the forms $\mathbf{y}\T A_k\mathbf{y}/\sigma^2$ are then independent noncentral chi-square variables. Each has degrees of freedom $\operatorname{rank}A_k$ and noncentrality $\boldsymbol\mu\T A_k\boldsymbol\mu/\sigma^2$. The noncentrality is zero exactly when that effect is absent.

Geometrically, the observation vector is split into orthogonal parts: the grand mean, the treatment effect and the residual (Linear Algebra, Chapters~11--12). The ratio of two independent scaled central chi-squares is $F$. This proves Theorem~\ref{thm:anova-F}.

\section{The Random Effects Model}

\subsection{Fixed Versus Random Factors}

\begin{itemize}
  \item \textbf{Fixed effects.} The $a$ treatments are the specific levels of interest. Conclusions refer to those levels only.
  \item \textbf{Random effects.} The $a$ levels are a random sample from a large population of levels (for example, $a$ operators, batches or looms chosen at random). We want to describe the whole population, in particular how much variability the factor contributes.
\end{itemize}

\subsection{ANOVA and Variance Components}

\begin{definition}[Random effects model]
$Y_{ij}=\mu+\tau_i+\varepsilon_{ij}$ with $\tau_i\sim\Normal(0,\sigma_\tau^2)$ and $\varepsilon_{ij}\sim\Normal(0,\sigma^2)$, all independent.
\begin{itemize}
  \item The \emph{variance components} are $\sigma_\tau^2$ and $\sigma^2$.
  \item $\Var(Y_{ij})=\sigma_\tau^2+\sigma^2$.
\end{itemize}
\end{definition}

The ANOVA table and sums of squares are exactly as before. The hypothesis and the interpretation change.
\begin{itemize}
  \item Test $H_0:\sigma_\tau^2=0$ against $H_1:\sigma_\tau^2>0$ with the same $F_0=MS_{\mathrm{Tr}}/MS_E\sim F_{a-1,\,N-a}$ under $H_0$.
  \item Testing individual effects $\tau_i=0$ is meaningless here. Also $\sum\tau_i=0$ does not apply.
\end{itemize}
For balanced data ($n_i=n$) the expected mean squares are
\[ \Exp[MS_E]=\sigma^2,\qquad\Exp[MS_{\mathrm{Tr}}]=\sigma^2+n\sigma_\tau^2. \]
So the method-of-moments (\emph{ANOVA}) estimators are
\[ \hat\sigma^2=MS_E,\qquad\hat\sigma_\tau^2=\frac{MS_{\mathrm{Tr}}-MS_E}n. \]

\begin{example}
Treat the settings of Example~\ref{ex:anovacalc} as a random sample of levels.
\begin{itemize}
  \item $\hat\sigma^2=1.667$ and $\hat\sigma_\tau^2=(17.33-1.667)/4=3.92$.
  \item The share of total variance due to the factor (the \emph{intraclass correlation}) is $3.92/(3.92+1.667)=0.70$.
\end{itemize}
\end{example}

\begin{example}[Looms]
A textile company samples four looms at random and makes four strength determinations per loom. The ANOVA shows significant loom-to-loom variability.
\begin{itemize}
  \item $\hat\sigma^2=1.90$ (within looms) and $\hat\sigma_\tau^2=6.96$ (between looms).
  \item $\widehat{V}(Y_{ij})=6.96+1.90=8.86$. The process standard deviation is $\sqrt{8.86}=2.98$ psi about a mean of $95.45$ psi. Most of the variability comes from differences between looms.
  \item If the loom-to-loom variability were removed, the standard deviation would fall to $\sqrt{1.90}=1.38$ psi. That would greatly reduce fallout below a lower specification limit of $90$ psi.
\end{itemize}
\end{example}

\begin{pitfall}{negative variance estimates}
If $MS_{\mathrm{Tr}}<MS_E$, the ANOVA estimate $\hat\sigma_\tau^2$ is negative. This suggests the true $\sigma_\tau^2$ is near or at zero. The textbook lists three courses of action:
\begin{itemize}
  \item Accept it as evidence that the true component is zero (this disturbs the properties of the other estimates). In practice, report $\hat\sigma_\tau^2=0$.
  \item Re-estimate with a method that always gives nonnegative values, such as maximum likelihood or REML.
  \item Treat it as a sign that the assumed linear model is wrong, and study the model.
\end{itemize}
\end{pitfall}

\subsection{Determining Sample Size in the Random Model (CD Only)}

Under $H_1$, the statistic $F_0$ has a \emph{central} $F$ distribution, scaled by $\lambda=\sqrt{1+n\sigma_\tau^2/\sigma^2}$:
\[ F_0/\lambda^2\sim F_{a-1,\,a(n-1)}. \]
So the power is $\Prob\big(F_{a-1,\,a(n-1)}>f_{\alpha}/\lambda^2\big)$. It grows with $n\sigma_\tau^2/\sigma^2$.

The levels are themselves sampled. So at a fixed total $N=an$, adding levels $a$ usually helps more than adding replication $n$.

\section{Randomized Complete Block Design}\label{sec:rcbd}

\subsection{Design and Statistical Analysis}

When units are different (different days, batches, operators), treatment comparisons in a CRD get mixed up with that variation. This inflates $MS_E$.

In a \emph{randomized complete block design} (RCBD):
\begin{itemize}
  \item Units are grouped into $b$ \emph{blocks} of $a$ similar units.
  \item Within each block, each of the $a$ treatments appears \emph{once}, in random order.
  \item Treatments are compared within blocks, where the units are alike.
\end{itemize}

\begin{definition}[RCBD model]
$Y_{ij}=\mu+\tau_i+\beta_j+\varepsilon_{ij}$, with $i=1,\dots,a$ (treatments) and $j=1,\dots,b$ (blocks). Also $\sum\tau_i=0$, $\sum\beta_j=0$, and $\varepsilon_{ij}\sim\Normal(0,\sigma^2)$ independent.
\end{definition}

\begin{theorem}[Partition for the RCBD]\label{thm:rcbd-ss}
\[ SS_T=SS_{\mathrm{Tr}}+SS_{\mathrm{Blocks}}+SS_E, \]
with
\begin{align*}
SS_{\mathrm{Tr}}&=b\sum_i(\bar{y}_{i\cdot}-\bar{y}_{\cdot\cdot})^2,\\
SS_{\mathrm{Blocks}}&=a\sum_j(\bar{y}_{\cdot j}-\bar{y}_{\cdot\cdot})^2,\\
SS_E&=\sum_{i,j}(y_{ij}-\bar{y}_{i\cdot}-\bar{y}_{\cdot j}+\bar{y}_{\cdot\cdot})^2,
\end{align*}
and degrees of freedom $ab-1=(a-1)+(b-1)+(a-1)(b-1)$.
\end{theorem}

\begin{theorem}[Tests for the RCBD]\label{thm:rcbd-F}
To test $H_0:\tau_1=\dots=\tau_a=0$, use
\[ F_0=\frac{MS_{\mathrm{Tr}}}{MS_E}\sim F_{a-1,\,(a-1)(b-1)}\ \text{ under }H_0. \]
Reject if $f_0>f_{\alpha,\,a-1,\,(a-1)(b-1)}$.

Blocks can be tested by $MS_{\mathrm{Blocks}}/MS_E$. This is of secondary interest: blocks were not randomised, so it is only a rough guide to whether blocking was worthwhile.
\end{theorem}

Computing formulas, with totals $y_{i\cdot}$ (treatments), $y_{\cdot j}$ (blocks) and $y_{\cdot\cdot}$:
\[ SS_{\mathrm{Tr}}=\frac1b\sum y_{i\cdot}^2-\frac{y_{\cdot\cdot}^2}{ab},\qquad SS_{\mathrm{Blocks}}=\frac1a\sum y_{\cdot j}^2-\frac{y_{\cdot\cdot}^2}{ab}. \]

\begin{example}\label{ex:rcbd}
Three treatments in four blocks (rows are treatments, columns are blocks):
\begin{center}
\begin{tabular}{l cccc|c}
\toprule
 & B1 & B2 & B3 & B4 & $y_{i\cdot}$\\
\midrule
T1 & 10 & 12 & 11 & 15 & 48\\
T2 & 13 & 14 & 13 & 18 & 58\\
T3 & 12 & 12 & 13 & 16 & 53\\
\midrule
$y_{\cdot j}$ & 35 & 38 & 37 & 49 & 159\\
\bottomrule
\end{tabular}
\end{center}
\begin{itemize}
  \item Correction term: $159^2/12=2106.75$. With $\sum y^2=2161$, $SS_T=54.25$.
  \item $SS_{\mathrm{Tr}}=(48^2+58^2+53^2)/4-2106.75=12.50$.
  \item $SS_{\mathrm{Blocks}}=(35^2+38^2+37^2+49^2)/3-2106.75=39.58$.
  \item $SS_E=54.25-12.50-39.58=2.17$.
  \item Degrees of freedom $2,3,6$: $MS_{\mathrm{Tr}}=6.25$, $MS_{\mathrm{Blocks}}=13.19$, $MS_E=0.361$.
  \item $f_0=17.3>f_{0.05,2,6}=5.14$, so the treatment means differ. The block ratio $F=36.5$ shows large block-to-block differences.
\end{itemize}
\end{example}

\subsection{Multiple Comparisons}

Each treatment mean is an average of $b$ observations. So:
\[ \text{Tukey: }T_\alpha=q_\alpha\big(a,(a-1)(b-1)\big)\sqrt{MS_E/b},\qquad\text{LSD: }t_{\alpha/2,\,(a-1)(b-1)}\sqrt{2MS_E/b}. \]

\begin{example}
For Example~\ref{ex:rcbd}: $q_{0.05}(3,6)=4.34$ and $T=4.34\sqrt{0.361/4}=1.30$.

The treatment means are $12.0$, $14.5$, $13.25$. T2 differs from T1 (difference $2.5$). T3 is not distinguished from either (differences $1.25$ and $1.25$).
\end{example}

\begin{example}[Fabric strength]
Four chemicals ($a=4$) are tested once each on five fabric samples ($b=5$, the blocks), in random order within each sample. The error has $(a-1)(b-1)=12$ degrees of freedom.
\begin{itemize}
  \item $f_0=75.13>f_{0.05,3,12}=5.95$, so the chemicals differ.
  \item Fisher's LSD (each average uses $b=5$ observations) is $0.39$. Any pair of averages differing by $0.39$ or more is significantly different.
  \item Chemical~4 differs from the other three. Chemicals $2$ and $3$ do not differ; neither do $1$ and $3$.
  \item Without blocking, the error would have $16$ degrees of freedom, so blocking cost only $b-1=4$. The block effect was large and would otherwise have inflated $MS_E$.
\end{itemize}
Rule of thumb: when in doubt whether block effects matter, block. The block $F$-ratio that software prints is of doubtful validity as a formal test, because randomisation was restricted to within blocks.
\end{example}

\begin{interview}{was blocking worth it?}
Compare with a CRD on the same units. The error mean square a CRD would have had is estimated by
\[ \widehat{MS}_E^{\mathrm{CRD}}=\frac{(b-1)MS_{\mathrm{Blocks}}+b(a-1)MS_E}{ab-1}. \]
The relative efficiency is $\widehat{MS}_E^{\mathrm{CRD}}/MS_E$.

In Example~\ref{ex:rcbd}: $\big(3(13.19)+4\cdot2(0.361)\big)/11=3.86$, and $3.86/0.361=10.7$. A CRD would need about ten times as many units for the same precision.

Blocking costs $b-1$ error degrees of freedom. It can hurt only if the blocks do not actually reduce variability.
\end{interview}

\subsection{Residual Analysis and Model Checking}

The residuals are $e_{ij}=y_{ij}-\bar{y}_{i\cdot}-\bar{y}_{\cdot j}+\bar{y}_{\cdot\cdot}$. Use the usual plots: normal probability, residuals against fitted values, against treatment, against block and against time.

The RCBD assumes \emph{no treatment-by-block interaction}: treatment effects are the same in every block.
\begin{itemize}
  \item Non-additivity shows up as a pattern in the residual plots. Tukey's one-degree-of-freedom test for nonadditivity can detect it.
  \item If the interaction is real, rethink the model, for example with replication within blocks (Chapter~8).
\end{itemize}

\subsection{Randomized Complete Block Design with Random Factors (CD Only)}

If the blocks (or the treatments) are a random sample, the model has random effects $\beta_j\sim\Normal(0,\sigma_\beta^2)$. The ANOVA table and the treatment $F$-test do not change. With fixed treatments the expected mean squares are
\[ \Exp[MS_E]=\sigma^2,\qquad\Exp[MS_{\mathrm{Blocks}}]=\sigma^2+a\sigma_\beta^2,\qquad\Exp[MS_{\mathrm{Tr}}]=\sigma^2+\frac{b\sum\tau_i^2}{a-1}. \]
So the block variance component is $\hat\sigma_\beta^2=(MS_{\mathrm{Blocks}}-MS_E)/a$.

For Example~\ref{ex:rcbd}: $\hat\sigma_\beta^2=(13.19-0.361)/3=4.28$ and $\hat\sigma^2=0.361$. Nearly all the variability is between blocks.

\begin{summarybox}
\begin{itemize}
  \item Basic design principles: replication, randomisation, blocking. CRD model $y_{ij}=\mu+\tau_i+\varepsilon_{ij}$; $SS_T=SS_{\mathrm{Tr}}+SS_E$; $F_0=MS_{\mathrm{Tr}}/MS_E\sim F_{a-1,N-a}$ (upper tail).
  \item After a significant $F$, use Tukey ($T_\alpha=q_\alpha(a,N-a)\sqrt{MS_E/n}$), Fisher's LSD, Bonferroni, Dunnett or Scheff\'e. Many separate $t$-tests inflate the error rate.
  \item Check normality, constant variance (Bartlett, Levene) and independence (run order) from residuals. Transform the response if needed.
  \item Random effects: same $F$-test, now of $H_0:\sigma_\tau^2=0$. $\hat\sigma_\tau^2=(MS_{\mathrm{Tr}}-MS_E)/n$.
  \item RCBD: $SS_T=SS_{\mathrm{Tr}}+SS_{\mathrm{Blocks}}+SS_E$ with $(a-1)(b-1)$ error df. Compare treatments within blocks. Efficiency against a CRD comes from $\widehat{MS}_E^{\mathrm{CRD}}$.
\end{itemize}
\end{summarybox}

\section{Exercises}

\begin{exercise}
Four methods of teaching are compared, with scores: method 1: $72,75,78,70,74$; method 2: $80,84,79,85,82$; method 3: $74,76,73,78,75$; method 4: $68,71,70,66,69$. Construct the ANOVA table, test equality of the means at $\alpha=0.05$, and use Tukey's method to find which methods differ ($q_{0.05}(4,16)=4.05$).
\end{exercise}

\begin{exercise}
For Exercise~1, compute the residuals, make a normal probability plot, and test equality of variances with Bartlett's test.
\end{exercise}

\begin{exercise}
Derive $\Exp[MS_E]=\sigma^2$ for the single-factor model by showing that $SS_E/\sigma^2$ is a sum of $a$ independent $\chi^2_{n_i-1}$ variables.
\end{exercise}

\begin{exercise}
Five operators are randomly selected from a large pool and each produces four items; $MS_{\mathrm{Operators}}=48$ and $MS_E=12$. Test $H_0:\sigma_\tau^2=0$ ($f_{0.05,4,15}=3.06$) and estimate $\sigma_\tau^2$, $\sigma^2$ and the proportion of variance due to operators.
\end{exercise}

\begin{exercise}
Four fertilisers are tested on plots in five fields (blocks); the ANOVA gives $SS_{\mathrm{Tr}}=120$, $SS_{\mathrm{Blocks}}=300$, $SS_E=90$. Complete the table, test treatments at $\alpha=0.05$, and estimate the relative efficiency of the RCBD to a CRD.
\end{exercise}

\begin{exercise}
Show that in a balanced design with $a=2$ treatments the ANOVA $F$-test is equivalent to the pooled two-sample $t$-test ($F_0=T_0^2$).
\end{exercise}

\begin{exercise}
An experimenter runs four treatments on $24$ units arranged in a CRD but discovers the units fall into six identifiable batches. Explain how the analysis should change, and why ignoring the batches affects the $F$-test.
\end{exercise}

\begin{exercise}
For $a=4$ treatments and $\sigma=2$, find (using $\Phi^2=nD^2/(2a\sigma^2)$) the replication $n$ needed to detect a range of means $D=4$ with high power; discuss how the answer would change with $\sigma=3$.
\end{exercise}

\chapter{Design of Experiments with Several Factors}

\section{Introduction}

Chapter~7 varied one factor. Real processes depend on many factors, and the effect of one factor often depends on the level of another.
\begin{itemize}
  \item The \emph{one-factor-at-a-time} strategy cannot detect this, and it uses the data inefficiently.
  \item In a \emph{factorial experiment}, all combinations of the levels of all factors are run. Each observation then helps estimate every effect.
\end{itemize}

\section{Some Applications of Designed Experiments (CD Only)}

Factorial and fractional factorial designs are used to:
\begin{itemize}
  \item \emph{screen} many factors for the few that matter;
  \item \emph{characterise} a process (find the settings and interactions that set the response);
  \item \emph{optimise} it (response surface methods);
  \item establish \emph{robustness} (find settings where the response is insensitive to noise factors such as temperature or operator).
\end{itemize}
Typical uses: chemical yield studies, tuning machining parameters, reliability testing, software configuration testing.

\section{Factorial Experiments}

\begin{definition}[Main effect, interaction]
\begin{itemize}
  \item The \emph{main effect} of a factor is the change in the mean response caused by changing that factor's level, averaged over the levels of the other factors.
  \item There is an \emph{interaction} between two factors if the effect of one depends on the level of the other.
\end{itemize}
\end{definition}

\begin{example}\label{ex:interaction}
Two factors A and B at two levels each. Mean responses: at B low, A low $=10$ and A high $=20$. At B high, A low $=12$ and A high $=14$.
\begin{itemize}
  \item The effect of A is $+10$ at low B, but only $+2$ at high B. So A and B interact.
  \item The average (main) effect of A is $6$. This misrepresents both situations.
  \item Plot the means against the levels of A, one line per level of B. The lines are \emph{not parallel}. Parallel lines mean no interaction.
\end{itemize}
\end{example}

For two factors at two levels, effects are differences of averages. Take the textbook's Table~14-1: $A_{\text{low}}$ gives $10$ ($B_{\text{low}}$) and $20$ ($B_{\text{high}}$); $A_{\text{high}}$ gives $30$ and $40$. Then
\begin{align*}
A&=\frac{30+40}2-\frac{10+20}2=20,\\
B&=\frac{20+40}2-\frac{10+30}2=10,\\
AB&=\frac{10+40}2-\frac{20+30}2=0.
\end{align*}
The interaction effect is the difference of the diagonal averages. The lines of response against $A$ are parallel.

Now change the $B_{\text{high}}$, $A_{\text{high}}$ response to $0$ (Table~14-2: $10,20$ and $30,0$).
\begin{itemize}
  \item The $A$ effect is $30-10=20$ at $B_{\text{low}}$ and $0-20=-20$ at $B_{\text{high}}$.
  \item The main effect of $A$ is $(30+0)/2-(10+20)/2=0$. This wrongly suggests $A$ does nothing.
  \item The interaction is large: $AB=(10+0)/2-(20+30)/2=-20$.
\end{itemize}
A significant interaction can hide the main effects. Then knowing the interaction is more useful than knowing the main effects.

\begin{nuance}{why not one factor at a time?}
A chemist fixes temperature at $155^\circ$F and varies time ($0.5$ to $2.5$ h). The best yield is near $1.7$ h. Then time is fixed at $1.7$ h and temperature varied ($140$ to $180^\circ$F). The best is near $155^\circ$F.

The conclusion ``$155^\circ$F and $1.7$ h'' can be at least $20$ yield points below the true optimum, which lies at shorter times and higher temperatures. The method fails because it cannot see the time--temperature interaction. A factorial experiment can. It also needs fewer runs.
\end{nuance}

\begin{figure}[h]
\centering
\begin{tikzpicture}[scale=0.9]
  \begin{scope}
    \draw[axis] (0,0)--(4.4,0) node[right,lab]{A}; \draw[axis] (0,0)--(0,3.4) node[above,lab]{response};
    \draw[thick,accent2] (0.8,0.8)--(3.4,1.5); \draw[thick,accent] (0.8,1.3)--(3.4,2.0);
    \node[lab] at (2.1,-0.6){no interaction (parallel)};
  \end{scope}
  \begin{scope}[xshift=6cm]
    \draw[axis] (0,0)--(4.4,0) node[right,lab]{A}; \draw[axis] (0,0)--(0,3.4) node[above,lab]{response};
    \draw[thick,accent2] (0.8,0.8)--(3.4,3.0); \draw[thick,accent] (0.8,1.3)--(3.4,1.7);
    \node[lab] at (2.1,-0.6){interaction (lines cross or diverge)};
  \end{scope}
\end{tikzpicture}
\caption{Interaction plots: one line per level of the second factor. Non-parallel lines indicate interaction.}
\end{figure}

\section{Two-Factor Factorial Experiments}

\subsection{Statistical Analysis of the Fixed-Effects Model}

Factor A has $a$ levels and factor B has $b$ levels. Each of the $ab$ combinations is replicated $n$ times, in random order.

\begin{definition}[Two-factor model with interaction]
\[ Y_{ijk}=\mu+\tau_i+\beta_j+(\tau\beta)_{ij}+\varepsilon_{ijk},\qquad i=1,\dots,a;\ j=1,\dots,b;\ k=1,\dots,n. \]
\begin{itemize}
  \item $\sum\tau_i=\sum\beta_j=0$.
  \item $\sum_i(\tau\beta)_{ij}=\sum_j(\tau\beta)_{ij}=0$.
  \item $\varepsilon_{ijk}\sim\Normal(0,\sigma^2)$, independent.
\end{itemize}
\end{definition}

\begin{theorem}[ANOVA partition]\label{thm:twofactor}
Let $\bar{y}_{ij\cdot}$ be cell means, $\bar{y}_{i\cdot\cdot}$ row means, $\bar{y}_{\cdot j\cdot}$ column means and $\bar{y}_{\cdot\cdot\cdot}$ the grand mean. Then $SS_T=SS_A+SS_B+SS_{AB}+SS_E$, where
\begin{align*}
SS_A&=bn\sum_i(\bar{y}_{i\cdot\cdot}-\bar{y}_{\cdot\cdot\cdot})^2,\\
SS_B&=an\sum_j(\bar{y}_{\cdot j\cdot}-\bar{y}_{\cdot\cdot\cdot})^2,\\
SS_{AB}&=n\sum_{i,j}(\bar{y}_{ij\cdot}-\bar{y}_{i\cdot\cdot}-\bar{y}_{\cdot j\cdot}+\bar{y}_{\cdot\cdot\cdot})^2,\\
SS_E&=\sum_{i,j,k}(y_{ijk}-\bar{y}_{ij\cdot})^2.
\end{align*}
The degrees of freedom add as $abn-1=(a-1)+(b-1)+(a-1)(b-1)+ab(n-1)$.
\end{theorem}

\begin{theorem}[$F$-tests (fixed effects)]\label{thm:twofactor-F}
Each ratio $F_A=MS_A/MS_E$, $F_B=MS_B/MS_E$, $F_{AB}=MS_{AB}/MS_E$ has an $F$ distribution under its null hypothesis.
\begin{itemize}
  \item Nulls: all $\tau_i=0$; all $\beta_j=0$; all $(\tau\beta)_{ij}=0$.
  \item Numerator degrees of freedom: $a-1$, $b-1$, $(a-1)(b-1)$.
  \item Denominator degrees of freedom: $ab(n-1)$.
\end{itemize}
Reject each null in the upper tail.
\end{theorem}

\begin{center}
\begin{tabular}{lcccc}
\toprule
Source & SS & df & MS & $F_0$\\
\midrule
A & $SS_A$ & $a-1$ & $MS_A$ & $MS_A/MS_E$\\
B & $SS_B$ & $b-1$ & $MS_B$ & $MS_B/MS_E$\\
AB & $SS_{AB}$ & $(a-1)(b-1)$ & $MS_{AB}$ & $MS_{AB}/MS_E$\\
Error & $SS_E$ & $ab(n-1)$ & $MS_E$ & \\
Total & $SS_T$ & $abn-1$ & & \\
\bottomrule
\end{tabular}
\end{center}

Computing formulas:
\begin{align*}
SS_T&=\sum y_{ijk}^2-\frac{y_{\cdots}^2}{abn},\\
SS_A&=\frac1{bn}\sum y_{i\cdot\cdot}^2-\frac{y_{\cdots}^2}{abn},\\
SS_B&=\frac1{an}\sum y_{\cdot j\cdot}^2-\frac{y_{\cdots}^2}{abn},\\
SS_{AB}&=\frac1n\sum\sum y_{ij\cdot}^2-\frac{y_{\cdots}^2}{abn}-SS_A-SS_B,\\
SS_E&=SS_T-SS_{AB}-SS_A-SS_B.
\end{align*}
Test the interaction first, then the main effects.

\begin{example}\label{ex:two-factor}
Factor A has $2$ levels, factor B has $3$, and $n=2$:
\begin{center}
\begin{tabular}{lccc}
\toprule
 & B1 & B2 & B3\\
\midrule
A1 & 10, 12 & 14, 16 & 20, 18\\
A2 & 11, 13 & 18, 20 & 28, 26\\
\bottomrule
\end{tabular}
\end{center}
\begin{itemize}
  \item Totals: A1 $=90$, A2 $=116$; B1 $=46$, B2 $=68$, B3 $=92$. Grand total $206$.
  \item Correction $206^2/12=3536.33$ and $\sum y^2=3894$, so $SS_T=357.67$.
  \item $SS_A=(90^2+116^2)/6-3536.33=56.33$.
  \item $SS_B=(46^2+68^2+92^2)/4-3536.33=264.67$.
  \item Cell totals $22,30,38,24,38,54$ give $SS_{\mathrm{cells}}=(22^2+30^2+38^2+24^2+38^2+54^2)/2-3536.33=345.67$.
  \item So $SS_{AB}=345.67-56.33-264.67=24.67$ and $SS_E=357.67-345.67=12.0$.
  \item Mean squares $56.33,\ 132.33,\ 12.33,\ 2.0$ on $1,2,2,6$ df.
  \item $F_A=28.2>f_{0.05,1,6}=5.99$. $F_B=66.2>f_{0.05,2,6}=5.14$. $F_{AB}=6.17>5.14$.
\end{itemize}
Both main effects and the interaction are significant.
\end{example}

\begin{example}[Primer paint adhesion]\label{ex:primer}
Three primer types and two application methods (dipping, spraying), $n=3$ specimens per cell. Cell totals: primer~1 dip $12.8$, spray $15.9$; primer~2 $15.9$, $18.2$; primer~3 $11.5$, $15.5$. The grand total is $89.8$ over $18$ observations.
\begin{itemize}
  \item $SS_T=10.72$, $SS_{\text{types}}=4.58$, $SS_{\text{methods}}=4.91$, $SS_{\text{interaction}}=0.24$. By subtraction $SS_E=0.99$, with $12$ degrees of freedom ($MS_E=0.08$).
  \item $f_{\text{types}}=2.29/0.08=28.63>f_{0.05,2,12}=3.89$. Significant.
  \item $f_{\text{methods}}=4.91/0.08=61.38>f_{0.05,1,12}=4.75$. Significant.
  \item $f_{AB}=0.12/0.08=1.50<3.89$ ($P=0.262$). No evidence of interaction. The interaction plot has nearly parallel lines.
\end{itemize}
Spraying gives higher adhesion, and primer~2 the highest. With no interaction, compare the primer row averages (for example by Fisher's LSD, as in Chapter~7). If the interaction were significant, compare the means of one factor with the other factor held at a fixed level.
\end{example}

\begin{pitfall}{interpret the interaction first}
If the interaction is significant, the main effects of A and B \emph{cannot be read separately}. The effect of A depends on B.
\begin{itemize}
  \item In Example~\ref{ex:two-factor}, the A effect grows from $+1$ at B1 to $+8$ at B3.
  \item Compare the cell means (an interaction plot), not the marginal means.
  \item Make multiple comparisons among cell means, at fixed levels of the other factor.
  \item If the interaction is not significant, report the main effects and compare marginal means with Tukey's method.
\end{itemize}
\end{pitfall}

\subsection{Model Adequacy Checking}

The residuals are $e_{ijk}=y_{ijk}-\bar{y}_{ij\cdot}$. Check them as in Chapter~7:
\begin{itemize}
  \item a normal probability plot;
  \item residuals against fitted values $\bar{y}_{ij\cdot}$ (the spread should be equal across cells);
  \item residuals against factor levels and against run order.
\end{itemize}
Unequal variances across cells are the most damaging violation. If the spread depends on the mean, transform the response.

\subsection{One Observation Per Cell}

With $n=1$ there is no replication. $SS_E$ in the model above is zero, and the interaction cannot be separated from error.

If it is reasonable to assume \emph{no interaction}:
\begin{itemize}
  \item use $SS_{AB}$ (with $(a-1)(b-1)$ df) as the error sum of squares;
  \item test with $F_A=MS_A/MS_{AB}$ and $F_B=MS_B/MS_{AB}$.
\end{itemize}
This is the randomized block analysis of Chapter~7, with the second factor playing the role of blocks. Tukey's one-degree-of-freedom test for nonadditivity checks the no-interaction assumption.

\subsection{Factorial Experiments with Random Factors: Overview}

If a factor is random (its levels are a random sample), its effects are random and the $F$-ratios change.

\textbf{Both A and B random:}
\begin{align*}
\Exp[MS_A]&=\sigma^2+n\sigma_{\tau\beta}^2+bn\sigma_\tau^2,\\
\Exp[MS_B]&=\sigma^2+n\sigma_{\tau\beta}^2+an\sigma_\beta^2,\\
\Exp[MS_{AB}]&=\sigma^2+n\sigma_{\tau\beta}^2,\\
\Exp[MS_E]&=\sigma^2.
\end{align*}
\begin{itemize}
  \item So $F_A=MS_A/MS_{AB}$ and $F_B=MS_B/MS_{AB}$. The interaction mean square is the denominator.
  \item $F_{AB}=MS_{AB}/MS_E$.
  \item Variance components come from solving these equations, for example $\hat\sigma_\tau^2=(MS_A-MS_{AB})/bn$.
\end{itemize}
\textbf{Mixed model} (A fixed, B random): the test for A uses $MS_{AB}$, and the test for B uses $MS_E$.

\emph{Rule:} form each $F$-ratio so that numerator and denominator have the same expected value under the null hypothesis.

\section{General Factorial Experiments}

\begin{definition}[Three-factor model]
For factors A, B, C with $a,b,c$ levels and $n$ replicates, the model has three main effects, three two-factor interactions $AB,AC,BC$, and the three-factor interaction $ABC$:
\[ Y_{ijkl}=\mu+\tau_i+\beta_j+\gamma_k+(\tau\beta)_{ij}+(\tau\gamma)_{ik}+(\beta\gamma)_{jk}+(\tau\beta\gamma)_{ijk}+\varepsilon_{ijkl}. \]
\end{definition}

The analysis extends the two-factor case.
\begin{itemize}
  \item Each sum of squares comes from the matching marginal totals.
  \item Degrees of freedom: $a-1$, $b-1$, $c-1$, $(a-1)(b-1)$, $(a-1)(c-1)$, $(b-1)(c-1)$, $(a-1)(b-1)(c-1)$. Error: $abc(n-1)$.
  \item For fixed effects, $F=MS_{\text{effect}}/MS_E$.
  \item Higher-order interactions are rarely large. Test the highest-order interaction first.
  \item At least two replicates ($n\ge2$) are needed to compute $SS_E$.
\end{itemize}

\begin{example}[Surface roughness]
Feed rate (A), depth of cut (B) and tool angle (C), each at two levels, with $n=2$. That gives $16$ observations and $8$ error degrees of freedom. The critical value for every effect is $f_{0.05,1,8}=5.32$.
\begin{itemize}
  \item Feed rate has a strong main effect.
  \item Depth of cut has $P=0.0710$. Only suggestive.
  \item The next largest effect is the $AB$ (feed $\times$ depth) interaction.
\end{itemize}
Both feed rate and depth of cut are probably important process variables.
\end{example}

The number of runs $abc\,n$ grows fast with the number of factors. This motivates the $2^k$ designs.

\section{Factorial Experiments with Random Factors (CD Only)}

For random factors, the expected mean squares in the general model decide the test statistics, as in the two-factor case.
\begin{itemize}
  \item Each effect is tested against the mean square whose expectation has the same variance components under the null.
  \item Sometimes no single mean square fits. Then use a \emph{linear combination} of mean squares (Satterthwaite's approximation), with approximate degrees of freedom.
  \item The ANOVA table uses the same sums of squares as for fixed effects. Only the denominators of the $F$-ratios change, and the aim becomes estimating variance components rather than comparing level means.
\end{itemize}

\section{$2^k$ Factorial Designs}

When each of $k$ factors has just two levels (``low'' $-$ and ``high'' $+$), the full factorial has $2^k$ runs per replicate. These designs are the workhorse of screening.

\subsection{The $2^2$ Design}

With factors A and B the four treatment combinations are:
\begin{itemize}
  \item $(1)$: both low;
  \item $a$: A high, B low;
  \item $b$: A low, B high;
  \item $ab$: both high.
\end{itemize}
The same symbol also stands for the total of the $n$ observations at that combination.

\begin{theorem}[Effects in the $2^2$ design]\label{thm:2k2}
\[ A=\frac1{2n}\big[ab+a-b-(1)\big],\quad B=\frac1{2n}\big[ab+b-a-(1)\big],\quad AB=\frac1{2n}\big[ab+(1)-a-b\big]. \]
The bracket is the effect's \emph{contrast} $C$. Each effect has one degree of freedom and $SS=C^2/(4n)$.

Also $SS_T=\sum\sum\sum y^2-y_{\cdot\cdot\cdot}^2/(4n)$ and $SS_E=SS_T-SS_A-SS_B-SS_{AB}$, with $4(n-1)$ degrees of freedom.
\end{theorem}

The coefficients of the contrasts form the \emph{sign table}. Each effect is the average response where its sign is $+$ minus the average where it is $-$.

\begin{center}
\begin{tabular}{l|cccc}
\toprule
Effect & $(1)$ & $a$ & $b$ & $ab$\\
\midrule
A & $-$ & $+$ & $-$ & $+$\\
B & $-$ & $-$ & $+$ & $+$\\
AB & $+$ & $-$ & $-$ & $+$\\
\bottomrule
\end{tabular}
\end{center}

\begin{example}\label{ex:22}
A $2^2$ design with $n=3$: $(1)$: $28,25,27$ (total $80$); $a$: $36,32,32$ ($100$); $b$: $18,19,23$ ($60$); $ab$: $31,30,29$ ($90$).
\begin{itemize}
  \item Contrasts: $A=ab+a-b-(1)=50$, $B=ab+b-a-(1)=-30$, $AB=ab+(1)-a-b=10$.
  \item Effects (divide by $2n=6$): $A=8.33$, $B=-5.0$, $AB=1.67$.
  \item $SS_A=50^2/12=208.3$, $SS_B=30^2/12=75.0$, $SS_{AB}=10^2/12=8.33$.
  \item $y_{\cdot\cdot\cdot}=330$ and $\sum y^2=9398$, so $SS_T=9398-330^2/12=323.0$. Then $SS_E=323.0-291.7=31.3$ on $8$ df, $MS_E=3.92$.
  \item $F_A=53.2$, $F_B=19.2$, $F_{AB}=2.13$, against $f_{0.05,1,8}=5.32$.
\end{itemize}
The main effects are significant. The interaction is not.
\end{example}

\begin{example}[Epitaxial layer]
A $2^2$ design with $n=4$: deposition time (A) and arsenic flow rate (B). The response is layer thickness in $\mu$m.
\begin{itemize}
  \item Changing deposition time from low to high changes the mean thickness by $0.836\,\mu$m. B and AB are small.
  \item The ANOVA confirms that only deposition time is significant.
  \item The fitted model is $\hat{y}=\bar{y}+\tfrac12(0.836)x_1$ with $x_1=\pm1$. Residuals are $y-\hat{y}$ at each run. One large residual flags a possible outlier.
\end{itemize}
\end{example}

Facts to remember:
\begin{itemize}
  \item In a $2^k$ design, the square of the software's $t$-statistic for an effect equals that effect's $F$-ratio ($t^2=F$, on $1$ and $2^k(n-1)$ df).
  \item The design is orthogonal, so sequential and adjusted sums of squares coincide.
  \item \textbf{Regression view.} Code the factors $x_A,x_B=\pm1$. Then $\hat{y}=\hat\beta_0+\hat\beta_Ax_A+\hat\beta_Bx_B+\hat\beta_{AB}x_Ax_B$, with $\hat\beta_0=\bar{y}$. Each coefficient is \emph{half} the effect, because $x$ moves by $2$ units from $-1$ to $+1$.
  \item In Example~\ref{ex:22}: $\hat{y}=27.5+4.17x_A-2.5x_B+0.83x_Ax_B$.
  \item The columns are orthogonal, so the estimates are uncorrelated and each has variance $\sigma^2/(4n)$.
\end{itemize}

\subsection{The $2^k$ Design for $k\ge3$ Factors}

For $k=3$ there are $8$ combinations: $(1),a,b,ab,c,ac,bc,abc$. The seven effects are $A,B,C,AB,AC,BC,ABC$.

\begin{theorem}[Effects and sums of squares in the $2^k$ design]\label{thm:2k}
Let $C$ be an effect's contrast: the signed sum of treatment totals, with signs from the sign table. (A letter gets $+$ if the factor is high. Interaction signs are products.) Then
\[ \text{effect}=\frac{C}{2^{k-1}n},\qquad SS=\frac{C^2}{2^kn},\qquad\text{each with one degree of freedom.} \]
The regression coefficient is half the effect. Total degrees of freedom: $2^kn-1=(2^k-1)+2^k(n-1)$.
\end{theorem}

\begin{center}
\begin{tabular}{l|cccccccc}
\toprule
 & $(1)$ & $a$ & $b$ & $ab$ & $c$ & $ac$ & $bc$ & $abc$\\
\midrule
A & $-$ & $+$ & $-$ & $+$ & $-$ & $+$ & $-$ & $+$\\
B & $-$ & $-$ & $+$ & $+$ & $-$ & $-$ & $+$ & $+$\\
C & $-$ & $-$ & $-$ & $-$ & $+$ & $+$ & $+$ & $+$\\
AB & $+$ & $-$ & $-$ & $+$ & $+$ & $-$ & $-$ & $+$\\
AC & $+$ & $-$ & $+$ & $-$ & $-$ & $+$ & $-$ & $+$\\
BC & $+$ & $+$ & $-$ & $-$ & $-$ & $-$ & $+$ & $+$\\
ABC & $-$ & $+$ & $+$ & $-$ & $+$ & $-$ & $-$ & $+$\\
\bottomrule
\end{tabular}
\end{center}

Example: the A effect in $2^3$ is $\dfrac1{4n}\big[a+ab+ac+abc-(1)-b-c-bc\big]$.

The columns of the sign table are orthogonal. So each effect is estimated independently, with equal precision, and the sums of squares add up to $SS_T-SS_E$.

\subsection{Single Replicate of the $2^k$ Design}

The number of runs grows as $2^k$. For $k\ge4$ one often runs a \emph{single replicate} ($n=1$). Then there are no degrees of freedom for error. Two fixes:
\begin{itemize}
  \item \textbf{Normal probability plot of the effects.} Effects of inactive factors are random draws from one normal distribution, so they fall on a straight line. Active effects stand off the line. No estimate of $\sigma$ is needed.
  \item \textbf{Pooling.} By the \emph{sparsity of effects principle}, high-order interactions (three-factor and above) are usually negligible. Pool their sums of squares into an error estimate and $F$-test the rest.
\end{itemize}
Then refit with only the significant terms and check the residuals.

\emph{Projection.} Dropping $h$ negligible factors from a $2^k$ design with $n$ replicates leaves a $2^{k-h}$ design with $n2^h$ replicates. Dropping factor $C$ collapses the cube onto a square, with each corner run twice as often.

Counting effects: a $2^4$ has $4$ main effects, $6$ two-factor, $4$ three-factor and $1$ four-factor interaction. A $2^6$ has $6,15,20,15,6,1$.

\subsection{Addition of Center Points to a $2^k$ Design (CD Only)}

Add $n_C$ replicated runs at the \emph{center point} ($x_i=0$ for all factors). They give:
\begin{itemize}
  \item an estimate of pure error from the center runs;
  \item a test for \emph{curvature}. If the true response is a plane, the center average equals the average of the factorial runs.
\end{itemize}
Let $\bar{y}_F$ be the mean of the $n_F$ factorial runs and $\bar{y}_C$ the mean of the $n_C$ center runs. Then
\[ SS_{\text{curvature}}=\frac{n_Fn_C(\bar{y}_F-\bar{y}_C)^2}{n_F+n_C},\qquad F_0=\frac{SS_{\text{curvature}}}{MS_E}\quad(1,\ n_C-1\text{ df}), \]
where $MS_E$ is the pure-error variance from the center runs. A significant value means the linear model is inadequate and quadratic terms are needed. Center points do not disturb the estimates of main effects and interactions.

\section{Blocking and Confounding in the $2^k$ Design}

A full replicate of $2^k$ may be too large for one homogeneous block (one batch of raw material, one day of production). \emph{Confounding} splits it into smaller blocks. It deliberately makes one effect, usually the highest-order interaction, indistinguishable from the block difference.
\begin{itemize}
  \item That effect is sacrificed.
  \item All other effects are estimated free of block differences.
\end{itemize}

\begin{definition}[Confounding scheme for two blocks]
To confound an effect, for example $ABC$ in a $2^3$, let $L=x_1+x_2+x_3$, where $x_i=0$ (low) or $1$ (high). Put runs with $L\equiv0\pmod2$ in one block, and runs with $L\equiv1\pmod2$ in the other.
\end{definition}

\begin{example}
$2^3$ confounded with $ABC$.
\begin{itemize}
  \item Block 1: $(1),\ ab,\ ac,\ bc$ (an even number of high factors).
  \item Block 2: $a,\ b,\ c,\ abc$.
  \item Every main-effect and two-factor contrast has two $+$ and two $-$ in each block, so the block difference cancels.
  \item The $ABC$ contrast is $+$ for all of block 2 and $-$ for all of block 1. It cannot be separated from the block effect.
\end{itemize}
If $ABC$ is negligible, as is usual, nothing of value is lost.
\end{example}

\textbf{More blocks.} Confound several effects.
\begin{itemize}
  \item With four blocks, choose two defining effects, say $AB$ and $AC$. Their product mod 2, $BC$, is confounded too (the \emph{generalised interaction}). Three effects are lost.
  \item In general, $2^p$ blocks need $p$ independent confounded effects and sacrifice $2^p-1$ effects.
  \item Construction: with $L=\alpha_1x_1+\dots+\alpha_kx_k\pmod2$ (where $x_i\in\{0,1\}$ is the level of factor $i$ and $\alpha_i$ its exponent in the confounded effect), runs with the same $L$ go in the same block.
  \item The block containing $(1)$ is the \emph{principal block}. Get the other by multiplying one of its elements into each element of the principal block, mod 2 on the exponents.
  \item For a $2^2$ with $AB$ confounded, the blocks are $\{(1),ab\}$ and $\{a,b\}$. $A$ and $B$ are unaffected, since each contrast has one plus and one minus run in each block. The block effect and $AB$ are identical.
\end{itemize}
In the textbook's missile example, a $2^4$ is run in two blocks (two gunners) with $ABCD$ confounded. The three-factor interactions are pooled as error. Target type (A), target range (D) and the $AC$ and $AD$ interactions are large.

\emph{Partial confounding.} Replicate the design with a \emph{different} effect confounded in each replicate. Then every effect is estimable from some replicates.

\section{Fractional Replication of the $2^k$ Design}

Even a single replicate may be too costly when $k$ is large ($2^7=128$ runs). High-order interactions are usually negligible, so we can run only a \emph{fraction}.

\subsection{One-Half Fraction of the $2^k$ Design}

Take the runs of the $2^3$ for which $ABC$ is $+$: $a,b,c,abc$. This is the \emph{principal fraction} of a $2^{3-1}$ design, with \emph{defining relation} $I=ABC$.

\begin{definition}[Alias]
Two effects are \emph{aliases} if their contrasts are identical in the fraction, so they cannot be separated. Find aliases by multiplying the effect by the defining relation, with $A^2=B^2=C^2=I$ (powers reduce mod 2).
\end{definition}

For $I=ABC$: $A\cdot ABC=BC$, so $A=BC$. Likewise $B=AC$ and $C=AB$.
\begin{itemize}
  \item The estimate of ``A'' is really an estimate of $A+BC$. If $BC$ is negligible, it estimates $A$.
  \item The \emph{complementary fraction} $I=-ABC$ (runs $(1),ab,ac,bc$) estimates $A-BC$, and so on.
  \item Running both recovers the full factorial. The sum and difference separate the aliases.
\end{itemize}

\begin{definition}[Resolution]
The \emph{resolution} of a fractional design is the length of the shortest word in its defining relation.
\begin{itemize}
  \item \textbf{Resolution III} (for example $2^{3-1}_{\mathrm{III}}$, $I=ABC$): main effects are aliased with two-factor interactions.
  \item \textbf{Resolution IV} (for example $2^{4-1}_{\mathrm{IV}}$, $I=ABCD$): main effects are clear of two-factor interactions. Two-factor interactions are aliased with each other ($AB=CD$).
  \item \textbf{Resolution V}: main effects and two-factor interactions are clear of each other. They are aliased only with three-factor and higher.
\end{itemize}
\end{definition}

\begin{example}
A $2^{4-1}_{\mathrm{IV}}$ design (eight runs) with $I=ABCD$.
\begin{itemize}
  \item Aliases: $A=BCD$, $B=ACD$, $C=ABD$, $D=ABC$, $AB=CD$, $AC=BD$, $AD=BC$.
  \item Construction: write the full $2^3$ in A, B, C and set $D=ABC$ (multiply the signs).
  \end{itemize}
Four factors are studied in half the runs. Main effects are estimable if three-factor interactions are negligible.
\end{example}

Textbook examples:
\begin{itemize}
  \item \textbf{Plasma etch.} The $2^{4-1}$ with $I=ABCD$ shows large $A$ (gap), $D$ (power) and $AD$ effects, agreeing with the full factorial. Dropping $B$ and $C$ leaves a replicated $2^2$ in $A$ and $D$.
  \item \textbf{Unresolved aliases.} Run the alternate fraction ($I=-ABC$) and combine. With $\ell$ and $\ell'$ the two estimates, $\tfrac12(\ell+\ell')$ estimates $A$ and $\tfrac12(\ell-\ell')$ estimates $BC$.
  \item \textbf{Bad generator.} Choosing $I=AB$ in a $2^{3-1}$ aliases two main effects and loses important information.
  \item \textbf{Injection molding.} A $2^{6-2}$ design (16 runs) shows large $A$ (mold temperature), $B$ (screw speed) and $AB$. Residuals against holding time $C$ show less scatter at low $C$. That points to a setting that reduces variability.
\end{itemize}

\begin{interview}{why fractions work}
Three principles justify fractions:
\begin{itemize}
  \item \emph{Sparsity of effects}: few effects are large.
  \item \emph{Hierarchy}: main effects matter more than two-factor interactions, which matter more than three-factor interactions.
  \item \emph{Projection}: a resolution $R$ design projects onto a full factorial in any $R-1$ factors.
\end{itemize}
If one factor is inactive, the $2^{4-1}$ collapses to a full $2^3$ in the others. No information is wasted.
\end{interview}

\subsection{Smaller Fractions: The $2^{k-p}$ Fractional Factorial}

A $2^{k-p}$ design has $2^{k-p}$ runs, a $1/2^p$ fraction.
\begin{itemize}
  \item Start from a full $2^{k-p}$ design in $k-p$ \emph{basic} factors.
  \item Assign the other $p$ factors to interactions of the basic factors, using $p$ independent \emph{generators}.
  \item The defining relation holds the $p$ generator words and all their $2^p-1$ products.
  \item Each effect has $2^p-1$ aliases. The resolution is the length of the shortest word in the complete defining relation.
\end{itemize}

\begin{example}[A $2^{6-2}$ design]
Start from a $2^4$ in $A,B,C,D$, with generators $E=ABC$ and $F=BCD$. So $I=ABCE=BCDF$. Multiplying, $ABCE\cdot BCDF=ADEF$. The full relation is $I=ABCE=BCDF=ADEF$.
\begin{itemize}
  \item The shortest word has $4$ letters, so this is resolution IV.
  \item Main effects are aliased with three-factor interactions, for example $A=BCE=DEF=ABCDF$.
  \item Two-factor interactions are aliased with each other: $AB=CE$, $AC=BE$, $AD=EF$, $AF=DE$, \dots
\end{itemize}
\end{example}

Other designs:
\begin{itemize}
  \item $2^{7-3}_{\mathrm{IV}}$ (16 runs): $E=ABC$, $F=BCD$, $G=ACD$, with $I=ABCE=BCDF=ADEF=ACDG=BDEG=CEFG=ABFG$.
  \item $2^{7-4}_{\mathrm{III}}$ (eight runs for seven factors): generators $D=AB$, $E=AC$, $F=BC$, $G=ABC$. It is \emph{saturated}: all seven degrees of freedom estimate main effects. Its $15$-word defining relation includes $ABD=ACE=BCF=ABCG$ and their products.
\end{itemize}

Practical advice:
\begin{itemize}
  \item Choose generators to maximise the resolution. Tables of best designs exist.
  \item Resolution III is for pure screening (assume interactions are negligible). Use resolution IV or V when interactions matter.
  \item \emph{Foldover} (run the design again with all signs reversed) breaks the aliasing of main effects with two-factor interactions. It turns resolution III into resolution IV.
  \item Use fractions sequentially: begin with a small screening fraction, then add runs to resolve ambiguities.
\end{itemize}

\section{Response Surface Methods and Designs (CD Only)}

Once the important factors are known, the goal is often to find the settings that optimise the response. \emph{Response surface methodology} (RSM) fits a low-order polynomial in the coded factors and moves toward the optimum.

\begin{itemize}
  \item \textbf{Steepest ascent.} Far from the optimum, a first-order model $\hat{y}=\hat\beta_0+\sum\hat\beta_ix_i$ (from a $2^k$ with center points) is adequate.
  \begin{itemize}
    \item The gradient $(\hat\beta_1,\dots,\hat\beta_k)$ gives the path of steepest ascent (descent, for a minimum).
    \item Run experiments along it until the response stops improving. Then repeat with a new factorial centred there.
  \end{itemize}
  \item \textbf{Second-order model.} Near the optimum, curvature appears (center-point test). Fit $\hat{y}=\hat\beta_0+\sum\hat\beta_ix_i+\sum\hat\beta_{ii}x_i^2+\sum_{i<j}\hat\beta_{ij}x_ix_j$. This needs at least three levels of each factor. The \emph{central composite design} (CCD) adds $2k$ axial points $(\pm\alpha,0,\dots,0),\dots$ and center runs to a $2^k$ factorial. The Box--Behnken design is an alternative.
  \item \textbf{Canonical analysis.} Write $\hat{y}=\hat\beta_0+\mathbf{x}\T\mathbf{b}+\mathbf{x}\T\hat{B}\mathbf{x}$. The stationary point is $\mathbf{x}_s=-\tfrac12\hat{B}^{-1}\mathbf{b}$. The eigenvalues of the symmetric matrix $\hat{B}$ classify it: all negative, a maximum; all positive, a minimum; mixed signs, a saddle (Linear Algebra, Chapter~16).
\end{itemize}

\begin{summarybox}
\begin{itemize}
  \item Factorial designs vary all factors together. They estimate main effects \emph{and} interactions efficiently. If an interaction is significant, interpret cell means, not marginal means.
  \item Two-factor ANOVA: $SS_T=SS_A+SS_B+SS_{AB}+SS_E$, with df $a-1$, $b-1$, $(a-1)(b-1)$, $ab(n-1)$. Fixed effects are tested against $MS_E$. One observation per cell requires assuming no interaction. For random factors, choose denominators from the expected mean squares.
  \item $2^k$: effect $=C/(2^{k-1}n)$, $SS=C^2/(2^kn)$, regression coefficient $=$ half the effect. Single replicates use normal plots of effects or pooled high-order interactions. Center points test curvature.
  \item Blocking by confounding sacrifices high-order interactions (for example $ABC$ in two blocks). Generalised interactions are confounded too.
  \item Fractions $2^{k-p}$ use defining relations. Find aliases by multiplying. Resolution III/IV/V says which effects are aliased. Foldover and sequential experiments resolve ambiguities.
  \item Response surface methods: steepest ascent, central composite designs, canonical analysis.
\end{itemize}
\end{summarybox}

\section{Exercises}

\begin{exercise}
Analyse the following two-factor experiment ($a=3$, $b=2$, $n=2$): A1B1: $8,10$; A1B2: $12,14$; A2B1: $9,11$; A2B2: $20,22$; A3B1: $10,12$; A3B2: $15,17$. Construct the ANOVA table, test for interaction and main effects at $\alpha=0.05$, and draw the interaction plot.
\end{exercise}

\begin{exercise}
In a $2^2$ design with $n=2$: $(1)=\{20,22\}$, $a=\{30,34\}$, $b=\{18,20\}$, $ab=\{28,32\}$. Compute the three effects, their sums of squares, the ANOVA table and the fitted regression equation.
\end{exercise}

\begin{exercise}
Write the contrast for the effect $AC$ in a $2^3$ design and verify that it is orthogonal to the contrast for $B$.
\end{exercise}

\begin{exercise}
Show that in a $2^k$ design with $n$ replicates each effect estimate has variance $\sigma^2/(2^{k-2}n)$, and that the regression coefficient has variance $\sigma^2/(2^kn)$.
\end{exercise}

\begin{exercise}
Design a $2^3$ experiment in two blocks of four runs confounding $ABC$, and one in four blocks of two runs confounding $AB$ and $AC$; list the blocks and the confounded effects.
\end{exercise}

\begin{exercise}
For a $2^{5-1}$ design with $I=ABCDE$, find the resolution, the aliases of $A$ and of $AB$, and explain what can be estimated if three-factor and higher interactions are negligible.
\end{exercise}

\begin{exercise}
Construct the eight runs of the $2^{4-1}$ design with $D=ABC$, list the aliases of every main effect and two-factor interaction, and say why it is called resolution IV.
\end{exercise}

\begin{exercise}
A first-order fit from a $2^2$ design with center points is $\hat{y}=50+3x_1-1.5x_2$ (coded units). Give the direction of steepest ascent and the coordinates of the first three points along it with step $\Delta x_1=1$.
\end{exercise}


\part{Nonparametric Statistics and Quality Control}
\chapter{Nonparametric Statistics}

\section{Introduction}

Most procedures in Chapters~2--10 assume the data come from a \emph{normal} population. These are \emph{parametric} methods: they rest on one parametric family.

The $t$- and $F$-procedures tolerate slight non-normality. When the population is clearly non-normal, or the data are ranks or categories, use \emph{nonparametric} (\emph{distribution-free}) methods. They:
\begin{itemize}
  \item assume little beyond a continuous distribution;
  \item work with signs and ranks;
  \item keep the stated significance level for many different distributions.
\end{itemize}

\begin{itemize}
  \item \textbf{Advantages:} fewer assumptions; usable for ranked or categorical data; quick to compute; robust to outliers.
  \item \textbf{Disadvantage:} they do not use all the information in the sample. When the population really is normal, they are less efficient than the parametric test and need somewhat larger samples for the same power.
\end{itemize}
Rule of thumb: if the parametric assumptions are justified, use the parametric test. If they are doubtful or cannot be checked, use the nonparametric one.

Many nonparametric methods analyse \emph{ranks}: the positions $1,2,\dots,N$ in the ordered sample.

\section{Sign Test}

\subsection{Description of the Test}

The \emph{sign test} tests the \emph{median} $\tilde\mu$ of a continuous distribution, the value with $\Prob(X\le\tilde\mu)=\Prob(X\ge\tilde\mu)=0.5$.
\begin{itemize}
  \item For a symmetric distribution (such as the normal) the median equals the mean. So the test can replace the $t$-test.
  \item It is valid for \emph{any} continuous distribution.
\end{itemize}

\textbf{Idea.} To test $H_0:\tilde\mu=\tilde\mu_0$, form the differences $X_i-\tilde\mu_0$. Under $H_0$, each difference is equally likely to be positive or negative. So the number of plus signs $R^+$ is binomial with $n$ trials and $p=\tfrac12$.

\begin{definition}[Sign test and its $P$-value]
Let $r^+$ be the number of positive differences and $r^-$ the number of negative ones. Differences equal to zero (\emph{ties}) are discarded, and $n$ is reduced, so $r^++r^-=n$. For $R\sim\mathrm{Bin}(n,\tfrac12)$:
\begin{itemize}
  \item $H_1:\tilde\mu<\tilde\mu_0$: $P=\Prob(R\le r^+)$.
  \item $H_1:\tilde\mu>\tilde\mu_0$: $P=\Prob(R\ge r^+)$.
  \item $H_1:\tilde\mu\ne\tilde\mu_0$: $P=2\Prob(R\ge r^+)$ if $r^+>n/2$, and $P=2\Prob(R\le r^+)$ if $r^+<n/2$.
\end{itemize}
Reject $H_0$ if $P\le\alpha$. Here $\Prob(R\ge r)=\sum_{k=r}^n\binom nk(0.5)^n$.
\end{definition}

\begin{example}\label{ex:sign}
A claim says the median breaking strength of a thread is $50$. Eleven non-tied measurements: $52,48,55,51,47,56,53,49,58,54,57$. (One measurement equal to $50$ was discarded.) Test $H_0:\tilde\mu=50$ against $H_1:\tilde\mu>50$ at $\alpha=0.05$.
\begin{enumerate}[label=\arabic*.]
  \item Parameter: the median breaking strength.
  \item $H_0:\tilde\mu=50$.
  \item $H_1:\tilde\mu>50$.
  \item $\alpha=0.05$.
  \item Statistic: the number of plus differences, $r^+$.
  \item Reject $H_0$ if $P\le0.05$.
  \item The signs are $+,-,+,+,-,+,+,-,+,+,+$, so $r^+=8$ and $n=11$. $P=\Prob(R\ge8)=(165+55+11+1)/2048=0.113$.
  \item Since $0.113>0.05$, do not reject $H_0$. There is no significant evidence that the median exceeds $50$.
\end{enumerate}
\end{example}

\subsubsection*{Critical values}

Tables avoid the binomial sum. Let $r=\min(r^+,r^-)$.
\begin{itemize}
  \item \textbf{Two-sided.} Reject $H_0$ if $r\le r^*_\alpha$. Here $r^*_\alpha$ is the largest value whose tail probability does not exceed $\alpha/2$ in each direction. The actual level is as close to $\alpha$ as the discrete distribution allows, without exceeding it.
  \item \textbf{One-sided.} Put all of $\alpha$ in one tail. For $H_1:\tilde\mu>\tilde\mu_0$, use $r^-$ (reject when the number of minus signs is small). For $H_1:\tilde\mu<\tilde\mu_0$, use $r^+$. Reject if that count is $\le r^*_{2\alpha}$ from the two-sided table. A two-sided level $2\alpha$ is a one-sided level $\alpha$.
\end{itemize}
Do not use the two-sided critical value at the same $\alpha$ for a one-sided test.

\begin{center}
\small
\begin{tabular}{c|cccccccccccccccc}
\toprule
$n$ & 6 & 7 & 8 & 9 & 10 & 11 & 12 & 13 & 14 & 15 & 16 & 17 & 18 & 19 & 20\\
\midrule
$r^*_{0.05}$ & 0 & 0 & 0 & 1 & 1 & 1 & 2 & 2 & 2 & 3 & 3 & 4 & 4 & 4 & 5\\
$r^*_{0.01}$ & -- & -- & 0 & 0 & 0 & 0 & 1 & 1 & 1 & 2 & 2 & 2 & 3 & 3 & 3\\
\bottomrule
\end{tabular}
\end{center}

\subsubsection*{The normal approximation}

For $n$ at least about $10$, $R$ is approximately normal with mean $0.5n$ and variance $0.25n$. So
\[ Z_0=\frac{R^+-0.5n}{0.5\sqrt n}\ \text{ is approximately }\Normal(0,1)\text{ under }H_0. \]
Reject if $\abs{z_0}>z_{\alpha/2}$ (two-sided), $z_0>z_\alpha$ ($H_1:\tilde\mu>\tilde\mu_0$) or $z_0<-z_\alpha$ ($H_1:\tilde\mu<\tilde\mu_0$).

For Example~\ref{ex:sign}: $z_0=(8-5.5)/(0.5\sqrt{11})=1.51$ and $P=0.066$. Same decision.

\begin{pitfall}{ties}
A difference of exactly zero has no sign. The population is assumed continuous, so ties should not occur. If rounding creates them, set them aside, reduce $n$, and apply the test to the rest.
\end{pitfall}

\begin{example}[Propellant shear strength]
Test $H_0:\tilde\mu=2000$ psi against $H_1:\tilde\mu\ne2000$ at $\alpha=0.05$, with $n=20$ motors. The data give $r^+=14$ and $r^-=6$.
\begin{itemize}
  \item \textbf{Exact.} Since $r^+>n/2=10$, $P=2\Prob(R\ge14)=2(0.0577)=0.1153>0.05$. Do not reject.
  \item \textbf{Table.} $r=\min(14,6)=6$. For $n=20$ and $\alpha=0.05$ (two-sided), $r^*=5$. Since $6>5$, do not reject.
  \item \textbf{Normal approximation.} $z_0=(14-10)/(0.5\sqrt{20})=1.789<z_{0.025}=1.96$. Do not reject.
\end{itemize}
All three agree.
\end{example}

\subsection{Sign Test for Paired Samples}

For pairs $(X_{1j},X_{2j})$ from two continuous populations, form $D_j=X_{1j}-X_{2j}$. Apply the sign test to the $n$ differences to test $H_0:\tilde\mu_D=0$ (the two populations have a common median).

\begin{example}
Ten workers are timed (minutes) on a task before and after training. The differences (before $-$ after) have $8$ positive and $2$ negative values.
\begin{itemize}
  \item Two-sided test: $r=\min(8,2)=2$. For $n=10$, $r^*_{0.05}=1$. Since $2>1$, do not reject $H_0$.
  \item Exact $P$-value: $2\Prob(R\le2)=2(1+10+45)/1024=0.109$.
\end{itemize}
\end{example}

\begin{example}[Fuel metering devices]
Twelve cars are each tested with two metering devices. Apply the test to the $12$ paired differences at $\alpha=0.05$. The table gives $r^*=2$. The data give $r=4>2$. So the devices are not shown to differ in median mileage.
\end{example}

\subsection{Type II Error for the Sign Test}

The sign test holds $\alpha$ for any continuous distribution. But $\beta$ depends on the shape of the true distribution. $\beta$ is a binomial probability whose success probability $p$ is the area of the true distribution on the relevant side of $\tilde\mu_0$.

\begin{example}
Take a normal population with $\sigma=1$. Test $H_0:\tilde\mu=0$ (two-sided, $n=12$, $r^*_{0.05}=2$). Suppose the true median is $1.5$.
\begin{itemize}
  \item A negative difference has probability $p=\Prob(X\le0)=\Phi(-1.5)=0.0668$.
  \item The test fails to reject when $R^-\ge3$. So
  \[ \beta=1-\sum_{x=0}^2\binom{12}x(0.0668)^x(0.9332)^{12-x}=1-(0.436+0.375+0.148)=0.042. \]
\end{itemize}
For a non-normal population (say exponential) with the same shift, $p$ and $\beta$ would be very different.
\end{example}

\subsection{Comparison to the $t$-Test}

\begin{itemize}
  \item For a normal population, the $t$-test has the smallest $\beta$ among all tests of the same size. It beats the sign test. The sign test's large-sample efficiency relative to the $t$-test is $2/\pi\approx0.64$.
  \item For symmetric non-normal populations, the $t$-test is still usually better unless the tails are very heavy.
  \item So use the sign test mainly as a quick test about a median, or when only signs are available. The Wilcoxon signed-rank test, which also uses magnitudes, is preferable.
\end{itemize}

\section{Wilcoxon Signed-Rank Test}

The sign test uses only the direction of each difference. The \emph{Wilcoxon signed-rank test} also uses the magnitudes. It applies to \emph{symmetric} continuous distributions, where mean $=$ median. So it tests the mean, $\mu=\mu_0$.

\subsection{Description of the Test}

To test $H_0:\mu=\mu_0$:
\begin{enumerate}
  \item Compute $X_i-\mu_0$.
  \item Rank the absolute differences $\abs{X_i-\mu_0}$ from $1$ (smallest) to $n$.
  \item Attach to each rank the sign of its difference.
\end{enumerate}

\begin{definition}[Signed-rank statistics]
\begin{itemize}
  \item $W^+$: the sum of the ranks with positive signs.
  \item $W^-$: the absolute value of the sum of the ranks with negative signs.
  \item $W=\min(W^+,W^-)$.
\end{itemize}
Since $W^++W^-=n(n+1)/2$, one determines the other.
\end{definition}

Under $H_0$, each rank is equally likely to carry $+$ or $-$. So $W^+$ and $W^-$ should be about equal. The test rejects when they differ a lot.
\begin{itemize}
  \item $H_1:\mu\ne\mu_0$: reject if $w\le w^*$ (table).
  \item $H_1:\mu>\mu_0$: reject if $w^-\le w^*$.
  \item $H_1:\mu<\mu_0$: reject if $w^+\le w^*$. (For one-sided tests use the one-sided significance levels.)
  \item Tied absolute differences get the average of the ranks they would have shared. Zero differences are discarded.
\end{itemize}

\begin{center}
\small
\begin{tabular}{c|cccccccccc}
\toprule
$n$ & 6 & 7 & 8 & 9 & 10 & 11 & 12 & 13 & 14 & 15\\
\midrule
$w^*_{0.05}$ & 0 & 2 & 3 & 5 & 8 & 10 & 13 & 17 & 21 & 25\\
$w^*_{0.10}$ & 2 & 3 & 5 & 8 & 10 & 13 & 17 & 21 & 25 & 30\\
\bottomrule
\end{tabular}
\end{center}
The table gives two-sided levels $0.05$ and $0.10$. One-sided levels are half of these ($0.025$ and $0.05$). So for a one-sided test at $\alpha=0.05$, use the column labelled two-sided $0.10$.

\begin{example}\label{ex:wilcoxon}
Eight pairs give the differences $d=1.2,\ -0.4,\ 2.0,\ 0.8,\ -0.6,\ 1.5,\ 0.9,\ 2.2$. Test $H_0:\mu_D=0$ against $H_1:\mu_D\ne0$ at $\alpha=0.05$.
\begin{itemize}
  \item Absolute values in order: $0.4,0.6,0.8,0.9,1.2,1.5,2.0,2.2$, with ranks $1,\dots,8$.
  \item The negative differences are $-0.4$ (rank $1$) and $-0.6$ (rank $2$). So $w^-=3$, $w^+=36-3=33$, and $w=3$.
  \item For $n=8$, $w^*_{0.05}=3$. Since $w=3\le3$, reject $H_0$. The median difference is not zero at the $5\%$ level.
\end{itemize}
\end{example}

\begin{example}[Textbook examples]
\begin{itemize}
  \item Propellant data ($n=20$, symmetric continuous distribution): $w^+=150$, $w^-=60$, so $w=60$. The two-sided critical value at $\alpha=0.05$ is $w^*=52$. Since $60>52$, do not reject. Same conclusion as the sign test.
  \item Paired fuel-metering data ($n=12$): $w=22.5$ against $w^*=13$. Do not reject.
\end{itemize}
\end{example}

\subsection{Large-Sample Approximation}

If $n\ge20$, $W^+$ (or $W^-$) is approximately normal with
\[ \mu_{W^+}=\frac{n(n+1)}4,\qquad\sigma^2_{W^+}=\frac{n(n+1)(2n+1)}{24}. \]
So $Z_0=\dfrac{W^+-n(n+1)/4}{\sqrt{n(n+1)(2n+1)/24}}$ is approximately standard normal under $H_0$. Use the usual critical regions.

\subsection{Paired Observations}

For pairs from two continuous populations that differ only in their means, rank $\abs{D_j}$ where $D_j=X_{1j}-X_{2j}$, attach signs, and proceed as above. This tests $H_0:\mu_1=\mu_2$ (equivalently $\mu_D=0$).
\begin{itemize}
  \item The \emph{difference} must be symmetric. $X_1$ and $X_2$ themselves need not be.
  \item Ties receive average ranks.
\end{itemize}

\subsection{Comparison to the $t$-Test}

For non-normal data, comparing $\beta$ is hard. So we use the \emph{asymptotic relative efficiency} (ARE): the limiting ratio of the sample sizes needed for the same error probabilities.
\begin{itemize}
  \item Normal populations: ARE of the signed-rank test against the $t$-test is about $0.95$.
  \item Non-normal populations: at least $0.86$, and often above $1$. (Heavy tails hurt the $t$-test because the sample mean is unstable.)
\end{itemize}
So the signed-rank test is never much worse than the $t$-test, and can be much better.

\section{Wilcoxon Rank-Sum Test}

Two independent random samples, of sizes $n_1\le n_2$, come from continuous populations with means $\mu_1,\mu_2$. The distributions are assumed to have the same shape and spread, differing at most in location. Test $H_0:\mu_1=\mu_2$.

\subsection{Description of the Test}

\begin{definition}[Rank-sum statistics]
Arrange all $n_1+n_2$ observations in ascending order and rank them $1,\dots,n_1+n_2$ (ties get average ranks). Let $W_1$ be the sum of the ranks in the \emph{smaller} sample (sample~1) and $W_2$ the sum in the other. Then
\[ W_1+W_2=\frac{(n_1+n_2)(n_1+n_2+1)}2. \]
\end{definition}

If the means are equal, the ranks of the two samples are well mixed, so $W_1$ and $W_2$ are in proportion to the sample sizes.
\begin{itemize}
  \item $H_1:\mu_1\ne\mu_2$: reject if $w_1$ or $w_2$ is at most the tabulated value $w_\alpha$ for $(n_1,n_2)$.
  \item $H_1:\mu_1>\mu_2$: reject if $w_2\le w_\alpha$.
  \item $H_1:\mu_1<\mu_2$: reject if $w_1\le w_\alpha$. (Tabulated one-sided values correspond to $\alpha=0.025$ and $0.005$.)
  \item Examples: $w_{0.05}=78$ for $n_1=n_2=10$. $w_{0.05}=18$ for $n_1=5$, $n_2=6$.
\end{itemize}

\begin{example}\label{ex:ranksum}
Machine~1: $14,18,12,20,17$ ($n_1=5$). Machine~2: $22,25,19,27,24,21$ ($n_2=6$). Test $H_0:\mu_1=\mu_2$ against $H_1:\mu_1\ne\mu_2$ at $\alpha=0.05$.
\begin{itemize}
  \item Ordered, with ranks: $12(1),14(2),17(3),18(4),19(5),20(6),21(7),22(8),24(9),25(10),27(11)$.
  \item Machine~1 has ranks $1,2,4,6,3$, so $w_1=16$ and $w_2=66-16=50$.
  \item The critical value for $(5,6)$ is $18$. Since $16\le18$, reject $H_0$. The machines differ in mean.
\end{itemize}
\end{example}

\begin{example}[Aluminum-lithium alloys]
Ten specimens of a traditional alloy and ten of a new alloy are tested for axial stress. Test $H_0:\mu_1=\mu_2$ at $\alpha=0.05$. For $n_1=n_2=10$, $w_{0.05}=78$ (two-sided). Neither $w_1$ nor $w_2$ is $\le78$, so do not reject. There is no evidence the alloys differ in mean axial stress.
\end{example}

\subsection{Large-Sample Approximation}

When both $n_1,n_2$ exceed $8$, $W_1$ is approximately normal with
\[ \mu_{W_1}=\frac{n_1(n_1+n_2+1)}2,\qquad\sigma^2_{W_1}=\frac{n_1n_2(n_1+n_2+1)}{12}. \]
So $Z_0=(W_1-\mu_{W_1})/\sigma_{W_1}$ is approximately $\Normal(0,1)$ under $H_0$. The critical regions are $\abs{z_0}>z_{\alpha/2}$, $z_0>z_\alpha$ or $z_0<-z_\alpha$, as appropriate.

For Example~\ref{ex:ranksum}: $\mu=30$, $\sigma=\sqrt{30}$, $z_0=(16-30)/5.48=-2.56$, $P=0.011$. (Here $n_1=5$ is below the recommended size, so this is only a rough check.)

\subsection{Comparison to the $t$-Test}

As in the one-sample case:
\begin{itemize}
  \item With normal populations, the rank-sum test is about $95\%$ as efficient as the two-sample $t$-test in large samples.
  \item Whatever the distributions, it is at least $86\%$ as efficient.
  \item It is usually more efficient when the tails are heavier than normal.
\end{itemize}

\section{Nonparametric Methods in the Analysis of Variance}

\subsection{Kruskal--Wallis Test}

The single-factor model of Chapter~7 is $Y_{ij}=\mu+\tau_i+\varepsilon_{ij}$. The \emph{Kruskal--Wallis test} is a nonparametric alternative to its $F$-test. It requires only that the $\varepsilon_{ij}$ have the same continuous distribution for all $a$ treatments.

\textbf{Setup.} Let $N=\sum n_i$. Rank all $N$ observations from $1$ to $N$ (average ranks for ties). Let $R_{i\cdot}$ be the total and $\bar{R}_{i\cdot}$ the average of the ranks in treatment $i$. Under $H_0:\mu_1=\dots=\mu_a$, the expected rank of each observation is $(N+1)/2$.

\begin{theorem}[Kruskal--Wallis statistic]\label{thm:kw}
\[ H=\frac{12}{N(N+1)}\sum_{i=1}^an_i\Big(\bar{R}_{i\cdot}-\frac{N+1}2\Big)^2=\frac{12}{N(N+1)}\sum_{i=1}^a\frac{R_{i\cdot}^2}{n_i}-3(N+1). \]
When $H_0$ is true, and either $a=3$ with all $n_i\ge6$, or $a>3$ with all $n_i\ge5$, $H$ is approximately $\chi^2_{a-1}$. Reject $H_0$ if $h>\chi^2_{\alpha,\,a-1}$.

With many ties, use the corrected statistic
\[ H'=\frac1{S^2}\left[\sum_i\frac{R_{i\cdot}^2}{n_i}-\frac{N(N+1)^2}4\right],\qquad S^2=\frac1{N-1}\left[\sum_{i,j}R_{ij}^2-\frac{N(N+1)^2}4\right]. \]
$S^2$ is the variance of the ranks.
\end{theorem}

\begin{example}\label{ex:kw}
Three treatments, $n_i=5$ each: A: $12,15,11,14,13$. B: $18,20,17,19,16$. C: $14,16,15,13,17$.
\begin{itemize}
  \item Ranking all $15$ values (average ranks for ties) gives rank totals $R_{A\cdot}=19.5$, $R_{B\cdot}=63.0$, $R_{C\cdot}=38.0$. They sum to $120=15\cdot16/2$.
  \item $H=\dfrac{12}{15\cdot16}\cdot\dfrac{19.5^2+63^2+38^2}5-3(16)=0.05(1158.65)-48=9.93$.
  \item With ties: $S^2=(1237.5-960)/14=19.82$ and $H'=(1158.65-960)/19.82=10.02$.
  \item With $a-1=2$ degrees of freedom, $\chi^2_{0.05,2}=5.99$. Reject $H_0$ ($P\approx0.007$). The treatments do not all have the same location.
\end{itemize}
\end{example}

\subsection{Rank Transformation}

A general trick: \emph{replace each observation by its rank} (over all the data, average ranks for ties). Then run the usual normal-theory procedure on the ranks.
\begin{itemize}
  \item For the single-factor design, the $F$-test of Chapter~7 on the ranks is the \emph{rank transformation} test. $H$ is a monotone function of that $F$ statistic, so the two are closely related. It is approximately valid for large samples whatever the distribution.
  \item The same idea gives nonparametric versions of other $F$-tests (blocked designs, regression). A $t$-test on the ranks approximately reproduces the rank-sum test.
\end{itemize}

\textbf{Textbook advice.} If you worry about normality or outliers, run the usual ANOVA on both the original data and the ranks.
\begin{itemize}
  \item If they agree, the ANOVA assumptions are probably fine.
  \item If they disagree, prefer the rank-based result. It is less distorted by non-normality and wild values. Then look into transformations and possible outliers.
\end{itemize}
In the textbook example, five cotton levels are compared at $\alpha=0.01$. There are many ties, so the tie-corrected statistic is used. The critical value is $\chi^2_{0.01,4}=13.28$. Kruskal--Wallis agrees with the ANOVA $F$-test in rejecting $H_0$.

\begin{summarybox}
\begin{itemize}
  \item Nonparametric tests need little beyond continuity and work with signs and ranks. They are slightly less efficient than parametric tests under normality, and often more powerful when tails are heavy.
  \item Sign test: $R^+\sim\mathrm{Bin}(n,\tfrac12)$ under $H_0:\tilde\mu=\tilde\mu_0$. Use exact $P$-values, the table with $r=\min(r^+,r^-)$, or $z_0=(r^+-0.5n)/(0.5\sqrt n)$. Discard ties. For pairs, apply it to the differences. $\beta$ depends on the distribution. The $t$-test is better under normality.
  \item Wilcoxon signed-rank: rank $\abs{X_i-\mu_0}$; $W^+$, $W^-$, $W=\min$; reject if $w\le w^*$. Large $n$: mean $n(n+1)/4$, variance $n(n+1)(2n+1)/24$. ARE $\approx0.95$ against $t$.
  \item Wilcoxon rank-sum: ranks in the combined sample; $W_1$ (smaller sample) and $W_2=(n_1+n_2)(n_1+n_2+1)/2-W_1$. Mean $n_1(n_1+n_2+1)/2$, variance $n_1n_2(n_1+n_2+1)/12$.
  \item Kruskal--Wallis: $H=\frac{12}{N(N+1)}\sum R_{i\cdot}^2/n_i-3(N+1)\approx\chi^2_{a-1}$. Use the tie-corrected $H'$ when there are many ties. The rank transformation turns any ANOVA into a nonparametric test.
\end{itemize}
\end{summarybox}

\section{Exercises}

\begin{exercise}
The following are measurements of the lifetime (hours) of $12$ bulbs: $945,\ 1010,\ 930,\ 1055,\ 990,\ 1100,\ 970,\ 1040,\ 1010,\ 960,\ 1075,\ 985$. Use the sign test to test $H_0:\tilde\mu=1000$ against $H_1:\tilde\mu\ne1000$ at $\alpha=0.05$, find the exact $P$-value and compare with the normal approximation.
\end{exercise}

\begin{exercise}
Eleven pairs of measurements give the paired differences $0.4,\ 0.9,\ -0.2,\ 1.1,\ 0.6,\ 0.3,\ 1.4,\ -0.1,\ 0.8,\ 0.5,\ 1.0$. Test $H_0:\mu_D=0$ with (a) the sign test and (b) the Wilcoxon signed-rank test at $\alpha=0.05$ and compare the conclusions.
\end{exercise}

\begin{exercise}
Compute $\beta$ for the sign test with $n=10$ ($r^*_{0.05}=1$) against a normal alternative whose median is $1.0$ standard deviation above the hypothesised median, assuming the test rejects only for small $r^-$. Compare qualitatively with the $t$-test.
\end{exercise}

\begin{exercise}
Two methods are compared: method~1: $31,\ 28,\ 35,\ 33,\ 29,\ 36,\ 30$ and method~2: $38,\ 41,\ 34,\ 43,\ 40,\ 37,\ 39,\ 42$. Use the Wilcoxon rank-sum test at $\alpha=0.05$ (use the normal approximation) to test for a difference in means, and compare with the two-sample $t$-test.
\end{exercise}

\begin{exercise}
Show that under $H_0$ the Wilcoxon rank-sum statistic $W_1$ has mean $n_1(n_1+n_2+1)/2$. (Hint: every observation has the same expected rank $(n_1+n_2+1)/2$.)
\end{exercise}

\begin{exercise}
Four treatments with five observations each have rank totals $R_{1\cdot}=30,\ R_{2\cdot}=45,\ R_{3\cdot}=75,\ R_{4\cdot}=60$. Compute $H$ and test equality of locations at $\alpha=0.05$.
\end{exercise}

\begin{exercise}
Verify that $W^++W^-=n(n+1)/2$ and that, under $H_0$, $\Exp[W^+]=n(n+1)/4$ (each rank enters $W^+$ with probability $1/2$), and compute the exact distribution of $W^+$ for $n=4$.
\end{exercise}

\begin{exercise}
Explain when you would prefer (a) the sign test, (b) the signed-rank test and (c) the $t$-test for a single sample of $n=15$ measurements, in terms of what is known about the population.
\end{exercise}

\chapter{Statistical Quality Control}

\section{Quality Improvement and Statistics}

\emph{Quality} means \emph{fitness for use}. It has two parts:
\begin{itemize}
  \item \emph{Quality of design}: the grade, performance and reliability deliberately built into the product.
  \item \emph{Quality of conformance}: how consistently units meet the design. This means steadily reducing variability and defects.
\end{itemize}
Quality improvement is the systematic removal of waste: scrap, rework, inspection, warranty cost and time spent doing things over.

Statistics helps throughout: designed experiments in development (Chapters~7--8), process-capability studies, control charts for monitoring, and life testing for reliability.

\section{Statistical Quality Control}

Statistical quality control (SQC) is the set of statistical and engineering methods for measuring, monitoring, controlling and improving quality.
\begin{itemize}
  \item W.\,A.\ Shewhart (Bell Laboratories) introduced the control chart in 1924.
  \item H.\,F.\ Dodge and H.\,G.\ Romig developed statistical sampling inspection.
  \item W.\,E.\ Deming and J.\,M.\ Juran spread the methods widely after the Second World War.
\end{itemize}

\section{Statistical Process Control}

Quality cannot be inspected into a product. The process must be right the first time. \emph{Statistical process control} (SPC) is the on-line use of statistics to make a process stable and to improve its capability by cutting variability.

The basic problem-solving tools are the \emph{seven tools}:
\begin{enumerate}[label=(\arabic*)]
  \item histogram (or stem-and-leaf plot);
  \item Pareto chart;
  \item cause-and-effect (fishbone) diagram;
  \item defect-concentration diagram;
  \item control chart;
  \item scatter diagram;
  \item check sheet.
\end{enumerate}
These are the technical side of SPC. The other essential part is attitude: a commitment at all levels to continuous improvement by reducing variability. The control chart is the most powerful tool.

\section{Introduction to Control Charts}

\subsection{Basic Principles}

\begin{itemize}
  \item Every process has natural variability from many small, unavoidable \emph{chance causes} (common causes). A process with only chance causes is \emph{in statistical control}.
  \item Other variability comes from badly adjusted machines, operator errors or defective material. It is usually large compared with the background noise. It signals an \emph{assignable cause} (special cause). A process with assignable causes is \emph{out of control}.
  \item A main aim of SPC is to find assignable causes quickly, so corrective action comes before many nonconforming units are made.
\end{itemize}

A \emph{control chart} plots a quality statistic from samples against sample number (or time).
\begin{itemize}
  \item The \emph{center line} (CL) is the average of the statistic when only chance causes act.
  \item The \emph{upper} and \emph{lower control limits} (UCL, LCL) are set so nearly all points fall between them while the process is in control.
  \item A point outside the limits, or a non-random pattern inside them, is evidence of an assignable cause. Investigate.
\end{itemize}

\begin{definition}[General model of a Shewhart control chart]
Let $W$ be a sample statistic measuring a quality characteristic, with mean $\mu_W$ and standard deviation $\sigma_W$. Then
\[ \mathrm{UCL}=\mu_W+k\sigma_W,\qquad\mathrm{CL}=\mu_W,\qquad\mathrm{LCL}=\mu_W-k\sigma_W. \]
$k$ is the distance of the limits from the center line, in standard deviations. The standard choice is $k=3$ (\emph{three-sigma limits}).
\end{definition}

\textbf{A control chart is a repeated hypothesis test.}
\begin{itemize}
  \item For normal $W$, three-sigma limits give $\alpha=0.0027$ per point. That is the chance a point falls outside when the process is in control.
  \item A point inside the limits is like \emph{failing to reject} $H_0$ (process in control). A point outside is like rejecting it.
\end{itemize}

Two types of chart:
\begin{itemize}
  \item \emph{Variables charts} monitor a continuous measurement (a mean chart plus a range or standard-deviation chart).
  \item \emph{Attributes charts} monitor counts of nonconforming units or of nonconformities (defects).
\end{itemize}

\subsection{Design of a Control Chart}

\begin{example}\label{ex:xbardesign}
A process fills containers with nominal mean $\mu=250$ g and known $\sigma=3$ g. Samples of $n=4$ are taken hourly and $\bar{x}$ is plotted.
\begin{itemize}
  \item $\sigma_{\bar{X}}=\sigma/\sqrt n=1.5$.
  \item $\mathrm{UCL}=250+3(1.5)=254.5$ and $\mathrm{LCL}=250-4.5=245.5$.
\end{itemize}
This chart tests $H_0:\mu=250$ at each time point.
\end{example}

For a given multiple of sigma, the limits narrow as $\sqrt n$ grows. The designer must choose sample size and sampling frequency.
\begin{itemize}
  \item Larger samples detect smaller shifts. If only large shifts matter, small samples are enough.
  \item You usually cannot sample both large and often. Current practice favours small, frequent samples.
\end{itemize}

\begin{example}[Piston rings]
Inside diameter has process mean $74$ mm and known $\sigma=0.01$ mm. Samples of $n=5$ are taken hourly.
\begin{itemize}
  \item $\sigma_{\bar{X}}=0.01/\sqrt5=0.0045$.
  \item $\mathrm{UCL}=74+3(0.0045)=74.0135$ and $\mathrm{LCL}=73.9865$ ($\alpha=0.0027$).
\end{itemize}
Choosing these limits is the same as the critical region for repeatedly testing $H_0:\mu=74$ against $H_1:\mu\ne74$ with $\sigma=0.01$ known.
\end{example}

\subsection{Rational Subgroups}

Choose samples (\emph{subgroups}) so that variation \emph{within} a subgroup is only chance variation, and assignable causes show up \emph{between} subgroups. Then the limits bound the chance variation, and assignable causes push points outside.

Two common approaches:
\begin{itemize}
  \item Each subgroup is units made at the same time (a ``snapshot''). This gives the best chance of detecting shifts.
  \item Each subgroup is a random sample of everything made since the last sample. This is used when the aim is to decide whether to accept the product.
\end{itemize}
Rules:
\begin{itemize}
  \item Subgroups must not straddle assignable changes (for example the end of one shift and the start of the next).
  \item Chart the output of several machines or heads \emph{separately}. Do not pool it.
\end{itemize}

\subsection{Analysis of Patterns on Control Charts}

A chart can signal trouble even when all points are inside the limits, if the points follow a non-random pattern: a long \emph{run} (for example all above the center line), a trend, cycles, or many points near the limits. Eight or more points on one side of the center line is very unlikely by chance.

The \emph{Western Electric rules} formalise this. The process is out of control if any of these happens:
\begin{enumerate}[label=\arabic*.]
  \item One point is outside the three-sigma limits.
  \item Two of three consecutive points are beyond a two-sigma limit (same side of the center line).
  \item Four of five consecutive points are at least one sigma from the center line (same side).
  \item Eight consecutive points are on one side of the center line.
\end{enumerate}
\begin{itemize}
  \item The one-, two- and three-sigma lines divide each side of the chart into zones C, B and A.
  \item Rules 2 and 3 apply to one side at a time. A point above the upper two-sigma line followed by one below the lower two-sigma line is not a signal.
  \item These rules add sensitivity but also raise the false-alarm rate. Interpreting patterns needs knowledge of the process.
\end{itemize}

\section{$\bar{X}$ and $R$ or $S$ Control Charts}

The $\bar{X}$ chart controls the average. The range chart ($R$) or standard-deviation chart ($S$) controls the variability.

If $\mu$ and $\sigma$ are known, the $\bar{X}$ chart has
\[ \mathrm{UCL}=\mu+\frac{3\sigma}{\sqrt n},\qquad\mathrm{CL}=\mu,\qquad\mathrm{LCL}=\mu-\frac{3\sigma}{\sqrt n}. \]
When they are unknown, estimate them from $m$ \emph{preliminary samples} of size $n$, taken while the process is believed to be in control. Use at least $20$ to $25$ samples, with $n$ typically $4$, $5$ or $6$. Let $\bar{x}_i$, $r_i$, $s_i$ be the mean, range and standard deviation of sample $i$. Define
\[ \bar{\bar{x}}=\frac1m\sum_{i=1}^m\bar{x}_i,\qquad\bar{r}=\frac1m\sum_{i=1}^mr_i,\qquad\bar{s}=\frac1m\sum_{i=1}^ms_i. \]

\subsubsection*{Estimating $\sigma$ from the range}

For a normal sample, the \emph{relative range} $W=R/\sigma$ has a distribution that depends only on $n$. Its mean is $d_2$ and its standard deviation is $d_3$ (tabulated). So $\Exp[R]=d_2\sigma$ and $\sigma_R=d_3\sigma$. Hence $\hat\sigma=\bar{R}/d_2$ is unbiased, and $\hat\sigma_R=d_3\bar{R}/d_2$.

\begin{theorem}[$\bar{X}$ and $R$ charts]\label{thm:xbarR}
The $\bar{X}$ chart (from $\bar{R}$):
\[ \mathrm{UCL}=\bar{\bar{x}}+A_2\bar{r},\qquad\mathrm{CL}=\bar{\bar{x}},\qquad\mathrm{LCL}=\bar{\bar{x}}-A_2\bar{r},\qquad A_2=\frac3{d_2\sqrt n}. \]
The $R$ chart:
\[ \mathrm{UCL}=D_4\bar{r},\qquad\mathrm{CL}=\bar{r},\qquad\mathrm{LCL}=D_3\bar{r},\qquad D_3=1-\frac{3d_3}{d_2},\quad D_4=1+\frac{3d_3}{d_2}. \]
If $D_3$ (the lower limit) comes out negative, set it to $0$.
\end{theorem}

\subsubsection*{Estimating $\sigma$ from the standard deviation: $S$ chart}

$S$ is a biased estimator: $\Exp[S]=c_4\sigma$, with $c_4$ close to (but below) $1$. Also $\sigma_S=\sigma\sqrt{1-c_4^2}$. With $\bar{s}$ the average sample standard deviation, $\hat\sigma=\bar{s}/c_4$ is unbiased.

\begin{theorem}[$\bar{X}$ and $S$ charts]\label{thm:xbarS}
The $S$ chart:
\[ \mathrm{UCL}=\bar{s}+\frac{3\bar{s}}{c_4}\sqrt{1-c_4^2},\qquad\mathrm{CL}=\bar{s},\qquad\mathrm{LCL}=\bar{s}-\frac{3\bar{s}}{c_4}\sqrt{1-c_4^2}\ (\text{set to }0\text{ if negative}). \]
The $\bar{X}$ chart (from $\bar{S}$):
\[ \mathrm{UCL}=\bar{\bar{x}}+\frac{3\bar{s}}{c_4\sqrt n},\qquad\mathrm{CL}=\bar{\bar{x}},\qquad\mathrm{LCL}=\bar{\bar{x}}-\frac{3\bar{s}}{c_4\sqrt n}. \]
\end{theorem}

\subsubsection*{Factors for control charts}

\begin{center}
\small
\begin{tabular}{c|ccccc|c}
\toprule
$n$ & $A_2$ & $d_2$ & $D_3$ & $D_4$ & $c_4$ & $A=3/\sqrt n$\\
\midrule
2 & 1.880 & 1.128 & 0 & 3.267 & 0.7979 & 2.121\\
3 & 1.023 & 1.693 & 0 & 2.575 & 0.8862 & 1.732\\
4 & 0.729 & 2.059 & 0 & 2.282 & 0.9213 & 1.500\\
5 & 0.577 & 2.326 & 0 & 2.115 & 0.9400 & 1.342\\
6 & 0.483 & 2.534 & 0 & 2.004 & 0.9515 & 1.225\\
7 & 0.419 & 2.704 & 0.076 & 1.924 & 0.9594 & 1.134\\
8 & 0.373 & 2.847 & 0.136 & 1.864 & 0.9650 & 1.061\\
9 & 0.337 & 2.970 & 0.184 & 1.816 & 0.9693 & 1.000\\
10 & 0.308 & 3.078 & 0.223 & 1.777 & 0.9727 & 0.949\\
\bottomrule
\end{tabular}
\end{center}

\begin{example}\label{ex:xbarR}
Ten preliminary samples of size $n=4$.
\begin{itemize}
  \item Means: $50.1,\ 49.8,\ 50.6,\ 50.3,\ 49.5,\ 50.9,\ 50.0,\ 50.4,\ 49.9,\ 50.5$.
  \item Ranges: $3.2,\ 2.8,\ 4.1,\ 3.5,\ 2.9,\ 3.8,\ 3.0,\ 3.3,\ 2.7,\ 3.7$.
  \item $\bar{\bar{x}}=502.0/10=50.20$ and $\bar{r}=33.0/10=3.30$. For $n=4$: $A_2=0.729$, $D_3=0$, $D_4=2.282$.
  \item $\bar{X}$ chart: $\mathrm{CL}=50.20$, $\mathrm{UCL}=50.20+0.729(3.30)=52.61$, $\mathrm{LCL}=50.20-0.729(3.30)=47.79$.
  \item $R$ chart: $\mathrm{CL}=3.30$, $\mathrm{UCL}=2.282(3.30)=7.53$, $\mathrm{LCL}=0$.
\end{itemize}
All means lie between $47.79$ and $52.61$, and all ranges are below $7.53$. The trial limits show the process is in control, so use them to monitor future production. The estimated process standard deviation is $\hat\sigma=\bar{r}/d_2=3.30/2.059=1.60$.
\end{example}

\begin{pitfall}{trial limits and points outside}
Limits from preliminary data are \emph{trial limits}.
\begin{itemize}
  \item Plot the $m$ preliminary means and ranges. Investigate any point outside the limits.
  \item If an assignable cause is found and removed, discard those samples and \emph{recompute} the limits from the rest.
  \item Study the $R$ (or $S$) chart first. If the variability is not stable, the $\bar{X}$ limits (which use $\bar{R}$) are misleading.
  \item Revise limits periodically, and whenever the process improves.
  \item Never put specification limits on a control chart. Specification limits and control limits are unrelated.
\end{itemize}
\end{pitfall}

\begin{example}[Vane opening]
Twenty samples of five parts each (coded data) give $\bar{\bar{x}}=33.3$ and $\bar{r}=5.8$. For $n=5$, $A_2=0.577$.
\begin{itemize}
  \item Trial $\bar{X}$ limits: $33.3\pm0.577(5.8)$, about $29.95$ and $36.65$.
  \item $R$ chart: $\mathrm{CL}=5.8$, $\mathrm{UCL}=D_4\bar{r}$, $\mathrm{LCL}=0$ (since $D_3=0$ for $n=5$).
  \item Samples $6,8,11,19$ are out of control on the $\bar{X}$ chart, and sample $9$ on the $R$ chart (rule 1).
  \item All causes traced to a defective tool in the wax-molding area. Discard these five samples and recompute the limits.
\end{itemize}
Limits from $\bar{s}$ are almost the same as from $\bar{r}$ for such small subgroups. These first $20$ samples are estimation data. The revised limits are then used to judge future production.
\end{example}

\section{Control Charts for Individual Measurements}

Sometimes each sample is a single unit ($n=1$). Every unit may be measured automatically, output may be slow, or only one measurement exists (a batch process).

The \emph{chart for individuals} estimates variability from the \emph{moving range} of two successive observations, $MR_i=\abs{x_i-x_{i-1}}$. Since $d_2=1.128$ for $n=2$, $\hat\sigma=\overline{MR}/1.128$.

\begin{theorem}[Individuals and moving-range charts]\label{thm:individuals}
Individuals chart:
\[ \mathrm{UCL}=\bar{x}+\frac{3\overline{mr}}{1.128},\qquad\mathrm{CL}=\bar{x},\qquad\mathrm{LCL}=\bar{x}-\frac{3\overline{mr}}{1.128}. \]
Moving-range chart ($n=2$, $D_3=0$, $D_4=3.267$): $\mathrm{UCL}=3.267\,\overline{mr}$, $\mathrm{CL}=\overline{mr}$, $\mathrm{LCL}=0$.
\end{theorem}

\begin{example}
Ten concentration readings: $98.2,\ 99.1,\ 97.8,\ 100.4,\ 99.5,\ 98.7,\ 101.0,\ 99.3,\ 98.0,\ 99.9$.
\begin{itemize}
  \item $\bar{x}=99.19$.
  \item The nine moving ranges are $0.9,1.3,2.6,0.9,0.8,2.3,1.7,1.3,1.9$, so $\overline{mr}=13.7/9=1.52$.
  \item $\hat\sigma=1.52/1.128=1.35$.
  \item Individuals chart: $\mathrm{UCL}=99.19+3(1.35)=103.24$, $\mathrm{LCL}=95.14$.
  \item Moving-range chart: $\mathrm{UCL}=3.267(1.52)=4.97$.
\end{itemize}
All points are in control.
\end{example}

Read the individuals chart like an $\bar{X}$ chart, with these cautions:
\begin{itemize}
  \item Moving ranges are correlated, so apparent runs or cycles on the $MR$ chart need not be real.
  \item The observations are assumed uncorrelated, and non-normality matters more.
  \item The chart is slow to detect small shifts. A one-sigma shift takes about $43.9$ points on average.
  \item Narrower limits are \emph{not} the answer: they greatly increase false alarms. Use a CUSUM or EWMA chart for small shifts (Section~\ref{sec:cusum}).
\end{itemize}

\begin{example}[Concentration]
Twenty hourly readings have $\bar{x}=99.1$ and average moving range $\overline{mr}=2.59$. With $D_3=0$, $D_4=3.267$, $d_2=1.128$:
\begin{itemize}
  \item Moving-range chart: $\mathrm{CL}=2.59$, $\mathrm{LCL}=0$, $\mathrm{UCL}=3.267(2.59)=8.46$.
  \item Individuals chart: $99.1\pm3(2.59)/1.128=99.1\pm6.89$.
  \item No point is outside, so the process is in control.
\end{itemize}
\end{example}

\section{Process Capability}

\emph{Process capability} is how a process performs when it is in control. Is its natural variability acceptable relative to the specifications? It is meaningful only for a \emph{stable} process.

Two graphical tools:
\begin{itemize}
  \item The \emph{tolerance chart}: individual measurements plotted with the specification limits LSL and USL. This is the only chart where specification limits belong.
  \item The \emph{histogram}: shows centering and spread.
\end{itemize}

\begin{definition}[Process capability ratio]
With specification limits USL and LSL and process standard deviation $\sigma$,
\[ \mathrm{PCR}=\frac{\mathrm{USL}-\mathrm{LSL}}{6\sigma}. \]
$6\sigma$ is the \emph{width of the process}. The limits $\mu\pm3\sigma$ are the \emph{natural tolerance limits}. The quantity $(1/\mathrm{PCR})\times100$ is the percentage of the specification width the process uses.
\end{definition}

How to read PCR:
\begin{itemize}
  \item $\mathrm{PCR}>1$: the natural limits lie inside the specifications. Few nonconforming units.
  \item $\mathrm{PCR}=1$, centered and normal: $0.27\%$ ($2700$ parts per million) fall outside.
  \item $\mathrm{PCR}<1$: the process is very yield-sensitive. Many nonconforming units.
\end{itemize}
PCR assumes the process is centered at the nominal value, so it measures \emph{potential} capability. For an off-center process, use the \emph{actual} capability:

\begin{definition}[PCR$_k$]
\[ \mathrm{PCR}_k=\min\left[\frac{\mathrm{USL}-\mu}{3\sigma},\ \frac{\mu-\mathrm{LSL}}{3\sigma}\right]. \]
This is a one-sided ratio, using the specification limit nearest the process mean. $\mathrm{PCR}_k=\mathrm{PCR}$ exactly when the process is centered. $\mathrm{PCR}_k<\mathrm{PCR}$ signals off-center operation.
\end{definition}

For a normal process in control, the nonconforming fractions are $\Prob(X<\mathrm{LSL})=\Phi\big((\mathrm{LSL}-\mu)/\sigma\big)$ and $\Prob(X>\mathrm{USL})=1-\Phi\big((\mathrm{USL}-\mu)/\sigma\big)$. These assume normality and control, and use estimated $\mu,\sigma$. Treat them as approximate guidelines.

\begin{example}
For the process of Example~\ref{ex:xbarR}, let the specification be $50\pm6$ ($\mathrm{LSL}=44$, $\mathrm{USL}=56$). With $\hat\sigma=1.603$ and $\bar{\bar{x}}=50.20$:
\begin{itemize}
  \item $\mathrm{PCR}=12/(6\cdot1.603)=1.25$.
  \item $\mathrm{PCR}_k=\min\big[(56-50.2)/(3\cdot1.603),\ (50.2-44)/(3\cdot1.603)\big]=\min(1.21,1.29)=1.21$.
  \item Fallout: $\Prob(X<44)+\Prob(X>56)=\Phi(-3.87)+\big(1-\Phi(3.62)\big)=0.00005+0.00015=0.0002$. About $200$ parts per million.
\end{itemize}
\end{example}

\begin{example}
A current has specifications $100\pm10$ mA. The process has $\mu=107.0$, $\sigma=1.5$.
\begin{itemize}
  \item $\mathrm{PCR}=20/(6\cdot1.5)=2.22$.
  \item $\mathrm{PCR}_k=(110-107)/(3\cdot1.5)=0.67$.
  \item The process is capable but badly off-center. $\Prob(X>110)=\Prob(Z>2)=0.023$: a $2.3\%$ fallout.
\end{itemize}
\end{example}

\emph{Guidelines.}
\begin{itemize}
  \item A common minimum target is $\mathrm{PCR}=1.33$. Use up to $1.66$ for strength, safety or critical characteristics. Some companies require $\mathrm{PCR}_k\ge2.0$.
  \item A process with $\mathrm{PCR}=\mathrm{PCR}_k=2$ is a \emph{six-sigma process}: the nearest specification is six standard deviations away. If its mean shifts by $1.5\sigma$, $\mathrm{PCR}_k$ falls to $1.5$ and fallout is only about $3.4$ parts per million.
  \item In the automotive industry, $C_p$ and $C_{pk}$ mean PCR and PCR$_k$.
\end{itemize}

\section{Attribute Control Charts}

When a unit is just classified as conforming or nonconforming (defective), or defects per unit are counted, use attribute charts. They usually need larger samples than variables charts.

\subsection{$P$ Chart (Control Chart for Proportions)}

Let $D$ be the number of nonconforming units in a random sample of $n$. $D$ is binomial$(n,p)$. The sample fraction nonconforming $\hat{P}=D/n$ has $\mu_{\hat{P}}=p$ and $\sigma^2_{\hat{P}}=p(1-p)/n$. Estimate $p$ by $\bar{p}=\dfrac1{mn}\sum_{i=1}^mD_i$ from $m$ preliminary samples.

\begin{theorem}[$P$ chart]\label{thm:pchart}
\[ \mathrm{UCL}=\bar{p}+3\sqrt{\frac{\bar{p}(1-\bar{p})}n},\qquad\mathrm{CL}=\bar{p},\qquad\mathrm{LCL}=\bar{p}-3\sqrt{\frac{\bar{p}(1-\bar{p})}n}\ (\text{set to }0\text{ if negative}). \]
\end{theorem}

\begin{itemize}
  \item The limits use the normal approximation to the binomial. For small $p$ it is poor; then take limits from binomial probabilities.
  \item The \emph{$NP$ chart} plots the number nonconforming $n\hat{p}$ itself. Multiply everything by $n$: $\mathrm{CL}=n\bar{p}$, limits $n\bar{p}\pm3\sqrt{n\bar{p}(1-\bar{p})}$.
\end{itemize}

\begin{example}
Twenty samples of $n=200$ contain $160$ nonconforming units in total.
\begin{itemize}
  \item $\bar{p}=160/4000=0.04$ and $\sqrt{0.04\cdot0.96/200}=0.01386$.
  \item $\mathrm{UCL}=0.04+0.0416=0.0816$. $\mathrm{LCL}=0.04-0.0416<0$, so set it to $0$.
  \item A sample with more than $16$ nonconforming units ($>0.0816\cdot200$) signals a problem.
\end{itemize}
\end{example}

\subsection{$U$ Chart (Control Chart for Defects per Unit)}

To monitor the number of \emph{defects} (nonconformities) per unit, such as flaws per yard of cloth, model the count as Poisson. If a sample of $n$ inspection units has $C$ defects in total, then $U=C/n$ is the average defects per unit, with mean $\lambda$ and variance $\lambda/n$. With $\bar{u}=\frac1m\sum U_i$:

\begin{theorem}[$U$ chart]\label{thm:uchart}
\[ \mathrm{UCL}=\bar{u}+3\sqrt{\frac{\bar{u}}n},\qquad\mathrm{CL}=\bar{u},\qquad\mathrm{LCL}=\bar{u}-3\sqrt{\frac{\bar{u}}n}\ (\text{set to }0\text{ if negative}). \]
\end{theorem}

The $C$ chart plots the total defects $C$ in a sample: $\mathrm{CL}=n\bar{u}$, limits $n\bar{u}\pm3\sqrt{n\bar{u}}$. These limits use the normal approximation to the Poisson and are poor when $\lambda$ is small.

\begin{example}
Twenty samples of $5$ circuit boards each show $120$ defects in total.
\begin{itemize}
  \item $\bar{u}=120/100=1.2$ defects per board.
  \item $\mathrm{UCL}=1.2+3\sqrt{1.2/5}=2.67$. $\mathrm{LCL}<0$, so set it to $0$.
\end{itemize}
A chart can show the process is in control while the defect level is still too high. $1.2$ defects per board needs process improvement: study the defect types, often with designed experiments.
\end{example}

\begin{example}[$P$ chart and $U$ chart from the textbook]
\emph{$P$ chart.} Twenty samples of $100$ ceramic substrates contain $800$ defectives ($2000$ units in total across samples), so $\bar{p}=0.40$. Limits: $0.40\pm3\sqrt{0.40(0.60)/100}=0.40\pm0.147$, that is, $0.253$ to $0.547$. All samples are in control, but a $40\%$ defective rate is very poor. Analyse the defects by type. An $NP$ chart plots $D$ directly, with limits $n$ times those of the $P$ chart.

\emph{$U$ chart.} Printed circuit boards: five boards inspected per hour, $20$ samples, about $1.6$ defects per board on average. $\bar{u}$ is the average of the $u_i=c_i/n$. Limits: $\bar{u}\pm3\sqrt{\bar{u}/n}$ (a negative $\mathrm{LCL}$ becomes $0$). The process is in control, but $8$ defects per five boards is too many. A $C$ chart plots total defects $c$, with limits $n$ times those of the $U$ chart.
\end{example}

\section{Control Chart Performance}

Moving the limits farther from the center line lowers the Type~I risk (a false alarm: a point outside when no assignable cause is present). But it raises the Type~II risk (failing to signal a real shift).

A measure of performance is the \emph{average run length} (ARL): the average number of points plotted before one signals. If $p$ is the probability that a single point exceeds the limits, the run length is geometric, and
\[ \mathrm{ARL}=\frac1p. \]

\begin{itemize}
  \item \textbf{In control.} For a three-sigma $\bar{X}$ chart $p=0.0027$, so $\mathrm{ARL}_0=1/0.0027=370$. A false alarm comes on average once every $370$ samples.
  \item \textbf{Out of control.} If the mean shifts, $p=1-\beta$ with $\beta=\Prob(\mathrm{LCL}\le\bar{X}\le\mathrm{UCL}\mid\text{new mean})$. If the mean shifts exactly to the upper control limit, $\beta=\Prob(-6\le Z\le0)=0.5$. So $p=0.5$ and $\mathrm{ARL}=2$: detection after two samples on average.
\end{itemize}
The average time to signal is $\mathrm{ARL}\times$ (time between samples). To shorten it, sample more often or increase $n$.

For a shift of $k$ standard deviations of the individual observations, the ARL of a three-sigma $\bar{X}$ chart is:

\begin{center}
\begin{tabular}{c|cc}
\toprule
Shift (multiple of $\sigma$) & ARL, $n=1$ & ARL, $n=4$\\
\midrule
0 & 370.4 & 370.4\\
0.5 & 155.2 & 43.9\\
1.0 & 43.9 & 6.3\\
1.5 & 15.0 & 2.0\\
2.0 & 6.3 & 1.2\\
3.0 & 2.0 & 1.0\\
\bottomrule
\end{tabular}
\end{center}

A Shewhart chart detects large shifts quickly but small shifts slowly. With $n=1$, a $1\sigma$ shift takes $43.9$ points on average. This motivates charts with memory.

\begin{example}[Piston-ring ARL]
The mean shifts to $74.0135$ mm and $n=5$.
\begin{itemize}
  \item The chance that $\bar{X}$ falls between the limits is $0.50$, so $p=0.50$ and $\mathrm{ARL}=1/0.5=2$ samples. On average two hours pass between the shift and its detection.
  \item To detect it faster, sample more often (every half hour gives one hour) or increase $n$. With $n=10$ the limits narrow to $73.9905$ and $74.0095$. The chance of falling between them is about $0.1$, so $p=0.9$ and $\mathrm{ARL}\approx1.11$. About twice as fast.
\end{itemize}
Western Electric rules would lower the in-control ARL below $370$.
\end{example}

\section{Cumulative Sum Control Chart}\label{sec:cusum}

\emph{Cumulative sum} (CUSUM) charts detect \emph{small} shifts better than Shewhart charts, and work well with $n=1$. With $\bar{x}_j$ the mean of sample $j$ and $\mu_0$ the target,
\[ S_i=\sum_{j=1}^i(\bar{x}_j-\mu_0). \]
On target, $S_i$ wanders about $0$. A persistent drift up or down means the mean has shifted. The plain plot has no limits. The older fix is a \emph{V-mask}. The \emph{tabular} CUSUM is preferred.

\begin{definition}[Tabular CUSUM]
With target $\mu_0$, \emph{reference value} $K$, and $S_H(0)=S_L(0)=0$:
\[ S_H(i)=\max\big[0,\ x_i-(\mu_0+K)+S_H(i-1)\big],\qquad S_L(i)=\max\big[0,\ (\mu_0-K)-x_i+S_L(i-1)\big]. \]
\begin{itemize}
  \item $S_H$ and $S_L$ are the \emph{upper} and \emph{lower} one-sided CUSUMs.
  \item They add up deviations beyond $K$ from the target, and reset to $0$ when they would go negative.
  \item The process is out of control if either exceeds the \emph{decision interval} $H$.
\end{itemize}
\end{definition}

After a signal:
\begin{itemize}
  \item If $S_H(i)>H$, estimate the new mean by $\hat\mu=\mu_0+K+S_H(i)/N_H$, where $N_H$ is the number of consecutive periods with $S_H$ nonzero.
  \item If $S_L(i)>H$, use $\hat\mu=\mu_0-K-S_L(i)/N_L$.
  \item Find the cause, then restart the CUSUM at zero.
\end{itemize}

\begin{theorem}[Design of the CUSUM]\label{thm:cusumdesign}
Choose $K$ about half the shift to detect: $K=\delta\sigma_X/2$ for a shift $\delta\sigma_X$. So $K=k\sigma_X$ with $k=\tfrac12$ for a one-sigma shift. Choose $H=h\sigma_X$ with $h=4$ or $5$. Here $\sigma_X$ is the standard deviation of the plotted quantity ($\sigma/\sqrt n$ for sample means). Performance for $k=\tfrac12$:
\begin{center}
\begin{tabular}{c|cc}
\toprule
Shift (multiple of $\sigma_X$) & ARL, $h=4$ & ARL, $h=5$\\
\midrule
0 & 168 & 465\\
0.25 & 74.2 & 139\\
0.50 & 26.6 & 38.0\\
0.75 & 13.3 & 17.0\\
1.00 & 8.38 & 10.4\\
1.50 & 4.75 & 5.75\\
2.00 & 3.34 & 4.01\\
2.50 & 2.62 & 3.11\\
3.00 & 2.19 & 2.57\\
4.00 & 1.71 & 2.01\\
\bottomrule
\end{tabular}
\end{center}
\end{theorem}

\begin{itemize}
  \item For a one-sigma shift the CUSUM needs $8.38$ or $10.4$ samples, against $43.9$ for the Shewhart chart.
  \item The in-control ARL is $168$ for $h=4$ and $465$ for $h=5$. So $h=5$ gives fewer false alarms.
  \item \emph{Standardised CUSUM}: use $y_i=(x_i-\mu_0)/\sigma_X$, with $S_H(i)=\max[0,y_i-k+S_H(i-1)]$ and $S_L(i)=\max[0,-k-y_i+S_L(i-1)]$.
  \item Do \emph{not} apply the Western Electric rules to CUSUM values. Successive values are highly correlated and the rules would give many false alarms.
\end{itemize}

\begin{example}\label{ex:cusum}
Target $\mu_0=50$, individual standard deviation $\sigma=1$. Design for a one-sigma shift with $k=\tfrac12$ and $h=5$: $K=0.5$, $H=5$. Readings: $50.5,\ 49.0,\ 51.8,\ 52.4,\ 50.9,\ 53.1$.
\begin{itemize}
  \item Upper CUSUM $S_H(i)=\max[0,\ x_i-50.5+S_H(i-1)]$: $0,\ 0,\ 1.3,\ 3.2,\ 3.6,\ 6.2$.
  \item Lower CUSUM $S_L(i)=\max[0,\ 49.5-x_i+S_L(i-1)]$: $0,\ 0.5,\ 0,\ 0,\ 0,\ 0$.
  \item At $i=6$, $S_H=6.2>H=5$. Signal an upward shift.
  \item $S_H$ has been nonzero for $N_H=4$ periods ($i=3,\dots,6$). So $\hat\mu=\mu_0+K+S_H/N_H=50+0.5+6.2/4=52.05$.
\end{itemize}
No single reading is far outside three-sigma limits (only $53.1$ would be flagged by $50\pm3$), yet the CUSUM has gathered evidence of a persistent shift.
\end{example}

\begin{example}[Tabular CUSUM from the textbook]
Chemical concentration data: target $\mu_0=99$, $\sigma=2$, $K=1$ ($k=\tfrac12$), $H=10$ ($h=5$). Use $S_H(i)=\max[0,\,x_i-(\mu_0+K)+S_H(i-1)]$ and $S_L(i)=\max[0,\,(\mu_0-K)-x_i+S_L(i-1)]$, both starting at $0$. Neither ever exceeds $H=10$, so the process is in control.
\begin{itemize}
  \item On a signal, find the cause, correct it, and restart both CUSUMs at zero.
  \item The V-mask method is equivalent but less convenient.
  \item Do not apply Western Electric rules to a CUSUM.
  \item The recommended $h=4$ or $5$, $k=\tfrac12$ detects a $1\sigma$ shift in $8.38$ or $10.4$ samples on average.
\end{itemize}
\end{example}

\section{Other SPC Problem-Solving Tools}

The control chart is one of seven basic tools. The others support problem-solving and teamwork.
\begin{itemize}
  \item \textbf{Pareto chart.} Bars of defect types in order of decreasing frequency, with a cumulative-percentage line. It shows the ``vital few, trivial many'' and where effort pays most.
  \item \textbf{Cause-and-effect (fishbone, Ishikawa) diagram.} A team's list of possible causes of a problem, grouped by category (machines, methods, materials, measurement, people, environment). It points to factors worth investigating, often by a designed experiment.
  \item \textbf{Defect-concentration diagram.} A sketch of the product with defect locations marked. Clusters show where, and hence suggest why, defects occur.
  \item \textbf{Histogram.} Shows the shape, center and spread of the data, and how they relate to the specifications.
  \item \textbf{Scatter diagram.} Shows the relationship between two variables (correlation, Chapter~5).
  \item \textbf{Check sheet.} A simple form for collecting data in an organised way at the point of origin.
\end{itemize}
A typical improvement project:
\begin{enumerate}
  \item Rank defect types by Pareto analysis.
  \item Build a fishbone diagram of candidate causes.
  \item Locate a physical source with a defect-concentration diagram.
  \item Use a designed experiment (Chapters~7--8) to find the settings that remove the problem.
\end{enumerate}

\section{Implementing SPC}

Statistical tools alone do not give quality improvement. Successful programmes share these elements:
\begin{enumerate}[label=(\arabic*)]
  \item management leadership and commitment (management is the role model);
  \item a team approach (the basic tools help build teams);
  \item education and training at all levels;
  \item continuous improvement (weekly, quarterly, annual), not a one-time programme;
  \item a way to recognise success.
\end{enumerate}
\begin{itemize}
  \item Introduce control charts early. Processes do not naturally run in control. Using the charts removes assignable causes, reduces variability and stabilises the process. Managing with facts replaces managing by judgement alone.
  \item An \emph{out-of-control-action plan} is essential. A control chart only detects assignable causes. Operators and engineers must find and remove the root cause, not apply a cosmetic fix.
\end{itemize}

\begin{summarybox}
\begin{itemize}
  \item Chance causes are inherent (in control). Assignable causes need action (out of control). A control chart has CL $=\mu_W$ and limits $\mu_W\pm3\sigma_W$ ($\alpha=0.0027$). It is a repeated hypothesis test.
  \item Western Electric rules: one point beyond $3\sigma$; two of three beyond $2\sigma$; four of five beyond $1\sigma$; eight in a row on one side.
  \item $\bar{X}$--$R$: $\bar{\bar{x}}\pm A_2\bar{r}$; $R$: $D_3\bar{r}$, $D_4\bar{r}$; $\hat\sigma=\bar{r}/d_2$. $\bar{X}$--$S$: $\bar{\bar{x}}\pm3\bar{s}/(c_4\sqrt n)$; $S$: $\bar{s}\pm3(\bar{s}/c_4)\sqrt{1-c_4^2}$. Individuals: $\bar{x}\pm3\overline{mr}/1.128$; $MR$: $3.267\overline{mr}$.
  \item Capability: $\mathrm{PCR}=(\mathrm{USL}-\mathrm{LSL})/6\sigma$; $\mathrm{PCR}_k=\min[(\mathrm{USL}-\mu)/3\sigma,(\mu-\mathrm{LSL})/3\sigma]$. Targets $1.33$, $1.66$, $2.0$. Never plot specifications on a control chart.
  \item Attributes: $P$ chart $\bar{p}\pm3\sqrt{\bar{p}(1-\bar{p})/n}$; $U$ chart $\bar{u}\pm3\sqrt{\bar{u}/n}$ (and the $NP$ and $C$ versions).
  \item $\mathrm{ARL}=1/p$. $\mathrm{ARL}_0=370$ for three-sigma limits. CUSUM: $S_H=\max[0,x-(\mu_0+K)+S_H]$, $S_L=\max[0,(\mu_0-K)-x+S_L]$, with $K=k\sigma_X$, $H=h\sigma_X$, $k=\frac12$, $h=4$ or $5$. Much better than Shewhart for small shifts.
\end{itemize}
\end{summarybox}

\section{Exercises}

\begin{exercise}
Twenty-five samples of size $n=5$ have $\sum\bar{x}_i=1250.0$ and $\sum r_i=60.0$. Compute trial control limits for $\bar{X}$ and $R$ charts, and estimate the process standard deviation. If the specifications are $50\pm5$, compute PCR and PCR$_k$.
\end{exercise}

\begin{exercise}
For $n=5$, derive $D_3$ and $D_4$ from $d_2=2.326$ and $d_3=0.864$, and $A_2$ from $d_2$. Show that $B_4=1+3\sqrt{1-c_4^2}/c_4=2.089$ for $n=5$.
\end{exercise}

\begin{exercise}
A $\bar{X}$ chart with three-sigma limits and $n=4$ is used. Compute the probability that a single point plots outside the limits when the process mean shifts by $1.5\sigma$, and the resulting ARL; compare with the table.
\end{exercise}

\begin{exercise}
Twenty-two successive observations on viscosity have $\bar{x}=502.4$ and $\overline{mr}=12.3$. Compute limits for the individuals chart and the moving-range chart, and say whether an observation of $543$ would signal.
\end{exercise}

\begin{exercise}
Twenty-five samples of $n=100$ items contain $175$ defective items in total. Construct a $P$ chart. If the next sample has $13$ defectives is the process in control? What sample size would make the lower control limit positive?
\end{exercise}

\begin{exercise}
A process has USL $=24$, LSL $=12$, $\mu=19$ and $\sigma=1.2$. Compute PCR and PCR$_k$, the fractions nonconforming above USL and below LSL, and comment on centering.
\end{exercise}

\begin{exercise}
Use the tabular CUSUM with $\mu_0=100$, $K=2$ and $H=8$ on the data $101,103,104,102,106,105,108,107,109$. When is a signal given, and what is the estimated new mean?
\end{exercise}

\begin{exercise}
Explain how a Pareto chart and a fishbone diagram would be used in a project to reduce solder defects on circuit boards, and where a control chart and a designed experiment would then enter.
\end{exercise}


\part{Supplementary Chapters}
\chapter{Supplement: Sampling Distributions and Descriptive Statistics}

\begin{nuance}{place in the course}
This is a supplementary chapter. It is not part of the textbook's Chapter~7. It collects what the textbook states without proof or leaves to earlier chapters:
\begin{itemize}
  \item descriptive statistics;
  \item the sampling distribution of the sample mean;
  \item the exact distributions of $S^2$, $t$ and $F$ for normal samples, with the proofs the textbook only cites.
\end{itemize}
Chapters~1--3 refer to it.
\end{nuance}

Probability starts from a model and asks what data it is likely to produce. Statistics runs the other way: it starts from data and asks which model could have produced it.

The bridge is the \emph{sampling distribution}. A statistic computed from random data is itself random. Its distribution tells us how far the statistic can be trusted.

This chapter does four things:
\begin{itemize}
  \item sets up the language of populations and samples;
  \item summarises data;
  \item derives the exact sampling distributions of the mean and variance of a normal sample;
  \item introduces the chi-square, $t$ and $F$ distributions, which carry almost all of classical inference.
\end{itemize}

\section{Populations, samples and statistics}

\begin{itemize}
  \item A \emph{population} is the full collection of units of interest (all bolts from a factory, all future customers).
  \item A \emph{parameter} is a fixed number describing it, such as the mean $\mu$ or a proportion $p$.
  \item We rarely observe the whole population, so we observe a \emph{sample}.
\end{itemize}

\begin{definition}[Random sample]
A \emph{random sample} of size $n$ from a distribution $F$ is a set $X_1,\dots,X_n$ of independent random variables, each with distribution $F$. We say the $X_i$ are \emph{i.i.d.} (independent and identically distributed). The observed values $x_1,\dots,x_n$ are the \emph{data}.
\end{definition}

\begin{definition}[Statistic, sampling distribution]
A \emph{statistic} is any function $T=T(X_1,\dots,X_n)$ of the sample that does not involve unknown parameters. Its \emph{sampling distribution} is the distribution of the values $T$ would take over repeated samples of the same size from the same population.
\end{definition}

\begin{intuition}
One data set gives one value of $T$. The sampling distribution is the whole cloud of values you would see if you could repeat the experiment many times.
\begin{itemize}
  \item Inference means using the one value you have to work out what the cloud looks like.
  \item A narrow cloud means a reliable statistic.
\end{itemize}
\end{intuition}

\begin{pitfall}{i.i.d. is a modelling assumption, not a fact}
Almost every formula in this book assumes independent, identically distributed observations. These break it:
\begin{itemize}
  \item clustered data (students in the same school);
  \item repeated measurements on the same unit;
  \item time series;
  \item convenience samples.
\end{itemize}
Using i.i.d. formulas on such data usually understates the uncertainty, sometimes by a lot. How the data were collected matters more than any calculation that follows.
\end{pitfall}

\section{Describing a data set}

Before fitting any model, look at the data. The standard summaries follow.

\begin{definition}[Sample mean, variance, standard deviation]
For data $x_1,\dots,x_n$:
\[ \bar{x}=\frac1n\sum_{i=1}^nx_i,\qquad s^2=\frac1{n-1}\sum_{i=1}^n(x_i-\bar{x})^2,\qquad s=\sqrt{s^2}. \]
\end{definition}

The divisor $n-1$ (not $n$) makes $s^2$ an unbiased estimator of the population variance. Chapter~1 (Section 1.2) explains why.

A handy identity:
\[ \sum_{i=1}^n(x_i-\bar{x})^2=\sum_{i=1}^nx_i^2-n\bar{x}^2. \]
For numerical safety, prefer a two-pass or Welford-type algorithm. The identity subtracts two large, nearly equal numbers.

\begin{definition}[Order statistics, median, quantiles]
Sort the data: $x_{(1)}\le x_{(2)}\le\dots\le x_{(n)}$.
\begin{itemize}
  \item The \emph{median} is $x_{((n+1)/2)}$ for odd $n$, and $\tfrac12(x_{(n/2)}+x_{(n/2+1)})$ for even $n$.
  \item The \emph{$p$-th sample quantile} has a fraction $p$ of the data at or below it. (There are several interpolation conventions.)
  \item The quartiles $Q_1,Q_2,Q_3$ sit at $p=0.25,0.5,0.75$.
  \item The \emph{interquartile range} is $\mathrm{IQR}=Q_3-Q_1$.
\end{itemize}
\end{definition}

\begin{example}
Ten shaft diameters (mm): $9.9,\ 10.1,\ 10.0,\ 10.2,\ 9.8,\ 10.0,\ 10.1,\ 9.9,\ 10.0,\ 12.0$.
\begin{itemize}
  \item Mean: $\bar{x}=10.2$. Standard deviation: $s\approx0.64$.
  \item Sorted: $9.8,9.9,9.9,10.0,10.0,10.0,10.1,10.1,10.2,12.0$. Median $=10.0$. IQR $\approx0.2$.
  \item The single outlier $12.0$ pulls the mean up to $10.2$ and inflates $s$. It barely moves the median or the IQR.
\end{itemize}
\end{example}

\begin{theorem}[Robustness comparison]\label{thm:robust}
The mean and standard deviation have \emph{breakdown point} $0$: one wild observation can push them arbitrarily far. The median has breakdown point $1/2$: nearly half the data can be corrupted without moving the median outside the range of the clean data.
\end{theorem}

\begin{proof}
Replace $x_n$ by $t\to\infty$. Then $\bar{x}\to\infty$ and $s\to\infty$. For the median, as long as fewer than $\lceil n/2\rceil$ points are altered, the middle order statistic stays between two clean order statistics.
\end{proof}

\begin{interview}{mean or median?}
\begin{itemize}
  \item Use the mean for symmetric, light-tailed data when efficiency matters.
  \item Use the median and IQR for skewed or contaminated data: incomes, latencies, response times.
  \item Report both when they differ strongly. The gap itself says something about skew.
  \item Related fact: the mean minimises $\sum(x_i-c)^2$ over $c$. The median minimises $\sum\abs{x_i-c}$.
\end{itemize}
\end{interview}

\subsection*{Graphical summaries}

\begin{itemize}
  \item A \emph{histogram} splits the range into bins and counts the data in each. Its shape shows skew, several modes and outliers.
  \item A \emph{box plot} shows the median, the quartiles, and points beyond $1.5\,\mathrm{IQR}$ from the box as possible outliers.
  \item A \emph{normal quantile (Q--Q) plot} graphs the ordered data $x_{(i)}$ against normal quantiles $\Phi^{-1}\big(\tfrac{i-0.5}{n}\big)$. Roughly normal data lie near a straight line. Systematic curvature means skew or heavy tails.
\end{itemize}

\begin{figure}[h]
\centering
\begin{tikzpicture}[scale=1.0]
  \draw[axis] (-3.4,0)--(3.6,0); \draw[axis] (0,-0.2)--(0,2.6);
  \draw[thick,accent2,domain=-3.2:3.2,samples=60] plot (\x,{1.2*exp(-\x*\x/2)});
  \draw[thick,accent,domain=-3.2:3.2,samples=80] plot (\x,{2.4*exp(-2*\x*\x)});
  \node[lab,accent2] at (2.3,0.9){one observation}; \node[lab,accent] at (-1.7,2.2){$\bar{X}$, four observations};
  \node[lab] at (0,-0.6){the distribution of $\bar{X}$ concentrates as $n$ grows};
\end{tikzpicture}
\caption{Two bell curves centred at the true mean: the spread of one observation, and the narrower spread of the average of several observations.}
\end{figure}

\section{The sampling distribution of the sample mean}

\begin{theorem}[Mean and variance of $\bar{X}$]\label{thm:xbar}
If $X_1,\dots,X_n$ are i.i.d. with mean $\mu$ and variance $\sigma^2$, then
\[ \Exp[\bar{X}]=\mu,\qquad\Var(\bar{X})=\frac{\sigma^2}n,\qquad\se(\bar{X})=\frac{\sigma}{\sqrt n}. \]
\end{theorem}

\begin{proof}
By linearity, $\Exp[\bar{X}]=\frac1n\sum\mu=\mu$. By independence the variances add: $\Var(\bar{X})=\frac1{n^2}\sum\sigma^2=\sigma^2/n$.
\end{proof}

The \emph{standard error} of a statistic is its standard deviation, $\se(T)=\sqrt{\Var[T]}$. The $1/\sqrt n$ law is the most important rate in statistics: to halve the uncertainty you need four times the data.

\begin{theorem}[Normal samples]\label{thm:normalmean}
If $X_1,\dots,X_n$ are i.i.d. $\Normal(\mu,\sigma^2)$, then $\bar{X}\sim\Normal(\mu,\sigma^2/n)$ exactly.
\end{theorem}

\begin{proof}
A linear combination of independent normals is normal. (Moment generating functions multiply: $M_{\bar{X}}(t)=\prod M_{X_i}(t/n)=e^{\mu t+\sigma^2t^2/(2n)}$.) The mean and variance are as computed above.
\end{proof}

Without normality, the central limit theorem gives an approximation.

\begin{theorem}[CLT for the sample mean]\label{thm:clt}
If the $X_i$ are i.i.d. with mean $\mu$ and finite variance $\sigma^2>0$, then $\sqrt n(\bar{X}-\mu)/\sigma\to\Normal(0,1)$ in distribution as $n\to\infty$.
\end{theorem}

The CLT is proved in the Probability book (Limit Theorems). In practice:
\begin{itemize}
  \item For moderately skewed data, $n\ge30$ is a common rule of thumb.
  \item For heavily skewed or heavy-tailed data, much larger samples may be needed.
  \item For distributions with no finite variance (the Cauchy), the theorem does not apply. The mean of $n$ Cauchy variables has the same distribution as one.
\end{itemize}

\section{The chi-square distribution and the sample variance}

\begin{definition}[Chi-square distribution]
If $Z_1,\dots,Z_\nu$ are i.i.d. $\Normal(0,1)$, then $V=Z_1^2+\dots+Z_\nu^2$ has the \emph{chi-square distribution with $\nu$ degrees of freedom}, $V\sim\chi^2_\nu$.

It is the gamma distribution with shape $\nu/2$ and scale $2$. Its density is $f(v)=\dfrac{v^{\nu/2-1}e^{-v/2}}{2^{\nu/2}\Gamma(\nu/2)}$ for $v>0$.
\end{definition}

\begin{proposition}[Properties of $\chi^2_\nu$]\label{prop:chisq}
$\Exp[V]=\nu$ and $\Var[V]=2\nu$. If $V_1\sim\chi^2_{\nu_1}$ and $V_2\sim\chi^2_{\nu_2}$ are independent, then $V_1+V_2\sim\chi^2_{\nu_1+\nu_2}$.
\end{proposition}

\begin{proof}
$\Exp[Z^2]=1$ and $\Var[Z^2]=\Exp[Z^4]-1=3-1=2$. Then $\nu$ independent terms add their means and variances. Additivity is clear from the definition as a sum of squares (or from the gamma moment generating function $(1-2t)^{-\nu/2}$).
\end{proof}

The next result is the cornerstone of exact normal-theory inference.

\begin{theorem}[Distribution of $\bar{X}$ and $S^2$ for a normal sample]\label{thm:cochran}
Let $X_1,\dots,X_n$ be i.i.d. $\Normal(\mu,\sigma^2)$, with $S^2=\frac1{n-1}\sum(X_i-\bar{X})^2$. Then
\begin{enumerate}[label=(\roman*)]
  \item $\bar{X}$ and $S^2$ are independent;
  \item $\dfrac{(n-1)S^2}{\sigma^2}\sim\chi^2_{n-1}$.
\end{enumerate}
\end{theorem}

\begin{proof}
Standardise: $Z_i=(X_i-\mu)/\sigma$ are i.i.d. $\Normal(0,1)$, so $\mathbf{Z}=(Z_1,\dots,Z_n)\T\sim\Normal(\mathbf{0},I_n)$.

Choose an orthogonal matrix $Q$ whose first row is $\mathbf{1}\T/\sqrt n$. (Extend that unit vector to an orthonormal basis; see Chapter~11 of the Linear Algebra book.) Put $\mathbf{Y}=Q\mathbf{Z}$. Since $Q$ is orthogonal, $\mathbf{Y}\sim\Normal(\mathbf{0},QQ\T)=\Normal(\mathbf{0},I_n)$. The $Y_i$ are again i.i.d. standard normal.

The first coordinate is $Y_1=\sqrt n\,\bar{Z}$. Orthogonal maps preserve length, so $\sum Y_i^2=\sum Z_i^2$. Therefore
\[ \sum_{i=1}^n(Z_i-\bar{Z})^2=\sum Z_i^2-n\bar{Z}^2=\sum_{i=1}^nY_i^2-Y_1^2=\sum_{i=2}^nY_i^2. \]
This is a sum of $n-1$ squared independent standard normals, so it is $\chi^2_{n-1}$. It depends only on $Y_2,\dots,Y_n$, while $\bar{Z}=Y_1/\sqrt n$ depends only on $Y_1$. So the two are independent.

Finally, $(n-1)S^2/\sigma^2=\sum(Z_i-\bar{Z})^2$ and $\bar{X}=\mu+\sigma\bar{Z}$.
\end{proof}

\begin{corollary}[Mean and variance of $S^2$]\label{cor:s2}
For normal samples, $\Exp[S^2]=\sigma^2$ and $\Var[S^2]=2\sigma^4/(n-1)$. (The unbiasedness $\Exp[S^2]=\sigma^2$ holds for every distribution with finite variance; see Chapter~1.)
\end{corollary}

\begin{proof}
$\Exp[\chi^2_{n-1}]=n-1$ and $\Var[\chi^2_{n-1}]=2(n-1)$. Then $S^2=\sigma^2V/(n-1)$ with $V\sim\chi^2_{n-1}$.
\end{proof}

\begin{nuance}{why $n-1$ degrees of freedom}
The $n$ deviations $X_i-\bar{X}$ must sum to zero. So only $n-1$ of them are free. They lie in an $(n-1)$-dimensional subspace: the orthogonal complement of $\mathbf{1}$.

That is what ``degrees of freedom'' means: the dimension of the subspace where the sum of squares lives. The same count reappears in regression and ANOVA (Chapters~5--8). Estimating $p$ parameters removes $p$ degrees of freedom.
\end{nuance}

\section{The $t$ and $F$ distributions}

The standardised mean $(\bar{X}-\mu)/(\sigma/\sqrt n)$ is standard normal. In practice $\sigma$ is unknown, and we replace it by $S$. The distribution changes.

\begin{definition}[Student's $t$ distribution]
If $Z\sim\Normal(0,1)$ and $V\sim\chi^2_\nu$ are independent, then $T=Z/\sqrt{V/\nu}$ has the \emph{$t$ distribution with $\nu$ degrees of freedom}, $T\sim t_\nu$.
\end{definition}

\begin{theorem}[The $t$ statistic]\label{thm:tstat}
For a normal sample, $T=\dfrac{\bar{X}-\mu}{S/\sqrt n}\sim t_{n-1}$.
\end{theorem}

\begin{proof}
Write
\[ T=\dfrac{(\bar{X}-\mu)/(\sigma/\sqrt n)}{\sqrt{S^2/\sigma^2}}=\dfrac{Z}{\sqrt{V/(n-1)}}, \]
with $Z=\sqrt n(\bar{X}-\mu)/\sigma\sim\Normal(0,1)$ and $V=(n-1)S^2/\sigma^2\sim\chi^2_{n-1}$. They are independent by Theorem~\ref{thm:cochran}(i).
\end{proof}

Properties of $t_\nu$:
\begin{itemize}
  \item Symmetric about $0$, with density $f(t)=\dfrac{\Gamma(\frac{\nu+1}2)}{\sqrt{\nu\pi}\,\Gamma(\frac\nu2)}\big(1+\tfrac{t^2}\nu\big)^{-(\nu+1)/2}$.
  \item Heavier tails than the normal. This is the price of estimating $\sigma$.
  \item $\Var[T]=\nu/(\nu-2)$ for $\nu>2$.
  \item $t_1$ is the Cauchy distribution.
  \item $t_\nu\to\Normal(0,1)$ as $\nu\to\infty$, because $V/\nu\to1$ by the law of large numbers.
\end{itemize}

\begin{figure}[h]
\centering
\begin{tikzpicture}[scale=1.0]
  \draw[axis] (-4.4,0)--(4.6,0); \draw[axis] (0,-0.2)--(0,2.8);
  \draw[thick,accent2,domain=-4:4,samples=70] plot (\x,{2.6*exp(-\x*\x/2)});
  \draw[thick,accent,domain=-4:4,samples=70] plot (\x,{2.6*0.9*(1+\x*\x/3)^(-2)});
  \node[lab,accent2] at (1.9,1.5){$\Normal(0,1)$}; \node[lab,accent] at (2.4,0.5){$t_3$: heavier tails};
\end{tikzpicture}
\caption{The $t_3$ density has a lower peak and heavier tails than the standard normal. Large observations of $T$ are much more likely when $\sigma$ must be estimated from a small sample.}
\end{figure}

\begin{definition}[$F$ distribution]
If $U\sim\chi^2_{a}$ and $V\sim\chi^2_{b}$ are independent, then $F=\dfrac{U/a}{V/b}$ has the \emph{$F$ distribution with $(a,b)$ degrees of freedom}, $F\sim F_{a,b}$.
\end{definition}

Two quick consequences:
\begin{itemize}
  \item $1/F\sim F_{b,a}$.
  \item If $T\sim t_\nu$, then $T^2\sim F_{1,\nu}$, because $T^2=Z^2/(V/\nu)$ with $Z^2\sim\chi^2_1$.
\end{itemize}

\begin{theorem}[Ratio of sample variances]\label{thm:fratio}
Let $X_1,\dots,X_m$ be i.i.d. $\Normal(\mu_1,\sigma_1^2)$, and independently $Y_1,\dots,Y_n$ i.i.d. $\Normal(\mu_2,\sigma_2^2)$, with sample variances $S_1^2,S_2^2$. Then
\[ \frac{S_1^2/\sigma_1^2}{S_2^2/\sigma_2^2}\sim F_{m-1,\,n-1}. \]
\end{theorem}

\begin{proof}
By Theorem~\ref{thm:cochran}, $(m-1)S_1^2/\sigma_1^2\sim\chi^2_{m-1}$ and $(n-1)S_2^2/\sigma_2^2\sim\chi^2_{n-1}$. They are independent because the samples are. Divide each by its degrees of freedom.
\end{proof}

\begin{interview}{why is normality so central here?}
Theorems~\ref{thm:cochran}--\ref{thm:fratio} are \emph{exact} only for normal data. In fact, independence of $\bar{X}$ and $S^2$ \emph{characterises} the normal distribution.
\begin{itemize}
  \item For non-normal data, the results for the mean become large-sample approximations through the CLT.
  \item The $\chi^2$ and $F$ results for variances are \emph{not} robust to non-normality, because they depend on the fourth moment.
  \item When normality is doubtful, use the bootstrap or permutation methods (Chapter~13).
\end{itemize}
\end{interview}

\section{Sampling distributions in practice}

\begin{example}[Using the $t$ distribution]
A sample of $n=9$ from a normal population has $\bar{x}=52$ and $s=6$.
\begin{itemize}
  \item Standard error: $s/\sqrt n=2$. The statistic $T=(\bar{X}-\mu)/2$ follows $t_8$.
  \item $t_{8,0.025}\approx2.306$. So with probability $0.95$, $\abs{\bar{X}-\mu}\le2.306\times2=4.61$.
  \item The normal approximation would give $1.96\times2=3.92$. Ignoring the $t$ correction would understate the uncertainty by about $18\%$.
\end{itemize}
\end{example}

\begin{example}[Probability about a sample variance]
A normal population has $\sigma^2=4$. For a sample of $n=10$:
\[ \Prob(S^2>6)=\Prob\big(\chi^2_9>9\cdot6/4=13.5\big)\approx0.14. \]
\end{example}

\begin{summarybox}
\begin{itemize}
  \item A statistic is a function of the sample. Its sampling distribution (what it does over repeated samples) is the basis of all inference.
  \item Descriptive statistics: mean, variance, median, quantiles, IQR, histograms, box plots, Q--Q plots. The median and IQR are robust. The mean and variance are not.
  \item $\Exp[\bar{X}]=\mu$ and $\Var[\bar{X}]=\sigma^2/n$. For normal data $\bar{X}\sim\Normal(\mu,\sigma^2/n)$ exactly, and approximately for large $n$ by the CLT.
  \item For a normal sample, $\bar{X}$ and $S^2$ are independent and $(n-1)S^2/\sigma^2\sim\chi^2_{n-1}$. The proof is an orthogonal change of variables.
  \item $t=Z/\sqrt{V/\nu}$ appears when $\sigma$ is estimated: $(\bar{X}-\mu)/(S/\sqrt n)\sim t_{n-1}$. The ratio of independent scaled $\chi^2$ variables is $F$, which governs variance comparisons.
\end{itemize}
\end{summarybox}

\section{Exercises}

\begin{exercise}
For the data $3,5,5,7,8,9,12$ compute the mean, median, variance, standard deviation, quartiles and IQR. What happens to each if $12$ is replaced by $120$?
\end{exercise}

\begin{exercise}
Prove that $\sum(x_i-\bar{x})=0$ and that the mean minimises $\sum(x_i-c)^2$. Prove that the median minimises $\sum\abs{x_i-c}$ for $n$ odd.
\end{exercise}

\begin{exercise}
Show that if $Y=aX+b$ then $\bar{y}=a\bar{x}+b$ and $s_y=\abs{a}s_x$. What do the quartiles and the IQR do under this transformation?
\end{exercise}

\begin{exercise}
A machine fills bottles with mean $500$ ml and standard deviation $4$ ml. What is the approximate probability that the mean of $64$ bottles exceeds $501$ ml? What assumptions are used?
\end{exercise}

\begin{exercise}
If $Z_1,Z_2,Z_3$ are i.i.d. standard normal, find the distribution of $(Z_1-Z_2)^2/2$ and of $(Z_1+Z_2+Z_3)^2/3$. Are they independent? Justify.
\end{exercise}

\begin{exercise}
For a sample of size $n=16$ from $\Normal(\mu,\sigma^2)$, find $c$ so that $\Prob(S^2\le c\sigma^2)=0.95$ (use $\chi^2_{15}$ quantiles from software).
\end{exercise}

\begin{exercise}
Prove that $T\sim t_\nu$ has $\Exp[T]=0$ for $\nu>1$ and $\Var[T]=\nu/(\nu-2)$ for $\nu>2$ by writing $T=Z/\sqrt{V/\nu}$ and using $\Exp(1/V)=1/(\nu-2)$.
\end{exercise}

\begin{exercise}
Simulate $10{,}000$ samples of size $5$ from a normal population and from an exponential population; compute the $t$ statistic for each and compare histograms with the $t_4$ density. What do you conclude about the effect of non-normality at small $n$?
\end{exercise}

\chapter{Supplement: Estimation Theory}

\begin{nuance}{place in the course}
This supplementary chapter proves what Chapter~1 states without proof:
\begin{itemize}
  \item why the sample mean is the MVUE for a normal mean;
  \item how much precision is possible at all (the Cram\'er--Rao bound);
  \item why the maximum likelihood estimator is asymptotically optimal;
  \item what it means for a statistic to carry all the information about a parameter (sufficiency).
\end{itemize}
It is for readers who want the mathematics behind Chapter~1. The textbook's problems do not need it.
\end{nuance}

\section{Score function and Fisher information}

Let $X$ have density (or mass function) $f(x;\theta)$ with a scalar parameter $\theta$. Assume the usual \emph{regularity conditions}:
\begin{itemize}
  \item the support of $f$ does not depend on $\theta$;
  \item $f$ is twice differentiable in $\theta$;
  \item differentiation and integration can be interchanged.
\end{itemize}

\begin{definition}[Score and Fisher information]
\begin{itemize}
  \item The \emph{score} is $U(\theta)=\dfrac{\partial}{\partial\theta}\ln f(X;\theta)$.
  \item The \emph{Fisher information} of one observation is $\Fish(\theta)=\Var\big(U(\theta)\big)$.
  \item For a random sample of size $n$ the information is $n\Fish(\theta)$.
\end{itemize}
\end{definition}

\begin{theorem}[Properties of the score]\label{thm:score}
Under regularity, $\Exp[U(\theta)]=0$ and
\[ \Fish(\theta)=\Exp\Big[\Big(\frac{\partial\ln f}{\partial\theta}\Big)^2\Big]=-\Exp\Big[\frac{\partial^2\ln f}{\partial\theta^2}\Big]. \]
\end{theorem}

\begin{proof}
Since $\int f(x;\theta)\,dx=1$, differentiating gives $0=\int\partial_\theta f\,dx=\int(\partial_\theta\ln f)f\,dx=\Exp[U]$.

Differentiate $\int(\partial_\theta\ln f)f\,dx=0$ once more:
\[ 0=\int(\partial_\theta^2\ln f)f\,dx+\int(\partial_\theta\ln f)\partial_\theta f\,dx=\Exp[\partial_\theta^2\ln f]+\Exp[U^2]. \]
Since $\Exp[U]=0$, $\Var[U]=\Exp[U^2]=-\Exp[\partial_\theta^2\ln f]$.
\end{proof}

The second form says: information is the expected \emph{curvature} of the log-likelihood. A sharply peaked likelihood pins down $\theta$ precisely.

\begin{example}[Standard cases]\label{ex:fisher}
\begin{itemize}
  \item Bernoulli$(p)$: $\ln f=x\ln p+(1-x)\ln(1-p)$ and $\partial^2\ln f/\partial p^2=-x/p^2-(1-x)/(1-p)^2$. So $\Fish(p)=1/p+1/(1-p)=\dfrac1{p(1-p)}$.
  \item Poisson$(\lambda)$: $\ln f=x\ln\lambda-\lambda-\ln x!$. $\Fish(\lambda)=\Exp[X]/\lambda^2=1/\lambda$.
  \item Exponential with rate $\lambda$: $\ln f=\ln\lambda-\lambda x$. $\Fish(\lambda)=1/\lambda^2$.
  \item $\Normal(\mu,\sigma^2)$: with $\sigma^2$ known, $\Fish(\mu)=1/\sigma^2$. With $\mu$ known, $\Fish(\sigma^2)=1/(2\sigma^4)$.
\end{itemize}
\end{example}

\section{The Cram\'er--Rao lower bound}

\begin{theorem}[Cram\'er--Rao]\label{thm:cramerrao}
Let $X_1,\dots,X_n$ be a random sample and $\hat\Theta$ an unbiased estimator of $\theta$. Under regularity,
\[ \Var(\hat\Theta)\ \ge\ \frac1{n\Fish(\theta)}. \]
More generally, if $\Exp[\hat\Theta]=\theta+b(\theta)$, then $\Var(\hat\Theta)\ge[1+b'(\theta)]^2/(n\Fish(\theta))$.
\end{theorem}

\begin{proof}
Let $U_n=\sum_i\partial_\theta\ln f(X_i;\theta)$ be the score of the whole sample. Then $\Exp[U_n]=0$ and $\Var[U_n]=n\Fish(\theta)$.

For unbiased $\hat\Theta$, $\theta=\Exp[\hat\Theta]=\int\hat\theta(\mathbf{x})f(\mathbf{x};\theta)\,d\mathbf{x}$. Differentiate in $\theta$:
\[ 1=\int\hat\theta\,\partial_\theta f\,d\mathbf{x}=\Exp[\hat\Theta U_n]=\Cov(\hat\Theta,U_n), \]
using $\Exp[U_n]=0$. By Cauchy--Schwarz for covariances,
\[ 1=\Cov(\hat\Theta,U_n)^2\le\Var(\hat\Theta)\Var(U_n)=\Var(\hat\Theta)\,n\Fish(\theta). \]
\end{proof}

\begin{definition}[Efficient estimator]
An unbiased estimator whose variance equals the Cram\'er--Rao bound is \emph{efficient}. The \emph{efficiency} of an unbiased estimator is the bound divided by its variance. It is at most $1$.
\end{definition}

\begin{theorem}[$\bar{X}$ is the MVUE for a normal mean]\label{thm:mvue-proof}
For $X_i\sim\Normal(\mu,\sigma^2)$ with $\sigma^2$ known, $\bar{X}$ is efficient: $\Var[\bar{X}]=\sigma^2/n=1/(n\Fish(\mu))$. Hence $\bar{X}$ is the MVUE of $\mu$ (Theorem~\ref{thm:mvue} of Chapter~1).
\end{theorem}

\begin{example}
\begin{itemize}
  \item Bernoulli sample: $\Var[\hat{P}]=p(1-p)/n=1/(n\Fish(p))$. The sample proportion is efficient.
  \item $\Normal(\mu,\sigma^2)$: $\Var[S^2]=2\sigma^4/(n-1)$, while the bound for unbiased estimators of $\sigma^2$ is $2\sigma^4/n$. So $S^2$ is not efficient, but its efficiency $(n-1)/n\to1$.
\end{itemize}
\end{example}

\begin{pitfall}{when the bound does not apply}
The regularity conditions fail when the support depends on the parameter.

For the uniform distribution on $(0,\theta)$, the unbiased estimator $\frac{n+1}n\max X_i$ has variance $\theta^2/(n(n+2))$, of order $1/n^2$. That is \emph{smaller} than the order-$1/n$ value the formal Cram\'er--Rao expression would give. There is no contradiction: the bound is not valid for this family.
\end{pitfall}

\section{Asymptotic properties of the maximum likelihood estimator}

\begin{theorem}[Asymptotic normality and efficiency of the MLE]\label{thm:mle-asym}
Under regularity, the MLE $\hat\Theta_n$ from $n$ i.i.d.\ observations is consistent, and
\[ \sqrt n\big(\hat\Theta_n-\theta\big)\xrightarrow{d}\Normal\Big(0,\frac1{\Fish(\theta)}\Big). \]
So for large $n$, $\hat\Theta_n\approx\Normal\big(\theta,\ 1/(n\Fish(\theta))\big)$. The MLE reaches the Cram\'er--Rao bound asymptotically.
\end{theorem}

\begin{proof}[Sketch]
Let $\ell(\theta)=\sum\ln f(X_i;\theta)$. Expand the likelihood equation $\ell'(\hat\theta)=0$ about the true $\theta$:
\[ 0=\ell'(\theta)+(\hat\theta-\theta)\ell''(\tilde\theta)\quad\text{for some }\tilde\theta\text{ between }\hat\theta\text{ and }\theta. \]
Hence
\[ \sqrt n(\hat\theta-\theta)=\frac{n^{-1/2}\ell'(\theta)}{-n^{-1}\ell''(\tilde\theta)}. \]
\begin{itemize}
  \item The numerator is $n^{-1/2}$ times a sum of i.i.d.\ scores with mean $0$ and variance $\Fish(\theta)$. By the CLT it tends to $\Normal(0,\Fish(\theta))$.
  \item The denominator tends in probability to $-\Exp[\partial^2\ln f]=\Fish(\theta)$, by the law of large numbers. (Consistency lets us replace $\tilde\theta$ by $\theta$.)
  \item By Slutsky's theorem the ratio tends to $\Normal(0,\Fish(\theta)/\Fish(\theta)^2)=\Normal(0,1/\Fish(\theta))$.
\end{itemize}
\end{proof}

\begin{corollary}[Large-sample interval and Wald test]\label{cor:wald}
\begin{itemize}
  \item Approximate two-sided $100(1-\alpha)\%$ interval: $\hat\theta\pm z_{\alpha/2}/\sqrt{n\Fish(\hat\theta)}$.
  \item One-sided $100(1-\alpha)\%$ bounds use $z_\alpha$ and one limit: $\theta\le\hat\theta+z_\alpha/\sqrt{n\Fish(\hat\theta)}$, or $\hat\theta-z_\alpha/\sqrt{n\Fish(\hat\theta)}\le\theta$.
  \item To test $H_0:\theta=\theta_0$, use $Z=(\hat\theta-\theta_0)\sqrt{n\Fish(\theta_0)}$ (or $\hat\theta$ in place of $\theta_0$ inside $\Fish$). It is approximately $\Normal(0,1)$ under $H_0$.
\end{itemize}
\end{corollary}

\begin{example}
$n=200$ exponential observations with $\bar{x}=4$.
\begin{itemize}
  \item $\hat\lambda=0.25$ and $\Fish(\lambda)=1/\lambda^2$.
  \item $\se(\hat\lambda)=\hat\lambda/\sqrt n=0.25/14.14=0.0177$.
  \item $95\%$ interval: $0.25\pm0.035$.
\end{itemize}
\end{example}

\begin{theorem}[Delta method]\label{thm:delta}
If $\sqrt n(\hat\theta-\theta)\to\Normal(0,v)$ and $g$ is differentiable with $g'(\theta)\ne0$, then
\[ \sqrt n\big(g(\hat\theta)-g(\theta)\big)\to\Normal\big(0,\ g'(\theta)^2v\big). \]
\end{theorem}

\begin{proof}
$g(\hat\theta)-g(\theta)=g'(\tilde\theta)(\hat\theta-\theta)$ for some $\tilde\theta$ between them. As $\tilde\theta\to\theta$, $g'(\tilde\theta)\to g'(\theta)$. Apply Slutsky.
\end{proof}

\begin{example}
The MLE of $\sigma$ in a normal sample is $\hat\sigma=\sqrt{\hat\sigma^2}$, where $\sqrt n(\hat\sigma^2-\sigma^2)\to\Normal(0,2\sigma^4)$. With $g(t)=\sqrt t$ we get $g'(\sigma^2)=1/(2\sigma)$. So $\sqrt n(\hat\sigma-\sigma)\to\Normal(0,\sigma^2/2)$ and $\se(\hat\sigma)\approx\sigma/\sqrt{2n}$.
\end{example}

\section{Sufficiency}

\begin{definition}[Sufficient statistic]
A statistic $T=T(X_1,\dots,X_n)$ is \emph{sufficient} for $\theta$ if the conditional distribution of the sample given $T$ does not depend on $\theta$. Once $T$ is known, the rest of the data carry no more information about $\theta$.
\end{definition}

\begin{theorem}[Fisher--Neyman factorisation]\label{thm:factor}
$T$ is sufficient for $\theta$ if and only if the joint density factors as
\[ f(x_1,\dots,x_n;\theta)=g\big(T(\mathbf{x});\theta\big)\,h(\mathbf{x}), \]
where $h$ does not depend on $\theta$.
\end{theorem}

\begin{example}[Standard sufficient statistics]
\begin{itemize}
  \item Bernoulli$(p)$: $f=p^{\sum x_i}(1-p)^{n-\sum x_i}$. So $T=\sum X_i$ (the number of successes).
  \item Poisson$(\lambda)$: $f=e^{-n\lambda}\lambda^{\sum x_i}/\prod x_i!$. So $T=\sum X_i$.
  \item Exponential$(\lambda)$: $f=\lambda^ne^{-\lambda\sum x_i}$. So $T=\sum X_i$.
  \item $\Normal(\mu,\sigma^2)$: $\sum(x_i-\mu)^2=\sum(x_i-\bar{x})^2+n(\bar{x}-\mu)^2$ shows $(\bar{X},S^2)$ is jointly sufficient for $(\mu,\sigma^2)$. For $\mu$ alone ($\sigma^2$ known), $\bar{X}$ is sufficient.
  \item Uniform$(0,\theta)$: $f=\theta^{-n}\mathbf{1}\{\max x_i\le\theta\}$. So $T=\max X_i$.
\end{itemize}
\end{example}

By the factorisation, the likelihood depends on the data only through $T$. So the MLE is a function of any sufficient statistic. In every example above, the MLE is a function of $T$.

\begin{theorem}[Rao--Blackwell]\label{thm:rb}
Let $\hat\Theta$ be an unbiased estimator of $\theta$ and $T$ a sufficient statistic. Then $\tilde\Theta=\Exp[\hat\Theta\mid T]$ is an estimator (sufficiency makes it free of $\theta$), it is unbiased, and
\[ \Var(\tilde\Theta)\le\Var(\hat\Theta), \]
with equality only if $\hat\Theta$ is already a function of $T$.
\end{theorem}

\begin{proof}
Unbiased: $\Exp[\tilde\Theta]=\Exp\big[\Exp[\hat\Theta\mid T]\big]=\Exp[\hat\Theta]=\theta$.

By the law of total variance, $\Var[\hat\Theta]=\Var(\Exp[\hat\Theta\mid T])+\Exp[\Var(\hat\Theta\mid T)]\ge\Var[\tilde\Theta]$.
\end{proof}

\begin{example}
For a Bernoulli sample, $\hat\Theta=X_1$ is an unbiased (and very poor) estimator of $p$. With $T=\sum X_i$, $\Exp[X_1\mid T]=T/n=\bar{X}$, the sample proportion. Its variance $p(1-p)/n$ is far below $p(1-p)$. Conditioning on a sufficient statistic always helps.
\end{example}

\begin{definition}[Complete statistic]
A statistic $T$ is \emph{complete} if $\Exp_\theta[\phi(T)]=0$ for all $\theta$ implies $\phi(T)=0$ with probability one. For one-parameter exponential families (normal with known variance, Bernoulli, Poisson, exponential), the natural sufficient statistic is complete.
\end{definition}

\begin{theorem}[Lehmann--Scheff\'e]\label{thm:lehmann}
If $T$ is complete and sufficient and $\phi(T)$ is an unbiased estimator of $\theta$, then $\phi(T)$ is the unique MVUE of $\theta$.
\end{theorem}

\begin{proof}[Idea]
Take any unbiased estimator $\hat\Theta$. By Rao--Blackwell, $\Exp[\hat\Theta\mid T]$ is a function of $T$ with no larger variance. Two unbiased functions of $T$ differ by a function with zero mean, which is zero by completeness. So the Rao--Blackwellised estimator is unique, and it has minimum variance.
\end{proof}

\begin{example}
For $X_i\sim\Normal(\mu,\sigma^2)$, the MVUE of $\mu$ is $\bar{X}$ and the MVUE of $\sigma^2$ is $S^2$. Both are unbiased functions of the complete sufficient statistic $(\bar{X},S^2)$. This settles the questions of Chapter~1. For the uniform$(0,\theta)$ family, $\frac{n+1}n\max X_i$ is the MVUE of $\theta$.
\end{example}

\section{Bayes estimators and loss functions}

A \emph{loss function} $L(\hat\theta,\theta)$ is the cost of estimating $\theta$ by $\hat\theta$. Standard choices: squared error $(\hat\theta-\theta)^2$ and absolute error $\abs{\hat\theta-\theta}$. The \emph{Bayes risk} of an estimator is the expected loss, averaged over the prior and the data.

\begin{theorem}[Bayes estimators]\label{thm:bayes}
With prior $\pi(\theta)$ and posterior $\pi(\theta\mid\mathbf{x})$, the estimator that minimises the Bayes risk is:
\begin{itemize}
  \item for squared-error loss, the \emph{posterior mean} $\Exp[\theta\mid\mathbf{x}]$;
  \item for absolute-error loss, the \emph{posterior median};
  \item for $0$--$1$ loss, the \emph{posterior mode} (MAP estimate).
\end{itemize}
\end{theorem}

\begin{proof}[Squared error]
Minimise $\Exp[(\theta-a)^2\mid\mathbf{x}]$ over $a$. Let $m=\Exp[\theta\mid\mathbf{x}]$ and write $(\theta-a)^2=(\theta-m)^2+2(\theta-m)(m-a)+(m-a)^2$. The cross term has zero posterior expectation. What remains is $\Var(\theta\mid\mathbf{x})+(m-a)^2$, minimised at $a=m$.
\end{proof}

\begin{example}[Conjugacy and shrinkage]
$X\sim\mathrm{Binomial}(n,p)$ with a $\mathrm{Beta}(a,b)$ prior gives the posterior $\mathrm{Beta}(a+x,b+n-x)$, with mean $\dfrac{a+x}{a+b+n}$ (Example~\ref{ex:betabin}). It is the weighted average
\[ w\cdot\frac xn+(1-w)\cdot\frac a{a+b},\qquad w=\frac n{n+a+b}. \]
\begin{itemize}
  \item The Bayes estimator shrinks the MLE toward the prior mean. The prior acts like $a+b$ extra pseudo-observations.
  \item As $n\to\infty$, $w\to1$, and the estimator approaches the MLE, with variance approaching the Cram\'er--Rao bound.
  \item For small $n$, shrinkage trades a little bias for a bigger drop in variance. This can lower the mean square error (compare Example~\ref{ex:mseexample}).
\end{itemize}
\end{example}

\begin{interview}{frequentist versus Bayesian in one paragraph}
A frequentist treats $\theta$ as a fixed unknown and judges an estimator by its sampling distribution (bias, variance, coverage of intervals). A Bayesian treats $\theta$ as random with a prior, and reports the posterior.

A $95\%$ \emph{credible interval} contains $\theta$ with posterior probability $0.95$. This is the statement people often wrongly read into a confidence interval.

For large samples with smooth priors the two approaches agree (Bernstein--von Mises: the posterior is approximately $\Normal(\hat\theta,1/(n\Fish(\hat\theta)))$). For small samples the prior matters.
\end{interview}

\begin{summarybox}
\begin{itemize}
  \item Fisher information $\Fish(\theta)=\Var(\partial_\theta\ln f)=-\Exp[\partial_\theta^2\ln f]$ is the curvature of the log-likelihood. It adds up: $n\Fish$ for a sample of size $n$.
  \item Cram\'er--Rao: every unbiased estimator has $\Var\ge1/(n\Fish(\theta))$. $\bar{X}$ (normal mean), $\hat{P}$ (proportion) and $\bar{X}$ (Poisson and exponential mean) reach it.
  \item The MLE is asymptotically $\Normal(\theta,1/(n\Fish))$: consistent, asymptotically unbiased and efficient. The delta method carries this to functions of $\theta$.
  \item Sufficient statistics (factorisation theorem) keep all the information. Rao--Blackwell improves any unbiased estimator by conditioning on one. Lehmann--Scheff\'e gives the unique MVUE from a complete sufficient statistic.
  \item Bayes estimators are posterior summaries (the mean, for squared error). Conjugate priors give closed forms and shrinkage toward the prior mean.
\end{itemize}
\end{summarybox}

\section{Exercises}

\begin{exercise}
Compute the Fisher information for (a) $\Normal(0,\sigma^2)$ with respect to $\sigma^2$, (b) the geometric distribution $f(x;p)=p(1-p)^{x-1}$, and give the Cram\'er--Rao bound for unbiased estimators of $p$ based on a sample of size $n$.
\end{exercise}

\begin{exercise}
Show that for a Poisson sample $\bar{X}$ attains the Cram\'er--Rao bound for estimating $\lambda$, and that $\bar{X}^2-\bar{X}/n$ is an unbiased estimator of $\lambda^2$. Is it efficient for $\lambda^2$?
\end{exercise}

\begin{exercise}
Prove the factorisation theorem in the discrete case (the easy direction: if the joint pmf factors as $g(T;\theta)h(\mathbf{x})$ then $T$ is sufficient).
\end{exercise}

\begin{exercise}
Show that $\sum X_i$ is sufficient for $\lambda$ in a Poisson sample, and use Rao--Blackwell to improve the estimator $\mathbf{1}\{X_1=0\}$ of $e^{-\lambda}$; find its improved form $\big(\frac{n-1}n\big)^{T}$.
\end{exercise}

\begin{exercise}
Using the delta method, find the approximate standard error of $\hat{p}/(1-\hat{p})$ (the odds) and of $\ln\big(\hat{p}/(1-\hat{p})\big)$ (the log-odds) for a binomial sample.
\end{exercise}

\begin{exercise}
In a Bayesian analysis of $n=50$ trials with $x=7$ successes and a $\mathrm{Beta}(1,1)$ (uniform) prior, find the posterior mean, the posterior mode, and compare with the MLE and with a Wald interval.
\end{exercise}

\begin{exercise}
For a $\Normal(\mu,\sigma^2)$ sample with known $\sigma^2$ and prior $\mu\sim\Normal(\mu_0,\tau^2)$, show that the Bayes estimator is $w\bar{x}+(1-w)\mu_0$ with $w=\dfrac{n/\sigma^2}{n/\sigma^2+1/\tau^2}$ and find its mean square error as a function of the true $\mu$. For which $\mu$ is it better than $\bar{X}$?
\end{exercise}

\begin{exercise}
Prove that the MLE of $\theta$ in the uniform$(0,\theta)$ family is biased and find the bias; show that the bias-corrected estimator has variance $\theta^2/(n(n+2))$.
\end{exercise}

\chapter{Supplement: Multiple Testing, Resampling and A/B Tests}

\begin{nuance}{place in the course}
This supplementary chapter covers topics that the textbook does not treat but that modern practice (and most interviews) require: what goes wrong when many hypotheses are tested, how to correct for it, how to do inference by resampling when no formula is available (the bootstrap and permutation tests), and how to design and read an A/B test. It builds on Chapters~2--4.
\end{nuance}

\section{The multiple testing problem}

A single test at level $\alpha=0.05$ has a $5\%$ chance of a false positive when $H_0$ is true. If we run $m$ independent tests of true null hypotheses, the probability of at least one false positive is
\[ 1-(1-\alpha)^m,\qquad\text{which is }0.40\text{ for }m=10\text{ and }0.994\text{ for }m=100. \]
Testing many hypotheses (many genes, many metrics, many subgroups, many models) and reporting only the significant ones is the main source of non-reproducible findings (``$p$-hacking'', or the ``garden of forking paths'').

\begin{definition}[Error rates for $m$ tests]
Let $V$ be the number of false positives (true nulls rejected) and $R$ the total number of rejections.
\begin{itemize}
  \item The \emph{family-wise error rate} is $\mathrm{FWER}=\Prob(V\ge1)$.
  \item The \emph{false discovery rate} is $\mathrm{FDR}=\Exp\Big[\dfrac V{\max(R,1)}\Big]$, the expected proportion of rejections that are false.
\end{itemize}
\end{definition}

\subsection{Controlling the family-wise error rate}

\begin{theorem}[Bonferroni correction]\label{thm:bonf}
Rejecting each hypothesis $H_{0i}$ whose $P$-value satisfies $p_i\le\alpha/m$ controls $\mathrm{FWER}\le\alpha$, whatever the dependence among the tests.
\end{theorem}

\begin{proof}
Let $I_0$ be the set of true nulls, $\abs{I_0}\le m$. By the union bound, $\Prob(V\ge1)=\Prob\big(\bigcup_{i\in I_0}\{p_i\le\alpha/m\}\big)\le\sum_{i\in I_0}\Prob(p_i\le\alpha/m)\le m\cdot\alpha/m=\alpha$, using that a $P$-value is (super)uniform under $H_0$.
\end{proof}

Bonferroni is simple but conservative. \emph{Holm's step-down procedure} is uniformly more powerful with the same guarantee: sort the $P$-values $p_{(1)}\le\dots\le p_{(m)}$; reject $H_{(1)},\dots,H_{(k-1)}$ where $k$ is the smallest index with $p_{(k)}>\alpha/(m-k+1)$ (and stop there).

\begin{example}
Five tests give $P$-values $0.003,\ 0.011,\ 0.020,\ 0.045,\ 0.30$. Bonferroni at $\alpha=0.05$ uses the threshold $0.01$ and rejects only the first. Holm compares $0.003$ with $0.05/5=0.010$ (reject), $0.011$ with $0.05/4=0.0125$ (reject), $0.020$ with $0.05/3=0.0167$ (fail, stop): it rejects the first two.
\end{example}

\subsection{Controlling the false discovery rate}

When thousands of hypotheses are screened, controlling FWER is too strict. The \emph{Benjamini--Hochberg} (BH) procedure controls the FDR.

\begin{theorem}[Benjamini--Hochberg]\label{thm:bh}
Order the $P$-values $p_{(1)}\le\dots\le p_{(m)}$ and let $k=\max\{i:p_{(i)}\le i\,q/m\}$. Rejecting $H_{(1)},\dots,H_{(k)}$ controls $\mathrm{FDR}\le q\,m_0/m\le q$ when the $P$-values of true nulls are independent (and under positive dependence), where $m_0$ is the number of true nulls.
\end{theorem}

For the five $P$-values above with $q=0.05$: the thresholds $iq/m=0.01,0.02,0.03,0.04,0.05$; $p_{(4)}=0.045>0.04$, $p_{(3)}=0.020\le0.03$, so $k=3$ and the first three hypotheses are rejected. BH is more powerful than Holm because it tolerates a controlled fraction of false discoveries.

\begin{interview}{adjusted $P$-values and what ``significant'' means}
A Bonferroni-adjusted $P$-value is $\min(1,mp_i)$; BH-adjusted values ($q$-values) are the smallest $q$ at which $H_i$ would be rejected. Use FWER control when a single false positive is costly (clinical confirmatory trials, a safety claim) and FDR control for exploratory screens where follow-up experiments will weed out false leads. Reporting only the best of many metrics without adjustment is a classic interview red flag.
\end{interview}

\section{The bootstrap}

The bootstrap approximates the sampling distribution of a statistic by resampling the observed data. Let $x_1,\dots,x_n$ be the data, $\hat F_n$ the empirical distribution (probability $1/n$ at each $x_i$) and $\hat\theta=T(x_1,\dots,x_n)$ the statistic of interest.

\begin{definition}[Nonparametric bootstrap]
A \emph{bootstrap sample} $x_1^*,\dots,x_n^*$ is drawn \emph{with replacement} from the data (i.e.\ i.i.d.\ from $\hat F_n$); $\hat\theta^*=T(x_1^*,\dots,x_n^*)$. Repeating $B$ times (typically $B=1000$--$10{,}000$) gives the bootstrap distribution of $\hat\theta^*$, used as an approximation of the unknown sampling distribution of $\hat\theta$.
\end{definition}

\begin{itemize}
  \item \textbf{Standard error:} $\se_{\mathrm{boot}}(\hat\theta)=$ the standard deviation of $\hat\theta^*_1,\dots,\hat\theta^*_B$ (Chapter~1).
  \item \textbf{Bias estimate:} $\overline{\hat\theta^*}-\hat\theta$.
  \item \textbf{Percentile interval:} the $\alpha/2$ and $1-\alpha/2$ quantiles of the bootstrap distribution (Chapter~2).
  \item \textbf{Basic (pivotal) interval:} $\big[2\hat\theta-q^*_{1-\alpha/2},\ 2\hat\theta-q^*_{\alpha/2}\big]$, which corrects for bias and skewness.
  \item \textbf{Bootstrap-$t$ interval:} resample the studentised statistic $(\hat\theta^*-\hat\theta)/\se^*$ and use its quantiles; it has better coverage than the percentile interval for the mean.
\end{itemize}

\begin{example}
For the data $4.1,\ 5.3,\ 3.9,\ 6.2,\ 4.8,\ 5.5,\ 4.4$ (median $4.8$), $B=2000$ bootstrap medians would take values among the data points themselves (the median of a bootstrap sample of seven values is always one of the seven), with the distribution concentrated on $4.4,\ 4.8,\ 5.3$. A percentile $95\%$ interval is read from the sorted bootstrap medians, for example $(4.1,\ 5.5)$. No formula for the standard error of a median was needed.
\end{example}

\begin{nuance}{when the bootstrap fails}
The bootstrap works for smooth statistics (means, regression coefficients, quantiles in the interior) with i.i.d.\ data. It fails or needs modification for extremes (the maximum of a sample, as in the uniform$(0,\theta)$ problem), for heavy-tailed data with infinite variance, for dependent data (use the block bootstrap) and for very small samples, where the data are a poor picture of the population. It also quantifies only sampling variability, not bias in the data collection.
\end{nuance}

\begin{theorem}[Why the bootstrap works for the mean]\label{thm:bootstrap-mean}
For the sample mean, conditional on the data, $\sqrt n(\bar{X}^*-\bar{x})$ has the same limiting distribution $\Normal(0,\sigma^2)$ as $\sqrt n(\bar{X}-\mu)$; moreover the bootstrap variance of $\bar{X}^*$ is exactly $\hat\sigma^2/n$ with $\hat\sigma^2=\frac1n\sum(x_i-\bar{x})^2$.
\end{theorem}

\begin{proof}[Idea]
Conditional on the data, the $x_i^*$ are i.i.d.\ with mean $\bar{x}$ and variance $\hat\sigma^2$, so $\Var(\bar{X}^*)=\hat\sigma^2/n$ exactly, and the CLT applies to them (with $\hat\sigma^2\to\sigma^2$).
\end{proof}

\section{Permutation tests}

A \emph{permutation test} gives an exact distribution-free test of the hypothesis that two groups have the \emph{same distribution}. Under $H_0$ the group labels are exchangeable, so all ways of reassigning the labels are equally likely.

\begin{definition}[Permutation test for two groups]
Choose a statistic $T$ (for instance the difference of means, $\bar{x}_1-\bar{x}_2$). Compute $T_{\mathrm{obs}}$. Repeatedly shuffle the pooled data into groups of the same sizes and recompute $T$. The $P$-value is the proportion of shuffles with $\abs{T}\ge\abs{T_{\mathrm{obs}}}$ (two-sided). With all $\binom{n_1+n_2}{n_1}$ shuffles the test is exact; in practice one uses a large random sample of shuffles ($P=(1+\#\{|T^*|\ge|T_{\mathrm{obs}}|\})/(B+1)$).
\end{definition}

\begin{example}
Groups: A: $12,15,11$ ($\bar{x}=12.67$) and B: $18,20,17$ ($\bar{x}=18.33$). Observed difference $T_{\mathrm{obs}}=-5.67$. There are $\binom63=20$ ways to choose which three of the six values form group A; only the two extreme assignments (A $=\{11,12,15\}$ and its mirror image A $=\{17,18,20\}$) give $\abs{T}\ge5.67$, so the exact two-sided $P=2/20=0.10$. A rank-based statistic gives the Wilcoxon rank-sum test (Chapter~9) as a permutation test on ranks.
\end{example}

Permutation tests make no distributional assumption and work for any statistic. They test exchangeability, which is stronger than equality of means; they are the natural tool for randomised experiments, where the randomisation itself justifies the reference distribution.

\section{A/B testing}

An A/B test is a randomised comparison of two versions (A control, B treatment) on a metric: conversion rate, revenue per user, latency. Statistically it is the two-sample problem of Chapter~4 (two proportions or two means), but practice adds several issues.

\subsection{Sample size and power}

For two proportions with baseline $p_1$ and a minimum detectable effect giving $p_2$, the required sample size \emph{per arm} for a two-sided test at level $\alpha$ with power $1-\beta$ is (Chapter~4)
\[ n\approx\frac{\Big(z_{\alpha/2}\sqrt{2\bar{p}\bar{q}}+z_\beta\sqrt{p_1q_1+p_2q_2}\Big)^2}{(p_2-p_1)^2},\qquad\bar{p}=\tfrac{p_1+p_2}2. \]
A useful rule of thumb for $\alpha=0.05$ and power $0.8$: $n\approx16\,\sigma^2/\Delta^2$ per arm for means with standard deviation $\sigma$ and effect $\Delta$, and $n\approx16\,\bar{p}\bar{q}/\Delta^2$ for proportions. Because $n$ scales as $1/\Delta^2$, detecting an effect half as large needs four times the traffic.

\begin{example}
Baseline conversion $p_1=10\%$; the smallest change worth detecting is $+1$ percentage point ($p_2=11\%$). The rule of thumb gives $n\approx16(0.105)(0.895)/0.01^2=15{,}000$ users per arm; the full formula gives about $14{,}750$.
\end{example}

\subsection{Pitfalls}

\begin{itemize}
  \item \textbf{Peeking and optional stopping.} Checking the $P$-value daily and stopping when it first falls below $0.05$ inflates the Type~I error far above $5\%$. Fix the sample size in advance, or use a \emph{sequential} design with adjusted boundaries (for example the Pocock or O'Brien--Fleming group-sequential boundaries) or an always-valid method.
  \item \textbf{Multiple metrics and segments.} Declare one primary metric in advance; correct the secondary ones (see the first section of this chapter).
  \item \textbf{Sample ratio mismatch.} If the observed split of users between arms differs from the planned split (tested with a chi-square test), the randomisation or logging is broken and the results cannot be trusted.
  \item \textbf{Novelty and seasonality effects,} interference between users (network effects), and non-independent units (several observations per user: randomise and analyse at the user level).
  \item \textbf{Statistical versus practical significance:} always report the estimated effect with a confidence interval, not just $P$.
\end{itemize}

\subsection{Variance reduction}

A pre-experiment covariate $X$ (for example the user's activity before the test) correlated with the outcome $Y$ can be used to reduce the variance: the \emph{CUPED} estimator replaces $Y$ by $Y-\theta(X-\bar{X})$ with $\theta=\Cov(X,Y)/\Var X$; the adjusted outcome has variance $(1-\rho^2)\Var Y$, where $\rho=\Corr(X,Y)$. With $\rho=0.7$ the variance halves, so half the sample size suffices. This is the regression-adjustment idea of Chapter~6 applied to the treatment effect.

\begin{summarybox}
\begin{itemize}
  \item $m$ true nulls tested at $\alpha$ give $\Prob(\text{some false positive})=1-(1-\alpha)^m$. Bonferroni ($\alpha/m$) and Holm control the FWER; Benjamini--Hochberg ($p_{(i)}\le iq/m$) controls the FDR.
  \item The bootstrap resamples the data with replacement to estimate standard errors, bias and intervals (percentile, basic, bootstrap-$t$) for any statistic; it fails for extremes and dependent data without modification.
  \item Permutation tests reshuffle group labels to get exact distribution-free $P$-values; they underpin randomisation inference.
  \item A/B tests: plan $n\approx16\sigma^2/\Delta^2$ per arm, avoid peeking, pre-specify the primary metric, check for sample ratio mismatch, report effect sizes with intervals, and use covariate adjustment (CUPED) to gain power.
\end{itemize}
\end{summarybox}

\section{Exercises}

\begin{exercise}
Twelve $P$-values are $0.001,\ 0.004,\ 0.008,\ 0.012,\ 0.019,\ 0.027,\ 0.041,\ 0.050,\ 0.12,\ 0.24,\ 0.51,\ 0.88$. Find the rejections under unadjusted testing at $0.05$, Bonferroni, Holm and Benjamini--Hochberg at $q=0.05$.
\end{exercise}

\begin{exercise}
Show that the probability of at least one false positive among $m$ independent tests is $1-(1-\alpha)^m$ and that Bonferroni's threshold $\alpha/m$ gives FWER at most $\alpha$ also when the tests are dependent.
\end{exercise}

\begin{exercise}
For the data $7,\ 9,\ 12,\ 14,\ 18$, perform a bootstrap by hand with $B=5$ resamples (use your own random draws) to estimate the standard error of the mean and compare with $s/\sqrt n$.
\end{exercise}

\begin{exercise}
Carry out an exact permutation test for the difference of means with groups $\{3,5,8\}$ and $\{9,11,14\}$ and give the two-sided $P$-value.
\end{exercise}

\begin{exercise}
An A/B test has baseline conversion $4\%$. How many users per arm are needed to detect a relative improvement of $10\%$ ($4\%\to4.4\%$) with power $0.8$ at $\alpha=0.05$? What if the improvement is $20\%$ relative?
\end{exercise}

\begin{exercise}
Simulate the effect of peeking: draw two groups from the same normal distribution, test after every $10$ observations per group up to $200$, and record how often some $P$-value falls below $0.05$. Explain your finding.
\end{exercise}

\begin{exercise}
Show that the CUPED adjustment $Y-\theta(X-\bar{X})$ is unbiased for the treatment effect when $X$ is measured before randomisation, and derive its variance reduction factor $1-\rho^2$.
\end{exercise}

\begin{exercise}
Explain why the bootstrap fails for the sample maximum of a uniform$(0,\theta)$ sample: what is the probability that the bootstrap maximum equals the observed maximum?
\end{exercise}

\chapter{Supplement: Generalised Linear Models and Model Assessment}

\begin{nuance}{place in the course}
This supplementary chapter extends the regression of Chapters~5--6 in two directions the textbook does not cover.
\begin{itemize}
  \item Responses that are not normal and continuous (yes/no outcomes, counts) are handled by \emph{generalised linear models}. Logistic regression is the most important.
  \item We ask how well a fitted model will predict \emph{new} data. This leads to cross-validation, information criteria and regularisation (lasso and ridge).
\end{itemize}
\end{nuance}

\section{Why ordinary regression is not enough}

Suppose the response is binary, $Y\in\{0,1\}$ (a customer converts or not). The linear model $\Exp[Y\mid x]=\beta_0+\beta_1x$ has three problems:
\begin{itemize}
  \item it can predict probabilities outside $[0,1]$;
  \item the errors have non-constant variance ($\Var[Y]=p(1-p)$ depends on the mean);
  \item the errors are far from normal.
\end{itemize}
Counts have the same trouble: negative predictions, and variance that grows with the mean.

The fix: model a \emph{transformation} of the mean as a linear function of the regressors, and use the correct distribution for the response.

\section{Generalised linear models}

\begin{definition}[Generalised linear model]
A generalised linear model (GLM) has three parts:
\begin{enumerate}[label=(\arabic*)]
  \item a \emph{random component}: independent $Y_i$ from an exponential-family distribution (normal, Bernoulli/binomial, Poisson, gamma, \dots) with mean $\mu_i=\Exp[Y_i]$;
  \item a \emph{linear predictor} $\eta_i=\beta_0+\beta_1x_{i1}+\dots+\beta_kx_{ik}=\mathbf{x}_i\T\boldsymbol\beta$;
  \item a \emph{link function} $g$ with $g(\mu_i)=\eta_i$.
\end{enumerate}
\end{definition}

\begin{center}
\begin{tabular}{llll}
\toprule
Response & Distribution & Canonical link $g(\mu)$ & Inverse link\\
\midrule
continuous & normal & identity, $\mu$ & $\eta$\\
binary & Bernoulli & logit, $\ln\frac\mu{1-\mu}$ & $\dfrac1{1+e^{-\eta}}$\\
counts & Poisson & log, $\ln\mu$ & $e^\eta$\\
positive continuous & gamma & inverse, $1/\mu$ & $1/\eta$\\
\bottomrule
\end{tabular}
\end{center}

Ordinary multiple regression (Chapter~6) is the GLM with a normal response and identity link.

\section{Logistic regression}

For $Y_i\sim\mathrm{Bernoulli}(p_i)$, the \emph{logistic regression} model is
\[ \ln\frac{p_i}{1-p_i}=\mathbf{x}_i\T\boldsymbol\beta,\qquad\text{that is,}\qquad p_i=\frac1{1+e^{-\mathbf{x}_i\T\boldsymbol\beta}}. \]
$p/(1-p)$ is the \emph{odds} of success. $\ln\frac p{1-p}$ is the \emph{log-odds} (logit).

\begin{theorem}[Interpretation of the coefficients]\label{thm:logit-interp}
A one-unit increase in $x_j$, with the other regressors fixed, adds $\beta_j$ to the log-odds. It \emph{multiplies the odds by} $e^{\beta_j}$, the \emph{odds ratio}. For a binary regressor, $e^{\beta_j}$ is the odds in the group with $x_j=1$ divided by the odds in the group with $x_j=0$.
\end{theorem}

\begin{example}
In a conversion model, $\hat\beta_1=0.70$ for the indicator ``saw the new design''.
\begin{itemize}
  \item Odds ratio: $e^{0.70}=2.01$. The odds of converting are about twice as high.
  \item If the baseline probability is $10\%$ (odds $1/9$), the new odds are $2/9$, so $p=2/11=18\%$.
  \item The effect on the \emph{probability} depends on the baseline. The effect on the odds is constant.
\end{itemize}
\end{example}

\subsection{Maximum likelihood estimation}

For independent Bernoulli observations, the log-likelihood is $\ell(\boldsymbol\beta)=\sum_i\big[y_i\eta_i-\ln(1+e^{\eta_i})\big]$ with $\eta_i=\mathbf{x}_i\T\boldsymbol\beta$.

\begin{theorem}[Score and Hessian for logistic regression]\label{thm:logit-score}
With $p_i=1/(1+e^{-\eta_i})$,
\[ \nabla\ell(\boldsymbol\beta)=X\T(\mathbf{y}-\mathbf{p}),\qquad\nabla^2\ell(\boldsymbol\beta)=-X\T WX,\qquad W=\diag\big(p_i(1-p_i)\big). \]
$X\T WX$ is positive semidefinite (positive definite when $X$ has full column rank and $0<p_i<1$). So $\ell$ is concave, and the MLE, when it exists, is unique.
\end{theorem}

\begin{proof}
$\partial\ell/\partial\beta_j=\sum_i(y_i-p_i)x_{ij}$, since $\partial\ln(1+e^{\eta})/\partial\eta=p$.

Differentiate again: $\partial p_i/\partial\beta_k=p_i(1-p_i)x_{ik}$. This gives the Hessian $-\sum_ip_i(1-p_i)\mathbf{x}_i\mathbf{x}_i\T=-X\T WX$.

For any $\mathbf{v}$: $\mathbf{v}\T X\T WX\mathbf{v}=\sum_iw_i(\mathbf{x}_i\T\mathbf{v})^2\ge0$.
\end{proof}

The likelihood equations $X\T(\mathbf{y}-\mathbf{p})=\mathbf{0}$ have no closed form. Solve them by \emph{Newton's method}, which here is called \emph{iteratively reweighted least squares} (IRLS):
\[ \boldsymbol\beta^{(t+1)}=\boldsymbol\beta^{(t)}+(X\T W^{(t)}X)^{-1}X\T(\mathbf{y}-\mathbf{p}^{(t)})=(X\T W^{(t)}X)^{-1}X\T W^{(t)}\mathbf{z}^{(t)}. \]
The \emph{working response} is $z_i=\eta_i+(y_i-p_i)/(p_i(1-p_i))$. Each iteration is a weighted least squares fit (Chapter~6), with weights $W^{(t)}$ updated from the current fit.

\begin{pitfall}{perfect separation}
Suppose some combination of the regressors separates the $0$'s from the $1$'s perfectly. Then the likelihood keeps increasing as the coefficients grow. The MLE does not exist: some $\hat\beta_j\to\pm\infty$, and software warns about convergence.

Remedies: penalised (ridge-type) logistic regression, a prior (Bayesian logistic regression), or dropping the offending regressor.
\end{pitfall}

\subsection{Inference}

By Theorem~\ref{thm:mle-asym}, in large samples $\hat{\boldsymbol\beta}\approx\Normal\big(\boldsymbol\beta,(X\T WX)^{-1}\big)$, with $W$ evaluated at $\hat{\boldsymbol\beta}$. So:
\begin{itemize}
  \item \textbf{Wald test} for $H_0:\beta_j=0$: $z=\hat\beta_j/\se(\hat\beta_j)$, with $\se(\hat\beta_j)=\sqrt{[(X\T WX)^{-1}]_{jj}}$, compared with $\Normal(0,1)$. A $95\%$ interval for the odds ratio is $\exp\big(\hat\beta_j\pm1.96\,\se(\hat\beta_j)\big)$.
  \item \textbf{Likelihood ratio test} for nested models: $-2\big(\ell_{\mathrm{reduced}}-\ell_{\mathrm{full}}\big)\approx\chi^2_r$ under $H_0$, where $r$ is the number of coefficients set to zero (Wilks' theorem, Chapter~3). This is the GLM version of the extra-sum-of-squares $F$-test.
\end{itemize}

\begin{definition}[Deviance]
The \emph{deviance} is $D=2\big(\ell_{\mathrm{saturated}}-\ell(\hat{\boldsymbol\beta})\big)$. It is twice the log-likelihood gap to a saturated model that fits each observation exactly. For ungrouped Bernoulli data, $D=-2\ell(\hat{\boldsymbol\beta})$. It generalises the residual sum of squares: the change in deviance between nested models is the likelihood ratio statistic.
\end{definition}

\begin{example}[Likelihood ratio test]
Adding two regressors lowers the deviance from $412.6$ to $405.1$. The statistic is $7.5$. Compare with $\chi^2_{0.05,2}=5.99$. The two regressors are jointly significant at $5\%$ ($P=0.023$).
\end{example}

\subsection{Classification and evaluation}

A fitted logistic model gives probabilities $\hat{p}_i$. A decision rule calls a case positive if $\hat{p}_i>c$. The usual choice is $c=0.5$, but the threshold should reflect the costs of the two kinds of error.

Ways to measure performance:
\begin{itemize}
  \item The confusion matrix: sensitivity (true positive rate) and specificity (true negative rate).
  \item The \emph{ROC curve}: true positive rate against false positive rate as $c$ varies.
  \item The area under it (\emph{AUC}): the probability that a random positive has a higher score than a random negative.
  \item \emph{Calibration}: among cases with $\hat{p}\approx0.3$, are about $30\%$ positive?
\end{itemize}

\section{Poisson regression}

For counts $Y_i\sim\mathrm{Poisson}(\mu_i)$ with $\ln\mu_i=\mathbf{x}_i\T\boldsymbol\beta$:
\begin{itemize}
  \item An \emph{offset} $\ln t_i$ (exposure time or population size) models a \emph{rate}.
  \item The score is $X\T(\mathbf{y}-\boldsymbol\mu)$. The information matrix is $X\T WX$ with $W=\diag(\mu_i)$.
  \item $e^{\beta_j}$ is the factor by which the expected count is multiplied for a unit increase in $x_j$ (a \emph{rate ratio}).
  \item The Poisson model forces $\Var[Y]=\mu$. If the data are \emph{overdispersed} ($\Var[Y]>\mu$, estimated by Pearson's statistic over its degrees of freedom), use a quasi-Poisson or negative binomial model. Otherwise standard errors are too small and tests too liberal.
\end{itemize}

\section{Model assessment: how well will the model predict?}

A model fitted to data fits those data better than it fits new data, because it adapts to the noise.

\begin{definition}[Training error and test error]
\begin{itemize}
  \item The \emph{training error} is the average loss on the data used to fit the model.
  \item The \emph{test (generalisation) error} is the expected loss on new observations from the same population.
\end{itemize}
For squared error and a fitted function $\hat{f}$, $\mathrm{Err}=\Exp[(Y_0-\hat{f}(\mathbf{x}_0))^2]$ for a new point $(\mathbf{x}_0,Y_0)$.
\end{definition}

\begin{theorem}[Bias--variance decomposition]\label{thm:biasvar-pred}
Let $Y_0=f(\mathbf{x}_0)+\varepsilon_0$ with $\Exp[\varepsilon_0]=0$ and $\Var[\varepsilon_0]=\sigma^2$, independent of the training data. Then
\[ \Exp\big[(Y_0-\hat{f}(\mathbf{x}_0))^2\big]=\underbrace{\sigma^2}_{\text{irreducible noise}}+\underbrace{\big(\Exp[\hat{f}(\mathbf{x}_0)]-f(\mathbf{x}_0)\big)^2}_{\text{bias}^2}+\underbrace{\Var[\hat{f}(\mathbf{x}_0)]}_{\text{variance}}. \]
\end{theorem}

\begin{proof}
Add and subtract $f(\mathbf{x}_0)$ and $\Exp[\hat{f}(\mathbf{x}_0)]$, and expand. The cross terms vanish because $\varepsilon_0$ has mean zero and is independent of $\hat{f}$, and $\Exp[\hat{f}-\Exp[\hat{f}]]=0$. This is the decomposition $\MSE=\Var+\bias^2$ of Chapter~1 applied to prediction, plus the unavoidable noise $\sigma^2$.
\end{proof}

\begin{itemize}
  \item More flexible models (more regressors, higher-degree polynomials, weaker regularisation) have lower bias and higher variance.
  \item Simple models have high bias and low variance.
  \item Test error is lowest at an intermediate complexity. This is the bias--variance trade-off.
  \item Training error keeps falling as complexity grows, so it is useless for choosing a model.
\end{itemize}

\begin{theorem}[Optimism of the training error]\label{thm:optimism}
For a least squares fit with $p$ parameters and error variance $\sigma^2$, the expected test error on new responses at the \emph{same} regressor values exceeds the expected training error by $2p\sigma^2/n$ per observation:
\[ \Exp[\mathrm{Err}_{\mathrm{in}}]=\Exp[\mathrm{err}_{\mathrm{train}}]+\frac{2p}n\sigma^2. \]
\end{theorem}

\begin{proof}
For the hat matrix $H$ (Chapter~6), $\hat{\mathbf{y}}=H\mathbf{y}$ and $\sum_i\Cov(\hat{y}_i,y_i)=\sigma^2\tr H=p\sigma^2$. The standard identity $\Exp[\mathrm{Err}_{\mathrm{in}}-\mathrm{err}_{\mathrm{train}}]=\frac2n\sum\Cov(\hat{y}_i,y_i)$ gives the result.
\end{proof}

This motivates Mallows' $C_p=\mathrm{err}_{\mathrm{train}}\cdot n/\hat\sigma^2+2p-n$ (Chapter~6) and the information criteria below.

\begin{definition}[AIC and BIC]
For a model with $p$ parameters and maximised log-likelihood $\hat\ell$ on $n$ observations,
\[ \mathrm{AIC}=-2\hat\ell+2p,\qquad\mathrm{BIC}=-2\hat\ell+p\ln n. \]
Smaller is better.
\begin{itemize}
  \item AIC aims at predictive accuracy (it estimates out-of-sample log-loss).
  \item BIC penalises complexity more for $n>7$. It is consistent for picking the true model among those considered, when one of them is true.
\end{itemize}
\end{definition}

\section{Cross-validation}

Information criteria rely on distributional assumptions. \emph{Cross-validation} (CV) estimates the test error directly, by repeatedly holding out part of the data.

\begin{definition}[$K$-fold cross-validation]
Split the data randomly into $K$ roughly equal folds. For $k=1,\dots,K$: fit the model on the other $K-1$ folds and compute the loss on fold $k$. The CV estimate is the average held-out loss over all $n$ observations. $K=5$ or $10$ is typical. $K=n$ is \emph{leave-one-out} cross-validation (LOOCV).
\end{definition}

\begin{theorem}[LOOCV shortcut for linear regression]\label{thm:loocv}
For least squares with hat-matrix diagonal $h_{ii}$ and residuals $e_i$, the leave-one-out prediction error for observation $i$ is $e_i/(1-h_{ii})$. So
\[ \mathrm{LOOCV}=\frac1n\sum_{i=1}^n\Big(\frac{e_i}{1-h_{ii}}\Big)^2=\frac{PRESS}n, \]
from a single fit.
\end{theorem}

This is the PRESS statistic of Chapter~6. Practical rules:
\begin{itemize}
  \item Do any data-dependent preprocessing, feature selection or tuning \emph{inside} each training fold. Otherwise the CV estimate is optimistic.
  \item For time-ordered data, use forward-chaining splits, not random folds.
  \item For grouped data (several rows per user), split by group.
\end{itemize}

\begin{interview}{training, validation and test sets}
A clean workflow:
\begin{itemize}
  \item Training data: fit models.
  \item Validation set (or CV): choose among models and tune hyperparameters.
  \item A final untouched \emph{test set}, used once: report the performance of the chosen model.
\end{itemize}
Reusing the test set to pick models turns it into a validation set and gives optimistic estimates.
\end{interview}

\section{Regularisation: ridge and lasso}

When $p$ is large relative to $n$, least squares overfits. Penalising the size of the coefficients trades a little bias for a big drop in variance (Theorem~\ref{thm:biasvar-pred}).

\begin{definition}[Ridge and lasso]
For $\lambda\ge0$ (regressors standardised, intercept unpenalised):
\begin{align*}
\hat{\boldsymbol\beta}^{\mathrm{ridge}}&=\arg\min_{\boldsymbol\beta}\ \tfrac12\norm{\mathbf{y}-X\boldsymbol\beta}^2+\tfrac\lambda2\norm{\boldsymbol\beta}_2^2,\\
\hat{\boldsymbol\beta}^{\mathrm{lasso}}&=\arg\min_{\boldsymbol\beta}\ \tfrac12\norm{\mathbf{y}-X\boldsymbol\beta}^2+\lambda\norm{\boldsymbol\beta}_1.
\end{align*}
\end{definition}

\begin{itemize}
  \item Ridge: $(X\T X+\lambda I)^{-1}X\T\mathbf{y}$ (Chapter~6; Linear Algebra, Chapter~17). Coefficients shrink smoothly but are never exactly zero.
  \item Lasso: no closed form in general. For an \emph{orthonormal} design ($X\T X=I$) it is explicit.
\end{itemize}

\begin{theorem}[Lasso with orthonormal design]\label{thm:lasso}
If $X\T X=I$, and $\hat\beta_j^{\mathrm{OLS}}$ are the least squares coefficients, then
\[ \hat\beta_j^{\mathrm{lasso}}=\sgn\big(\hat\beta_j^{\mathrm{OLS}}\big)\big(\abs{\hat\beta_j^{\mathrm{OLS}}}-\lambda\big)_+\quad\text{(soft thresholding)},\qquad\hat\beta_j^{\mathrm{ridge}}=\frac{\hat\beta_j^{\mathrm{OLS}}}{1+\lambda}. \]
\end{theorem}

\begin{proof}
With $X\T X=I$, $\tfrac12\norm{\mathbf{y}-X\boldsymbol\beta}^2=\text{const}+\tfrac12\sum_j(\beta_j-\hat\beta_j^{\mathrm{OLS}})^2$. The problem splits into $p$ scalar problems $\min_\beta\tfrac12(\beta-b)^2+\lambda\abs{\beta}$.

Their solution is $\sgn(b)(\abs{b}-\lambda)_+$. For $\beta>0$ the derivative $\beta-b+\lambda$ vanishes at $b-\lambda$, valid when $b>\lambda$. And $\beta=0$ is optimal exactly when $\abs{b}\le\lambda$.

The ridge case differentiates $\tfrac12(\beta-b)^2+\tfrac\lambda2\beta^2$.
\end{proof}

So:
\begin{itemize}
  \item The lasso sets small coefficients exactly to zero (\emph{variable selection}) and shrinks the others by the constant $\lambda$.
  \item Ridge shrinks every coefficient by the same factor $1/(1+\lambda)$.
  \item Geometrically, the contours of the squared error first touch the diamond $\norm{\boldsymbol\beta}_1\le t$ at a corner, where some coordinates vanish. The round ball $\norm{\boldsymbol\beta}_2\le t$ has no corners (Linear Algebra, Chapter~1).
  \item Choose $\lambda$ by cross-validation. It is the complexity knob of the bias--variance trade-off.
  \item The same penalties work for logistic and other GLMs: add them to the negative log-likelihood.
\end{itemize}

\begin{summarybox}
\begin{itemize}
  \item A GLM is an exponential-family response, a linear predictor $\eta=\mathbf{x}\T\boldsymbol\beta$, and a link $g(\mu)=\eta$. Logistic regression: logit link, $e^{\beta_j}$ is an odds ratio. Poisson regression: log link, $e^{\beta_j}$ is a rate ratio.
  \item Logistic MLE: score $X\T(\mathbf{y}-\mathbf{p})$, Hessian $-X\T WX$ (concave), solved by IRLS. Use Wald and likelihood-ratio tests. Deviance generalises the residual sum of squares. Beware perfect separation and overdispersion.
  \item Test error $=\sigma^2+\text{bias}^2+\text{variance}$. Training error is optimistic by $2p\sigma^2/n$. Compare models by AIC, BIC or cross-validation ($K$-fold; LOOCV $=PRESS/n$ for linear regression).
  \item Ridge shrinks smoothly. Lasso soft-thresholds and selects variables. Tune $\lambda$ by cross-validation, with all preprocessing inside the folds and a final untouched test set.
\end{itemize}
\end{summarybox}

\section{Exercises}

\begin{exercise}
In a logistic regression, $\hat\beta_0=-2.0$ and $\hat\beta_1=0.5$ for a regressor $x$. Compute the estimated probability of success at $x=0,2,4$ and the odds ratio for a unit increase in $x$.
\end{exercise}

\begin{exercise}
Derive the score and the Hessian of the Poisson regression log-likelihood $\ell(\boldsymbol\beta)=\sum_i[y_i\mathbf{x}_i\T\boldsymbol\beta-e^{\mathbf{x}_i\T\boldsymbol\beta}-\ln y_i!]$, and write one IRLS update.
\end{exercise}

\begin{exercise}
A $2\times2$ table gives: exposed, $30$ cases of $100$; unexposed, $15$ cases of $100$. Compute the odds ratio and show that it equals $e^{\hat\beta_1}$ for the logistic regression on the exposure indicator. Give a $95\%$ confidence interval using the standard error $\sqrt{1/a+1/b+1/c+1/d}$ of $\ln\widehat{OR}$.
\end{exercise}

\begin{exercise}
Two models for $n=200$ observations have maximised log-likelihoods $-310.2$ ($p=3$) and $-305.9$ ($p=6$). Compute AIC and BIC for each; which would each criterion choose, and does the likelihood ratio test agree?
\end{exercise}

\begin{exercise}
Prove Theorem~\ref{thm:loocv} for simple linear regression through the origin ($\hat\beta=\sum x_iy_i/\sum x_i^2$) and check it numerically for the data $(1,2),(2,3),(3,5)$.
\end{exercise}

\begin{exercise}
Show that with orthonormal design the lasso solution is zero exactly when $\abs{\hat\beta_j^{\mathrm{OLS}}}\le\lambda$, and compute the lasso and ridge coefficients for $\hat\beta^{\mathrm{OLS}}=(3.0,\ 0.8,\ -1.5,\ 0.2)$ and $\lambda=1$.
\end{exercise}

\begin{exercise}
Explain why the training error always decreases as polynomial degree increases while the cross-validation error eventually increases, and sketch both curves. What polynomial degree would you choose from a CV curve that is flat beyond degree $3$?
\end{exercise}

\begin{exercise}
A classifier scores $5$ positives $0.9,0.8,0.55,0.4,0.3$ and $5$ negatives $0.7,0.5,0.35,0.2,0.1$. Compute the AUC as the probability that a random positive outscores a random negative, and the sensitivity and specificity at threshold $0.5$.
\end{exercise}


\end{document}
